% \iffalse meta-comment
%%
%% ----------------------------------------------------------------
%% Copyright (C) 2003-2021 by Dr. Uwe Kern <xcolor at ukern dot de>
%%               2021-2024    LaTeX Project https://github.com/latex3/xcolor/issues
%% ----------------------------------------------------------------
%%
%<*package>
%% Please send error reports and suggestions for
%% improvements to https://github.com/latex3/xcolor/issues.
%%
%% This package contains code copied from
%%
%% color.dtx [1999/02/16 v1.0i Standard LaTeX Color (DPC)]
%% Copyright (C) 1994--1999 David Carlisle
%%
%% which is part of the Standard LaTeX `Graphics Bundle'.
%%
%</package>
%<*extract|driver|test1|test2|test3|test4>
\def\XCfileversion{v3.02}%
\def\XCfiledate{2024/09/29}%
%</extract|driver|test1|test2|test3|test4>
%
% This file may be distributed and/or modified under the
% conditions of the LaTeX Project Public License, either version 1.3c
% of this license or (at your option) any later version.
% The latest version of this license is in:
%
%    http://www.latex-project.org/lppl.txt
%
% and version 1.3c or later is part of all distributions of LaTeX
% version 2008-05-04 or later.
%
% \fi
%
%
% \DoNotIndex{\def,\gdef,\edef,\xdef,\long,\let}
% \DoNotIndex{\ifnum,\ifdim,\iftrue,\iffalse,\ifx,\ifcase,\else,\or,\fi}
% \DoNotIndex{\bgroup,\egroup,\begingroup,\endgroup,\begin,\end}
% \DoNotIndex{\relax,\endinput}
% \DoNotIndex{\csname,\endcsname,\string,\the,\noexpand,\expandafter}
% \DoNotIndex{\advance,\count,\dimen}
% \DoNotIndex{\m@ne,\z@,\@ne,\tw@,\p@,\@@tmp}
%
% \CharacterTable
%  {Upper-case    \A\B\C\D\E\F\G\H\I\J\K\L\M\N\O\P\Q\R\S\T\U\V\W\X\Y\Z
%   Lower-case    \a\b\c\d\e\f\g\h\i\j\k\l\m\n\o\p\q\r\s\t\u\v\w\x\y\z
%   Digits        \0\1\2\3\4\5\6\7\8\9
%   Exclamation   \!     Double quote  \"     Hash (number) \#
%   Dollar        \$     Percent       \%     Ampersand     \&
%   Acute accent  \'     Left paren    \(     Right paren   \)
%   Asterisk      \*     Plus          \+     Comma         \,
%   Minus         \-     Point         \.     Solidus       \/
%   Colon         \:     Semicolon     \;     Less than     \<
%   Equals        \=     Greater than  \>     Question mark \?
%   Commercial at \@     Left bracket  \[     Backslash     \\
%   Right bracket \]     Circumflex    \^     Underscore    \_
%   Grave accent  \`     Left brace    \{     Vertical bar  \|
%   Right brace   \}     Tilde         \~}
%
%
% \iffalse
%<*batchfile>
\IfFileExists{xcolor.ins}{\let\uknext=\relax}{\let\uknext=\stop}
\begin{filecontents*}{xcolor.ins}
%% Package `xcolor' by Dr. Uwe Kern
%%
%% Extract all package files:
%%   `latex xcolor.ins'
%%
%% Put the generated files to their respective locations
%% within the TeX installation:
%%   *.sty -> /tex/latex/xcolor/
%%   *.def -> /tex/latex/xcolor/
%%   *.pro -> /dvips/xcolor/
%%
%% Create the documentation:
%%   `latex xcolor.dtx'
%%   `latex xcolor.dtx'
%%   `makeindex -s gind.ist xcolor.idx'
%%   `latex xcolor.dtx'
%%   `latex xcolor.dtx'
%%
%% Run the test files (N=1,2,...):
%%   `latex xcolorN.tex'
%%   `latex xcolorN.tex'
%%
\def\batchfile{xcolor.ins}
\input docstrip.tex
\askforoverwritefalse
\generate{\file{xcolor.lox}{\from{xcolor.dtx}{extract}}}
\batchinput{xcolor.lox}
\endbatchfile
\end{filecontents*}
%</batchfile>
%
%<*driver>
\listfiles
\ProvidesFile{xcolor.dtx}[\XCfiledate\space\XCfileversion\space
                          `xcolor' source and documentation (UK)]
\PassOptionsToPackage
 {
   hyperindex=false,bookmarks,bookmarksopen,bookmarksopenlevel=1,%
   pdftitle={xcolor \XCfileversion{} (\XCfiledate)},
   pdfauthor={Dr. Uwe Kern},%
   pdfsubject={Color extensions for LaTeX and pdfLaTeX},%
   pdfkeywords={xcolor,color,colour,model,tint,tone,shade,harmony,spot,latex,pdftex,dvips,%
   conversion,blend,mix,mask,separation,rgb,cmy,cmyk,hsb,gray,html,wave,thsb,wheel}
 }{hyperref}
\documentclass{ltxdoc}
\usepackage[UTF8,scheme=plain,fontset=none]{ctex}
\setCJKmainfont{Source Han Serif SC}
\setCJKsansfont{Source Han Sans SC}
\setCJKmonofont{Source Han Sans SC}
\xeCJKsetup{CJKmath=true}
\usepackage[table,showerrors,dvipsnames,hyperref]{xcolor}[2005/06/06]
\usepackage{chngpage,varioref,tabularx,amsmath,amsthm}
\OnlyDescription
\EnableCrossrefs
\CodelineIndex
\makeatletter
\def\Meaning#1{\texttt{\expandafter\expandafter\expandafter\strip@prefix\expandafter\meaning\csname #1\endcsname}}
\makeatother
\pagestyle{myheadings}
\markright{\hskip-2.5cm xcolor \XCfileversion\space(\XCfiledate)}
\begin{document}
 \DocInput{xcolor-CN.dtx}
 \let\Section\section\def\section*#1{\Section*{#1}\addcontentsline{toc}{section}{#1}}
 \begin{adjustwidth}{-2.5cm}{}
 \def\markboth#1#2{}\PrintIndex
 \end{adjustwidth}
\end{document}
%</driver>
%
%<*extract>
\def\ukfile#1#2#3#4%
 {\def\ukhead
   {\ifnum#1>2 \perCent!^^J\fi
    \defaultpreamble^^J\MetaPrefix^^J%
    \ifcase#1
      \string\ProvidesPackage{#2}^^J\or
      \string\ProvidesFile{#2}^^J\or
      \string\ProvidesFile{#2.def}^^J\or
      \MetaPrefix\space #2.pro (PostScript/dvips header file)^^J\MetaPrefix\space
    \fi
    \space[\XCfiledate\space\XCfileversion\space#4\space(UK)]}
  \usepreamble\ukhead
  \file{#2.\ifcase#1sty\or tex\or def\or pro\fi}{\from{xcolor.dtx}{#3}}}
\generate{
\ukfile0{xcolor}{package}{LaTeX color extensions}
\ukfile2{svgnam}{svgnames}{Predefined colors according to SVG 1.1}
\ukfile2{x11nam}{x11names}{Predefined colors according to Unix/X11}
\ukfile1{xcolor1}{test1}{Color conversion test}
\ukfile1{xcolor2}{test2}{Pstricks color test}
\ukfile1{xcolor3}{test3}{Color logging test}
\ukfile1{xcolor4}{test4}{Color driver test}
\def\PStail{\MetaPrefix\space End of file `\outFileName'.}
\usepostamble\PStail
\ukfile3{xcolor}{dvips}{LaTeX color extensions}
}
%</extract>
% \fi
%
%
%
% \iffalse
%-------------------------------------------------------------------------------
% \fi
%
%
% \makeatletter
%
% \def\testclr#1#{\@testclr{#1}}
% \def\@testclr#1#2{{\fboxsep\z@\fbox{\colorbox#1{#2}{\phantom{XX}}}}}
% \def\Testclr#1#{\@Testclr{#1}}
% \def\@Testclr#1#2#3{\testclr#1{#2}~\rlap{\Color[-]{#3}}\\}
% \def\TestClr#1#{\@TestClr{#1}}
% \def\@TestClr#1#2#3{\testclr#1{#2}~\rlap{\Color[+]{#3}}\\}
% \newif\if@Meta
% \DeclareRobustCommand*\Meta{\@ifstar{\@Metatrue\@Meta}{\@Metafalse\@Meta}}
% \newcommand*\@Meta[2][]{\meta{#2}\@ifxempty{#1}{}{\ensuremath{\mskip-1mu\rule[-.4ex]{0pt}{1ex}_{#1}}}\if@Meta\index{arguments\levelchar\meta{#2}\string|usage}\fi}
% \def\DescribeArg#1{\leavevmode\@bsphack\marginpar{\raggedleft\Meta*{#1}}\@esphack\ignorespaces}
% \def\DescribeFunc#1{\leavevmode\@bsphack\marginpar{\raggedleft\Function{#1}\strut}\@esphack\ignorespaces}
%
% \makeatother
%
% \def\Index#1{#1\index{#1\string|usage}}
% \def\Package#1{\textsf{#1}\index{packages\levelchar\textsf{#1}\string|usage}}
% \def\Program#1{\textsf{\itshape#1}\index{programs\levelchar\textsf{\itshape#1}\string|usage}}
% \def\Model#1{\textsf{\bfseries#1}\index{color models\levelchar\textsf{\bfseries#1}\string|usage}}
% \def\nModel#1{`#1'\index{color models\levelchar`#1'\string|usage}}
% \def\pModel#1{\textbf{\itshape #1}}
% \def\People#1#2{#1 #2\index{people\levelchar#2, #1\string|usage}}
% \newcommand*\Color[2][+]{\textsl{#2}\ifx#1+\index{color names\levelchar\textsl{#2}\string|usage}\fi}
% \def\Option#1{\texttt{#1}\index{package options\levelchar\texttt{#1}\string|usage}}
% \def\oOption#1{\texttt{#1}\index{package options (obsolete)\levelchar\texttt{#1}\string|usage}}
% \def\Key#1{\textit{#1}\index{keys\levelchar\textit{#1}\string|usage}}
% \def\File#1{\texttt{#1}\index{files\levelchar\texttt{#1}\string|usage}}
% \def\Function#1{\textit{#1}\index{color functions\levelchar\textit{#1}\string|usage}}
% \def\Op#1#2{#1\nolimits_{\,\Model{#2}}}
% \def\med{\mathop{\textrm{med}}}
% \def\mix{\mathop{\textrm{mix}}}
% \def\rd{\mathop{\textrm{round}}}
% \def\tint{\mathop{\textrm{tint}}}
% \def\shade{\mathop{\textrm{shade}}}
% \def\tone{\mathop{\textrm{tone}}}
% \def\LC{\Package{color}}
% \def\LCP{\Package{color} package}
% \def\XC{\Package{xcolor}}
% \def\XCP{\Package{xcolor} package}
% \def\XCW{\texttt{\href{http://www.ukern.de/tex/xcolor.html}{www.ukern.de/tex/xcolor.html}}}
% \def\URL#1{\texttt{\url{#1}}}
% \def\CTAN#1#2{\texttt{\href{http://www.ctan.org/tex-archive/#1}{CTAN/#1#2}}}
% \newtheorem{lemma}{Lemma}
% \def\RtoI#1{\Gamma_{\!#1}}
% \def\ItoR#1{\Delta_{#1}}
% \def\bbbN{\textrm{I\kern-.333ex N}}
% \def\bbbR{\textrm{I\kern-.333ex R}}
% \def\Circ{\textsuperscript{$\circ$}}
%
% \newdimen\XCdim
%
% \parindent0pt
% \extrarowheight1pt
% \def\textfraction{0.01}
% \def\topfraction{0.99}
% \def\bottomfraction{0.99}
% \def\floatpagefraction{0.70}
% \setlength\abovecaptionskip{12pt}
% \setlength\belowcaptionskip{10pt}
% \colorlet{tableheadcolor}{gray!25}
% \colorlet{tablerowcolor}{gray!12.5}
%
%
%
% \title{扩展 \LaTeX{} 的颜色功能：\XCP}
% \author{\href{mailto:xcolor@ukern.de}{\fboxrule0pt\fboxsep2pt\fbox{Dr. Uwe Kern（尤尔根·克恩 博士）}}}
% \date{\XCfileversion{} (\XCfiledate)
%       \thanks{本宏包可从 \CTAN{macros/latex/contrib/xcolor/}{} 下载。
%               另有 \textsf{xcolor} 主页：\XCW。
%               如有错误报告或改进建议，请发送至：
%               \url{https://github.com/latex3/xcolor/issues}.}}
% \maketitle
%
% \begin{abstract}
% \noindent
% \XC{} 通过类似 |\color{red!50!green!20!blue}| 的颜色表达式，提供了与驱动无关的便捷方式，可访问多种颜色、色调（tint）、暗色（shade）、灰度调和（tone）以及任意颜色的混合。
% 它允许选择文档全局的目标颜色模型，并提供自动配色方案、十二种颜色模型之间的转换、表格行交替着色、颜色混合与蒙版、颜色分色以及色轮计算等工具。
% \end{abstract}
%
% \phantomsection\pdfbookmark[1]{\contentsname}{toc}
% \tableofcontents\listoftables\listoffigures\newpage
%
%
%
%
% \clearpage
% \section{引言}
%
% \subsection{本宏包的用途}
%
% \LCP{} 提供了一个强大而稳定的工具，用于在 (pdf)\LaTeX{} 中以一致且与驱动无关的方式处理颜色，并支持多种颜色模型（在驱动无关性方面略有折扣）。
%
% 尽管如此，它有时用起来还是略显笨拙，尤其是在涉及轻微颜色变化、颜色混合或颜色转换的场合：
% 这通常意味着需要借助另一个程序计算出必要的参数，然后再把这些参数复制到 \LaTeX{} 的 |\definecolor| 命令中。
% 很多时候，还得借助计算器来处理以下这类问题：
% \begin{itemize}
% \item 我的公司定义了一种企业色，印刷厂告诉我在新宣传册中使用超过两种颜色会很贵，但使用该颜色的各种色调（例如 75\% 浓度）则无需额外费用。
% 可是如何在 \LaTeX{} 中访问这些颜色变化呢？\\
% （答案：|\color{CorporateColor!75}| 等）
% \item 我的朋友用了一种很好看的颜色，我想在自己的文档中使用；不幸的是，它是用 \Model{hsb} 模型定义的，而我最常用的 pdf\LaTeX{} 不支持该模型。
% 现在怎么办？\\
% （答案：直接使用 \Model{hsb} 定义即可，\XC{} 会完成必要的计算）
% \item 40\% 的 \Color{green} 与 60\% 的 \Color{yellow} 混合后是什么样子？\\
% （答案：40\% \testclr{green} $+$ 60\% \testclr{yellow} $=$ \testclr{green!40!yellow}，例如 |\color{green!40!yellow}|）
% \item 那它的补色又是什么样子？\\
% （答案：\testclr{-green!40!yellow}，可通过 |\color{-green!40!yellow}| 访问）
% \item 现在我想把上面这种颜色取三份、其补色取两份、\Color{red} 取一份混合起来。
% 效果如何？\\
% （答案：$3\times$ \testclr{-green!40!yellow} $+ 2\times$ \testclr{green!40!yellow} $+ 1\times$ \testclr{red} $=$ \testclr{rgb:-green!40!yellow,3;green!40!yellow,2;red,1}，后一种颜色可通过 |\color{rgb:-green!40!yellow,3;green!40!yellow,2;red,1}| 访问）
% \item 我知道 485\,nm 的光波处于可见光范围内。但它代表什么颜色？\\
% （答案：大约是 \testclr[wave]{485}，通过 |\color[wave]{485}|）
% \item 我的印刷厂希望我文档中所有的颜色定义都转换为 \Model{cmyk} 模型。
% 如何高效完成这些计算？\\
% （答案：|\usepackage[cmyk]{xcolor}| 或 |\selectcolormodel{cmyk}|）
% \item 我有一张 50 行的表格。
% 怎样才能让整行交替着色，而不必复制 50 个 |\rowcolor| 命令？而且交替方案要从第 3 行开始。\\
% （答案：类似 |\rowcolors{3}{OddColor}{EvenColor}|）
% \end{itemize}
% 以上就是 \XCP{} 所要解决的部分问题。
% 其宗旨可以概括为：在保持 \LC{} 原有特性的同时，（希望能以）易用的接口提供额外的功能与灵活性。
%
%
%
% \subsection{色调、暗色、灰度调和与补色}
%
% 按照 \cite{smith} 的定义，我们有如下术语：
% \begin{itemize}
% \item \textbf{\Index{tint}（色调）}：加入 \Color{white} 后得到的颜色，
% \item \textbf{\Index{shade}（暗色）}：加入 \Color{black} 后得到的颜色，
% \item \textbf{\Index{tone}（灰度调和）}：加入 \Color{gray} 后得到的颜色。
% \end{itemize}
% 它们都是一般函数 $\mix(C,C',p)$ 的特例，该函数由 $p$ 份颜色 $C$ 与 $1-p$ 份颜色 $C'$ 构造出新颜色，其中 $0\le p\le1$。
% 因此我们定义：
% \begin{align}
% \tint(C,p) &:= \mix(C,|white|,p)\\
% \shade(C,p) &:= \mix(C,|black|,p)\\
% \tone(C,p) &:= \mix(C,|gray|,p)
% \end{align}
% 其中 |white|、|black|、|gray| 是与模型相关的常量，见表 \vref{tab.const}。
% 进一步定义：
% \begin{itemize}
% \item \textbf{补色（complement）}：与原色 $C$ 叠加后得到 \Color{white} 的颜色 $C^*$，
% \end{itemize}
% 须知补色还有其他不同的概念（例如 \emph{色轮} 上的对侧颜色）。
% 计算细节见第 \vref{sec.cnv} 节，色轮的相关说明见第 \vref{sec.harmony} 节。
%
%
%
% \subsection{颜色模型}
%
% 颜色模型是一种工具，用于以适合目标设备（如屏幕或打印机）的方式来描述或表示某一组颜色。
% 存在一些专有模型（如 \Index{Pantone} 或 \Index{HKS}），它们提供有限的颜色集合（通常称为 \emph{专色（spot color）}），用户只需从中选择而无需关心参数；
% 另一方面，也有参数驱动的模型，如 \Model{gray}、\Model{rgb}、\Model{cmyk}，它们旨在表示数量庞大乃至（理论上）无限的颜色集合，这些颜色基于很少的几种基色以及从基色构造其他颜色的规则。
% 例如，通过基色 \Color{red}、\Color{green}、\Color{blue} 的线性组合，可以构造出大范围的颜色。
% 另一方面，通常 \Index{spot color} 只能通过 \Model{cmyk} 或 \Model{rgb} 等模型中的参数来 \emph{近似}；原始颜色是物理混合而成的，甚至与所用纸张类型有关。
% 最后，还有一些颜色如 \Color{gold}、\Color{silver}，在标准墨水或激光打印机上几乎无法用任何参数驱动的颜色模型复现。
%
%
%
% \subsection{色轮与色彩和谐}\label{sec.harmony}
%
% 将颜色（色相）排列在色轮上来讨论理论或实践中的色彩问题，已有悠久历史（如 \People{Isaac}{Newton}、\People{Johann Wolfgang von}{Goethe}）。
% 原因之一或许是：圆本身就是展示共性关系与对立属性的天然工具。
%
% 如今，相关术语存在一定的混淆，因为与色彩相关的两大领域——艺术与图形设计，以及科学色彩理论——倾向于使用相同的词汇来描述颜色属性，但有时所指却大相径庭！
% 因此，不仅色轮的外观各不相同，"原色"或"补色"等概念的含义也存在差异。
%
%
% \paragraph{典型色轮的构造}
% 首先，将三种 \emph{原色} 放在 0\Circ、120\Circ、240\Circ 的位置（艺术家常选 \Color{red}、\Color{yellow}、\Color{blue}，而色彩科学家可能更倾向于 \Color{red}、\Color{green}、\Color{blue}）。
% 接着，将三种 \emph{间色} 放在 60\Circ、180\Circ、300\Circ 的位置。
% 然后，可在其间再放置六种 \emph{复色}（30\Circ、90\Circ 等位置）。
% 这就是为什么色轮常被描述为由十二种等距颜色组成，尽管该算法实际上可以任意次重复。
%
%
% \paragraph{由色轮导出的色彩和谐}
% 我们从任意一个色轮出发：
% \begin{itemize}
% \item \textbf{补色（complementary colors）}：在色轮上相距 180\Circ 的两种颜色，
% \item \textbf{三色组（color triads）}：每隔 120\Circ 取一种颜色，共三种，
% \item \textbf{四色组（color tetrads）}：每隔 90\Circ 取一种颜色，共四种。
% \end{itemize}
% 现在假设色轮被等分为 $2n$ 个扇形区域：
% \begin{itemize}
% \item 给定颜色的 \textbf{分裂补色（split complementary colors）} 是其补色的两个紧邻邻色，位置为 $\frac{n\pm1}{2n}\cdot360^\circ$，
% \item 给定颜色的 \textbf{类似色（analogous colors）} 是它的两个或四个邻色，位置为 $\pm\frac{1}{2n}\cdot360^\circ$ 和 $\pm\frac{2}{2n}\cdot360^\circ$。
% \end{itemize}
% 有了这些生成"色彩和谐"的方法之后，我们可以得出结论：结果在很大程度上取决于色轮是如何构造的！
% 此外，$n$ 的选择也会影响视觉效果。
% 示例见图 \vref{fig.harmony}。
%
%
%
%
% \clearpage
% \section{用户接口}
%
% \subsection{准备工作}
%
% \subsubsection{宏包安装}
%
% 首先，将 \File{xcolor.sty} 和所有 \File{.def} 文件放到 (pdf)\LaTeX{} 能够找到的位置。
% 按照 \emph{\TeX\ 目录结构}（TDS）的惯例，典型位置是 |texmf/tex/latex/xcolor| 目录，其中 |texmf| 表示你的 \TeX\ 安装主目录。
% 另外，将 \File{xcolor.pro} 放到 \Program{dvips} 能够找到的位置，通常是 |texmf/dvips/xcolor|。
% 通常，你需要运行某种文件名数据库更新命令，以便 \TeX\ 系统能够识别并快速搜索这些文件。
% 之后，在文档中使用 \XC{}（替代 \LC）即可。
% 因此，一般命令是在文档导言区使用 |\usepackage|\oarg{options}|{xcolor}|。
% 表 \vref{tab.package} 说明了加载宏包顺序时需要注意的事项。
%
%
% \subsubsection{宏包选项}
%
% 一般来说，选项分为以下几类：
% \begin{itemize}
% \item 决定颜色驱动的选项，如 \cite{dpc} 和 \cite{color} 所述，目前包括：\Option{dvips}, \Option{xdvi}, \Option{dvipdf}, \Option{dvipdfm}, \Option{dvipdfmx}, \Option{luatex}, \Option{pdftex}, \Option{dvipsone}, \Option{dviwindo}, \Option{emtex}, \Option{dviwin}, \Option{oztex}, \Option{textures}, \Option{pctexps}, \Option{pctexwin}, \Option{pctexhp}, \Option{pctex32}, \Option{truetex}, \Option{tcidvi}, \Option{vtex}, \Option{xetex}, \Option{dvisvgm},
% \item 决定目标颜色模型的选项\footnote{第 \vref{sec.tgt} 节说明了如何在文档任意位置覆盖此设置。}（\Option{natural}, \Option{rgb}, \Option{cmy}, \Option{cmyk}, \Option{hsb}, \Option{gray}, \Option{RGB}, \Option{HTML}, \Option{HSB}, \Option{Gray}），或者关闭彩色输出（\Option{monochrome}），
% \item 控制是否以及如何加载某些预定义颜色集合的选项：\Option{dvipsnames}, \Option{dvipsnames*}, \Option{svgnames}, \Option{svgnames*}, \Option{x11names}, \Option{x11names*},
% \item 决定加载或支持哪些其他宏包的选项：\Option{table},
% \item 影响其他命令行为的选项：\Option{prologue}, \Option{kernelfbox}, \Option{xcdraw}, \Option{noxcdraw}, \Option{fixinclude}, \Option{showerrors}, \Option{hideerrors},
% \item 已废弃的选项：
% \oOption{pst},
% \oOption{override}, \oOption{usenames}, \oOption{nodvipsnames}, \oOption{fixpdftex}, \oOption{hyperref}。
% \end{itemize}
% 所有可用的宏包选项（驱动选择和已废弃选项除外）列于表 \vref{tab.opt}。
% \DescribeMacro\GetGinDriver
% \DescribeMacro\GinDriver
% 为了便于与 \Package{hyperref} 宏包协作，提供了命令 |\GetGinDriver|\footnote{如果使用宏包选项 \Option{hyperref}，此命令会自动执行。}，它会抓取实际使用的驱动并将其存入命令 |\GinDriver|。
% 之后可在 \Package{hyperref}（或其他宏包）中使用该命令，参见第 \pageref{ex.range} 页的代码示例。
% 如果没有对应的 \Package{hyperref} 选项，则默认使用 \Option{hypertex}。
%
% 警告：\XC{} 与 \LC{} 在处理 \Option{dvips} 选项的方式上存在重大差异。
% \LCP{} 会在选择 |dvips|、|oztex|、|xdvi| 驱动之一时，隐式启用 \Option{dvipsnames} 选项。
% 这会降低文档的可移植性，因为当使用这些颜色而未显式指定 \Option{dvipsnames} 选项时，|pdftex| 等其他驱动会因未知颜色而报错。
% 因此，\XC{} 要求显式指定 \Option{dvipsnames} 选项才能使用这些颜色名称——这样便可在所有驱动下工作。
%
%
% \begin{table}[p]\caption{宏包选项}\label{tab.opt}
% \extrarowheight2pt
% \begin{adjustwidth}{-2.5cm}{}
% \advance\textwidth2.5cm
% \rowcolors{3}{tablerowcolor}{}
% \begin{tabularx}{\linewidth}{lX}
% \hline
% \rowcolor{tableheadcolor}
% \em 选项 & \em 说明\\
% \Option{natural}& （默认。）保留所有颜色的原始模型，但 \Model{RGB} 会转换为 \Model{rgb}，\Model{HSB} 会转换为 \Model{hsb}，\Model{Gray} 会转换为 \Model{gray}。\\
% \Option{rgb}& 将所有颜色转换为 \Model{rgb} 模型。\\
% \Option{cmy}& 将所有颜色转换为 \Model{cmy} 模型。\\
% \Option{cmyk}& 将所有颜色转换为 \Model{cmyk} 模型。\\
% \Option{hsb}& 将所有颜色转换为 \Model{hsb} 模型。\\
% \Option{gray}& 将所有颜色转换为 \Model{gray} 模型。对于模拟黑白打印机的输出效果特别有用。\\
% \Option{RGB}& 将所有颜色转换为 \Model{RGB} 模型（之后再转换为 \Model{rgb}）。\\
% \Option{HTML}& 将所有颜色转换为 \Model{HTML} 模型（之后再转换为 \Model{rgb}）。\\
% \Option{HSB}& 将所有颜色转换为 \Model{HSB} 模型（之后再转换为 \Model{hsb}）。\\
% \Option{Gray}& 将所有颜色转换为 \Model{Gray} 模型（之后再转换为 \Model{gray}）。\\
% \hline
% \Option{dvipsnames}, \Option{dvipsnames*}& 加载一组预定义颜色。\footnotemark[1]\\
% \Option{svgnames}, \Option{svgnames*}& 加载符合 \Index{SVG} 1.1 标准的一组预定义颜色。\footnotemark[1]\\
% \Option{x11names}, \Option{x11names*}& 加载符合 \Index{Unix}/\Index{X11} 标准的一组预定义颜色。\footnotemark[1]\\
% \hline
% \Option{table}& 加载 \Package{colortbl} 宏包，以便使用表格中行、列、单元格着色工具。\\
% \Option{fixpdftex}& 此选项已废弃，不做任何事情。\\
% \Option{hyperref}& 通过定义额外的键值（见第 \vref{sec.hyperref} 节），在颜色表达式方面支持 \Package{hyperref} 宏包。\\
% \hline
% \Option{prologue}& 为每个颜色定义写入 prologue 信息到 \File{.xcp} 文件（如第 \vref{sec.prologue} 节所述）。\\
% \hline
% \Option{kernelfbox}& 使用 \LaTeX{} 内核方法绘制 \texttt{\string\f(rame)box} 边框\footnotemark[2]。\\
% \Option{xcdraw}& 使用驱动特定命令绘制边框和颜色框。\footnotemark[2]\\
% \Option{noxcdraw}& （默认。）使用通用代码绘制边框和颜色框。\footnotemark[2]\\
% \Option{fixinclude}& 在嵌入 \File{.eps} 文件之前阻止 |dvips| 重置颜色（见第 \vref{sec.fixinclude} 节）。\\
% \hline
% \Option{showerrors}& （默认。）使用未定义颜色时显示错误信息（与原 \LCP 行为相同）。\\
% \Option{hideerrors}& 使用未定义颜色时仅显示警告，并将该颜色替换为 \Color{black}。\\
% \hline\hline
% \multicolumn2{c}{\footnotemark[1]\,见第 \vref{sec.additionalcolors} 节。\quad \footnotemark[2]\,见第 \vref{sec.fbox} 节。}\\
% \hline
% \end{tabularx}
% \end{adjustwidth}
% \end{table}
%
%
% \begin{table}[htbp]\caption{宏包加载顺序}\label{tab.package}
% \begin{adjustwidth}{-2.5cm}{}
% \raggedleft
% \extrarowheight2pt
% \rowcolors{3}{tablerowcolor}{}
% \begin{tabular}{l*{5}{c}}
% \hline
% \rowcolor{tableheadcolor}
% \em 操作/宏包    & \Package{colortbl}&  \Package{hyperref}& \Package{pstricks}& \LC& \Package{pstcol}\\
% 在 \XC 之前加载       & 允许 &  允许 & 允许\rlap{\footnotemark[1]} & 不可 & 不可  \\
% 通过 \XC{} 选项加载& \Option{table} &  --- & --- & --- & --- \\
% 在 \XC 之后加载        & 允许 &  允许 & 允许 & 不可 & 不可 \\
% \hline\hline
% \multicolumn5{c}{\footnotemark[1]\,\Package{pstricks} 默认会加载 \XC}\\
% \hline
% \end{tabular}
% \end{adjustwidth}
% \end{table}
%
%
% \begin{table}[htbp]\caption{支持的颜色模型}\label{tab.mod}
% \raggedleft
% \extrarowheight2pt
% \rowcolors{3}{tablerowcolor}{}
% \begin{tabular}{l>{\em}l>{$}l<{$}rl}
% \hline
% \rowcolor{tableheadcolor}
% \emph{名称}  & 基色/概念 & \emph{参数范围}& \multicolumn2{>{\columncolor{tableheadcolor}}l}{\emph{默认}}\\
% \Model{rgb} & 红、绿、蓝 & \left[0,1\right]^3&& \\
% \Model{cmy} & 青、品红、黄 & \left[0,1\right]^3&& \\
% \Model{cmyk} & 青、品红、黄、黑 & \left[0,1\right]^4&& \\
% \Model{hsb} & 色相、饱和度、亮度 & \left[0,1\right]^3&& \\
% \Model{Hsb} & 色相\Circ\!、饱和度、亮度 & \left[0,H\right]\times\left[0,1\right]^2&$H=$& $\mkern-15mu\rangeHsb$\\
% \Model{tHsb} & 色相\Circ\!、饱和度、亮度 & \left[0,H\right]\times\left[0,1\right]^2&$H=$& $\mkern-15mu\rangeHsb$\\
% \Model{gray} & 灰度 & \left[0,1\right]&& \\
% \Model{RGB} & 红、绿、蓝 & \{0,1,\dots,L\}^3& $L\,=$& $\mkern-15mu\rangeRGB$\\
% \Model{HTML} & RRGGBB & \{\texttt{000000},\dots,\texttt{FFFFFF}\}&& \\
% \Model{HSB} & 色相、饱和度、亮度 & \{0,1,\dots,M\}^3& $M=$& $\mkern-15mu\rangeHSB$\\
% \Model{Gray} & 灰度 & \{0,1,\dots,N\}& $N\,=$& $\mkern-15mu\rangeGray$\\
% \Model{wave} & 波长 \emph{(nm)}& \left[363,814\right]&& \\
% \hline\hline
% \multicolumn5{c}{$L,M,N$ 为正整数；$H$ 为正实数}\\
% \hline
% \end{tabular}
% \end{table}
%
%
% \subsubsection{执行额外的初始化命令}
%
% \DescribeMacro\xcolorcmd
% 这里提供了一个简单接口，用于传递应在 \XCP{} 末尾（紧接初始化命令 |\color{black}| 之前）执行的命令。
% 只需在加载 \XC{} 之前的某个位置写 |\def\xcolorcmd|\marg{commands} 即可。
%
% 示例：假设 |a.tex| 是一个完整的 \LaTeX{} 文档，在控制台执行如下命令：
% `|latex \def\xcolorcmd{\colorlet{black}{red}}\input{a}|'
% 会生成一个 |a.dvi| 文件，其中所有 \Color{black} 出现的地方都被替换为 \Color{red}，而无需修改源文件本身。
% （控制台命令的具体拼写可能因操作系统和 \TeX{} 发行版而异。）
%
%
%
% \subsection{颜色模型}
%
% \subsubsection{支持的颜色模型}
%
% 支持的颜色模型及其参数范围列于表 \vref{tab.mod}。
% 需要强调的是，这种颜色支持与所选驱动无关。
%
% "颜色模型支持"还意味着可以直接使用参数指定颜色，例如写 |\textcolor[cmy]{0.7,0.5,0.3}{foo}|（\textcolor[cmy]{0.7,0.5,0.3}{foo}）或 |\textcolor[HTML]{AFFE90}{foo}|（\textcolor[HTML]{AFFE90}{foo}）。
%
% \paragraph{\Model{rgb}, \Model{cmyk}, \Model{hsb}, \Model{gray}}
% 这些是 \Index{PostScript} 直接支持的模型。
% 其性质和相互关系请参见 \cite{plrm}。
% 还有一个特殊命令
% \DescribeMacro\adjustUCRBG
% 可用于精细调节转换到 \Model{cmyk} 模型时的 \emph{底色去除（undercolor-removal）}和 \emph{黑版生成（black-generation）}机制，详见第 \vref{sec.cmy2cmyk} 节。
%
% \paragraph{\Model{cmy}}
% 这主要是一个用于中间计算步骤的模型。
% 就此而言，它是 \Model{rgb} 的简单补集。
% 在可视化方面，\Model{cmy} 被视为 $k=0$ 的 \Model{cmyk}。
%
% \paragraph{\Model{HTML}}
% 这是从 \Model{rgb} 派生出来的模型，用于支持从网页或 CSS 文件中输入颜色参数。
% 因此，它本身并非真正独立的颜色模型，而是一个为方便使用而设的用户接口。
% 值得一提的是，\Model{HTML} 接受字符 |0|--|9|、|A|--|F|、|a|--|f| 的任意组合，只要字符串长度恰好为 6 个字符。
% 不过，转换为 \Model{HTML} 的输出将始终由数字和\emph{大写}字母组成。
%
% \paragraph{\Model{Hsb}, \Model{tHsb}}
% 首先，\Model{Hsb} 是一个"用户接口"模型，它将 $\emph{hue}\in[0,1]$ 映射到 $\emph{hue\Circ}\in[0,H]$，其中 $H$ 由 |\def\rangeHsb|\marg{H} 给出。
% \DescribeMacro\rangeHsb
% 因此，如果 $H=360$，我们可以把它想象成一个圆或色轮，用于指定 \emph{hue\Circ} 参数。
% 其次，\Model{Hsb} 是 \Model{tHsb}（又称 \pModel{tuned} \Model{Hsb}）的基础，它允许用户通过在圆上前后移动选定的 \emph{hue\Circ} 值，对 \emph{hue\Circ} 应用分段线性变换。
% 该变换由 |\def\rangetHsb{|$x_1,y_1;x_2,y_2;\dots$|}| 定义
% \DescribeMacro\rangetHsb
% 它指定 \Model{tHsb} 中的 $\emph{hue\Circ}=x_1$ 对应于 \Model{Hsb} 中的 $\emph{hue\Circ}=y_1$，依此类推。
% 例如，在 \Model{Hsb} 圆上，\Color{yellow} 位于 60\Circ（\Color{red} 位于 0\Circ），但在艺术家中使用的大多数色轮上，\Color{yellow} 位于 120\Circ。
% 因此，如果我们打算通过 \Model{tHsb} 复现艺术家的色轮，设置"|120,60|"将是合理的。
% 确切的变换公式和更多限制见第 \vref{sec.tHsb2Hsb} 节，色轮与色彩和谐见第 \vref{sec.harmony} 节。
% 图 \vref{fig.tHsb} 可供对比参考。
%
% 示例：`|\def\rangetHsb{|\texttt{\rangetHsb}|}|' 实际上就是 \XC 的默认设置。
%
% \paragraph{\Model{wave}}
% 通过这个模型，我们尝试将波长转换为标准颜色模型，以近似光波的视觉外观。
% 可见光谱范围约为 400--750\,nm，而 \XC{} 中的实现通常允许所有绝对值小于 \csname strip@pt\endcsname\maxdimen{}（\TeX{} 能表示的最大长度值）的实数"波长"$\lambda$。
% 不过，在 $\lambda\in[363,814]$ 范围之外得到任何非黑色的概率恰好为零。
% 因此，图 \vref{fig.spectrum} 仅展示上述波长区间。
% 注意：不可能将其他模型合理地转换为 \Model{wave}，因为后者仅覆盖有限的颜色范围。
%
% \paragraph{\Model{RGB}, \Model{HSB}, \Model{Gray}}
% 这些是派生模型，将 \Model{rgb}、\Model{hsb}、\Model{gray} 的连续 $[0,1]$ 参数范围转换为有限整数集合；因此我们称之为\emph{整数模型}。
% 表 \ref{tab.mod} 中的常量 $L, M, N$ 通过以下命令定义：|\def\rangeRGB|\marg{L}，
% \DescribeMacro\rangeRGB
% |\def\rangeHSB|\marg{M}，
% \DescribeMacro\rangeHSB
% 以及 |\def\rangeGray|\marg{N}。
% \DescribeMacro\rangeGray
% 这些常量可以在 \XCP{} 加载\emph{之前}或\emph{之后}修改，例如：
% \label{ex.range}\begin{verbatim}
% \documentclass{article}
% ...
% \def\rangeRGB{15}
% \usepackage[dvips]{xcolor}
% ...
% \GetGinDriver
% \usepackage[\GinDriver]{hyperref}
% ...
% \begin{document}
% ...
% \def\rangeRGB{63}
% ...
% \end{verbatim}
%
%
% \subsubsection{替换个别颜色模型}\label{sec.subst}
%
% \DescribeMacro\substitutecolormodel
%   \marg{source model}\marg{target model-list}\\
% 用 \Meta{target model-list} 中实际存在的第一个模型替换 \Meta{source model}。
% 仅允许 \Meta{num model} 类型的颜色模型；
% 所有更改仅作用于当前分组，但前置 |\xglobal| 仍会生效。
%
% 示例：假设当前驱动对 \Model{hsb} 的实现有问题，而 \Model{rgb} 表现正常。
% 那么 |\substitutecolormodel{hsb}{rgb}| 可能是一个好的选择，因为它从此处开始将所有 \Model{hsb} 颜色的定义通过 \XC 的算法转换为 \Model{rgb} 规格，而不影响其他模型。
%
%
% \subsubsection{在文档中途更改目标颜色模型}\label{sec.tgt}
%
% \DescribeMacro\selectcolormodel
%   \marg{num model}\\
% 将目标模型设置为 \Meta{num model}，后者是允许作为宏包选项的模型名称之一（即 \Option{natural}, \Option{rgb}, \Option{cmy}, \Option{cmyk}, \Option{hsb}, \Option{gray}, \Option{RGB}, \Option{HTML}, \Option{HSB}, \Option{Gray}），示例见图 \vref{fig.tgt}。
% 转换为目标模型可以在两个时机进行：
% \begin{itemize}
% \item \DescribeMacro\ifconvertcolorsD 在颜色\emph{定义}时\footnote{这意味着所有\emph{新}定义的颜色会先转换为目标模型，然后再保存。}（即在 |\definecolor| 及其类似命令内部）；由此开关 |\ifconvertcolorsD| 控制；
% \item \DescribeMacro\ifconvertcolorsU 在颜色\emph{使用}时（紧接颜色显示之前，因此涵盖在其他模型中定义的颜色以及直接指定的颜色，如 |\color[rgb]{.1,.2,.3}|）；由此开关 |\ifconvertcolorsU| 控制。
% \end{itemize}
% 选择除 \Option{natural} 以外的任何模型都会将两个开关都设为"真"，而 \Option{natural} 将它们设为"假"。
% 无论是通过宏包选项还是通过 |\selectcolormodel| 选择，均适用此规则。
% 为什么不在使用时才简单地转换所有颜色呢？
% 如果涉及大量颜色，在颜色定义阶段就完成所有转换可以节省处理时间。
% 最佳性能可通过 |\usepackage[rgb,...]{xcolor}\convertcolorsUfalse| 获得，这实际上就是 \XC{} 1.07 及更早版本的工作方式。
%
%
% \begin{table}[p]\caption{参数与术语}\label{tab.term}
% \def\RS{\ensuremath{\to}\space}
% \footnotesize\raggedleft
% \extrarowheight2pt
% \begin{adjustwidth}{-2.5cm}{}
% \advance\textwidth2.5cm
% \tabcolsep.25em
% \rowcolors{3}{tablerowcolor}{}
% \begin{tabularx}{\textwidth}{lX}
% \hline
% \rowcolor{tableheadcolor}
% \em 元素 & \em 替换串\\
% \Meta*{empty}& \RS empty string `'\\
% \Meta*{minus}& \RS non-empty string consisting of one or more minus signs `|-|'\\
% \Meta*{plus}& \RS non-empty string consisting of one or more plus signs `|+|'\\
% \hline
% \Meta*{int}& \RS 整数\hfill(\emph{整数})\\
% \Meta*{num}& \RS 非负整数\hfill(\emph{数})\\
% \Meta*{dec}& \RS 实数\hfill(\emph{小数})\\
% \Meta*{div}& \RS 非零实数\hfill(\emph{除数})\\
% \Meta*{pct}& \RS 区间 $[0,100]$ 内的实数\hfill(\emph{百分比})\\
% \hline
% \Meta*{id}& \RS 由字母和数字组成的非空串\hfill(\emph{标识符})\\
% \Meta*{ext id}& \RS \Meta{id}\newline
%                 \RS \Meta[1]{id}|=|\Meta[2]{id}\\
% \Meta*{id-list}& \RS \Meta[1]{ext id}|,|\Meta[2]{ext id}|,...,|\Meta[l]{ext id}\\
% \Meta*{name}& \RS \Meta{id}\hfill(\emph{显式名称})\newline
%               \RS `|.|'\hfill(\emph{隐式名称})\\
% \Meta*{core model}& \RS `|rgb|', `|cmy|', `|cmyk|', `|hsb|', `|gray|'\hfill(\emph{核心模型})\\
% \Meta*{num model}& \RS \Meta{core model}\newline
%                    \RS `|RGB|', `|HTML|', `|HSB|', `|Gray|'\hfill(\emph{整数模型})\newline
%                    \RS `|Hsb|', `|tHsb|', `|wave|'\hfill(\emph{小数模型})\\
% \Meta*{model}& \RS \Meta{num model}\hfill(\emph{数值模型})\newline
%                \RS `|named|'\hfill(\emph{伪模型})\\
% \Meta*{model-list}& \RS \Meta[1]{model}|/|\Meta[2]{model}|/.../|\Meta[m]{model} \hfill(\emph{多模型})\newline
%                     \RS \Meta{core model}|:|\Meta[1]{model}|/|\Meta[2]{model}|/.../|\Meta[m]{model}\\
% \Meta*{spec}& \RS 逗号分隔的数值列表\hfill(\emph{显式指定})\newline
%               \RS 空格分隔的数值列表\hfill(\emph{显式指定})\newline
%               \RS 某个 `named' 颜色的名称\hfill(\emph{隐式指定})\\
% \Meta*{spec-list}& \RS \Meta[1]{spec}|/|\Meta[2]{spec}|/.../|\Meta[m]{spec} \hfill(\emph{多重指定})\\
% \Meta*{type}& \RS \Meta{empty}\newline
%               \RS `|named|', `|ps|'\\
% \hline
% \Meta*{expr}& \RS \Meta{prefix}\Meta{name}\Meta{mix expr}\Meta{postfix}\hfill(\emph{标准颜色表达式})\\
% \Meta*{prefix}& \RS \Meta{empty}\newline
%                 \RS \Meta{minus}\hfill(\emph{补色指示})\\
% \Meta*{mix expr}& \RS |!|\Meta[1]{pct}|!|\Meta[1]{name}|!|\Meta[2]{pct}|!|\Meta[2]{name}|!...!|\Meta[n]{pct}|!|\Meta[n]{name} \hfill(\emph{完整混合表达式})\newline
%                   \RS |!|\Meta[1]{pct}|!|\Meta[1]{name}|!|\Meta[2]{pct}|!|\Meta[2]{name}|!...!|\Meta[n]{pct} \hfill(\emph{不完整混合表达式})\\
% \Meta*{postfix}& \RS \Meta{empty}\newline
%                  \RS |!!|\Meta{plus}\hfill(\emph{系列步进})\newline
%                  \RS |!![|\Meta{num}|]|\hfill(\emph{系列访问})\\
% \hline
% \Meta*{ext expr}& \RS \Meta{core model}|,|\Meta{div}|:|\Meta[1]{expr}|,|\Meta[1]{dec}|;|\Meta[2]{expr}|,|\Meta[2]{dec}|;...;|\Meta[k]{expr}|,|\Meta[k]{dec}\newline
%                   \RS \Meta{core model}|:|\Meta[1]{expr}|,|\Meta[1]{dec}|;|\Meta[2]{expr}|,|\Meta[2]{dec}|;...;|\Meta[k]{expr}|,|\Meta[k]{dec}\\
% \hline
% \Meta*{func expr}& \RS |>|\Meta{function}|,|\Meta[1]{arg}|,|\Meta[2]{arg}|,...,|\Meta[j]{arg}\hfill(\emph{颜色函数表达式})\\
% \Meta*{function}& \RS `|wheel|', `|twheel|'\hfill(\emph{颜色函数})\\
% \hline
% \Meta*{color}& \RS \Meta{color expr}\Meta[1]{func expr}\Meta[2]{func expr}|...|\Meta[i]{func expr}\\
% \Meta*{color expr}& \RS \Meta{name}\newline
%                     \RS \Meta{expr}\newline
%                     \RS \Meta{ext expr}\\
% \hline\hline
% \rowcolor{white}
% 说明：& 每个 \RS 表示左列元素的一种可能替换串；
%           但可能还有进一步的上下文相关限制。
%           详见正文。
%           字符串 `|foo|' 始终不包含引号。
%           $i,j,k,l,m,n$ 为非负整数，$k,l,m,n>0$，$m\le8$。\\
% \hline
% \end{tabularx}
% \end{adjustwidth}
% \end{table}
%
%
%
% \subsection{参数与术语}\label{sec.term}
%
% 在详细描述 \XC 的颜色相关命令之前，我们先定义一些在这些命令参数中反复出现的元素或标识符。
% 表 \vref{tab.term} 给出了通用语法概览。
%
%
% \subsubsection{参数的附加说明与限制}
%
% \paragraph{Basic strings and numbers}
% \DescribeArg{empty}
% \DescribeArg{minus}
% \DescribeArg{plus}
% \DescribeArg{int}
% \DescribeArg{num}
% \DescribeArg{dec}
% \DescribeArg{pct}
% \DescribeArg{div}
% 这些参数无需过多解释。
% 不过就数值而言，值得注意的是，(La)\TeX{} 中的实数——只要用于长度、维度或间距的上下文——其绝对值最大限制为 $<16384$。
% 当然，在数值计算链中，每个中间结果也必须遵守此限制，这通常意味着更严格的范围限制。
% 由于 \XC{} 在大多数计算中大量使用 \TeX 的内部维度寄存器，因此在使用 \Meta{ext expr} 表达式时应牢记这一点。
%
%
% \paragraph{颜色名称}
% \DescribeArg{name}
% \Meta{name} 表示\emph{颜色}或\emph{颜色系列}的已声明名称（或待声明名称）；
% 它可以通过以下命令之一\emph{显式}声明：|\definecolor|, |\providecolor|, |\colorlet|, |\definecolorset|, |\providecolorset|, |\definecolorseries|, |\definecolors|, |\providecolors|。
% 另一方面，保留颜色名 `|.|' 是\emph{隐式}声明的，表示\emph{当前颜色}。
% 实际上，除了字母和数字，某些其他字符也可用于 \Meta{name} 声明，但上述限制可避免误解并确保与 \XC 未来扩展的兼容性。
%
% 示例：`|red|', `|MySpecialGreen1980|', `|.|'.
%
%
% \paragraph{颜色模型}
% \DescribeArg{core model}
% \DescribeArg{num model}
% \DescribeArg{model}
% 区分\emph{核心模型}（\Model{rgb}, \Model{cmy}, \Model{cmyk}, \Model{hsb}, \Model{gray}）、\emph{整数模型}（\Model{RGB}, \Model{HTML}, \Model{HSB}, \Model{Gray}）、\emph{小数模型}（\Model{Hsb}, \Model{tHsb}, \Model{wave}）和\emph{伪模型}（目前为 \nModel{named}, \nModel{ps}）有一个简单原因：以单位区间 $[0,1]$ 为参数范围的核心模型最适合各种计算，而整数模型的目的主要是便于参数输入，之后再转换为某种核心模型。
% 最后，小数模型 \Model{Hsb} 和 \Model{tHsb} 是 \Model{hsb} 的特殊用途版本，而 \Model{wave} 和伪模型 \nModel{named} 具有特殊地位，因为它们"不擅计算"：通常只能将这种颜色转换为其他模型之一，而不能反向转换。
% 伪模型 \nModel{ps} 更是如此：由于这类颜色包含 \Index{PostScript} 代码，对 \TeX 来说完全是不透明的。
%
%
% \paragraph{颜色规格}
% \DescribeArg{spec}
% \Meta{spec} 参数用于指定颜色的参数，显然取决于底层颜色模型。
% 我们区分\emph{显式}指定和\emph{隐式}指定：前者指表 \vref{tab.mod} 中说明的数值参数，后者（理想情况下）指驱动提供的名称。
%
% 示例：`|.1,.2,.3|', `|.1 .2 .3|', `|0.56789|', `|89ABCD|', `|ForestGreen|'.
%
%
% \paragraph{多模型与多规格}
% \DescribeArg{model-list}
% \DescribeArg{spec-list}
% 这些参数总是在以下颜色定义命令中成对（显式或隐式）出现：
% |\definecolor|, |\providecolor|, |\definecolorset|, |\providecolorset|。
% 首先，将 \Meta{model-spec} 与当前目标模型（由宏包选项或 |\selectcolormodel| 命令设置）协调；如果没有完全匹配，则选择列表中的第一个模型。
% 然后，从 \Meta{spec-list} 中选择对应的颜色规格，从而得到一对合适的 $(\Meta{model}, \Meta{spec})$。
% 因此，在实际执行的颜色定义中不再存在歧义。
% 扩展形式 \Meta{core model}|:|\Meta[1]{model}|/|\Meta[2]{model}|/.../|\Meta[m]{model} 会立即将相关 \Meta{spec} 转换为 \Meta{core model}；未知的目标模型会被静默忽略。
%
% 示例：`|rgb/cmyk/named/gray|', `|0,0,0/0,0,0,1/Black/0|', `|rgb:cmy/hsb|'.
%
%
% \paragraph{type 参数}
% \DescribeArg{type}
% 仅在颜色定义命令的上下文中使用，参见 |\definecolor| 及其类似命令的说明。
%
%
% \paragraph{标准颜色表达式}
% \DescribeArg{expr}
% \DescribeArg{prefix}
% \DescribeArg{mix expr}
% \DescribeArg{postfix}
% 这些表达式是一种工具，用于方便地指定某种级联颜色混合形式，详见第~\ref{sec.expr-meaning} 节。
% \Meta{prefix} 参数控制其后的颜色还是其补色生效：
% 奇数个减号表示剩余表达式得到的颜色需要转换为其补色。
% \emph{不完整混合表达式} 只是 \emph{完整混合表达式} 在 $\Meta[n]{name} = |white|$ 时的缩写，以便在使用色调时减少按键次数。
% \Meta{postfix} 字符串通常为空，但在\emph{颜色系列}的情况下提供额外功能：
% 非空情形要求：
% \begin{itemize}
% \item \Meta{name} 表示\emph{颜色系列}的名称，
% \item \Meta{mix expr} 是\emph{完整的}混合表达式。
% \end{itemize}
% 示例：`|red|', `|-red|', `|--red!50!green!12.345|', `|red!50!green!20!blue|', `|foo!!+|', `|foo!![7]|', `|foo!25!red!!+++|', `|foo!25!red!70!green!![7]|'.
%
%
% \paragraph{扩展颜色表达式}
% \DescribeArg{ext expr}
% 这些表达式提供了另一种颜色混合方法，详见第 \vref{sec.extexpr-meaning} 节。
% 较短形式
% \begin{align*}
% \Meta{core model}&|:|\Meta[1]{expr}|,|\Meta[1]{dec}|;|\Meta[2]{expr}|,|\Meta[2]{dec}|;...;|\Meta[k]{expr}|!|\Meta[k]{dec}\\
% \intertext{是以下特殊情况（可能是最常用的）的缩写：}
% \Meta{core model}|,|\Meta{div}&|:|\Meta[1]{expr}|,|\Meta[1]{dec}|;|\Meta[2]{expr}|,|\Meta[2]{dec}|;...;|\Meta[k]{expr}|!|\Meta[k]{dec}\\
% \intertext{其定义如下（要求所有 \Meta[\kappa]{dec} 系数之和非零）：}
% \Meta{div} &:= \Meta[1]{dec} + \Meta[2]{dec} + \dots + \Meta[k]{dec} \neq0.
% \end{align*}
% 示例：`|rgb:red,1|', `|cmyk:red,1;-green!25!blue!60,11.25;blue,-2|'.
%
%
% \paragraph{函数表达式}
% \DescribeArg{func expression}
% \DescribeArg{function}
% 这些表达式通过对标准或扩展表达式的结果进行额外计算，扩展了其功能。
% 参数数量因函数而异，详见第 \vref{sec.functions} 节。
%
% 示例：`|>wheel,30|', `|>wheel,30,|', `|>twheel,1,12|', `|>twheel,-11,12|'.
%
%
%
% \paragraph{颜色}
% \DescribeArg{color}
% \DescribeArg{color expr}
% 最后，\Meta{color} 是一个"伞形"参数，涵盖指定颜色的各种概念。
% 这意味着，只要出现 \Meta{color} 参数，就可以使用上述全部名称和表达式。
%
%
% \subsubsection{标准颜色表达式的含义}\label{sec.expr-meaning}
%
% 我们现在解释表达式
% \begin{gather*}
% \Meta{prefix}\Meta{name}|!|\Meta[1]{pct}|!|\Meta[1]{name}|!|\Meta[2]{pct}|!|\dots|!|\Meta[n]{pct}|!|\Meta[n]{name}\Meta{postfix}
% \end{gather*}
% 是如何被解释和处理的：
% \begin{enumerate}
% \item 首先，提取 \Meta{name} 的模型和颜色参数，定义一个临时颜色 \Meta{temp}。
% 如果 \Meta{postfix} 具有 `|!![|\Meta{num}|]|' 的形式，则 \Meta{temp} 将是系列 \Meta{name} 中对应的（直接访问的）第 \Meta{num} 个颜色。
% \item 然后计算颜色混合：\Meta[1]{pct}\% 的颜色 \Meta{temp} 与 $(100-\Meta[1]{pct})$\% 的颜色 \Meta[1]{name}；这就是新的临时颜色 \Meta{temp}。
% \item 对剩余所有参数对 (\Meta[2]{pct},\Meta[2]{name}), \dots, (\Meta[n]{pct},\Meta[n]{name}) 重复上一步。
% \item 如果 \Meta{prefix} 由奇数个减号 `-` 组成，则将 \Meta{temp} 转换为其补色。
% \item 如果 \Meta{postfix} 具有 `|!!+|'、`|!!++|'、`|!!+++|' 等形式，则对底层颜色系列 \Meta{name} 执行相应数量的步进命令（即 `|+|' 号的数量）。
% 这对颜色 \Meta{temp} 没有影响。
% \item 现在，颜色 \Meta{temp} 被显示，或作为其他操作的输入，取决于调用命令。
% \end{enumerate}
% 注意，在典型的第 2 步表达式 \Meta{temp}|!|\Meta[\nu]{pct}|!|\Meta[\nu]{name} 中，如果 \Meta[\nu]{pct}=100 或 \Meta[\nu]{pct}=0，则直接使用颜色 \Meta{temp} 或 \Meta[\nu]{name}，不再进行变换。
% 在真正的混合情形 $0<$\Meta[\nu]{pct}$<100$ 下，涉及的两种颜色可能定义在不同的颜色模型中，例如 |\definecolor{foo}{rgb}{...}| 和 |\definecolor{bar}{cmyk}{...}|。
% 一般规则是，第二个颜色 \Meta[\nu]{name} 先转换为第一个颜色 \Meta{temp} 的模型，然后在该模型内计算混合。\,\footnote{例外：为避免奇怪结果，如果 \Meta{temp} 源自 \Model{gray} 模型，则此规则反转——将其转换为 \Meta[\nu]{name} 的底层模型。}
% 因此，\Meta{temp}|!|\Meta[\nu]{pct}|!|\Meta[\nu]{name} 与 \Meta[\nu]{name}|!|\Meta[\nu]{$100-$pct}|!|\Meta{temp} 理论上等价，但视觉上不一定产生完全相同的结果。
%
% 图 \vrefrange{fig.expr}{fig.ex} 展示了颜色和表达式的一些初步应用。
% 更多示例见图 \vref{fig.progress}。
% 此外，大量颜色示例可参见 \cite{chroma}。
%
%
% \subsubsection{扩展颜色表达式的含义}\label{sec.extexpr-meaning}
%
% \emph{扩展颜色表达式}
% \begin{gather*}
% \Meta{core model}|:|\Meta[1]{expr}|,|\Meta[1]{dec}|;|\Meta[2]{expr}|,|\Meta[2]{dec}|;...;|\Meta[k]{expr}|,|\Meta[k]{dec}
% \end{gather*}
% 模拟画家调色：指定一个颜色列表，每个颜色附带一个 \Meta{dec} 系数。
% 对于这样的 \Meta{ext expr}，每个标准颜色表达式 \Meta[\kappa]{expr} 都会转换为 \Meta{core model}，然后将所得向量乘以 $\Meta[\kappa]{dec}/\Meta{div}$，其中
% \begin{gather*}
% \Meta{div} := \Meta[1]{dec} + \Meta[2]{dec} + \dots + \Meta[k]{dec}.
% \end{gather*}
% 之后计算所有这些向量之和。
%
% 示例：将 4 份 \testclr{red} \Color{red}、2 份 \testclr{green} \Color{green}、1 份 \testclr{yellow} \Color{yellow} 混合，通过 |\color{rgb:red,4;green,2;yellow,1}| 得到 \testclr{rgb:red,4;green,2;yellow,1}。
% 如果改用 $-1$ 份 \Color{yellow}，则得到 \testclr{rgb:red,4;green,2;yellow,-1}。
% 注意，此机制也可用于以某种颜色模型显示单个颜色（表达式）：|\color{rgb:yellow,1}| 就会产生这样的转换。
% 一般形式
% \begin{gather*}
% \Meta{core model}|,|\Meta{div}|:|\Meta[1]{expr}|,|\Meta[1]{dec}|;|\Meta[2]{expr}|,|\Meta[2]{dec}|;...;|\Meta[k]{expr}|,|\Meta[k]{dec}
% \end{gather*}
% 执行相同操作，唯一区别是除数 \Meta{div} 是显式指定而非计算得出。
% 在上面的示例中，通过 |\color{rgb,9:red,4;green,2;yellow,1}| 得到一个加深版本 \testclr{rgb,9:red,4;green,2;yellow,1}。
% 注意：指定小于所有 \Meta[\kappa]{dec} 之和的 \Meta{div} 参数并不被禁止，这可能导致最终颜色规格参数超出区间 $[0,1]$。
% 不过，公式 (\ref{eq.U}) 的映射会处理这种情况。
%
%
% \subsubsection{颜色函数}\label{sec.functions}
%
% 颜色函数接受逗号分隔的参数列表，用于将\emph{给定颜色}（即函数调用前所有计算的结果）转换为新颜色。
%
% \paragraph{色轮计算}
% \DescribeFunc{wheel}
% \DescribeFunc{twheel}
% 参数：\Meta{angle} 或 \Meta{angle}|,|\Meta{full circle}，前者是 \Meta{angle}|,\rangeHsb| 的缩写。
% 这些函数允许根据色轮上的和谐关系计算相关颜色（参见第 \vref{sec.harmony} 节）。
% 第二个参数 \Meta{full circle} 声明一个完整圆由多少个单位组成，第一个参数说明给定颜色需要旋转多少个单位。
% 为此，给定颜色首先转换为 \Model{Hsb}（对于 \Function{wheel}），分别得到\emph{hue\Circ}、\emph{saturation}、\emph{brightness}。
% 然后
% \begin{gather}
% \displaystyle\emph{hue\Circ} := \emph{hue\Circ} + \frac{\Meta{angle}}{\Meta{full circle}}\cdot H, \qquad \emph{hue} := u\Bigl( \frac{\emph{hue\Circ}}{H} \Bigr)
% \end{gather}
% 其中 $u$ 是公式~(\ref{eq.U}) 的范围归约函数，$H = \texttt{\string\rangeHsb}$。
% \emph{saturation} 和 \emph{brightness} 保持不变，最终模型为 \Model{hsb}。
% \Function{twheel} 函数类似，但其参数对应于 \Model{tHsb} 而非 \Model{Hsb}。
% 示例见图 \vref{fig.harmony}。
%
%
% \subsection{预定义颜色}
%
% \subsubsection{始终可用的颜色}
%
% \def\test#1{\testclr{#1}~\Color{#1}}
% 在 \File{xcolor.sty} 中定义了以下颜色名称：
% \test{red}, \test{green}, \test{blue}, \test{cyan}, \test{magenta}, \test{yellow}, \test{black}, \test{gray}, \test{white}, \test{darkgray}, \test{lightgray}, \test{brown}, \test{lime}, \test{olive}, \test{orange}, \test{pink}, \test{purple}, \test{teal}, \test{violet}。
%
% 这组基础颜色可在各种颜色表达式中无限制使用，如第 \vref{sec.term} 节所述。
%
%
% \subsubsection{附加颜色集合}\label{sec.additionalcolors}
%
% 还有一些颜色名称集合可通过宏包选项加载，有两种变体：
% "普通"版本（如 \Option{dvipsnames}）和"带星号"版本（如 \Option{dvipsnames*}）。
% 第一种变体简单地\emph{立即}定义所有颜色，而第二种采用\emph{延迟}定义机制。
% 在后一种情况下，各颜色名称必须通过 |\definecolors| 或 |\providecolors| 命令激活，如第 \vref{sec.deferred} 节所述，之后才能在文档中使用。
% \begin{itemize}
% \item \Option{dvipsnames}/\Option{dvipsnames*} 加载 |dvips| 驱动中定义的 68 种 \Model{cmyk} 颜色。不过，这些颜色可在所有支持的驱动中使用。
% \item \Option{svgnames}/\Option{svgnames*} 加载一组 151 个 \Model{rgb} 颜色名称\footnote{实际上这些名称代表 141 种不同颜色。}，符合 \Index{SVG}~1.1 规范 \cite{svg}\footnote{实际上，所引用的规范只列出了小写名称，原始定义以 \Model{RGB} 参数给出，由作者转换为 \Model{rgb}。}，并补充了 4 个取自 \File{rgb.txt} 文件的名称（该文件属于 \Index{Unix}/\Index{X11} 发行版）。注意 \Index{HTML4} 接受 16 个颜色关键字的子集（使用相同规格），见 \cite{HTML4} 和第 \vref{sec.names} 节。
% \item \Option{x11names}/\Option{x11names*} 加载一组 317 个 \Model{rgb} 颜色名称\footnote{这些名称代表 315 种不同颜色。}，它们基本上是前述 \Index{SVG} 集合子集的变体，依据 \Index{Unix}/\Index{X11} 发行版中的 \File{rgb.txt} 文件\footnote{同样，原始定义以 \Model{RGB} 参数给出，由作者转换为 \Model{rgb}。}。
% 我们现在说明如何轻松访问 \File{rgb.txt} 中的全部 752 个颜色名称：
% \begin{itemize}
% \item 同时加载 \Option{x11names} 和 \Option{svgnames}。
% \item 首字母大写并去掉空格：用 \Color[-]{DarkSlateGray} 代替 \Color[-]{dark slate gray}。
% \item 不含数字的 \Index{X11} 名称与对应的 \Index{SVG} 颜色相同，但有 5 个例外：
% 应使用 \Color{Gray0}, \Color{Grey0}, \Color{Green0}, \Color{Maroon0}, \Color{Purple0} 而不是 \Color[-]{Gray}, \Color[-]{Grey}, \Color[-]{Green}, \Color[-]{Maroon}, \Color[-]{Purple}，以获得原始 \Index{X11} 颜色。
% \item 对于 $N=0,1,\dots,100$，使用 `|[gray]{|$N/100$|}|' 或 `|black!|$100-N$' 代替 \Color[-]{grayN} 或 \Color[-]{greyN}。
% \end{itemize}
% \end{itemize}
% 颜色名称及对应色块列于第 \vref{sec.names} 节。
% 第 \vref{sec.clash} 节说明了在同一文档中同时使用 \Option{svgnames} 和 \Option{dvipsnames} 时如何处理名称冲突。
% 另见 \cite{chroma} 中系统的颜色与混合示例。
%
%
% \subsection{颜色定义}\label{sec.def}
%
% \subsubsection{普通颜色与命名颜色}\label{sec.prologue}
%
% 在 \LCP{} 中，区分"颜色"（由 |\definecolor| 命令定义）和"命名颜色"（由 |\DefineNamedColor| 定义，后者仅允许在导言区使用）。
% 每当文档中使用普通颜色时，它会被转换为 |\special| 命令，其中包含颜色的——驱动相关的——数值描述，写入 \File{.dvi} 文件。
% 另一方面，命名颜色提供了在集中位置存储数值的机会，而在使用时颜色可通过名称识别，从而在输出设备需要时进行后处理。
%
% 标准 \Package{graphics} 宏包附带的所有驱动都支持定义和调用"命名颜色"的\emph{形式}。
% 不过，对其背后\emph{概念}（即用名称代替参数）的实际支持程度，从"完全没有"到"完整"不等。
% 我们以三种不同驱动为例说明当前情况：
% \begin{itemize}
% \item \Option{dvips} 对"命名"概念有很好的支持；由 \Option{dvipsnames} 定义的颜色名称对应的 \Index{PostScript} 等价物会被 \Program{dvips} 自动加载——除非关闭。
% 不过，额外的名称需要通过某种头文件告知 \Index{PostScript} 解释器。
% 从 2.01 版开始，\XC{} 为此提供了集成解决方案：启用宏包选项 \Option{prologue} 后，\Program{dvips} 会加载 \Index{PostScript} 头文件 \File{xcolor.pro}。
% 此外，在此选项下，每个颜色定义命令\footnote{不仅导言区如此，正文也一样。}（|\definecolor|, |\colorlet| 等）都会生成一些 \Index{PostScript} 代码，写入扩展名为 \File{.xcp} 的辅助文件（\textbf{xc}olor \textbf prologue 的缩写）。
% 该文件同样由 \Program{dvips} 作为 prologue 加载，从而使所有颜色名称对 \Index{PostScript} 解释器可用。
% 当然，\File{.xcp} 文件可以在运行 \Program{dvips} 之前编辑，便于在集中位置修改设备相关的颜色参数。
% 注意，\Index{PostScript} 代码的设计与 \File{color.pro} 类似：只定义\emph{新}名称。
% 这允许预加载其他包含颜色定义的 prologue 文件，而不会被 \XC 破坏。
% 另一方面，这要求用户自行注意颜色名称重定义的问题。\\
% 示例：|\colorlet{foo}{red}\colorlet{foo}{blue}\color{foo}| 在通常的 \XC{} 逻辑下会切换为 \Color{blue}，但 \File{.ps} 文件会显示 \Color{red}（除非 \Color{foo} 之前已被不同方式定义）。\\
% 需要强调，此机制仅由 \Option{prologue} 选项启用。
% 没有该选项时，由 \Option{dvipsnames} 选项激活的预定义"命名"颜色（不使用任何色调、暗色、颜色表达式等）可以这种方式使用，其他所有"命名"颜色对 \Index{PostScript} 都是未知的。
% \item \Option{dvipdfm} 仅支持标准 \Option{dvipsnames} 颜色，因为这些颜色硬编码在 \Program{dvipdfm} 程序内部；似乎没有办法加载任何用户自定义的 prologue 文件。
% \item \Option{pdftex} 不提供概念上的支持，所有"命名"颜色都会立即转换为其数值表示。因此它允许无限制地定义和使用命名颜色（尽管这样做并没有带来额外价值）。
% \end{itemize}
%
% 通常，\File{.dvi} 查看器在显示用户自定义"命名"颜色时会有困难。
% 例如，\Index{MiK\TeX} 的查看器 \Program{Yap} 目前仅显示 \Option{dvipsnames} 集合中的"命名"颜色。
% 因此，每当同时使用 \Option{prologue} 选项和 \Option{dvips} 时，\emph{所有}其他颜色都会显示为黑色。
% 不过，经过 \Program{dvips} 处理后，\Index{PostScript} 查看器应该能显示正确的颜色。
%
%
% \subsubsection{\XC 中的颜色定义}
%
% \DescribeMacro\definecolor
%   \oarg{type}\marg{name}\marg{model-list}\marg{spec-list}\footnote{2.00 版之前，此命令名为 \cs{xdefinecolor}，后一名称出于兼容性仍可用。}\\
% 这是用于将 \Meta{name} 赋予特定颜色的命令之一。
% 之后，该颜色对系统（当前分组内）已知，可在\emph{\Index{color expression}} 中使用，如第 \vref{sec.term} 节所述。
% 它合并了 \LC 的 |\DefineNamedColor| 和 |\definecolor|。
% 注意，已存在的颜色 \Meta{name} 会被覆盖。
% 变量 |\tracingcolors| 控制是否记录此类覆盖（详见第 \vref{sec.info} 节）。
% 参数说明见第 \vref{sec.term} 节。
% 因此，合法的颜色定义表达式包括：
% \begin{itemize}
% \item |\definecolor{red}{rgb}{1,0,0}|,
% \item |\definecolor{red}{rgb/cmyk}{1,0,0/0,1,1,0}|,
% \item |\definecolor{red}{hsb:rgb/cmyk}{1,0,0/0,1,1,0}|,
% \item |\definecolor[named]{Black}{cmyk}{0,0,0,1}|,
% \item |\definecolor{myblack}{named}{Black}|,
% \end{itemize}
% 其中最后一条命令等价于 |\colorlet{myblack}{Black}|（见下文）；第二条命令根据当前\emph{目标模型}的设置，在 \Model{rgb} 或 \Model{cmyk} 模型中定义 \Color{red}，而第三条命令会在保存前将颜色额外转换为 \Model{hsb}。
% 注意，存在一个与 \Package{pstricks} 相关的特殊版本，见第 \vref{sec.pstricks} 节。
%
% \medskip
%
% \DescribeMacro\providecolor
%   \oarg{type}\marg{name}\marg{model-list}\marg{spec-list}\\
% 与 |\definecolor| 类似，但仅当颜色 \Meta{name} 尚不存在时才定义。
%
% \medskip
%
% \DescribeMacro\colorlet
%   \oarg{type}\marg{name}\oarg{num model}\marg{color}\\
% 将 \Meta{color} 计算得到的实际颜色复制到 \Meta{name}。
% 如果 \Meta{num model} 非空，则先将 \Meta{color} 转换为指定模型，再定义 \Meta{name}。
% 伪模型 \nModel{named} 在此处\emph{不}允许使用，但可以在 \Meta{type} 参数中指定。
% 注意，已存在的颜色 \Meta{name} 会被覆盖。
%
% 示例：我们在本文档导言区写了 |\colorlet{tableheadcolor}{gray!25}|。
% 在大多数表格中，我们随后用 |\rowcolor{tableheadcolor}| 命令格式化第一行。
%
%
% \subsubsection{定义颜色集合}
%
% \DescribeMacro\definecolorset
%   \oarg{type}\marg{model-list}\marg{head}\marg{tail}\marg{set spec}\\
% 此命令便于构造\emph{\Index{color set}}，即（可能很大的）一组具有共同底层 \Meta{model-list} 和 \Meta{type} 的独立颜色。
% 其中 \Meta{set spec} = \Meta[1]{name},\Meta[1]{spec-list};\dots;\Meta[l]{name},\Meta[l]{spec-list}（$l\ge1$ 个 名称/规格列表 对）。
% 各颜色由单独的
% \begin{center}
% |\definecolor|\oarg{type}|{|\Meta{head}\Meta[\lambda]{name}\Meta{tail}|}|\marg{model-list}|{|\Meta[\lambda]{spec-list}|}|
% \end{center}
% 命令构造，$\lambda=1,\dots,l$。
% 例如：
% \begin{itemize}
% \item |\definecolorset{rgb}{}{}{red,1,0,0;green,0,1,0;blue,0,0,1}|\\
% 可用于定义基色 \Color{red}, \Color{green}, \Color{blue}；\footnote{实际上，\XC{} 使用了更复杂的变体来为不同底层模型提供基色（完整命令见源码）：\\
% \cs{definecolorset\string{rgb/hsb/cmyk/gray\string}\string{\string}\string{\string}\string{red,1,0,0/0,1,1/0,1,1,0/.3;green,...\string}}.}
% \item |\definecolorset{rgb}{x}{10}{red,1,0,0;green,0,1,0;blue,0,0,1}|\\
% 将定义颜色 \Color[-]{xred10}, \Color[-]{xgreen10}, \Color[-]{xblue10}。
% \end{itemize}
%
% \DescribeMacro\providecolorset
%   \oarg{type}\marg{model-list}\marg{head}\marg{tail}\marg{set spec}\\
% 与 |\definecolorset| 类似，但基于 |\providecolor|，因此各颜色仅在不存在时才定义。
%
%
% \subsubsection{立即定义与延迟定义}\label{sec.deferred}
%
% 传统上，上述颜色定义会立即构造一个命令，其中至少包含驱动显示所需颜色的信息。
% 因此，定义 300 种颜色（例如加载一大组预定义颜色）会产生 300 个新命令，而其中大多数——除了展示颜色列表外——几乎不会在文档中使用。
% 随着计算机内存的发展——容量增大、价格降低——现代 \TeX{} 实现增加了可供字符串、控制序列等使用的内部内存栈。
% 不过，由于内存终究有限，掌握一种能略微降低内存需求的方法仍然是有用的（有时甚至是必要的）。
% 这就是\emph{延迟颜色定义}的用武之地。
% 其原理很简单：对于此类定义（例如通过 |\preparecolor|），所有必要信息保存在一个特定的全局\emph{\Index{definition stack}} 上，之后可从中取出（例如通过 |\definecolors|）以构造实际的颜色命令。
%
% 注意，以下命令只能在文档导言区使用，因为延迟定义的颜色定义栈会在文档正文开始时被清空——以节省内存。
%
% \medskip
%
% \DescribeMacro\preparecolor
%   \oarg{type}\marg{name}\marg{model-list}\marg{spec-list}\\
% 与 |\definecolor| 类似，但颜色 \Meta{name} 暂不定义：立即求值 \Meta{model-list} 和 \Meta{spec-list} 参数，然后将所有必要参数（即 \Meta{type}, \Meta{name}, \Meta{model}, \Meta{spec}）放入\emph{\Index{definition stack}} 供以后使用。
%
% \medskip
%
% \DescribeMacro\preparecolorset
%   \oarg{type}\marg{model-list}\marg{head}\marg{tail}\marg{set spec}
% \DescribeMacro\ifdefinecolors\\
% 与 |\definecolorset| 类似，但取决于 |\ifdefinecolors| 开关：若为"真"，对集合中每个元素应用 |\definecolor| 命令（即立即定义）；若为"假"，应用 |\preparecolor|（即延迟定义）。
% 例如，宏包选项 \Option{svgnames} 执行类似 |\definecolorstrue\preparecolorset| 的操作，而 \Option{svgnames*} 类似 |\definecolorsfalse\preparecolorset|。
% 两个选项最后都会设置 |\definecolorstrue|，以便为其他颜色集合提供正确的起点。
%
% \medskip
%
% \DescribeMacro\DefineNamedColor
%   \marg{type}\marg{name}\marg{model-list}\marg{spec-list}
% 主要出于兼容性考虑而提供，特别是为了支持 \File{dvipsnam.def} 中的预定义颜色。
% 它等价于 \Meta{cmd}\oarg{type}\marg{name}\marg{model}\marg{spec}，其中 \Meta{cmd} 根据 |\ifdefinecolors| 的状态为 |\definecolor| 或 |\preparecolor|。
% 注意，\LC 中仅允许在导言区使用 |\DefineNamedColor| 的限制在 \XC 中已取消。
%
% \medskip
%
% \DescribeMacro\definecolors
%   \marg{id-list}\\
% 回顾一下，\Meta{id-list} 的形式为 \Meta[1]{ext id}|,...,|\Meta[l]{ext id}，其中每个 \Meta[\lambda]{ext id} 要么是标识符 \Meta[\lambda]{id}，要么是赋值 \Meta[\lambda']{id}|=|\Meta[\lambda]{id}。
% 我们把第一种情况视为 \Meta[\lambda]{id}|=|\Meta[\lambda]{id} 的缩写，并说明一般情形：
% 在\emph{\Index{definition stack}} 中搜索名称 \Meta[\lambda]{id} 及其对应的颜色参数；若没有匹配，则什么也不做；若名称 \Meta[\lambda]{id} 在栈上，且其颜色参数为 \Meta[\lambda]{type}, \Meta[\lambda]{model}, \Meta[\lambda]{spec}，则执行命令 |\definecolor[|\Meta[\lambda]{type}|]{|\Meta[\lambda']{id}|}{|\Meta[\lambda]{model}|}{|\Meta[\lambda]{spec}|}|。
% 因此，用户可以控制文档中使用\emph{已准备好的}颜色时所采用的名称。
% 注意，栈中的条目 \Meta[\lambda]{id} 不会被移除，因此可多次使用（甚至在同一个 |\definecolors| 命令内）。
%
% \medskip
%
% \DescribeMacro\providecolors
%   \marg{id-list}\\
% 与 |\definecolors| 类似，但基于 |\providecolor|，因此各颜色仅在不存在时才定义。
%
%
% \subsubsection{全局颜色定义}
%
% \DescribeMacro\ifglobalcolors
% 默认情况下，通过 |\definecolor|、|\providecolor| 等的定义仅在当前分组内有效。
% 设置 |\globalcolorstrue| 后，所有此类定义将全局有效——直到当前分组结束。\footnote{该开关也可在导言区设置，以控制整个文档。}
% \DescribeMacro\xglobal
% 另一种指定单个颜色定义为全局的方法是在其前加 |\xglobal|，例如 |\xglobal\definecolor{foo}...|。
%
%
%
% \subsection{颜色应用}
%
% \subsubsection{标准颜色命令}
%
% 以下是 \LCP 中已知的用户级颜色命令列表，但颜色语法已扩展，支持表达式等：\\
% \DescribeMacro\color\marg{color}\\
% \oarg{model-list}\marg{spec-list}\\
% 切换到由名称/表达式或由模型/规格给出的颜色。
% 该颜色在当前 \TeX\ 分组结束前持续有效。\\
% \DescribeMacro\textcolor\marg{color}\marg{text}\\
% \oarg{model-list}\marg{spec-list}\marg{text}\\
% 只是 |\color| 的另一种语法，分组由隐式添加。
% 因此 \meta{text} 以指定颜色显示，但之后颜色恢复为之前的值。
% 此外，它会调用 |\leavevmode| 以确保进入水平模式。\\
% \DescribeMacro\pagecolor\marg{color}\\
% \oarg{model-list}\marg{spec-list}\\
% 指定当前页及后续所有页的背景颜色。
% 这是全局声明，不受 \TeX 分组影响。
%
% 注意：除 |\color| 外，所有这些命令的 \Meta{color} 或 \Meta{spec} 参数都必须放在花括号 |{}| 中，即使它们藏在宏内部。
%
% 例如，定义 |\def\foo{red}| 后，可以写 |\color\foo|，但应始终写 |\textcolor{\foo}{bar}| 而不是 |\textcolor\foo{bar}|，以避免奇怪结果。
%
% 注意，其他宏包的颜色相关命令如果直接面对颜色表达式，可能产生意外结果（例如 \Package{soul} 的 |\sethlcolor| 等）。
% 不过，可以通过 |\colorlet| 将表达式转换为名称，然后尝试使用该名称。
%
% \DescribeMacro\nopagecolor
% 移除当前指定的页面颜色，恢复默认透明背景。
% 并非所有驱动选项都支持此命令，因此如果驱动文件中没有定义，会产生警告。
%
%
% \subsubsection{颜色框}\label{sec.fbox}
%
% \DescribeMacro\colorbox\marg{color}\marg{text}\\
% \oarg{model-list}\marg{spec-list}\marg{text}\\
% 参数形式与 |\textcolor| 相同，但颜色指定的是框的\emph{背景}色。\\
% \DescribeMacro\fcolorbox\marg{frame color}\marg{background color}\marg{text}\\
% \oarg{model-list}\marg{frame spec-list}\marg{background spec-list}\marg{text}\\
% \oarg{fr.~model-list}\marg{fr.~spec-list}\oarg{backgr.~model-list}\marg{backgr.~spec-list}\marg{text}\\
% \marg{frame color}\oarg{background model-list}\marg{background spec-list}\marg{text}\\
% 用第一种颜色画边框，围绕一个以第二种颜色为背景的框。
% 如果只给出第一个可选参数，则它同时指定两种颜色的颜色模型。
% 除了可以将颜色\emph{表达式}作为参数外，|\fcolorbox| 现在的参数比 \LC{} 版本更灵活：
% \begin{itemize}
% \item \fcolorbox{gray}{yellow}{test} |\fcolorbox{gray}{yellow}{test}|,
% \item \fcolorbox[cmyk]{0,0,0,0.5}{0,0,1,0}{test} |\fcolorbox[cmyk]{0,0,0,0.5}{0,0,1,0}{test}|,
% \item \fcolorbox[gray]{0.5}[wave]{580}{test} |\fcolorbox[gray]{0.5}[wave]{580}{test}|,
% \item \fcolorbox{gray}[wave]{580}{test} |\fcolorbox{gray}[wave]{580}{test}|.
% \end{itemize}
% 此外，|\fcolorbox| 采用了一种新的边框绘制方法，这是 \People{Donald}{Arseneau} 在 bug 报告 latex/3655 \cite{fbox} 中建议的扩展。
% 与 \LaTeX 实现的主要区别在于，框构造和边框绘制被拆分为独立操作，因此边框在框内容构造\emph{之后}绘制。
% 这确保边框始终在框的最上层。
% \People{Donald}{Arseneau} 还改进了此方法的速度和内存需求。
% 此外，引入了一个新宏：
%
% \DescribeMacro\boxframe\marg{width}\marg{height}\marg{depth}\\
% 绘制线宽为 |\fboxrule| 的边框。
% 返回一个外层尺寸为 \meta{width}, \meta{height}, \meta{depth} 的 |\hbox|。
% 通过这种方式，驱动文件也可以提供边框图元，以利用驱动特定的绘制设施（见下文）。
% 同样，此宏由 \People{Donald}{Arseneau} 优化。
%
% 新的边框方法用于 |\fcolorbox| 以及 \LaTeX 的 |\fbox| 和 |\framebox| 命令，除非指定 \Option{kernelfbox} 选项，后者恢复 \LaTeX 对 |\f(rame)box| 的原始定义。
%
% 选项 \Option{xcdraw} 在 \Option{dvips} 驱动下使用 \Index{PostScript} 命令绘制边框和颜色框，在 \Option{pdftex} 和 \Option{dvipdfm} 驱动下使用 \Index{PDF} 代码绘制边框。
% 这仍是实验性代码，可能会使 \File{.dvi} 查看器混淆。
% 相反的选项 \Option{noxcdraw} 强制使用通用（与驱动无关）代码。
%
% \subsubsection{数学中的颜色}
%
% 虽然可以使用 |\color| 命令给公式的部分内容上色，但这种方法相当繁琐。
% 因此，\cs{mathcolor} 命令被添加到 \Package{color} 和 \Package{xcolor} 宏包中，
% 其语法与 \cs{textcolor} 相同。完整文档和使用示例见 \LaTeX 文档中的 \texttt{mathcolor.pdf}。
%
% \subsubsection{使用当前颜色}
%
% 在颜色表达式中，`.` 作为当前颜色的占位符。
% 示例见图 \vref{fig.cc}。
%
% 也可以保存当前颜色供以后使用，例如通过命令 |\colorlet{foo}{.}|。
%
% 注意，在某些情况下当前颜色的用途相当有限，例如构造 |\fcolorbox| 时，在求值 \Meta{background color} 的时刻，当前颜色等于 \Meta{frame color}；此时 `|.|' 并不指框\emph{外部}的当前颜色。
%
%
% \subsubsection{颜色测试}
%
% \DescribeEnv{testcolors}
%   \oarg{num models}\\
% 这是一个简单的表格环境，用于测试（显示）不同模型中的颜色，同时展示视觉结果和模型相关参数。
% 可选参数 \Meta{num models} 是逗号分隔的\emph{数值}颜色模型列表（通常无空格），构成表格列；默认列表为
% {\makeatletter\ttfamily\let\@protected@testopt\@gobbletwo\edef\tmp{\testcolors}\expandafter\@firstofone\tmp}。
%
% \DescribeMacro\testcolor
%   \marg{color}\\
%   \oarg{model-list}\marg{spec-list}\\
% 每个 |\testcolor| 命令生成一行表格，包含一个显示样本以及各模型对应的参数。
% 如果列模型与所讨论颜色的模型匹配，则其参数加下划线。
% 注意，此命令仅在 |testcolors| 环境内可用。
%
% 应用示例见图 \vref{fig.testcolors} 以及图 \ref{fig.tHsb}, \ref{fig.harmony}。
%
%
%
% \subsection{颜色混合}
%
% \emph{颜色混合（color blending）}的目的是为后续所有显式颜色命令添加某种混合颜色（表达式）。
% 这样就可以对许多颜色执行此类混合操作，而无需逐个修改命令。
%
% \medskip
%
% \DescribeMacro\blendcolors
%   \marg{mix expr}\\
% \DescribeMacro{\blendcolors*}
%   \marg{mix expr}\\
% 初始化颜色混合所需的所有参数。
% 实际（完整的）颜色混合表达式存储在 |\colorblend| 中。
% 带星号版本会将参数追加到先前定义的混合表达式之后。
% 空的 \Meta{mix expr} 参数会关闭混合。
%
% 示例：执行 |\blendcolors{!50!yellow}| 后，颜色 \testclr{red} \testclr{green} \testclr{blue} 变为 {\blendcolors{!50!yellow}\testclr{red} \testclr{green} \testclr{blue}，再执行 |\blendcolors*{!50}| 得到 \blendcolors*{!50}\testclr{red} \testclr{green} \testclr{blue}}。
% \DescribeMacro\xglobal
% 为实现全局作用域，|\blendcolors| 前可加 |\xglobal|。
%
% 注意：颜色混合仅应用于\emph{显式}颜色命令，即 |\color|、|\fcolorbox| 等。
% 上例中的边框没有被混合，因为其颜色由驱动内部命令设置（切换回"当前颜色"）。
% 因此，要影响这些\emph{隐式}颜色，必须在混合\emph{之后}设置当前颜色：
% |\blendcolors{!50!yellow}\color{black}| 得到 {\blendcolors{!50!yellow}\color{black}\testclr{red} \testclr{green} \testclr{blue}，再执行 |\blendcolors*{!50}\color{black}| 得到 \blendcolors*{!50}\color{black}\testclr{red} \testclr{green} \testclr{blue}}。
%
%
%
% \subsection{颜色蒙版与分色}
%
% \emph{颜色分色（color separation）}的目的是将文档中出现的所有颜色表示为有限基色子集及其色调的组合。
% 最常见的是 \Model{cmyk} 分色，基色为 \Color{cyan}, \Color{magenta}, \Color{yellow}, \Color{black}，这是打印机所要求的。
% 这可以通过选择宏包选项 \Option{cmyk} 来实现，使所有颜色都转换为此模型，然后后处理输出文件。
% 我们现在介绍另一种——更通用的——解决方案：\emph{颜色蒙版（color masking）}。
% 它如何工作？
% 颜色蒙版基于指定的颜色模型 \Meta{m-model} 和参数向量 \Meta{m-spec}。
% 每当文档中要显示某种颜色时，先将其转换为 \Meta{m-model}，然后将所得颜色向量的每个分量乘以 \Meta{m-spec} 的对应分量。
% 例如，假设 \Meta{m-model} 为 |cmyk|，\Meta{m-spec} 为 $(\mu_c,\mu_m,\mu_y,\mu_k)$。
% 那么任意颜色 \Color{foo} 将按以下方式变换：
% \begin{gather}
% \Color{foo} \mapsto (c,m,y,k) \mapsto (\mu_c\cdot c,\mu_m\cdot m,\mu_y\cdot y,\mu_k\cdot k)
% \end{gather}
% 显然，颜色分色是用向量 $(1,0,0,0)$、$(0,1,0,0)$ 等进行蒙版的特例。
% 一个有趣的应用是在 \Model{rgb} 或 \Model{cmy} 模型中用 $(x,x,x)$ 蒙版来加深或提亮所有颜色，见图 \vref{fig.mask} 的最后两行。
%
% \medskip
%
% \DescribeMacro\maskcolors
%   \oarg{num model}\marg{color}\\
% 初始化颜色蒙版所需的所有参数：
% 如果未指定 \Meta{num model}（或为空），则 \Meta{m-model} 设为 \Meta{color} 的自然模型，否则设为 \Meta{num model}；
% 提取 \Meta{color} 的颜色规格以定义 \Meta{m-spec}。
% 此外，会执行 |\maskcolorstrue|。
% \DescribeMacro\ifmaskcolors
% 颜色蒙版可通过 |\maskcolorsfalse| 临时关闭，或——更彻底地——通过 |\maskcolors{}| 关闭，后者还会清除初始化参数。
% \DescribeMacro\xglobal
% 一般来说，|\maskcolors| 的作用域是当前分组（除非前面加了 |\xglobal| 命令），但也可在文档导言区使用。
% 颜色混合一节的最后一条注意事项同样适用于此。
%
% 现在，无需修改源代码即可轻松分色整个文档：
% |latex| |\def\xcolorcmd{\maskcolors[cmyk]{cyan}}\input{a}| 会完成 |a.tex| 的 \Color{cyan} 分色部分。
%
% \DescribeMacro\colormask
% 注意：\XC{} 无法感知通过 |\includegraphics| 命令包含的文件中的颜色，例如 \File{.eps}、\File{.pdf}、\File{.jpg} 或 \File{.png} 类型的图像。
% 这些文件必须单独分色。
% 尽管如此，\XC{} 提供了一些基本支持：将蒙版颜色存储在 |\colormask| 中，可用于决定包含哪个文件：
% \begin{verbatim}
% \def\temp{cyan}\ifx\colormask\temp \includegraphics{foo_c}\else
% \def\temp{magenta}\ifx\colormask\temp \includegraphics{foo_m}\else
% ...
% \fi\fi
% \end{verbatim}
%
%
%
% \subsection{颜色系列}\label{sec.series}
%
% 在图形或图表应用中，自动着色可能很有用，因为这类应用需要——可能数量庞大且未指定的——颜色，而用户不想或无法逐一指定每种颜色。
% 因此，我们引入\emph{颜色系列（color series）}的概念，它由一个基色和一个方案组成，用于从当前颜色构造下一种颜色。
%
% 实际应用包括三部分：定义颜色系列（通常在文档中一次）、初始化系列（可能多次）、以及应用——带或不带步进——系列的当前颜色（可能很多次）。
%
%
% \subsubsection{颜色系列的定义}
%
% \DescribeMacro\definecolorseries
%   \marg{name}\marg{core model}\marg{method}\oarg{b-model}\marg{b-spec}\oarg{s-model}\marg{s-spec}\\
% 定义名为 \Meta{name} 的颜色系列，其计算在颜色模型 \Meta{core model} 内进行，其中 \Meta{method} 选择算法（|step|、|grad|、|last| 之一，见下文）。
% 方法细节由其余参数决定：
% \begin{itemize}
% \item \oarg{b-model}\marg{b-spec} 指定算法中的\emph{基色}（即第一个颜色），可直接指定，例如 |[rgb]{1,0.5,0.5}|；如果省略可选参数，也可作为 \Meta{color} 指定，例如 |{-yellow!50}|。
% \item \oarg{s-model}\marg{s-spec} 根据所选 \Meta{method} 指定算法中\emph{步进向量}的计算方式：
%   \begin{itemize}
%   \item |step|、|grad|：可选参数无意义，\Meta{s-spec} 是参数向量，其维数由 \Meta{core model} 决定，例如 |rgb|、|cmy| 或 |hsb| 时为 |{0.1,-0.2,0.3}|。
%   \item |last|：最后一个颜色可直接指定，例如 |[rgb]{1,0.5,0.5}|；如果省略可选参数，也可作为 \Meta{color} 指定，例如 |{-yellow!50}|。
%   \end{itemize}
% \end{itemize}
% 这是一般方案：
% \begin{gather}
% \emph{color}_1 := \emph{base},\qquad
% \emph{color}_{n+1} := U\bigl(\emph{color}_{n} + \emph{step}\bigr)
% \end{gather}
% 其中 $n=1,2,\dots$，$U$ 将任意实数 $m$ 维向量映射到单位 $m$ 维立方体：
% \begin{gather}
% U(x_1,\dots,x_m) = (u(x_1),\dots,u(x_m)),\qquad
% u(x) =
% \begin{cases}
%  1 & \text{if $x=1$}\\
%  x-[x]& \text{if $x\neq1$}\\
% \end{cases}\label{eq.U}
% \end{gather}
% 因此，算法每一步都产生一个参数在区间 $[0,1]$ 内的合法颜色。
%
% \medskip
% 不同方法使用不同方案计算\emph{步进}向量：
% \begin{itemize}
% \item |step|、|grad|：最后一个参数 \marg{s-spec} 定义方向向量 \emph{grad}。
% \item |last|：\marg{s-spec} 或 \oarg{s-model}\marg{s-spec} 定义颜色参数向量 \emph{last}。
% \end{itemize}
% 然后，在 |\resetcolorseries| 时，计算实际的\emph{步进}向量：
% \begin{gather}
% \emph{step} :=
% \begin{cases}
% \emph{grad} & \text{if $\Meta{method} = |step|$}\\
% \frac1{\Meta{div}}\cdot\emph{grad} & \text{if $\Meta{method} = |grad|$}\\
% \frac1{\Meta{div}}\cdot(\emph{last}-\emph{base}) & \text{if $\Meta{method} = |last|$}\\
% \end{cases}\label{eq.cs}
% \end{gather}
% Please note that it is also possible to use the current color placeholder `.' within the definition of color series.
% Thus, |\definecolorseries{foo}{rgb}{last}{.}{-.}| will set up a series that starts with the current color and ends with its complement.
% 当然，与 \TeX 的 |\let| 原语类似，使用的是执行时刻当前颜色的\emph{当前}定义，与文档后续阶段的当前颜色无关。
%
%
% \subsubsection{颜色系列的初始化}
%
% \DescribeMacro\resetcolorseries
%   \oarg{div}\marg{name}\\
% 此命令至少需要执行一次，才能使用颜色系列 \Meta{name}。
% 它将系列的当前颜色重置为基色，并根据所选 \Meta{div}（一个非零实数）为 |grad| 和 |last| 方法计算实际步进向量，见公式 (\ref{eq.cs})。
% \DescribeMacro\colorseriescycle
% 如果可选参数为空，则使用宏 |\colorseriescycle| 中存储的值。
% 其默认值为 \colorseriescycle，可通过 |\def\colorseriescycle|\marg{div} 修改，但必须在 \XCP{} 加载\emph{之前}应用（与 |\rangeRGB| 等类似）。
% 对于 |step| 方法，可选参数被忽略。
%
%
% \subsubsection{颜色系列的应用}
%
% 显示颜色系列当前颜色有两种方式：
% 在 |\color|、|\textcolor| 等命令中使用第 \vref{sec.term} 节中的任意\emph{颜色表达式}，将按此类表达式的通常语法显示该颜色。
% 不过，当 \Meta{postfix} 为 `|!!+|' 时，|\color{|\Meta{name}|!!+}| 等不仅会显示颜色，还会执行一次步进操作。
% 因此，这种情况下系列的当前颜色会被改变。
% 表达式 |\color{|\Meta{name}|!![|\Meta{num}|]}| 可直接访问系列中的某个元素，其中 \Meta{num} = $0,1,2,\dots$，基色从 0 开始。
% 不同方法的演示见图 \vref{fig.csx}。
%
%
% \subsubsection{颜色与颜色系列的区别}
%
% 虽然在颜色表达式中使用时行为相似，但 |\definecolor| 和 |\definecolorseries| 定义的对象在作用域/可用性上有根本区别：
% 与 \LC 原始 |\definecolor| 命令一样，|\definecolor| 生成\emph{局部}颜色，而 |\definecolorseries| 生成\emph{全局}对象（否则无法方便地在表格或图形中使用步进机制）。
% 例如，假设 |bar| 是未定义颜色，那么在执行
% \begin{verbatim}
% \begingroup
% \definecolorseries{foo}{rgb}{last}{red}{blue}
% \resetcolorseries[10]{foo}
% \definecolor{bar}{rgb}{.6,.5,.4}
% \endgroup
% \end{verbatim}
% 之后，|\color{foo}| 或 |\color{foo!!+}| 等命令可无限制使用，而 |\color{bar}| 会报错。
% 不过，可以使用 |\colorlet{bar}{foo}| 或 |\colorlet{bar}{foo!!+}| 来在本地保存系列的当前颜色——带或不带步进。
%
%
%
%
% \subsection{超链接的边框颜色}\label{sec.hyperref}
%
% \textbf{2022 年变更}：本节描述的选项已从代码中移除，因为不再相关。
% \Package{hyperref} 现在可以在边框颜色键中直接处理颜色名称。
%
% 这影响宏包选项 \Option{hyperref}，以及键
% \Key{xcitebordercolor}, \Key{xfilebordercolor}, \Key{xlinkbordercolor}, \Key{xmenubordercolor},
% \Key{xpagebordercolor}, \Key{xrunbordercolor}, \Key{xurlbordercolor} 和 \Key{xpdfborder}。
%
%
% \Package{hyperref} 宏包为超链接、pdf 书签等提供了各种支持。
% 使超链接可见有两种标准方式（关于如何设置这些功能的更多信息，见宏包文档 \cite{hyperref}）：
% \begin{itemize}
% \item 用与正文不同的颜色打印超链接，使用键 \Key{citecolor}, \Key{filecolor}, \Key{linkcolor}, \Key{menucolor}, \Key{pagecolor}, \Key{runcolor}, \Key{urlcolor} 并配合颜色表达式，例如 |\hypersetup{urlcolor=-green!50}|；
% \item 在超链接周围显示彩色边框，使用键 \Key{citebordercolor}, \Key{filebordercolor}, \Key{linkbordercolor}, \Key{menubordercolor}, \Key{pagebordercolor}, \Key{runbordercolor}, \Key{urlbordercolor} 并配合显式数值 \Model{rgb} 参数，例如 |\hypersetup{urlbordercolor={1 0.5 0.25}}|。
% \end{itemize}
% 显然，第二种方法不太方便，因为它不允许使用颜色名称甚至颜色表达式。
% 因此，\XC{} 通过宏包选项 \Option{hyperref} 提供了一组扩展键 \Key{xcitebordercolor}, \Key{xfilebordercolor}, \Key{xlinkbordercolor}, \Key{xmenubordercolor}, \Key{xpagebordercolor}, \Key{xrunbordercolor}, \Key{xurlbordercolor}，可与颜色表达式配合使用，例如 |\hypersetup{xurlbordercolor=-green!50}|。
%
% 另一个新键 \Key{xpdfborder} 提供了处理 \Program{dvips} 相关问题的方法：
% 对于大多数驱动，|pdfborder={0 0 1}| 这样的设置会以点为单位确定超链接周围绘制的边框宽度。
% 但在 \Program{dvips} 情况下，数值参数是相对于将 \File{.dvi} 文件处理为 \File{.ps} 文件时所选输出分辨率来解释的。
% 不幸的是，在构造 \File{.dvi} 时，没有人知道之后是否以及以何种分辨率转换为 \File{.ps}。
% 因此，\Key{pdfborder} 的任何默认值可能有用也可能没用。
% 在 \Package{hyperref} 中，\Program{dvips} 的默认值是 |pdfborder={0 0 12}|，这在 600 或 1200 dpi 分辨率下工作良好，但在 8000 dpi（由命令行开关 |-Ppdf| 决定）下会产生不可见边框。
% 另一方面，设置 |pdfborder={0 0 80}| 在 8000 dpi 的 \Program{dvips} 下工作良好，但会使文档不可移植，因为其他驱动（甚至低分辨率下的 |dvips|）在这种情况下会画出非常粗的框。
% 这就是 \Key{xpdfborder} 键派上用场的地方：它对 |dvips| 的情况将参数缩放 80 倍（适配 8000 dpi），对其他驱动保持不变。
% 因此可以以与驱动无关的方式写 |xpdfborder={0 0 1}|。
%
%
%
% \subsection{\Package{pstricks} 环境中的附加颜色指定}\label{sec.pstricks}
%
% 对于 \Package{pstricks} 用户，在命令选项键中调用颜色有多种方式：
% \begin{itemize}
% \item |\psset{linecolor=green!50}|
% \item |\psset{linecolor=[rgb]{0.5,1,0.5}}|
% \item |\psframebox[linecolor={[rgb]{0.5,1,0.5}}]{foo}|
% \end{itemize}
% 注意最后一种情况中额外的花括号；没有它们，|\psframebox| 的可选参数会过早结束。
%
% \DescribeMacro\definecolor
%   |[ps]|\marg{name}\marg{core model-list}\marg{code}\\
% 将 \Index{PostScript} \Meta{code}（不应包含斜杠 `|/|' 字符）存储在颜色中。
% 示例：
% 执行 |\definecolor[ps]{foo}{rgb}{bar}| 后，\Package{pstricks} 命令 |\psline[linecolor=foo]...| 会在需要 linecolor 信息的地方插入 `|bar setrgbcolor|'——至少在 \Option{dvips} 驱动下如此。
% 另见 \href{xcolor2.pdf}{\File{xcolor2.tex}} 中的示例应用。
%
%
%
% \subsection{表格中的颜色}
% 本节中的命令
% 在以前版本中由 \Package{xcolor} 宏包作为
% \Package{colortbl} 命令的扩展提供。现在它们已
% 正式合并入 \Package{colortbl} 宏包。该宏包可通过
% \Option{table} 选项加载，也可通过 |\usepackage| 独立加载。
% 命令说明仅出于历史原因保留。如有疑问，
% 请查阅 \Package{colortbl} 的文档。
%
% \DescribeMacro\rowcolors
%   \oarg{commands}\marg{row}\marg{odd-row color}\marg{even-row color}\\
% \DescribeMacro{\rowcolors*}
%   \oarg{commands}\marg{row}\marg{odd-row color}\marg{even-row color}\\
% 这些命令之一必须在表格开始\emph{之前}执行。
% \Meta{row} 指定第一行应按 \Meta{odd-row color} 和 \Meta{even-row color} 方案着色的行号。
% 每个颜色参数也可以留空（=不着色）。
% 在带星号版本中，\Meta{commands} 在"rowcolors 状态"未激活的行中被忽略（见下文）；而在不带星号版本中，\Meta{commands} 应用于表格的每一行。
% 此类可选命令可以是 |\hline| 或 |\noalign|\marg{stuff}。
%
% \DescribeMacro\showrowcolors
% \DescribeMacro\hiderowcolors
% "rowcolors 状态"默认激活（即使用着色方案），也可通过 |\showrowcolors| 激活；通过 |\hiderowcolors| 命令可停用（即忽略着色方案）。
% \DescribeMacro\rownum
% 计数器 |\rownum| 可在此类表格中访问当前行号。
% 示例见图 \vref{fig.row}。
%
% 注意，表格着色可与颜色系列结合使用。
% 图 \vref{fig.csx} 中的示例就是用这种方法构造的。
%
%
%
% \subsection{颜色信息}\label{sec.info}
%
% \DescribeMacro\extractcolorspec
%   \marg{color}\marg{cmd}\\
% 提取 \Meta{color} 的颜色规格并存入 \Meta{cmd}；等价于 |\def\cmd{|\marg{model}\marg{spec}|}|。
%
% \DescribeMacro\extractcolorspecs
%   \marg{color}\marg{model-cmd}\marg{color-cmd}\\
% 提取 \Meta{color} 的颜色规格，分别存入 \Meta{model-cmd} 和 \Meta{color-cmd}。
%
% \medskip
% \DescribeMacro\tracingcolors
%   |=|\Meta{int}\\
% 控制写入 |log| 文件的信息量：
% \begin{itemize}
% \item \mbox{$\Meta{int}\le0$：}~不记录特定颜色信息。
% \item \mbox{$\Meta{int}\ge1$：}~记录因 |\providecolor| 而被忽略的颜色定义。
% \item \mbox{$\Meta{int}\ge2$：}~记录重复（即被覆盖的）颜色定义。
% \item \mbox{$\Meta{int}\ge3$：}~记录每个定义颜色的命令。
% \item \mbox{$\Meta{int}\ge4$：}~记录每个设置颜色的命令。
% \end{itemize}
% 与 \TeX 的 |\tracing...| 命令一样，此命令可全局使用（在文档导言区）或局部/分组使用。
% 宏包默认设置 |\tracingcolors=0|。
% 注意：由于寄存器有限且宝贵，此功能不浪费任何计数器。
%
% 注意，每当使用通过 \LC 的 |\definecolor| 命令（而非 \XC 新的 |\definecolor| 等）定义的颜色时，会发出警告"Incompatible color definition"。\footnote{这不应该发生，因为通常没有理由在加载 \XC{} 的同时并行加载 \LC{}。}
%
%
%
% \subsection{颜色转换}
%
% \DescribeMacro\convertcolorspec
%   \marg{model}\marg{spec}\marg{target model}\marg{cmd}\\
% 将由 \Meta{model} 中 \Meta{spec} 给出的颜色转换为 \Meta{target model}，并将新颜色规格存入 |\cmd|。
% \Meta{target model} 必须是 \Meta{num model} 类型，而 \Meta{model} 也可以是 `named'，此时 \Meta{spec} 就是颜色名称。
%
% 示例：|\convertcolorspec{cmyk}{0.81,1,0,0.07}{HTML}\tmp| 等价于 |\def| |\tmp{1F00ED}|。
%
%
%
% \subsection{问题与解决方案}
%
% \subsubsection{\Option{dvipsnames} 与 \Option{svgnames} 之间的名称冲突}\label{sec.clash}
%
% 由于固定的选项处理顺序（与选项在 |\usepackage| 命令中指定的顺序无关），\Option{svgnames} 颜色总是会覆盖同名的 \Option{dvipsnames} 颜色。
% 如果同时使用这两个选项，可能导致不希望的结果。
% 例如，\Color{Fuchsia} 在 \Option{dvipsnames} 下为 \testclr[cmyk]{0.47,0.91,0,0.08}，而在 \Option{svgnames} 下为 \testclr[rgb]{1,0,1}。
% 不过，有一个简单技巧——基于\emph{延迟颜色定义}——可以让我们按期望的方式使用两个集合中的颜色：
% \begin{verbatim}
% \usepackage[dvipsnames*,svgnames]{xcolor}
% \definecolors{Fuchsia}
% \end{verbatim}
% 现在 \Index{SVG} 集合中的所有颜色都可用（除了 \Color{Fuchsia}），外加另一集合中的 \Color{Fuchsia}。
%
%
% \subsubsection{更改嵌入的 \File{.eps} 文件颜色}\label{sec.fixinclude}
%
% 一般来说，\XC{} 无法更改通过 \Package{graphics} 或 \Package{graphicx} 宏包的 |\includegraphics| 命令嵌入的图像颜色。
% 不过，在有限范围内可以影响嵌入的 \Index{PostScript} 文件的当前颜色。
% 考虑以下文件 |foo.eps|，它绘制一个带边框的灰色框：
% \begin{verbatim}
% %!PS-Adobe-3.0 EPSF-3.0
% %%BoundingBox: 0 0 60 12
% 0 0 60 12 rectfill
% 0.75 setgray
% 2 2 56 8 rectfill
% \end{verbatim}
% 现在通过 \LaTeX{} 和 \Program{dvips} 运行以下代码：
% \begin{verbatim}
% \documentclass{minimal}
% \usepackage[fixinclude]{xcolor}
% \usepackage{graphics}
% \begin{document}
% \includegraphics{foo} \textcolor{red}{\includegraphics{foo}}
% \end{document}
% \end{verbatim}
% 生成的 \File{.ps} 文件会显示两个灰色框：第一个带黑色边框，第二个带红色边框。
% 如果省略 \Option{fixinclude} 选项，第二个框也会显示黑色边框。
% 这是因为 \Program{dvips} 通常在嵌入 \File{.eps} 文件之前立即将当前颜色重置为黑色。
%
%
%
%
% \clearpage
% \begin{adjustwidth}{-2.5cm}{}
% \section{示例}
%
%
% \begin{figure}[htbp]\caption{光谱}\label{fig.spectrum}
% \small
% \newcount\WL \unitlength.75pt
% \begin{picture}(460,60)(355,-10)
% \sffamily \tiny \linethickness{1.25\unitlength} \WL=360
% \multiput(360,0)(1,0){456}%
%  {{\color[wave]{\the\WL}\line(0,1){50}}\global\advance\WL1}
% \linethickness{0.25\unitlength}\WL=360
% \multiput(360,0)(20,0){23}%
%  {\picture(0,0)
%   \line(0,-1){5} \multiput(5,0)(5,0){3}{\line(0,-1){2.5}}
%   \put(0,-10){\makebox(0,0){\the\WL}}\global\advance\WL20
%   \endpicture}
% \end{picture}
% \vspace*{3ex}
% \begin{verbatim}
% \newcount\WL \unitlength.75pt
% \begin{picture}(460,60)(355,-10)
% \sffamily \tiny \linethickness{1.25\unitlength} \WL=360
% \multiput(360,0)(1,0){456}%
%  {{\color[wave]{\the\WL}\line(0,1){50}}\global\advance\WL1}
% \linethickness{0.25\unitlength}\WL=360
% \multiput(360,0)(20,0){23}%
%  {\picture(0,0)
%   \line(0,-1){5} \multiput(5,0)(5,0){3}{\line(0,-1){2.5}}
%   \put(0,-10){\makebox(0,0){\the\WL}}\global\advance\WL20
%   \endpicture}
% \end{picture}
% \end{verbatim}
% \end{figure}
%
%
% \begin{figure}[htbp]\caption{颜色测试}\label{fig.testcolors}
% \small
% \begin{adjustwidth}{-2.5cm}{}
% \raggedleft
% \sffamily
% \begin{testcolors}[rgb,cmyk,hsb,HTML,gray]
% \testcolor{olive}
% \testcolor{red!50!green}
% \testcolor{-cyan!50!magenta}
% \testcolor[cmyk]{0,0,1,0.5}
% \testcolor[cmyk]{0,0,.5,.5}
% \testcolor[rgb:cmyk]{0,0,.5,.5}
% \end{testcolors}
% \end{adjustwidth}
% \vspace*{3ex}
% \begin{adjustwidth}{2.5cm}{}
% \begin{verbatim}
% \sffamily
% \begin{testcolors}[rgb,cmyk,hsb,HTML,gray]
% \testcolor{olive}
% \testcolor{red!50!green}
% \testcolor{-cyan!50!magenta}
% \testcolor[cmyk]{0,0,1,0.5}
% \testcolor[cmyk]{0,0,.5,.5}
% \testcolor[rgb:cmyk]{0,0,.5,.5}
% \end{testcolors}
% \end{verbatim}
% \end{adjustwidth}
% \end{figure}
%
%
% \def\ukColor{MyGreen}
% \definecolor{\ukColor}{cmyk}{0.92,0,0.87,0.09}
% \colorlet{\ukColor-rgb}[rgb]{\ukColor}
% \colorlet{\ukColor-cmy}[cmy]{\ukColor}
% \colorlet{\ukColor-hsb}[hsb]{\ukColor}
% \colorlet{\ukColor-gray}[gray]{\ukColor}
%
% \begin{figure}[htbp]\caption{从一种颜色渐变到另一种颜色}\label{fig.progress}
% \small
% \begin{adjustwidth}{-2.5cm}{}
% \makeatletter
% \def\ColorLoop#1#2{\@tfor\@@tmp:={}{-rgb}{-cmy}{-hsb}{-gray}\do
%   {\ifcase#1
%    \expandafter\Colorloop\expandafter{\expandafter\ukColor\@@tmp}{#2}\or
%    \expandafter\Colorshow\expandafter{\expandafter\ukColor\@@tmp}\fi}\\[-2ex]\ignorespaces}
% \def\Colorloop#1#2{\@ifxempty{#1}{}{\texttt{#1!}}&
%   \count@105 \@whilenum\count@>\z@\do
%    {\advance\count@-5
%     \@ifxempty{#1}{\hb@xt@\XCdim{\hss\textsf{\the\count@}\hss}}{\testclr{#1!\the\count@!#2}}\hskip-\fboxrule}\hskip\fboxrule&
%   \@ifxempty{#1}{}{\texttt{!#2}}\\}
% \def\Colorshow#1{\texttt{#1}&
%   \multicolumn2{l@{}}{\expandafter\let\expandafter\@@tmp\csname\string\color@#1\endcsname
%   \edef\@@tmp{\expandafter\@gobbletwo\@@tmp}\Meaning{@@tmp}}\\}
% \makeatother
% \sbox0{\testclr{red}}\XCdim\wd0
% \begin{tabular}{@{}lll@{}}
% \Colorloop{}{}\\[-2ex]
% \ColorLoop0{white}
% \ColorLoop0{gray}
% \ColorLoop0{black}
% \ColorLoop0{red}
% \ColorLoop0{blue}
% \ColorLoop0{yellow}
% \\\noalign{\vspace*{-2.5ex}}\hline\noalign{\vspace*{1.5ex}}
% \itshape 颜色&\itshape 定义/表示（\textup{\Option{\csname XC@@drv\endcsname}} 驱动）\\
% \ColorLoop1{}
% \end{tabular}
% \end{adjustwidth}
% \end{figure}
%
%
% \begin{figure}[htbp]\caption{目标颜色模型}\label{fig.tgt}
% \small
% \begin{adjustwidth}{-2.5cm}{}
% \raggedleft
% \extrarowheight2pt
% \makeatletter
% \def\X#1%
%  {\def\tmp{{#1}}\dots\Meaning{tmp}&
%   \selectcolormodel{#1}\testclr{red} \testclr{green} \testclr{blue} \testclr{cyan} \testclr{magenta} \testclr{yellow} \testclr{brown} \testclr{lime} \testclr{olive} \testclr{orange} \testclr{pink} \testclr{purple} \testclr{teal} \testclr{violet} \testclr{black} \testclr{darkgray} \testclr{gray} \testclr{lightgray} \testclr{white}\\}
% \makeatother
% \begin{tabular}{@{}>{\ttfamily}l@{ }l@{}}
% \string\selectcolormodel&\\
% \X{natural}
% \X{rgb}
% \X{cmy}
% \X{cmyk}
% \X{hsb}
% \X{gray}
% \end{tabular}
% \end{adjustwidth}
% \end{figure}
%
%
% \begin{figure}[htbp]\caption{标准颜色表达式}\label{fig.expr}
% \centering\small
% \def\test#1{\testclr{#1}& \texttt{#1}}
% \def\X#1{\test{#1}& \test{-#1}}
% \begin{tabular}{@{}l@{ \ }ll@{ \ }l@{}}
% \X{red}\\
% \X{red!75}\\
% \X{red!75!green}\\
% \X{red!75!green!50}\\
% \X{red!75!green!50!blue}\\
% \X{red!75!green!50!blue!25}\\
% \X{red!75!green!50!blue!25!gray}\\
% \end{tabular}
% \end{figure}
%
%
% \begin{figure}[htbp]\caption{标准颜色表达式}\label{fig.ex}
% \begin{minipage}{\textwidth}\small
% \begin{verbatim}
% \fboxrule6pt
% \fcolorbox
%  {red!70!green}% outer frame
%  {yellow!30!blue}% outer background
%  {\fcolorbox
%    {-yellow!30!blue}% inner frame
%    {-red!70!green}% inner background
%    {Test\textcolor{red!72.75}{Test}\color{-green}Test}}
% \end{verbatim}
% \end{minipage}\relax
% \llap{\fboxrule6pt \fcolorbox
%   {red!70!green}{yellow!30!blue}{\fcolorbox
%     {-yellow!30!blue}{-red!70!green}{Test\textcolor{red!72.75}{Test}\color{-green}Test}}}
% \end{figure}
%
%
% \begin{figure}[htbp]\caption{当前颜色}\label{fig.cc}
% \begin{minipage}{\textwidth}\small
% \begin{verbatim}
% \def\test{current, \textcolor{.!50}{50\%},
%           \textcolor{-.}{complement},
%           \textcolor{yellow!50!.}{mix}}
% \textcolor{blue}{\test}\\
%  and \textcolor{red}{\test}\\
% \def\Test{\color{.!80}Test}
% \textcolor{blue}{\Test\Test\Test\Test\Test}\\
% and \textcolor{red}{\Test\Test\Test\Test\Test}
% \end{verbatim}
% \end{minipage}\relax
% \llap{\begin{minipage}{.5\textwidth}\small\raggedleft
% \def\test{current, \textcolor{.!50}{50\%},
%           \textcolor{-.}{complement},
%           \textcolor{yellow!50!.}{mix}}
% \textcolor{blue}{\test}\\
%  and \textcolor{red}{\test}\\
% \def\Test{\color{.!80}Test}
% \textcolor{blue}{\Test\Test\Test\Test\Test}\\
% and \textcolor{red}{\Test\Test\Test\Test\Test}
% \end{minipage}}
% \end{figure}
%
%
% \begin{figure}[htbp]\caption{颜色系列}\label{fig.csx}
% \centering
% \def\tabulartest#1{\definecolorseries{test}#1\resetcolorseries[12]{test}%
%  \rowcolors[\hline]{1}{test!!+}{test!!+}%
%  \begin{tabular}{c}
%  \number\rownum\\ \number\rownum\\ \number\rownum\\ \number\rownum\\
%  \number\rownum\\ \number\rownum\\ \number\rownum\\ \number\rownum\\
%  \number\rownum\\ \number\rownum\\ \number\rownum\\ \number\rownum\\
%  \number\rownum\\ \number\rownum\\ \number\rownum\\ \number\rownum\\
%  \end{tabular}}
%
% \begin{tabular}{@{}*{9}{c}@{}}
% $S_1$ & $S_2$ & $G_1$ & $G_2$ & $L_1$ & $L_2$ & $L_3$ & $L_4$ & $L_5$ \\[3pt]
% \tabulartest{{rgb}{step}[rgb]{.95,.85,.55}{.17,.47,.37}}&
% \tabulartest{{hsb}{step}[hsb]{.575,1,1}{.11,-.05,0}}&
% \tabulartest{{rgb}{grad}[rgb]{.95,.85,.55}{3,11,17}}&
% \tabulartest{{hsb}{grad}[hsb]{.575,1,1}{.987,-.234,0}}&
% \tabulartest{{rgb}{last}[rgb]{.95,.85,.55}[rgb]{.05,.15,.55}}&
% \tabulartest{{hsb}{last}[hsb]{.575,1,1}[hsb]{-.425,.15,1}}&
% \tabulartest{{rgb}{last}{yellow!50}{blue}}&
% \tabulartest{{hsb}{last}{yellow!50}{blue}}&
% \tabulartest{{cmy}{last}{yellow!50}{blue}}\\
% \end{tabular}
%
% \bigskip
%
% \footnotesize
% \rowcolors{3}{tablerowcolor}{}
% \begin{tabular}{>{$}l<{$}l}
% \hline
% \multicolumn2{>{\columncolor{tableheadcolor}}c}{\em 单独定义}\\
% S_1& |\definecolorseries{test}{rgb}{step}[rgb]{.95,.85,.55}{.17,.47,.37}|\\
% S_2& |\definecolorseries{test}{hsb}{step}[hsb]{.575,1,1}{.11,-.05,0}|\\
% G_1& |\definecolorseries{test}{rgb}{grad}[rgb]{.95,.85,.55}{3,11,17}|\\
% G_2& |\definecolorseries{test}{hsb}{grad}[hsb]{.575,1,1}{.987,-.234,0}|\\
% L_1& |\definecolorseries{test}{rgb}{last}[rgb]{.95,.85,.55}[rgb]{.05,.15,.55}|\\
% L_2& |\definecolorseries{test}{hsb}{last}[hsb]{.575,1,1}[hsb]{-.425,.15,1}|\\
% L_3& |\definecolorseries{test}{rgb}{last}{yellow!50}{blue}|\\
% L_4& |\definecolorseries{test}{hsb}{last}{yellow!50}{blue}|\\
% L_5& |\definecolorseries{test}{cmy}{last}{yellow!50}{blue}|\\
% \hline
% \multicolumn2{>{\columncolor{tableheadcolor}}c}{\em 公共定义}\\
% \hiderowcolors
% &|\resetcolorseries[12]{test}|\\
% &|\rowcolors[\hline]{1}{test!!+}{test!!+}|\\
% &|\begin{tabular}{c}|\\
% &|\number\rownum\\ \number\rownum\\ \number\rownum\\ \number\rownum\\|\\
% &|\number\rownum\\ \number\rownum\\ \number\rownum\\ \number\rownum\\|\\
% &|\number\rownum\\ \number\rownum\\ \number\rownum\\ \number\rownum\\|\\
% &|\number\rownum\\ \number\rownum\\ \number\rownum\\ \number\rownum\\|\\
% &|\end{tabular}|\\
% \hline
% \end{tabular}
% \end{figure}
%
%
% \begin{figure}[htbp]\caption{颜色蒙版}\label{fig.mask}
% \begin{adjustwidth}{-2.5cm}{}
% \small\raggedleft
% \extrarowheight2pt
% \makeatletter
% \def\X#1#2{\edef\tmp{\@ifxempty{#1}{}{[#1]}{#2}}\dots\Meaning{tmp}&
%   \maskcolors[#1]{#2}\testclr{red} \testclr{green} \testclr{blue} \testclr{cyan} \testclr{magenta} \testclr{yellow} \testclr{brown} \testclr{lime} \testclr{olive} \testclr{orange} \testclr{pink} \testclr{purple} \testclr{teal} \testclr{violet} \testclr{black} \testclr{darkgray} \testclr{gray} \testclr{lightgray} \testclr{white}\\}\relax
% \makeatother
% \begin{tabular}{@{}>{\ttfamily}l@{ }l@{}}
% \string\maskcolors&\\
% \hline
% \X{}{}
% \hline
% \X{cmyk}{cyan}
% \X{cmyk}{magenta}
% \X{cmyk}{yellow}
% \X{cmyk}{black}
% \hline
% \X{cmyk}{red}
% \X{cmyk}{green}
% \X{cmyk}{blue}
% \hline
% \X{rgb}{red}
% \X{rgb}{green}
% \X{rgb}{blue}
% \hline
% \X{hsb}{red}
% \X{hsb}{green}
% \X{hsb}{blue}
% \hline
% \X{rgb}{gray}
% \X{cmy}{gray}
% \end{tabular}
% \end{adjustwidth}
% \end{figure}
%
%
% \begin{figure}[htbp]\caption{表格中交替行颜色：\cmd\rowcolors{} 与 \cmd\rowcolors\texttt* 对比}\label{fig.row}
% \centering
% \begin{minipage}{\textwidth}
% \begin{verbatim}
% \rowcolors[\hline]{3}{green!25}{yellow!50} \arrayrulecolor{red!75!gray}
% \begin{tabular}{ll}
% test & row \number\rownum\\
% test & row \number\rownum\\
% test & row \number\rownum\\
% test & row \number\rownum\\
% \arrayrulecolor{black}
% test & row \number\rownum\\
% test & row \number\rownum\\
% \rowcolor{blue!25}
% test & row \number\rownum\\
% test & row \number\rownum\\
% \hiderowcolors
% test & row \number\rownum\\
% test & row \number\rownum\\
% \showrowcolors
% test & row \number\rownum\\
% test & row \number\rownum\\
% \multicolumn{1}%
%  {>{\columncolor{red!12}}l}{test} & row \number\rownum\\
% \end{tabular}
% \end{verbatim}
% \end{minipage}
% \hskip-.5\textwidth
% \rowcolors[\hline]{3}{green!25}{yellow!50} \arrayrulecolor{red!75!gray}
% \begin{tabular}{ll}
% test & row \number\rownum\\
% test & row \number\rownum\\
% test & row \number\rownum\\
% test & row \number\rownum\\
% \arrayrulecolor{black}
% test & row \number\rownum\\
% test & row \number\rownum\\
% \rowcolor{blue!25}
% test & row \number\rownum\\
% test & row \number\rownum\\
% \hiderowcolors
% test & row \number\rownum\\
% test & row \number\rownum\\
% \showrowcolors
% test & row \number\rownum\\
% test & row \number\rownum\\
% \multicolumn{1}%
%  {>{\columncolor{red!12}}l}{test} & row \number\rownum\\
% \end{tabular}
% \qquad
% \rowcolors*[\hline]{3}{green!25}{yellow!50}\arrayrulecolor{red!75!gray}
% \begin{tabular}{ll}
% test & row \number\rownum\\
% test & row \number\rownum\\
% test & row \number\rownum\\
% test & row \number\rownum\\
% \arrayrulecolor{black}
% test & row \number\rownum\\
% test & row \number\rownum\\
% \rowcolor{blue!25}
% test & row \number\rownum\\
% test & row \number\rownum\\
% \hiderowcolors
% test & row \number\rownum\\
% test & row \number\rownum\\
% \showrowcolors
% test & row \number\rownum\\
% test & row \number\rownum\\
% \multicolumn{1}%
%  {>{\columncolor{red!12}}l}{test} & row \number\rownum\\
% \end{tabular}
% \arrayrulecolor{black}
% \end{figure}
%
%
% \begin{figure}[htbp]\caption{\Model{Hsb} 与 \Model{tHsb}：\emph{hue\Circ} 按 $15^\circ$ 步进}\label{fig.tHsb}
% \begin{adjustwidth}{-2.5cm}{}
% \footnotesize\sffamily\raggedleft
% \extrarowheight.5pt
% \begin{testcolors}[rgb,cmyk,hsb,Hsb,tHsb]
% \testcolor[Hsb]{0,1,1}
% \testcolor[Hsb]{15,1,1}
% \testcolor[Hsb]{30,1,1}
% \testcolor[Hsb]{45,1,1}
% \testcolor[Hsb]{60,1,1}
% \testcolor[Hsb]{75,1,1}
% \testcolor[Hsb]{90,1,1}
% \testcolor[Hsb]{105,1,1}
% \testcolor[Hsb]{120,1,1}
% \testcolor[Hsb]{135,1,1}
% \testcolor[Hsb]{150,1,1}
% \testcolor[Hsb]{165,1,1}
% \testcolor[Hsb]{180,1,1}
% \testcolor[Hsb]{195,1,1}
% \testcolor[Hsb]{210,1,1}
% \testcolor[Hsb]{225,1,1}
% \testcolor[Hsb]{240,1,1}
% \testcolor[Hsb]{255,1,1}
% \testcolor[Hsb]{270,1,1}
% \testcolor[Hsb]{285,1,1}
% \testcolor[Hsb]{300,1,1}
% \testcolor[Hsb]{315,1,1}
% \testcolor[Hsb]{330,1,1}
% \testcolor[Hsb]{345,1,1}
% \testcolor[Hsb]{360,1,1}
% \noalign{\medskip}
% \testcolor[tHsb]{0,1,1}
% \testcolor[tHsb]{15,1,1}
% \testcolor[tHsb]{30,1,1}
% \testcolor[tHsb]{45,1,1}
% \testcolor[tHsb]{60,1,1}
% \testcolor[tHsb]{75,1,1}
% \testcolor[tHsb]{90,1,1}
% \testcolor[tHsb]{105,1,1}
% \testcolor[tHsb]{120,1,1}
% \testcolor[tHsb]{135,1,1}
% \testcolor[tHsb]{150,1,1}
% \testcolor[tHsb]{165,1,1}
% \testcolor[tHsb]{180,1,1}
% \testcolor[tHsb]{195,1,1}
% \testcolor[tHsb]{210,1,1}
% \testcolor[tHsb]{225,1,1}
% \testcolor[tHsb]{240,1,1}
% \testcolor[tHsb]{255,1,1}
% \testcolor[tHsb]{270,1,1}
% \testcolor[tHsb]{285,1,1}
% \testcolor[tHsb]{300,1,1}
% \testcolor[tHsb]{315,1,1}
% \testcolor[tHsb]{330,1,1}
% \testcolor[tHsb]{345,1,1}
% \testcolor[tHsb]{360,1,1}
% \end{testcolors}
% \end{adjustwidth}
% \end{figure}
%
%
% \begin{figure}[htbp]\caption{色彩和谐}\label{fig.harmony}
% \begin{adjustwidth}{-2.5cm}{}
% \small\sffamily\raggedleft
% \extrarowheight1.5pt
% \begin{testcolors}[rgb,cmyk,Hsb,tHsb]
% \noalign{\smallskip}
% \rlap{\em 补色（双色和谐）：}\\
% \testcolor{yellow>wheel,1,2}
% \testcolor{yellow}
% \testcolor{yellow>twheel,1,2}
% \noalign{\medskip}
% \rlap{\em 三色组（三色和谐）：}\\
% \testcolor{yellow>wheel,2,3}
% \testcolor{yellow>wheel,1,3}
% \testcolor{yellow}
% \testcolor{yellow>twheel,1,3}
% \testcolor{yellow>twheel,2,3}
% \noalign{\medskip}
% \rlap{\em 四色组（四色和谐）：}\\
% \testcolor{yellow>wheel,3,4}
% \testcolor{yellow>wheel,2,4}
% \testcolor{yellow>wheel,1,4}
% \testcolor{yellow}
% \testcolor{yellow>twheel,1,4}
% \testcolor{yellow>twheel,2,4}
% \testcolor{yellow>twheel,3,4}
% \noalign{\medskip}
% \rlap{\em 分裂补色：}\\
% \testcolor{yellow>wheel,7,12}
% \testcolor{yellow>wheel,5,12}
% \testcolor{yellow}
% \testcolor{yellow>twheel,5,12}
% \testcolor{yellow>twheel,7,12}
% \noalign{\medskip}
% \rlap{\em 类似色（邻接色）：}\\
% \testcolor{yellow>wheel,11,12}
% \testcolor{yellow>wheel,10,12}
% \testcolor{yellow>wheel,2,12}
% \testcolor{yellow>wheel,1,12}
% \testcolor{yellow}
% \testcolor{yellow>twheel,1,12}
% \testcolor{yellow>twheel,2,12}
% \testcolor{yellow>twheel,10,12}
% \testcolor{yellow>twheel,11,12}
% \end{testcolors}
% \end{adjustwidth}
% \end{figure}
%
% \end{adjustwidth}
%
%
% \clearpage
%
% \begin{adjustwidth}{-2.5cm}{}
% \raggedcolumns\columnsep2.5em
%
% \section{按名称列出的颜色}\label{sec.names}
%
% \begin{multicols}{5}[\subsection{基色（始终可用）}\label{sec.basenam}]
% \footnotesize\def\0#1{\Testclr{#1}{#1}}
% \0{black}
% \0{blue}
% \0{brown}
% \0{cyan}
% \0{darkgray}
% \0{gray}
% \0{green}
% \0{lightgray}
% \0{lime}
% \0{magenta}
% \0{olive}
% \0{orange}
% \0{pink}
% \0{purple}
% \0{red}
% \0{teal}
% \0{violet}
% \0{white}
% \0{yellow}
% \end{multicols}
%
%
% \begin{multicols}{5}[\subsection{通过 \Option{dvipsnames} 选项加载的颜色}\label{sec.dvipsnam}]
% \footnotesize\def\0#1{\Testclr{#1}{#1}}
% \0{Apricot}
% \0{Aquamarine}
% \0{Bittersweet}
% \0{Black}
% \0{Blue}
% \0{BlueGreen}
% \0{BlueViolet}
% \0{BrickRed}
% \0{Brown}
% \0{BurntOrange}
% \0{CadetBlue}
% \0{CarnationPink}
% \0{Cerulean}
% \0{CornflowerBlue}
% \0{Cyan}
% \0{Dandelion}
% \0{DarkOrchid}
% \0{Emerald}
% \0{ForestGreen}
% \0{Fuchsia}
% \0{Goldenrod}
% \0{Gray}
% \0{Green}
% \0{GreenYellow}
% \0{JungleGreen}
% \0{Lavender}
% \0{LimeGreen}
% \0{Magenta}
% \0{Mahogany}
% \0{Maroon}
% \0{Melon}
% \0{MidnightBlue}
% \0{Mulberry}
% \0{NavyBlue}
% \0{OliveGreen}
% \0{Orange}
% \0{OrangeRed}
% \0{Orchid}
% \0{Peach}
% \0{Periwinkle}
% \0{PineGreen}
% \0{Plum}
% \0{ProcessBlue}
% \0{Purple}
% \0{RawSienna}
% \0{Red}
% \0{RedOrange}
% \0{RedViolet}
% \0{Rhodamine}
% \0{RoyalBlue}
% \0{RoyalPurple}
% \0{RubineRed}
% \0{Salmon}
% \0{SeaGreen}
% \0{Sepia}
% \0{SkyBlue}
% \0{SpringGreen}
% \0{Tan}
% \0{TealBlue}
% \0{Thistle}
% \0{Turquoise}
% \0{Violet}
% \0{VioletRed}
% \0{White}
% \0{WildStrawberry}
% \0{Yellow}
% \0{YellowGreen}
% \0{YellowOrange}
% \end{multicols}
%
%
% \makeatletter
% \def\ColorList#1#2#3#4{\begin{multicols}{#1}[\subsection{通过 \Option{#2names} 选项加载的颜色}\label{sec.#2nam}]
%   \footnotesize
%   \def\preparecolorset
%    {\def\XC@@cmd[####1]####2####3####4{\definecolor[####1]{@test}{####3}{####4}#4{@test}{####2}}\@testopt{\XC@definecolorset}{}}
%   \edef\restore{\catcode`\noexpand\%\the\catcode`\%\relax}
%   \catcode`\%14\relax\input{#3nam.def}\restore
%   \end{multicols}}
% \makeatother
%
% \ColorList4{svg}{svg}\Testclr
% \medskip
% \begin{adjustwidth}{1cm}{1cm}
% \footnotesize\raggedright 重复颜色：
% \Color[-]{Aqua} = \Color[-]{Cyan},
% \Color[-]{Fuchsia} = \Color[-]{Magenta};
% \Color[-]{Navy} = \Color[-]{NavyBlue};
% \Color[-]{Gray} = \Color[-]{Grey},
% \Color[-]{DarkGray} = \Color[-]{DarkGrey},
% \Color[-]{LightGray} = \Color[-]{LightGrey},
% \Color[-]{SlateGray} = \Color[-]{SlateGrey},
% \Color[-]{DarkSlateGray} = \Color[-]{DarkSlateGrey},
% \Color[-]{LightSlateGray} = \Color[-]{LightSlateGrey},
% \Color[-]{DimGray} = \Color[-]{DimGrey}.
% \par\smallskip
% \Index{HTML4} 颜色关键字子集：
% \Color[-]{Aqua}, \Color[-]{Black}, \Color[-]{Blue}, \Color[-]{Fuchsia},
% \Color[-]{Gray}, \Color[-]{Green}, \Color[-]{Lime}, \Color[-]{Maroon},
% \Color[-]{Navy}, \Color[-]{Olive}, \Color[-]{Purple}, \Color[-]{Red},
% \Color[-]{Silver}, \Color[-]{Teal}, \Color[-]{White}, \Color[-]{Yellow}.
% \par\smallskip
% 取自 \Index{Unix}/\Index{X11} 的颜色：
% \Color[-]{LightGoldenrod},
% \Color[-]{LightSlateBlue},
% \Color[-]{NavyBlue},
% \Color[-]{VioletRed}.
% \end{adjustwidth}
%
% \ColorList4{x11}{x11}\Testclr
% \medskip
% \begin{adjustwidth}{1cm}{1cm}
% \footnotesize\raggedright 重复颜色：
% \Color[-]{Gray0} = \Color[-]{Grey0},
% \Color[-]{Green0} = \Color[-]{Green1}.
% \end{adjustwidth}
%
% \end{adjustwidth}
%
%
%
%
% \clearpage
% \section{技术补充}
%
% \subsection{驱动支持的颜色模型}
%
% 由于某些驱动只是假装支持 \Model{hsb} 模型，我们加入了一些代码来绕过这种行为。
% 实际由 \XC{} 添加的模型会显示在日志文件中。
% 表 \vref{tab.drv} 列出了当前 MiK\TeX{} \cite{mik} 发行版中包含的主要驱动及其颜色模型支持情况。
% 其他发行版的行为可能类似。
%
% \begin{table}[htbp]\caption{驱动与颜色模型}\label{tab.drv}
% \centering\small
% \extrarowheight1pt
% \tabcolsep3pt
% \rowcolors{3}{tablerowcolor}{}
% \begin{tabular}{ll*{9}{c}}
% \rowcolor{tableheadcolor}
% \hline
% \em 驱动&\em 版本&\Model{rgb}&\Model{cmy}&\Model{cmyk}&\Model{hsb}&\Model{gray}&\Model{RGB}&\Model{HTML}&\Model{HSB}&\Model{Gray}\\
% \Option{dvipdf}&2015/12/30 v3.0k&d&n&d&n&d&i&n&n&n\\
% \Option{dvips}&2015/12/30 v3.0k&d&n&d&d&d&i&n&n&n\\
% \Option{dvipsone}&2015/12/30 v3.0k&d&n&d&d&d&i&n&n&n\\
% \Option{pctex32}&2015/12/30 v3.0k&d&n&d&d&d&i&n&n&n\\
% \Option{pctexps}&2015/12/30 v3.0k&d&n&d&d&d&i&n&n&n\\
% \hline
% \Option{pdftex}&2011/05/27 v0.06d&d&n&d&n&d&i&n&n&n\\
% \Option{luatex}&2016/01/23 v0.01b&d&n&d&n&d&i&n&n&n\\
% \hline
% \Option{dvipdfm}&1999/9/6 vx.x&d&n&d&n&d&i&n&n&n\\
% \Option{dvipdfmx}&2016/04/06 v4.08&d&n&d&?&d&i&n&n&n\\
% \hline
% \Option{textures}&1997/5/28 v0.3&d&n&d&?&i&n&n&n&n\\
% \hline
% \Option{vtex}&1999/01/14 v6.3&d&n&d&n&i&i&n&n&n\\
% \hline
% \Option{xetex}&2016/04/06 v4.08&d&n&d&d&d&i&n&n&n\\
% \hline
% \Option{tcidvi}&2015/12/30 v3.0k&i&n&i&n&i&d&n&n&n\\
% \Option{truetex}&2015/12/30 v3.0k&i&n&i&n&i&d&n&n&n\\
% \hline
% \Option{dviwin}&2015/12/30 v3.0k&n&n&n&n&n&n&n&n&n\\
% \Option{emtex}&2015/12/30 v3.0k&n&n&n&n&n&n&n&n&n\\
% \Option{pctexhp}&2015/12/30 v3.0k&n&n&n&n&n&n&n&n&n\\
% \Option{pctexwin}&2015/12/30 v3.0k&n&n&n&n&n&n&n&n&n\\
% \hline\hline
% \multicolumn{11}c{\Option{dviwindo} = \Option{dvipsone}; \Option{oztex} = \Option{dvips}; \Option{xdvi} = \Option{dvips} + \Option{monochrome}}\\
% \hline
% \multicolumn{11}c{驱动颜色模型支持：d\,=\,直接支持，i\,=\,间接支持，n\,=\,不支持}\\
% \hline
% \end{tabular}
% \end{table}
%
%
%
% \subsection{\XC{} 如何处理驱动特定的颜色模型}
%
% 虽然各种驱动实现颜色可视化的方式各不相同，但它们都具有原 \LCP 定义的一些共同特征。
% 其中之一是：任何颜色模型 `foo' 都需要一个 |\color@foo|\marg{cmd}\marg{spec} 命令，将与 `foo' 相关的颜色 \Meta{spec} 翻译为存储在 \Meta{cmd} 中的驱动特定代码。
% 因此，\XC{} 通常通过 |\color@foo| 是否存在来检测驱动对 `foo' 模型的支持。
%
% 通过这种机制，\XC{} 还可以在不修改驱动文件本身的情况下改变某些模型的行为。
% 一个很好的例子是 |\substitutecolormodel| 命令，它在宏包初始化过程中用于提供对实际驱动未覆盖模型（如 \Option{pdftex} 的 \Model{hsb}）或实现不正确的模型（如 \Option{dvipdfm} 的 \Model{hsb}）的支持。
%
%
% \subsection{幕后：内部颜色表示}
%
% 每个按名称访问的颜色定义都需要颜色的内部表示，即一个包含驱动正确显示颜色所需信息的宏。
%
% \LC 的 |\definecolor{foo}{...}{...}| 生成一个命令 |\\color@foo|\footnote{双反斜杠是有意的。}，其中以驱动相关的方式包含颜色定义；因此之后访问颜色模型和参数虽然可行但并不简单（解决方案见 \Package{colorinfo} 宏包 \cite{colorinfo}）。
%
% \LC 的 |\DefineNamedColor{named}{foo}{...}{...}| 生成 |\col@foo|\footnote{单反斜杠是有意的。}，其中同样包含一些驱动相关信息。
% 在这种情况下，只有当宏包选项 \Option{usecolors} 激活时，才会额外定义 |\\color@foo|。
%
% \XC 的 |\definecolor{foo}{...}{...}| 也生成\footnote{这是 1.10 版引入的；在此之前生成的是语法不同的命令 \cs{\string\xcolor@foo}。} 一个命令 |\\color@foo|，它合并了前两个命令的特性，同时包含驱动相关和驱动无关的信息，从而可以标准化地访问相关参数。
% 虽然语法不同了，但 |\\color@foo| 展开后与原命令表达式相同。
% 另一方面，|\col@foo| 命令不再需要，因此在 `named' 情况下不再生成：\XC{} 使用单一的颜色数据结构（如上所述）。
%
% 表 \vref{tab.internal} 展示了两种最主要驱动的示例。
% 另见图 \vref{fig.progress}，它显示了处理本文档所用驱动的相关定义。
%
%
% \begin{table}[htb]\caption{驱动相关的内部颜色表示}\label{tab.internal}
% \centering\small
% \def\X{\hspace*{-.845\textwidth}\\}
% \extrarowheight1pt
% \tabcolsep8pt
% \begin{tabular}{l>{\color{blue}(}r<{)}l}
% \hline
% \multicolumn3{>{\columncolor{tableheadcolor}}l}{\Option{dvips} driver}\\
% \hline
% |\\color@Plum=macro:| & |\definecolor{Plum}{rgb}{.5,0,1}| & \LC\\
% |->rgb .5 0 1.|\X
% \hline
% |\\color@Plum=macro:| & |\definecolor{Plum}{rgb}{.5,0,1}|& \XC\\
% |->\xcolor@ {}{rgb 0.5 0 1}{rgb}{0.5,0,1}.|\X
% \hline
% |\col@Plum=macro:| & |\DefineNamedColor{Plum}{rgb}{.5,0,1}|& \LC\\
% |->\@nil .|\X
% |\\color@Plum=macro:| & with option \oOption{usenames}\\
% |-> Plum.|\X
% \hline
% |\\color@Plum=macro:| & |\definecolor[named]{Plum}{rgb}{.5,0,1}|& \XC\\
% |->\xcolor@ {named}{ Plum}{rgb}{0.5,0,1}.|\X
% \hline
% \multicolumn3{>{\columncolor{tableheadcolor}}l}{\Option{pdftex} driver}\\
% \hline
% |\\color@Plum=macro:| & |\definecolor{Plum}{rgb}{.5,0,1}|& \LC\\
% |->.5 0 1 rg .5 0 1 RG.|\X
% \hline
% |\\color@Plum=macro:| & |\definecolor{Plum}{rgb}{.5,0,1}|& \XC\\
% |->\xcolor@ {}{0.5 0 1 rg 0.5 0 1 RG}{rgb}{0.5,0,1}.|\X
% \hline
% |\col@Plum=macro:| & |\DefineNamedColor{Plum}{rgb}{.5,0,1}|& \LC\\
% |->.5 0 1 rg .5 0 1 RG.|\X
% |\\color@Plum=macro:| & with option \oOption{usenames}\\
% |->.5 0 1 rg .5 0 1 RG.|\X
% \hline
% |\\color@Plum=macro:| & |\definecolor[named]{Plum}{rgb}{.5,0,1}|& \XC\\
% |->\xcolor@ {}{0.5 0 1 rg 0.5 0 1 RG}{rgb}{0.5,0,1}.|\X
% \hline
% \end{tabular}
% \end{table}
%
%
%
% \subsection{关于精度的说明}
%
% 由于此处展示的宏需要进行一些计算，我们特别努力确保在 \TeX 有限的数值能力范围内，转换和混合公式具有最高精度。\footnote{例如，对区间 00000\dots 99999 内的所有五位数 \Meta{int} 应用"变换" \cs{dimen0=0.}\Meta{int}\texttt{pt} \cs{the\string\dimen0}，100000 个数中恰好有 34464 个不能保持不变。这里我们还没有讨论被吞掉的末尾零……}
% 我们决定开发并包含一小套命令来提高除法和乘法结果的质量，而不是加载 \Package{realcalc} 或 \Package{fp} 这类提供多位运算和更多功能的宏包。
% 后者对我们的用途来说边际贡献似乎不足以证明其使用价值。
% 因此，我们保持在一种定点运算框架内，通过 \TeX 的维度寄存器最多提供 5 位小数。
%
%
%
%
% \clearpage
% \section{公式}
%
% \subsection{颜色混合}
%
% 一般来说，我们使用线性插值进行颜色混合：
% \begin{align}
% \mix(C, C', p) = p \cdot C + (1-p) \cdot C'\label{eq.mix}
% \end{align}
% 注意，在 \Model{hsb} 情况下有一个特殊情形：
% 如果 $\emph{saturation}=0$，则该颜色等于一个亮度为 \emph{brightness} 的灰色，与 \emph{hue} 值无关。
% 因此，为了实现任意颜色到特定灰色（如白色或黑色）的平滑过渡，我们实际使用以下公式：
% \begin{align}
% \Op\tint{hsb}(C, p) &= p \cdot C + (1-p) \cdot \bigl(\emph{hue},0,1\bigr)\\
% \Op\shade{hsb}(C, p) &= p \cdot C + (1-p) \cdot \bigl(\emph{hue},0,0\bigr)\\
% \Op\tone{hsb}(C, p) &\textstyle= p \cdot C + (1-p) \cdot \bigl(\emph{hue},0,\frac12\bigr)
% \end{align}
% 其中 $C=(\emph{hue},\emph{saturation},\emph{brightness})$。
%
% 由公式 (\ref{eq.mix}) 以及第 \vref{sec.term} 节所述颜色表达式的解释方式，通过归纳法容易验证：颜色表达式
% \begin{gather}
% C_0!P_1!C_1!P_2!\dots!P_n!C_n
% \end{gather}
% 其中 $n\in\{0,1,2,\dots\}$，颜色 $C_0,C_1,\dots,C_n$，百分比 $P_1,\dots,P_n\in [0,100]$，将得到参数向量
% \begin{equation}
% \begin{split}
% C &= \sum_{\nu=0}^n \Biggl(\prod_{\mu=\nu+1}^n p_\mu \Biggr) (1-p_\nu)\cdot C_\nu\\
%   &= p_n\cdots p_1\cdot C_0\\
%   &\quad + p_n\cdots p_2(1-p_1)\cdot C_1\\
%   &\quad + p_n\cdots p_3(1-p_2)\cdot C_2\\
%   &\quad + \dots\\
%   &\quad + p_n(1-p_{n-1})\cdot C_{n-1}\\
%   &\quad + (1-p_n)\cdot C_n
% \end{split}
% \end{equation}
% 其中 $p_0 := 0$，$p_\nu := P_\nu/100$（$\nu=1,\dots,n$）。
% 我们还注意到一个拆分公式：
% \begin{equation}
% \begin{split}
% C_0!P_1!C_1!\dots!P_{n+k}!C_{n+k}
%  &= p_{n+k}\cdots p_{n+1}\cdot C_0!P_1!C_1!\dots!P_n!C_n\\
%  &\quad - p_{n+k}\cdots p_{n+1}\cdot C_n\\
%  &\quad + C_n!P_{n+1}!C_{n+1}!\dots!P_{n+k}!C_{n+k}
% \end{split}
% \end{equation}
%
%
% \begin{table}[p]\caption{颜色常量}\label{tab.const}
% \centering
% \extrarowheight2pt
% \rowcolors{3}{tablerowcolor}{}
% \begin{tabular}{*4{c}}
% \hline
% \rowcolor{tableheadcolor}
% \emph{模型/常量} &
%       \ttfamily white  &  \ttfamily black & \ttfamily gray\\[2pt]
% \Model{rgb}&  $(1,1,1)$   & $(0,0,0)$   & $(\frac12,\frac12,\frac12)$ \\[2pt]
% \Model{cmy}&  $(0,0,0)$   & $(1,1,1)$   & $(\frac12,\frac12,\frac12)$ \\[2pt]
% \Model{cmyk}& $(0,0,0,0)$ & $(0,0,0,1)$ & $(0,0,0,\frac12)$ \\[2pt]
% \Model{hsb}&  $(h,0,1)$   & $(h,0,0)$   & $(h,0,\frac12)$ \\[2pt]
% \Model{Hsb}&  $(h^\circ,0,1)$ & $(h^\circ,0,0)$ & $(h^\circ,0,\frac12)$ \\[2pt]
% \Model{tHsb}& $(h^\circ,0,1)$ & $(h^\circ,0,0)$ & $(h^\circ,0,\frac12)$ \\[2pt]
% \Model{gray}& $1$ & $0$ & $\frac12$ \\[2pt]
% \Model{RGB}&  $(L,L,L)$ & $(0,0,0)$ & $(\lfloor\frac {L+1}2\rfloor,\lfloor\frac {L+1}2\rfloor,\lfloor\frac {L+1}2\rfloor)$ \\[2pt]
% \Model{HTML}& \texttt{FFFFFF}& \texttt{000000}& \texttt{808080}\\[2pt]
% \Model{HSB}&  $(H,0,M)$ & $(H,0,0)$ & $(H,0,\lfloor\frac {M+1}2\rfloor)$ \\[2pt]
% \Model{Gray}& $N$ & $0$ & $\lfloor\frac {N+1}2\rfloor$ \\[2pt]
% \hline
% \end{tabular}
% \end{table}
%
%
% \begin{table}[p]\caption{颜色转换对}\label{tab.cnv}
% \begin{adjustwidth}{-2.5cm}{}
% \small\raggedleft
% \def\Em#1{\Model{#1}}
% \def\sp{\hbox to3em{\hfil$\ast$\hfil}}
% \def\id{\textrm{id}}
% \extrarowheight1pt
% \tabcolsep3pt
% \rowcolors{3}{tablerowcolor}{}
% \begin{tabular}{c*{11}{>{\sffamily}c}}
% \rowcolor{tableheadcolor}
% \hline
% \emph{从/到}&
%            \Em{rgb}& \Em{cmy}& \Em{cmyk}& \Em{hsb}& \Em{Hsb} & \Em{tHsb} & \Em{gray}& \Em{RGB}& \Em{HTML}& \Em{HSB}& \Em{Gray}\\
% \Em{rgb}&  \id     & \sp     & (cmy)    & \sp     & (hsb)    & (hsb)     & \sp      & \sp     & \sp      & (hsb)   & (gray)\\
% \Em{cmy}&  \sp     & \id     & \sp      & (rgb)   & (rgb)    & (rgb)     & \sp      & (rgb)   & (rgb)    & (rgb)   & (gray)\\
% \Em{cmyk}& (cmy)   & \sp     & \id      & (cmy)   & (cmy)    & (cmy)     & \sp      & (cmy)   & (cmy)    & (cmy)   & (gray)\\
% \Em{hsb}&  \sp     & (rgb)   & (rgb)    & \id     & \sp      & (Hsb)     & (rgb)    & (rgb)   & (rgb)    & \sp     & (rgb)\\
% \Em{Hsb}&  (hsb)   & (hsb)   & (hsb)    & \sp     & \id      & \sp       & (hsb)    & (hsb)   & (hsb)    & (hsb)   & (hsb)\\
% \Em{tHsb}& (Hsb)   & (Hsb)   & (Hsb)    & (Hsb)   & \sp      & \id       & (Hsb)    & (Hsb)   & (Hsb)    & (Hsb)   & (Hsb)\\
% \Em{gray}& \sp     & \sp     & \sp      & \sp     & \sp      & \sp       & \id      & \sp     & \sp      & \sp     & \sp\\
% \Em{RGB}&  \sp     & (rgb)   & (rgb)    & (rgb)   & (rgb)    & (rgb)     & (rgb)    & \id     & (rgb)    & (rgb)   & (rgb)\\
% \Em{HTML}& \sp     & (rgb)   & (rgb)    & (rgb)   & (rgb)    & (rgb)     & (rgb)    & (rgb)   & \id      & (rgb)   & (rgb)\\
% \Em{HSB}&  (hsb)   & (hsb)   & (hsb)    & \sp     & (hsb)    & (hsb)     & (hsb)    & (hsb)   & (hsb)    & \id     & (hsb)\\
% \Em{Gray}& (gray)  & (gray)  & (gray)   & (gray)  & (gray)   & (gray)    & \sp      & (gray)  & (gray)   & (gray)  & \id\\
% \Em{wave}& (hsb)   & (hsb)   & (hsb)    & \sp     & (hsb)    & (hsb)     & (hsb)    & (hsb)   & (hsb)    & (hsb)   & (hsb)\\
% \hline\hline
% \multicolumn{12}{c}{\id{} = 恒等函数；$\ast$ = 特定转换函数；}\\
% \multicolumn{12}{c}{\textsf{(model)} = 通过指定模型转换}\\
% \hline
% \end{tabular}
% \end{adjustwidth}
% \end{table}
%
%
%
% \subsection{整数模型与实数模型之间的转换}
%
% 我们固定一个正整数 $n$，定义集合 $\mathcal I_n:=\{0,1,\dots,n\}$ 和 $\mathcal R:= [0,1]$。
% $\nu\in\mathcal I_n$ 的补数为 $n-\nu$，$x\in\mathcal R$ 的补数为 $1-x$。
%
% \subsubsection{实数到整数的转换}\label{sec.real2int}
%
% 这种情况下直接的映射是
% \begin{gather}
% \textstyle\RtoI n : \mathcal R \to \mathcal I_n,\ x \mapsto \rd(n\cdot x,0) = \left\lfloor \frac12+n\cdot x\right\rfloor
% \end{gather}
% 其中 $\rd(r,d)$ 将实数 $r$ 四舍五入到 $d\ge0$ 位小数。
% 如下述引理所示，该映射几乎总是保持补数。
% \begin{lemma}[补数的保持]
% 对于 $x\in\mathcal R$，
% \begin{gather}
% \textstyle\RtoI n(x) + \RtoI n(1-x) = n \iff x \notin \mathcal R_n^\circ := \left\{ \frac1n\bigl(\nu - \frac12\bigr) \bigm\vert \nu=1,2,\dots,n\right\}.
% \end{gather}
% \end{lemma}
% \begin{proof}
% 令 $\nu:=\RtoI n(x)$，则由 $-\frac12 \le \eta := n\cdot x - \nu < \frac12$ 可得
% \begin{gather*}
% \textstyle \RtoI n(1-x) = \rd(n(1-x),0) = \rd(n-\nu-\eta,0) =
% \begin{cases}
%  n-\nu   & \text{若 $\eta\neq-\frac12$}\\
%  n-\nu+1 & \text{若 $\eta=-\frac12$}
% \end{cases}
% \end{gather*}
% 而 $\eta = -\frac12 \iff x = \frac1n\bigl(\nu-\frac12\bigr) \iff x\in\mathcal I_n'$。
% \end{proof}
% 注：集合 $\mathcal R_n^\circ$ 显然就是 $\RtoI n$ 不连续的点集。
%
%
% \subsubsection{整数到实数的转换}\label{sec.int2real}
%
% 这种情况下直接的函数是
% \begin{gather}
% \ItoR n^*: \mathcal I_n \to \mathcal R,\ \nu \mapsto \frac \nu n.
% \end{gather}
% 然而，这只是众多解中的一个：任何满足以下条件的函数 $\ItoR n:\mathcal I_n\to\mathcal R$
% \begin{gather}
% \nu\in\mathcal I_n \Rightarrow \Gamma_n\bigl(\ItoR n(\nu)\bigr) = \nu
% \end{gather}
% 等价于
% \begin{gather}
% \nu\in\mathcal I_n \Rightarrow \nu + \frac12 > n\cdot\ItoR n(\nu) \ge \nu - \frac12\label{eq.round}
% \end{gather}
% 至少能保证所有整数 $\nu$ 都可以通过将 $\ItoR n(\nu)$ 乘以 $n$ 并四舍五入到最近整数来重建。
% 补数的保持现在意味着
% \begin{gather}
% \nu\in\mathcal I_n \Rightarrow \ItoR n(\nu) + \ItoR n(n-\nu) = 1
% \end{gather}
% 显然 $\ItoR n = \ItoR n^*$ 满足此条件。
% 如果更一般地考虑变换
% \begin{gather}
% \ItoR n(\nu) = \frac{\nu + \alpha}{n + \beta}
% \end{gather}
% 其中 $\beta\neq-n$，则神奇不等式 (\ref{eq.round}) 等价于
% \begin{gather}
% \frac12 > \frac{\alpha n - \beta\nu}{n + \beta} \ge -\frac12
% \end{gather}
% 函数
% \begin{gather}
% \ItoR n':\mathcal I_n\to\mathcal R,\ \nu \mapsto
% \begin{cases}
% \frac\nu{n+1}& \text{若 $\nu\le\frac{n+1}2$}\\[0.25ex]
% \frac{\nu+1}{n+1}& \text{若 $\nu>\frac{n+1}2$}
% \end{cases}
% \end{gather}
% 满足此条件，且具有当 $n$ 为奇数时 $\ItoR n'\bigl(\frac{n+1}2\bigr) = \frac12$ 的优良特性。
% \begin{lemma}[补数的保持]
% 对于奇数 $n$ 和每个 $\nu\in\mathcal I_n$，
% \begin{gather}
% \textstyle\ItoR n'(\nu) + \ItoR n'(n-\nu) = 1 \iff \nu \notin \mathcal I_n^\circ := \left\{ \frac{n-1}2, \frac{n+1}2 \right\}.
% \end{gather}
% \end{lemma}
% \begin{proof}
% 该结论由以下论证得出：
% \begin{itemize}
% \item $\nu<\frac{n-1}2 \iff n-\nu>\frac{n+1}2$ \ 且 \ $\frac{n-1}2 + \frac{n+1}2 = n$；
% \item $\nu<\frac{n-1}2 \,\Rightarrow\, \ItoR n'(\nu) + \ItoR n'(n-\nu) = \frac\nu{n+1} + \frac{n-\nu+1}{n+1} = 1$；
% \item $\nu=\frac{n-1}2 \,\Rightarrow\, \ItoR n'(\nu) + \ItoR n'(n-\nu) = \frac{n-1}{2(n+1)} + \frac12 = \frac n{n+1} \neq1$。\qedhere
% \end{itemize}
% \end{proof}
%
% 目前，我们选择 \fbox{$\ItoR n := \ItoR n^*$} 作为默认变换函数。
%
% \medskip
% 另一种变体——对于大规模实时计算可能太慢——可用于构造预定义颜色集合。
% 基本思想是在保持对原始分辨率某种不变性的同时，最小化表示中的小数位数：
% \begin{gather}
% \textstyle\ItoR n'':\mathcal I_n\to\mathcal R,\ \nu \mapsto \rd\bigl(\frac\nu n, d_n(\frac\nu n)\bigr)
% \end{gather}
% 其中
% \begin{gather}
% d_n:[0,1]\to\bbbN,\ x \mapsto \min\bigl\{ d\in\bbbN \bigm\vert \RtoI n\bigl(\rd(\ItoR n^*(\RtoI n(x)),d)\bigr) = \RtoI n(x)\bigr\}
% \end{gather}
% 在最常见的情况 $n=255$ 下，最终最多保留 3 位小数；补数保持仅在 $\nu\in\{25,26,76,77,127,128,178,179,229,230\}$ 时被破坏，对应的十进制数集合为 $\{0.098,0.1,0.298,0.3,0.498,0.5,0.698,0.7,0.898,0.9\}$。
%
%
%
% \subsection{颜色转换与补数}\label{sec.cnv}
%
% 我们在此收集支持的颜色模型之间的具体转换公式。
% 表 \vref{tab.cnv} 概述了每种转换对的处理方式。
% 一般来说，\Index{PostScript}（如 \cite{plrm} 所述）被用作大多数计算的基础，因为它原生支持 \Model{rgb}、\Model{cmyk}、\Model{hsb} 和 \Model{gray} 颜色模型。
% 此外，\People{Alvy Ray}{Smith} 的论文 \cite{smith} 在 \cite{plrm} 中被引用为 \Model{hsb} 相关公式的参考。
%
% 首先，我们定义一个在所有转换公式中使用的常量：
% \begin{gather}
% E := (1,1,1)
% \end{gather}
%
%
% \subsubsection{\Model{rgb} 模型}
%
% \paragraph{\Model{rgb} 到 \Model{cmy} 的转换}\label{sec.rgb2cmy}
% 来源：\cite{plrm}，第 475 页。
% \begin{gather}
% (\emph{cyan}, \emph{magenta}, \emph{yellow}) := E-(\emph{red}, \emph{green}, \emph{blue})
% \end{gather}
%
%
% \paragraph{\Model{rgb} 到 \Model{hsb} 的转换 (1)}
% 我们设
% \begin{align}
% x & := \max\{\emph{red}, \emph{green}, \emph{blue}\}\\
% y & := \med\{\emph{red}, \emph{green}, \emph{blue}\}\\
% z & := \min\{\emph{red}, \emph{green}, \emph{blue}\}%\\
% \end{align}
% 其中 `$\med$' 表示这些值的中位数。
% 然后，
% \begin{align}
% \emph{brightness} & := x\\
% \intertext{情形 $x=z$：}
% \emph{saturation} & := 0\\
% \emph{hue} & := 0\\
% \intertext{情形 $x\neq z$：}
% \emph{saturation} & := \frac{x - z}x\\
% f & := \frac{x - y}{x - z}
% \end{align}
% \begin{gather}
% \emph{hue} := \frac 16\cdot
%  \begin{cases}
%  1-f & \text{若 $x=\emph{red}\ge\emph{green}\ge\emph{blue}=z$}\\
%  1+f & \text{若 $x=\emph{green}\ge\emph{red}\ge\emph{blue}=z$}\\
%  3-f & \text{若 $x=\emph{green}\ge\emph{blue}\ge\emph{red}=z$}\\
%  3+f & \text{若 $x=\emph{blue}\ge\emph{green}\ge\emph{red}=z$}\\
%  5-f & \text{若 $x=\emph{blue}\ge\emph{red}\ge\emph{green}=z$}\\
%  5+f & \text{若 $x=\emph{red}\ge\emph{blue}>\emph{green}=z$}
%  \end{cases}\label{eq.rgb2hsb1}
% \end{gather}
%
% \medskip
%
% 这基于 \cite{smith} 的 \emph{RGB to HSV Algorithm (Hexcone Model)}，其内容（稍作改写）为：
% \begin{gather}
% r := \frac{x - \emph{red}}{x - z},\qquad
% g := \frac{x - \emph{green}}{x - z},\qquad
% b := \frac{x - \emph{blue}}{x - z}\\
% \emph{hue} := \frac 16\cdot
%  \begin{cases}
%  5+b & \text{若 $\emph{red}=x$ 且 $\emph{green}=z$}\\
%  1-g & \text{若 $\emph{red}=x$ 且 $\emph{green}>z$}\\
%  1+r & \text{若 $\emph{green}=x$ 且 $\emph{blue}=z$}\\
%  3-b & \text{若 $\emph{green}=x$ 且 $\emph{blue}>z$}\\
%  3+g & \text{若 $\emph{blue}=x$ 且 $\emph{red}=z$}\\
%  5-r & \text{若 $\emph{blue}=x$ 且 $\emph{red}>z$}
%  \end{cases}
% \end{gather}
% 注意，特殊情形 $x=z$ 在 Smith 的原始算法中未完全覆盖；我们这里遵循 \Index{PostScript} 的实际行为。
%
% 由于需要对三个数排序来计算 $x,y,z$，该算法涉及多次比较。
% 我们现在介绍第二种更适合 \TeX 的方法。
%
%
% \paragraph{\Model{rgb} 到 \Model{hsb} 的转换 (2)}
% 令 $\beta$ 为一个以布尔表达式为参数的函数，表达式为真时返回 1，否则返回 0；设
% \begin{gather}
% i := 4\cdot\beta(\emph{red}\ge\emph{green}) + 2\cdot\beta(\emph{green}\ge\emph{blue}) + \beta(\emph{blue}\ge\emph{red}),
% \intertext{且}
% (\emph{hue},\emph{saturation},\emph{brightness}) :=
% \begin{cases}
% \Phi(\emph{blue},\emph{green},\emph{red},3,1) & \text{若 $i=1$}\\
% \Phi(\emph{green},\emph{red},\emph{blue},1,1) & \text{若 $i=2$}\\
% \Phi(\emph{green},\emph{blue},\emph{red},3,-1) & \text{若 $i=3$}\\
% \Phi(\emph{red},\emph{blue},\emph{green},5,1) & \text{若 $i=4$}\\
% \Phi(\emph{blue},\emph{red},\emph{green},5,-1) & \text{若 $i=5$}\\
% \Phi(\emph{red},\emph{green},\emph{blue},1,-1) & \text{若 $i=6$}\\
% (0,0,\emph{blue}) &  \text{若 $i=7$}
% \end{cases}\label{eq.rgb2hsb2}
% \intertext{其中}
% \Phi(x,y,z,u,v) := \Bigl( \frac{u\cdot(x-z)+v\cdot(x-y)}{6(x-z)}, \frac{x-z}x, x\Bigr)
% \end{gather}
%
% 特殊情形 $x=z$（等价于 $\emph{red}=\emph{green}=\emph{blue}$）由 $i=7$ 覆盖。
%
% \medskip
%
% 不难看出此算法是前一种方法的重新表述。
% 下表说明了从公式 (\ref{eq.rgb2hsb1}) 到公式 (\ref{eq.rgb2hsb2}) 的转换过程：
%
% \medskip
%
% \begin{tabular*}{\textwidth}{@{\extracolsep\fill}*5{>{$}c<{$}}l@{}}
% 6\cdot\emph{hue}& \emph{条件}& \emph{red}\ge\emph{green}& \emph{green}\ge\emph{blue}& \emph{blue}\ge\emph{red}& $i$\\\hline
%  1-f & \emph{red}\ge\emph{green}\ge\emph{blue}& 1& 1& *& \textbf6/7\\
%  1+f & \emph{green}\ge\emph{red}\ge\emph{blue}& *& 1& *& \textbf2/3/6/7\\
%  3-f & \emph{green}\ge\emph{blue}\ge\emph{red}& *& 1& 1& \textbf3/7\\
%  3+f & \emph{blue}\ge\emph{green}\ge\emph{red}& *& *& 1& \textbf1/3/5/7\\
%  5-f & \emph{blue}\ge\emph{red}\ge\emph{green}& 1& *& 1& \textbf5/7\\
%  5+f & \emph{red}\ge\emph{blue}\ge\emph{green}& 1& *& *& \textbf4/5/6/7\\
% \end{tabular*}
%
% \medskip
%
% 这里 $*$ 表示可能的 0 或 1 值。
% 粗体 $i$ 值标记该行所有 $*$ 值均为零时的主要情形。
% 最后一个不等式与公式 (\ref{eq.rgb2hsb1}) 的细微差异是有意为之，并无妨碍。
%
%
% \paragraph{\Model{rgb} 到 \Model{gray} 的转换}
% 来源：\cite{plrm}，第 474 页。
% \begin{align}
% \emph{gray} & := 0.3\cdot\emph{red} + 0.59\cdot\emph{green} + 0.11\cdot\emph{blue}
% \end{align}
%
%
% \paragraph{\Model{rgb} 到 \Model{RGB} 的转换}
% 如第 \vref{sec.real2int} 节所述。
% \begin{gather}
% (\emph{Red},\emph{Green},\emph{Blue}) := \bigl(\RtoI L(\emph{red}), \RtoI L(\emph{green}), \RtoI L(\emph{blue})\bigr)
% \end{gather}
%
%
% \paragraph{\Model{rgb} 到 \Model{HTML} 的转换}
% 如第 \vref{sec.real2int} 节所述。
% 之后转换为 6 位十六进制数。
% 当然，乘法和求和可以简单地用 2 位十六进制数的文本拼接来代替。
% \begin{gather}
% \emph{RRGGBB} := \bigl( 65536\cdot\RtoI L(\emph{red}) + 256\cdot \RtoI L(\emph{green}) + \RtoI L(\emph{blue}) \bigr)_{\emph{hex}}
% \end{gather}
%
%
% \paragraph{\Model{rgb} 颜色的补色}
% 我们直接取补向量：
% \begin{gather}
% (\emph{red}^*, \emph{green}^*, \emph{blue}^*) := E - (\emph{red}, \emph{green}, \emph{blue})
% \end{gather}
%
%
% \subsubsection{\Model{cmy} 模型}
%
% \paragraph{\Model{cmy} 到 \Model{rgb} 的转换}
% 这只是 $\Model{rgb}\to\Model{cmy}$ 的逆运算，参见第 \vref{sec.rgb2cmy} 节。
% \begin{gather}
% (\emph{red}, \emph{green}, \emph{blue}) := E - (\emph{cyan}, \emph{magenta}, \emph{yellow})
% \end{gather}
%
%
% \paragraph{\Model{cmy} 到 \Model{cmyk} 的转换}\label{sec.cmy2cmyk}
% 这可能是我们转换任务中最难的：许多资料强调，由于设备依赖性，这种情况不存在通用的转换算法。
% 以下算法是 \cite{plrm} 第 476 页所给算法的扩展版本。
% \begin{align}
% k & := \min\{\emph{cyan},\emph{magenta},\emph{yellow}\}\\
% \emph{cyan} & := \min\{1,\max\{0,\emph{cyan} - \emph{UCR}_c(k)\}\}\\
% \emph{magenta} & := \min\{1,\max\{0,\emph{magenta} - \emph{UCR}_m(k)\}\}\\
% \emph{yellow} & := \min\{1,\max\{0,\emph{yellow} - \emph{UCR}_y(k)\}\}\\
% \emph{black} & := \emph{BG}(k)
% \end{align}
% 这里需要四个额外函数：
% \begin{align}
% \emph{UCR}_c, \emph{UCR}_m, \emph{UCR}_y & : [0,1] \to [-1,1]\tag*{\em 底色去除}\\
% \emph{BG} & : [0,1] \to [0,1]\tag*{\em 黑版生成}
% \end{align}
% 这些函数是设备相关的，参见 \cite{plrm} 中的说明。
% 虽然有迹象表明它们应选为非线性函数，但在对目标设备没有进一步了解的情况下，我们将其定义为线性函数：
% \begin{align}
% \emph{UCR}_c(k) &:= \beta_c\cdot k\\
% \emph{UCR}_m(k) &:= \beta_m\cdot k\\
% \emph{UCR}_y(k) &:= \beta_y\cdot k\\
% \emph{BG}(k) &:= \beta_k\cdot k
% \end{align}
% \DescribeMacro\adjustUCRBG
% 其中参数可在文档任意位置通过 |\def\adjustUCRBG{|\Meta{$\beta_c$},\Meta{$\beta_m$},\Meta{$\beta_y$},\Meta{$\beta_k$}|}| 指定，默认为 |{|$\adjustUCRBG$|}|。
%
%
% \paragraph{\Model{cmy} 到 \Model{gray} 的转换}
% 这由转换链 $\Model{cmy}\to\Model{rgb}\to\Model{gray}$ 推导而来。
% \begin{align}
% \emph{gray} & := 1 - (0.3\cdot\emph{cyan} + 0.59\cdot\emph{magenta} + 0.11\cdot\emph{yellow})
% \end{align}
%
%
% \paragraph{\Model{cmy} 颜色的补色}
% 我们直接取补向量：
% \begin{gather}
% (\emph{cyan}^*, \emph{magenta}^*, \emph{yellow}^*) := E - (\emph{cyan}, \emph{magenta}, \emph{yellow})
% \end{gather}
%
%
% \subsubsection{\Model{cmyk} 模型}
%
% \paragraph{\Model{cmyk} 到 \Model{cmy} 的转换}\label{sec.cmyk2cmy}
% 基于 \cite{plrm} 第 477 页，结合 $\Model{rgb}\to\Model{cmy}$ 转换。
% \begin{align}
% \emph{cyan} & := \min\{1,\emph{cyan} + \emph{black}\}\\
% \emph{magenta} & := \min\{1,\emph{magenta} + \emph{black}\}\\
% \emph{yellow} & := \min\{1,\emph{yellow} + \emph{black}\}
% \end{align}
%
%
% \paragraph{\Model{cmyk} 到 \Model{gray} 的转换}
% 来源：\cite{plrm}，第 475 页。
% \begin{align}
% \emph{gray} & := 1 - \min\{1,0.3\cdot\emph{cyan} + 0.59\cdot\emph{magenta} + 0.11\cdot\emph{yellow} + black\}
% \end{align}
%
%
% \paragraph{\Model{cmyk} 颜色的补色}
% 简单的向量补数不能产生有用的结果。
% 因此，我们先将其转换到 \Model{cmy} 模型，在那里计算补色，再转换回 \Model{cmyk}。
%
%
% \subsubsection{\Model{hsb} 模型}
%
% \paragraph{\Model{hsb} 到 \Model{rgb} 的转换}
% \begin{gather}
% (\emph{red},\emph{green},\emph{blue}) :=
%  \emph{brightness}\cdot( E - \emph{saturation}\cdot F)\\
% \intertext{with}
% i := \left\lfloor 6\cdot\emph{hue}\right\rfloor,\qquad f := 6\cdot\emph{hue} - i\\
% \intertext{and}
% F :=
%  \begin{cases}
%   (0,1-f,1) & \text{if $i=0$}\\
%   (f,0,1) & \text{if $i=1$}\\
%   (1,0,1-f) & \text{if $i=2$}\\
%   (1,f,0) & \text{if $i=3$}\\
%   (1-f,1,0) & \text{if $i=4$}\\
%   (0,1,f) & \text{if $i=5$}\\
%   (0,1,1) & \text{if $i=6$}
%  \end{cases}\label{eq.hsb2rgb2}
% \end{gather}
%
% \medskip
%
% 这基于 \cite{smith} 的 \emph{HSV to RGB Algorithm (Hexcone Model)}，其内容（稍作改写）为：
% \begin{align}
% m & := 1-\emph{saturation}\\
% n & := 1-f\cdot\emph{saturation}\\
% k & := 1-(1-f)\cdot\emph{saturation}\\
% (\emph{red},\emph{green},\emph{blue}) & := \emph{brightness}\cdot
%  \begin{cases}
%   (1,k,m) & \text{若 $i=0,6$}\\
%   (n,1,m) & \text{若 $i=1$}\\
%   (m,1,k) & \text{若 $i=2$}\\
%   (m,n,1) & \text{若 $i=3$}\\
%   (k,m,1) & \text{若 $i=4$}\\
%   (1,m,n) & \text{若 $i=5$}
%  \end{cases}\label{eq.hsb2rgb1}
% \end{align}
% 注意，$i=6$ 的情形（由 $\emph{hue}=1$ 产生）在 Smith 的算法中缺失。
% 由于
% \begin{align}
% \lim_{f\to1} (0,1,f) = (0,1,1) = \lim_{f\to0} (0,1-f,1)
% \end{align}
% 显然，为了得到连续函数，$i=6$ 时 $F$ 只有一种定义方式，如公式 (\ref{eq.hsb2rgb2}) 所示。
% 这已被转换回公式 (\ref{eq.hsb2rgb1})。
% 类似的论证表明，$F$ 确实是 \emph{hue} 在整个区间 $[0,1]$ 上的连续函数。
%
%
% \paragraph{\Model{hsb} 到 \Model{Hsb} 的转换}
% 只需改变第一个分量。
% \begin{gather}
% (\emph{hue\Circ}, \emph{saturation}, \emph{brightness}) := (H\cdot\emph{hue}, \emph{saturation}, \emph{brightness})
% \end{gather}
%
%
% \paragraph{\Model{hsb} 到 \Model{HSB} 的转换}
% 如第 \vref{sec.real2int} 节所述。
% \begin{gather}
% (\emph{Hue}, \emph{Saturation}, \emph{Brightness}) := \bigl(\RtoI M(\emph{hue}), \RtoI M(\emph{saturation}), \RtoI M(\emph{brightness})\bigr)
% \end{gather}
%
%
% \paragraph{\Model{hsb} 颜色的补色}
% 我们在文献中未找到公式，因此在此给出简短证明。
% \begin{lemma}
% \textup{\Model{hsb}} 补色可由以下公式计算：\em
% \begin{align}
% \emph{hue}^* & :=
% \begin{cases}
% \emph{hue} + \frac12 & \text{若 $\emph{hue}<\frac12$}\\[1ex]
% \emph{hue} - \frac12 & \text{若 $\emph{hue}\ge\frac12$}\\
% \end{cases}\\[1ex]
% \emph{brightness}^* & := 1 - \emph{brightness}\cdot(1-\emph{saturation})\\[1ex]
% \emph{saturation}^* & :=
% \begin{cases}
% 0 & \text{若 $\emph{brightness}^*=0$}\\[1ex]
% \displaystyle\frac{\emph{brightness}\cdot\emph{saturation}}{\emph{brightness}^*} & \text{若 $\emph{brightness}^*\ne0$}
% \end{cases}
% \end{align}
% \end{lemma}
%
% \begin{proof}
% 从原始颜色 $C=(h,s,b)$ 出发，我们通过给定公式定义颜色 $C^*=(h^*,s^*,b^*)$，将 $C$ 和 $C^*$ 都转换到 \Model{rgb} 模型，并证明
% \begin{gather}
% C_{\Model{rgb}} + C_{\Model{rgb}}^* = b\cdot(E - s\cdot F) + b^*\cdot(E - s'\cdot F^*) \stackrel!= E,\label{eq.hsbcompl1}
% \end{gather}
% 即 $C_{\Model{rgb}}$ 是 $C_{\Model{rgb}}^*$ 的补色。
% 首先注意 $C^*$ 的参数在合法范围 $[0,1]$ 内。
% 这对 $h^*,b^*$ 是显然的。
% 由
% $b^* = 1-b\cdot(1-s) = 1 - b + b\cdot s$
% 可得 $b\cdot s = b^* - (1-b) \le b^*$，因此 $s^*\in[0,1]$，且
% \begin{gather*}
% b^* = 0 \Leftrightarrow s=0 \text{ 且 } b=1.
% \end{gather*}
% 因此公式 (\ref{eq.hsbcompl1}) 在 $b^*=0$ 时成立。
% 现在假设 $b^*\ne0$，则
% \begin{align*}
% C_{\Model{rgb}} + C_{\Model{rgb}}^*
% &= b\cdot(E - s\cdot F) + b^*\cdot\Bigl(E - \frac{b\cdot s}{b^*}\cdot F^*\Bigr) \\
% &= b\cdot E - b\cdot s\cdot F + b^*\cdot E - b\cdot s\cdot F^*\\[1ex]
% &= E - b\cdot s\cdot(F+F^*-E)
% \end{align*}
% 因为 $b^*=1-b+bs$。
% 因此只需证明
% \begin{gather}
% F+F^*=E.\label{eq.hsbcompl2}
% \end{gather}
% 由
% \begin{gather*}\textstyle
% h<\frac12 \Rightarrow h^*=h+\frac12 \Rightarrow 6h^* = 6h+3 \Rightarrow i^*=i+3 \text{ 且 } f^*=f
% \end{gather*}
% 从 (\ref{eq.hsb2rgb2}) 容易看出公式 (\ref{eq.hsbcompl2}) 在 $i=0,1,2$ 时成立。
% 类似地，
% \begin{gather*}\textstyle
% h\ge\frac12 \Rightarrow h^*=h-\frac12 \Rightarrow 6h^* = 6h-3 \Rightarrow i^*=i-3 \text{ 且 } f^*=f
% \end{gather*}
% 再由 (\ref{eq.hsb2rgb2}) 可得公式 (\ref{eq.hsbcompl2}) 在 $i=3,4,5$ 时成立。
% 最后，若 $i=6$，则 $f=0$，$F+F^* = (0,1,1) + (1,0,0) = E$。
% \end{proof}
%
%
% \subsubsection{\Model{Hsb} 模型}
%
% \paragraph{\Model{Hsb} 到 \Model{hsb} 的转换}
% 只需改变第一个分量。
% \begin{gather}
% (\emph{hue}, \emph{saturation}, \emph{brightness}) := (\emph{hue\Circ}/H, \emph{saturation}, \emph{brightness})
% \end{gather}
%
%
% \paragraph{\Model{Hsb} 到 \Model{tHsb} 的转换}
% 在公式 (\ref{eq.tHsb1})--(\ref{eq.tHsb3}) 的设置下，我们只需在公式 (\ref{eq.tHsb4}) 中交换字母 $x$ 和 $y$，即可得到逆变换：
% \begin{gather}
% \displaystyle \emph{hue\Circ}\in\left[y_{\eta-1},y_\eta\right] \Rightarrow
% \emph{hue\Circ} := x_{\eta-1} + \frac{x_\eta - x_{\eta-1}}{y_\eta - y_{\eta-1}}\cdot(\emph{hue\Circ} - y_{\eta-1})
% \end{gather}
% 而 \emph{saturation} 和 \emph{brightness} 保持不变。
%
%
% \subsubsection{\Model{tHsb} 模型}
%
% \paragraph{\Model{tHsb} 到 \Model{Hsb} 的转换}\label{sec.tHsb2Hsb}
% \DescribeMacro\rangeHsb
% \DescribeMacro\rangetHsb
% 我们假设 $\texttt{\string\rangeHsb}=H$，且 |\rangetHsb| 展开为
% \begin{gather}
% x_1,y_1;x_2,y_2;\dots;x_{h-1},y_{h-1}\label{eq.tHsb1}
% \intertext{其中}
% x_0:=0<x_1<x_2<\dots<x_{h-1}<x_h:=H\label{eq.tHsb2}\\
% y_0:=0<y_1<y_2<\dots<y_{h-1}<y_h:=H\label{eq.tHsb3}
% \end{gather}
% 其中 $h>0$ 为整数。
% 现在 $x$ 和 $y$ 值确定了一个分段线性变换：
% \begin{gather}
% \displaystyle \emph{hue\Circ}\in\left[x_{\eta-1},x_\eta\right] \Rightarrow
% \emph{hue\Circ} := y_{\eta-1} + \frac{y_\eta - y_{\eta-1}}{x_\eta - x_{\eta-1}}\cdot(\emph{hue\Circ} - x_{\eta-1})
% \label{eq.tHsb4}
% \end{gather}
% 而 \emph{saturation} 和 \emph{brightness} 保持不变。
%
%
% \subsubsection{\Model{gray} 模型}
%
% \paragraph{\Model{gray} 到 \Model{rgb} 的转换}
% 来源：\cite{plrm}，第 474 页。
% \begin{gather}
% (\emph{red}, \emph{green}, \emph{blue}) := \emph{gray}\cdot E
% \end{gather}
%
%
% \paragraph{\Model{gray} 到 \Model{cmy} 的转换}
% 这由转换链 $\Model{gray}\to\Model{rgb}\to\Model{cmy}$ 推导而来。
% \begin{gather}
% (\emph{cyan}, \emph{magenta}, \emph{yellow}) := (1-\emph{gray})\cdot E
% \end{gather}
%
%
% \paragraph{\Model{gray} 到 \Model{cmyk} 的转换}
% 来源：\cite{plrm}，第 475 页。
% \begin{gather}
% (\emph{cyan}, \emph{magenta}, \emph{yellow}, \emph{black}) := (0,0,0,1-\emph{gray})
% \end{gather}
%
%
% \paragraph{\Model{gray} 到 \Model{hsb} 的转换}
% 这由转换链 $\Model{gray}\to\Model{rgb}\to\Model{hsb}$ 推导而来。
% \begin{gather}
% (\emph{hue}, \emph{saturation}, \emph{brightness}) := (0,0,\emph{gray})
% \end{gather}
%
%
% \paragraph{\Model{gray} 到 \Model{Hsb}/\Model{tHsb} 的转换}
% 这由转换链 $\Model{gray}\to\Model{hsb}\to\Model{Hsb}$ 推导而来，如适用再经过 $\Model{Hsb}\to\Model{tHsb}$。
% \begin{gather}
% (\emph{hue\Circ}, \emph{saturation}, \emph{brightness}) := (0,0,\emph{gray})
% \end{gather}
%
%
% \paragraph{\Model{gray} 到 \Model{Gray} 的转换}
% 如第 \vref{sec.real2int} 节所述。
% \begin{gather}
% \emph{Gray} := \RtoI N(\emph{gray})
% \end{gather}
%
%
% \paragraph{\Model{gray} 颜色的补色}
% 这与 \Model{rgb} 的情形类似：
% \begin{gather}
% \emph{gray}^* := 1 - \emph{gray}
% \end{gather}
%
%
% \subsubsection{\Model{RGB} 模型}
%
% \paragraph{\Model{RGB} 到 \Model{rgb} 的转换}
% 如第 \vref{sec.int2real} 节所述。
% \begin{gather}
% (\emph{red}, \emph{green}, \emph{blue}) := \bigl(\ItoR L(\emph{Red}), \ItoR L(\emph{Green}), \ItoR L(\emph{Blue})\bigr)
% \end{gather}
%
%
% \subsubsection{\Model{HTML} 模型}
%
% \paragraph{\Model{HTML} 到 \Model{rgb} 的转换}
% 如第 \vref{sec.int2real} 节所述：从 \emph{RRGGBB} 出发，设
% \begin{gather}
% (\emph{red}, \emph{green}, \emph{blue}) := \bigl(\ItoR{255}(\emph{RR}_{\emph{dec}}), \ItoR{255}(\emph{GG}_{\emph{dec}}), \ItoR{255}(\emph{BB}_{\emph{dec}})\bigr)
% \end{gather}
%
%
% \subsubsection{\Model{HSB} 模型}
%
% \paragraph{\Model{HSB} 到 \Model{hsb} 的转换}
% 如第 \vref{sec.int2real} 节所述。
% \begin{gather}
% (\emph{hue}, \emph{saturation}, \emph{brightness}) := \bigl(\ItoR M(\emph{Hue}), \ItoR M(\emph{Saturation}), \ItoR M(\emph{Brightness})\bigr)
% \end{gather}
%
%
% \subsubsection{\Model{Gray} 模型}
%
% \paragraph{\Model{Gray} 到 \Model{gray} 的转换}
% 如第 \vref{sec.int2real} 节所述。
% \begin{gather}
% \emph{gray} := \ItoR N(\emph{Gray})
% \end{gather}
%
%
% \subsubsection{\Model{wave} 模型}
%
% \paragraph{\Model{wave} 到 \Model{rgb} 的转换}
% 来源：基于 \People{Dan}{Bruton} 的算法 \cite{Bruton}。
% 设 $\lambda$ 为可见光波长，单位为纳米（nm），即 $\lambda\in\left[380,780\right]$。
% 我们进一步假设 $\gamma>0$ 为固定常数（\cite{Bruton} 中取 $\gamma=0.8$）。
% 首先设
% \begin{align}
% (r,g,b) & :=
%  \begin{cases}
%   \displaystyle \Bigl(\frac{440-\lambda}{440-380},0,1\Bigr) & \text{若 $\lambda\in[380,440[$}\\[1.75ex]
%   \displaystyle \Bigl(0,\frac{\lambda-440}{490-440},1\Bigr) & \text{若 $\lambda\in[440,490[$}\\[1.75ex]
%   \displaystyle \Bigl(0,1,\frac{510-\lambda}{510-490}\Bigr) & \text{若 $\lambda\in[490,510[$}\\[1.75ex]
%   \displaystyle \Bigl(\frac{\lambda-510}{580-510},1,0\Bigr) & \text{若 $\lambda\in[510,580[$}\\[1.75ex]
%   \displaystyle \Bigl(1,\frac{645-\lambda}{645-580},0\Bigr) & \text{若 $\lambda\in[580,645[$}\\[2.5ex]
%   \displaystyle (1,0,0) & \text{若 $\lambda\in[645,780]$}\\
%  \end{cases}\label{eq.gamma2rgb1}\\
% \intertext{然后，为了让强度在视觉极限附近衰减，}
% f & :=
%  \begin{cases}
%   \displaystyle 0.3 + 0.7\cdot\frac{\lambda-380}{420-380} & \text{若 $\lambda\in\left[380,420\right[$}\\[1.25ex]
%   \displaystyle 1 & \text{若 $\lambda\in[420,700]$}\\[1ex]
%   \displaystyle 0.3 + 0.7\cdot\frac{780-\lambda}{780-700} & \text{若 $\lambda\in\left]700,780\right]$}\\
%  \end{cases}\label{eq.gamma2rgb2}\\
% \intertext{最后}
% (\emph{red},\emph{green},\emph{blue}) & := \bigl((f\cdot r)^\gamma, (f\cdot g)^\gamma, (f\cdot b)^\gamma\bigr)
% \end{align}
% 公式 (\ref{eq.gamma2rgb1}) 区间边界处的中间颜色 $(r,g,b)$ 是众所周知的：
% 当 $\lambda=380,440,490,510,580,645$ 时，分别得到 \Color{magenta}、\Color{blue}、\Color{cyan}、\Color{green}、\Color{yellow}、\Color{red}。
% 这些颜色在 \Model{hsb} 模型中表示为 $\emph{hue}=\frac56,\frac46,\frac36,\frac26,\frac16,\frac06$，而 $\emph{saturation}=\emph{brightness}=1$（这 6 种颜色均如此）。
% 此外，这些 \Model{hsb} 表示与实际的 $\gamma$ 值无关。
% 在此模型框架内，我们观察到视觉极限附近的强度衰减——如公式 (\ref{eq.gamma2rgb2}) 所示——转化为向边缘方向递减的 \emph{brightness} 参数。
% 简单计算表明，算法边缘 $\lambda=380,780$ 分别产生颜色 |magenta!|$0.3^\gamma$|!black|、|red!|$0.3^\gamma$|!black|。
% 我们认为没有理由不以类似方式扩展这些边缘，以使两端最终达到真正的 \Color{black}。
% 现在我们准备好将所有内容转换为另一种更自然的算法。
%
%
% \paragraph{\Model{wave} 到 \Model{hsb} 的转换}
% 设 $\lambda>0$ 为波长，单位为纳米（nm），并设
% \begin{gather}
% \varrho : \bbbR \to \left[0,1\right],\ x \mapsto \bigl( \min\{ 1, \max\{ 0,x\} \} )^\gamma
% \end{gather}
% 其中 $\gamma>0$ 为固定修正数。
% 则
% \begin{align}
% \emph{hue} & := \displaystyle \frac16\cdot
%  \begin{cases}
%   \displaystyle 4 + \varrho\Bigl(\frac{\lambda-440}{380-440}\Bigr) & \text{若 $\lambda<440$}\\[1.75ex]
%   \displaystyle 4 - \varrho\Bigl(\frac{\lambda-440}{490-440}\Bigr) & \text{若 $\lambda\in[440,490[$}\\[1.75ex]
%   \displaystyle 2 + \varrho\Bigl(\frac{\lambda-510}{490-510}\Bigr) & \text{若 $\lambda\in[490,510[$}\\[1.75ex]
%   \displaystyle 2 - \varrho\Bigl(\frac{\lambda-510}{580-510}\Bigr) & \text{若 $\lambda\in[510,580[$}\\[1.75ex]
%   \displaystyle 0 + \varrho\Bigl(\frac{\lambda-645}{580-645}\Bigr) & \text{若 $\lambda\in[580,645[$}\\[2.5ex]
%   \displaystyle 0 & \text{若 $\lambda\ge645$}\\
%  \end{cases}\label{eq.gamma2hsb1}\\[1ex]
% \emph{saturation} & := 1\\[1ex]
% \emph{brightness} & :=
%  \begin{cases}
%   \displaystyle \varrho\Bigl(0.3 + 0.7\cdot\frac{\lambda-380}{420-380}\Bigr) & \text{若 $\lambda<420$}\\[1.5ex]
%   \displaystyle 1 & \text{若 $\lambda\in[420,700]$}\\[1ex]
%   \displaystyle \varrho\Bigl(0.3 + 0.7\cdot\frac{\lambda-780}{700-780}\Bigr) & \text{若 $\lambda>700$}\\
%  \end{cases}\label{eq.gamma2hsb2}
% \end{align}
% 为完整起见我们注意到，与 $\gamma$ 无关，
% \begin{gather*}
% (\emph{hue},\emph{saturation},\emph{brightness}) =
%  \begin{cases}
%   \bigl(\frac56,1,0\bigr) & \text{若 $\lambda\le 380 - \frac{3\cdot(420-380)}7 = 362.857\dots$}\\[1.5ex]
%   \bigl(0,1,0\bigr) & \text{若 $\lambda\ge 780 + \frac{3\cdot(780-700)}7 = 814.285\dots$}\\
%  \end{cases}\label{eq.gamma2hsb3}
% \end{gather*}
% $\gamma$ 的最佳（或至少是一个好的）值是多少？
% 在原始算法 \cite{Bruton} 中，选择了 $\gamma=0.8$。
% 然而，我们未能发现 $\gamma=0.8$ 和 $\gamma=1$ 之间有明显的可见差异。
% 因此，目前 \XC 的实现使用后者，即纯线性方法。
% 在 \Package{pstricks} 示例文件 \href{xcolor2.pdf}{\File{xcolor2.tex}} 中，有不同 $\gamma$ 值的演示。
%
% \clearpage
%
%
%
% \begin{adjustwidth}{-2.5cm}{}
% \begingroup\raggedright
% \begin{thebibliography}{00}
% \addcontentsline{toc}{section}{参考文献}
%
% \bibitem{plrm}%
%  {Adobe Systems Incorporated: ``PostScript Language Reference Manual''. Addison-Wesley, third edition, 1999.
%   \URL{http://www.adobe.com/products/postscript/pdfs/PLRM.pdf}}
%
% \bibitem{fbox}%
%  {\People{Donald}{Arseneau}: ``Patch so \cmd\fbox{} draws frame on top of text''. \LaTeX{} bug report, latex/3655, 2004/03/18.
%   \URL{http://www.latex-project.org/cgi-bin/ltxbugs2html?pr=latex/3655}}
%
% \bibitem{url}%
%  {\People{Donald}{Arseneau}: \Package{url} package, ``2005/06/27  ver 3.2  Verb mode for urls, etc.''.
%   \CTAN{macros/latex/contrib/misc/}{url.sty}}
%
% \bibitem{Bruton}%
%  {\People{Dan}{Bruton}: ``Approximate RGB values for Visible Wavelengths'', 1996.
%   \URL{http://www.physics.sfasu.edu/astro/color/spectra.html}}
%
% \bibitem{dpc}%
%  {\People{David P.}{Carlisle}: ``Packages in the `graphics' bundle'', 2014.
%   \CTAN{macros/latex/required/graphics/}{grfguide.*}}
%
% \bibitem{color}%
%  {\People{David P.}{Carlisle}: \LCP, ``2016/01/03 v1.1b Standard LaTeX Color (DPC)''.
%   \CTAN{macros/latex/required/graphics/}{color.dtx}}
%
% \bibitem{colortbl}%
%  {\People{David P.}{Carlisle}: \Package{colortbl} package, ``2001/02/13 v0.1j Color table columns''.
%   \CTAN{macros/latex/contrib/colortbl/}{}}
%
% \bibitem{pstcol}%
%  {\People{David P.}{Carlisle}, \People{Herbert}{Vo\ss}, \People{Rolf}{Niepraschk}: \Package{pstcol} package, ``2005/11/16 v1.2 LaTeX wrapper for `PSTricks'{}''.
%   \CTAN{macros/graphics/pstricks/latex/}{pstcol.sty}}
%
% \bibitem{chroma}%
%  {\People{Uwe}{Kern}: ``Chroma: a reference book of \LaTeX{} colors''.
%   \CTAN{info/colour/chroma/}{} and
%   \URL{http://www.ukern.de/tex/chroma.html}}
%
% \bibitem{xcolor}%
%  {\People{Uwe}{Kern}: \Package{xcolor} package, ``\LaTeX{} color extensions''.
%   \CTAN{macros/latex/contrib/xcolor/}{} and
%   \URL{http://www.ukern.de/tex/xcolor.html}}
%
% \bibitem{mik}%
%  {MiK\TeX{} Project:
%   \URL{http://www.miktex.org/}}
%
% \bibitem{colorinfo}%
%  {\People{Rolf}{Niepraschk}: \Package{colorinfo} package, ``2003/05/04 v0.3c Info from defined colors''.
%   \CTAN{macros/latex/contrib/colorinfo/}{}}
%
% \bibitem{pdfcolmk}%
%  {\People{Heiko}{Oberdiek}: \Package{pdfcolmk} package, ``2006/02/20 v0.8 PDFtex COLor MarK''.
%   \CTAN{macros/latex/contrib/oberdiek/}{pdfcolmk.*}}
%
% \bibitem{hyperref}%
%  {\People{Sebastian}{Rahtz}, \People{Heiko}{Oberdiek}: \Package{hyperref} package, ``2006/09/06 v6.75e Hypertext links for \LaTeX''.
%   \CTAN{macros/latex/contrib/hyperref/}{}}
%
% \bibitem{smith}%
%  {\People{Alvy Ray}{Smith}: ``Color Gamut Transform Pairs''. \emph{Computer Graphics} (ACM SIGGRAPH), Volume 12, Number 3, August 1978.
%   \URL{http://alvyray.com/Papers/PapersCG.htm}}
%
% \bibitem{HTML4}%
%  {World Wide Web Consortium: ``HTML4 color keywords''.
%   \URL{http://www.w3.org/TR/css3-color/\#html4}}
%
% \bibitem{svg}%
%  {World Wide Web Consortium: ``Scalable Vector Graphics (\Index{SVG}) 1.1 Specification --- Basic Data Types and Interfaces''.
%   \URL{http://www.w3.org/TR/SVG11/types.html\#ColorKeywords}}
%
% \end{thebibliography}
% \endgroup
% \clearpage
%
% \columnsep2.25em
% \begin{multicols}{2}[\section*{附录}%
% \addcontentsline{toc}{section}{附录}]
%
% \subsection*{致谢}
% \addcontentsline{toc}{subsection}{致谢}
% 本宏包基于 \cite{color}（版权所有 (C) 1994--1999 \People{David P.}{Carlisle}）并包含其复制的代码，后者是标准 \LaTeX `Graphics Bundle' 的一部分。
% 虽然 \XC 中添加了许多命令和功能，且大多数原始 \LC 命令已被重写或改造，但没有 \LC 就不会有本宏包。
% 因此，作者感谢 \People{David P.}{Carlisle} 创建了 \LC 及其相关文件。
%
% \subsection*{商标}
% \addcontentsline{toc}{subsection}{商标}
% 本文档中出现的商标均未标注商标符号；它们归各自商标所有者所有。
% 此处无意侵权，其使用旨在惠及商标所有者。
%
% \raggedright
% \leftmargini1em
% \leftmarginii1em
% \leftmarginiii1em
%
% \subsection*{已知问题}
% \addcontentsline{toc}{subsection}{已知问题}
% \begin{itemize}
% \item |\rowcolors[\hline]...| 不能与 \Package{longtable} 配合使用。
% \end{itemize}
%
% \subsection*{历史}
% \addcontentsline{toc}{subsection}{历史}
% \subsubsection*{2023/11/15 v3.01}
% \begin{itemize}
% \item 使用过时选项 usenames 时仅发出警告而不报错（issue \#31）
% \end{itemize}
% \subsubsection*{2023/11/11 v3.00}
% \begin{itemize}
% \item 若检测到早于 2022-06-01 的 \LaTeX，则使用冻结版本。
% 所有新功能仅在检测到 LaTeX 2022-06-01 或更新版本时生效。
% 冻结版本也可通过 \verb+\usepackage{xcolor}[=v2]+ 请求。
% \item 宏包选项使用新的 LaTeX 键值处理程序。
% 这意味着选项冲突错误将减少（issue \#29）。
% \item |\rowcolors| 命令已移至 |colortbl| 宏包，不再由 |xcolor| 定义。
% \item fixpdftex 选项已过时，不再起任何作用（issue \#20）
% \item 扩展 |\normalcolor| 以同时设置点颜色，issue \#19
% \item 修复 Gray 颜色模型，issue \#23
% \end{itemize}
%
% \subsubsection*{2022/06/12 v2.14}
% \begin{itemize}
% \item 从 \LaTeX 加载代码，在可用时定义 |\mathcolor|。
% \item 错误修复：
%   \begin{itemize}
%   \item 修改 |\XC@info|，避免 |\colorlet{foo}{}| 产生虚假错误，issue \#12
%   \item 将 |\normalcolor|、|\colorbox|、|\fcolorbox|、|\textcolor|、|\pagecolor|、|\nopagecolor|、|\definecolor|、|\DefineNamedColor| 改为与 |color.sty| 相同的保护形式，issue \#17
%   \end{itemize}
% \end{itemize}

% \subsubsection*{2021/10/31 v2.13}
% \begin{itemize}
% \item 错误修复：
%   \begin{itemize}
%   \item 避免使用 |\string| 生成 csname
%   \item 声明 |dvisvgm| 选项，并使 |dvipdfm| 成为 |dvipdfmx| 的别名，以与 graphics 和 color 宏包保持一致。
%   \item 修改 |\XC@info|，避免 |\colorlet{foo}{.!50}| 产生多余文本
%   \end{itemize}
% \end{itemize}
%
% \subsubsection*{2016/05/11 v2.12}
% \begin{itemize}
% \item 新功能：
%   \begin{itemize}
%   \item |\nopagecolor| 命令，随 \LC v1.1a 引入（示例见 \File{xcolor3.tex}）；
%   \item \Option{luatex} 驱动选项（代码由 DPC 提供），修复因新版 Lua\TeX 变更导致的不兼容。
%   \end{itemize}
% \item 错误修复：
%   \begin{itemize}
%   \item |\XC@ifxcase| 调用可能导致的名称冲突；
%   \item 内部宏 |\@hex@@Hex| 不正确。
%   \end{itemize}
% \end{itemize}
%
% \subsubsection*{2007/01/21 v2.11}
% \begin{itemize}
% \item 新功能：
%   \begin{itemize}
%   \item 预定义颜色集合中添加颜色名 \Color{lime} 和 \Color{teal}。
%   \end{itemize}
% \item 错误修复：
%   \begin{itemize}
%   \item \Package{hyperref} 相关选项中 |\XC@strip@comma| 调用不正确。
%   \end{itemize}
% \end{itemize}
%
% \subsubsection*{2006/11/28 v2.10}
% \begin{itemize}
% \item 新功能：
%   \begin{itemize}
%   \item \Option{fixinclude} 选项防止 \Program{dvips} 在通过 |\color{red}\includegraphics{foo}| 实际插入 \File{.eps} 文件之前，显式将当前颜色重置为 \Color{black}。
%   \end{itemize}
% \item 变更：
%   \begin{itemize}
%   \item |\colorbox| 和 |\fcolorbox| 改为健壮命令；
%   \item 移除过时的宏包选项 \oOption{pst}；
%   \item 内部宏多处修改。
%   \end{itemize}
% \item 错误修复：
%   \begin{itemize}
%   \item \Model{cmyk} 类型当前颜色 `.'. 的处理不正确。
%   \end{itemize}
% \end{itemize}
%
% \subsubsection*{2005/12/21 v2.09}
% \begin{itemize}
% \item 新功能：
%   \begin{itemize}
%   \item |\definecolor| 和 |\color| 现接受空格分隔的颜色规范，例如 |\color| |[rgb]{1 .5 0}|；
%   \item 实验性 \Option{xcdraw} 选项扩展至 \Option{pdftex} 和 \Option{dvipdfm} 驱动。
%   \end{itemize}
% \item 变更：
%   \begin{itemize}
%   \item 测试文件 \File{xcolor2.tex} 与 \Package{pstricks} 的最新变更兼容；
%   \item 测试文件 \File{xcolor3.tex} 扩展；
%   \item 驱动测试文件 \File{xcolor4.tex} 扩展，演示不同的框线绘制方式；
%   \item 驱动相关代码实现更高效。
%   \end{itemize}
% \end{itemize}
%
% \subsubsection*{2005/11/25 v2.08}
% \begin{itemize}
% \item 新功能：
%   \begin{itemize}
%   \item |\fcolorbox| 参数更灵活，例如 |\fcolorbox| |[gray]{0.5}[wave]{580}{test}|；
%   \item |\boxframe| 返回给定尺寸的框；
%   \item |\f(rame)box| 和 |\fcolorbox| 的新实现（扩展错误报告 latex/3655），以减少输出设备上的像素定位误差；
%   \item \Option{kernelfbox} 选项供偏好旧版 |\f(rame)box| 方式的用户使用；
%   \item 实验性 \Option{xcdraw} 选项在 \Option{dvips} 情况下使用 \Index{PostScript} 命令绘制框和颜色框。
%   \end{itemize}
% \item 错误修复：
%   \begin{itemize}
%   \item |\colorlet| 内表达式类型检测不足；
%   \item 负整数的单位区间归约计算错误（影响颜色系列和扩展颜色表达式）。
%   \end{itemize}
% \end{itemize}
%
% \subsubsection*{2005/11/12 v2.07}
% \begin{itemize}
% \item 新功能：
%   \begin{itemize}
%   \item \Model{Hsb} 颜色模型允许以度为单位指定 \emph{hue}；
%   \item \Model{tHsb}（\pModel{tuned} \Model{Hsb}）颜色模型，用于在色轮上用户自定义 \emph{hue} 配置；
%   \item 轻松生成派生自 \Model{Hsb} 或 \Model{tHsb} 色轮的色彩和谐，例如 |\color{red>wheel,1,12}| 在 12 辐色轮上生成 \Color{red} 的"类似色"；
%   \item 根据 \File{rgb.txt} 新增 317 个预定义颜色名，该文件是 \Index{Unix}/\Index{X11} 发行版的一部分；
%   \item \Option{svgnames} 选项从 \File{rgb.txt} 扩展了 4 个颜色；
%   \item 增强的即时转换语法，例如 |\definecolor| |{foo}{rgb:gray}{0.3}| 或 |\color| |[rgb:wave]{478}|；
%   \item |\@ifundefinedcolor| 和 |\@ifundefinedmodel| 命令；
%   \end{itemize}
% \item 变更：
%   \begin{itemize}
%   \item 文档增强；
%   \item 内部宏多处修改。
%   \end{itemize}
% \item 错误修复：
%   \begin{itemize}
%   \item 负步进参数某些情况下颜色系列分量计算错误。
%   \end{itemize}
% \end{itemize}
%
% \subsubsection*{2005/10/15 v2.06}
% \begin{itemize}
% \item 新功能：
%   \begin{itemize}
%   \item \Model{wave} 颜色模型，用于（近似）可视化光波长，仍 somewhat experimental；
%   \item 伪模型 \nModel{ps}，用于配合 \Package{pstricks} 和 \Option{dvips} 以字面 \Index{PostScript} 代码定义颜色；\href{xcolor2.pdf}{\File{xcolor2.tex}} 中给出了 $\gamma$ 校正方法的示例；
%   \item |\substitutecolormodel| 命令，用于替换缺失或有缺陷的驱动特定颜色模型；
%   \item 改进驱动特定颜色模型的检测和处理；
%   \item \Option{dvipdfmx} 和 \Option{xetex} 选项以支持这些驱动；
%   \item 通用驱动测试文件 \File{xcolor4.tex}。
%   \end{itemize}
% \item 变更：
%   \begin{itemize}
%   \item |\XC@strip@comma| 不再生成尾随空格，这也改善了 |testcolors| 环境的输出。
%   \end{itemize}
% \end{itemize}
%
% \subsubsection*{2005/09/30 v2.05}
% \begin{itemize}
% \item 新功能：
%   \begin{itemize}
%   \item |testcolors| 环境帮助测试不同模型中的颜色，同时显示视觉结果和模型特定参数；
%   \item |\extractcolorspecs| 将模型/颜色规范放入两个独立命令，与 |\extractcolorspec| 相对；
%   \item 预定义颜色集合中添加颜色名 \Color{pink} 和 \Color{olive}。
%   \end{itemize}
% \item 错误修复：
%   \begin{itemize}
%   \item |\definecolor{foo}{named}{bar}| 在 v2.04 中不工作。
%   \end{itemize}
% \end{itemize}
%
% \subsubsection*{2005/09/23 v2.04}
% \begin{itemize}
% \item 新功能：
%   \begin{itemize}
%   \item 为使用额外的——驱动提供的——颜色模型做准备；
%   \item \Package{pstricks} 用户现在可以在 |\psset| 及相关命令中指定显式颜色参数，例如 |\psset{linecolor=[rgb]{1,0,0}}|；示例见 \href{xcolor2.pdf}{\File{xcolor2.tex}}。
%   \end{itemize}
% \item 变更：
%   \begin{itemize}
%   \item 整个宏包中颜色模型名被规范化（即转为 catcode~12）；
%   \item |\@namelet| 命令因与 \Package{memoir} 名称冲突而弃用——请改用 |\XC@let@cc|（还有更多 |\XC@let@..| 命令可用）；
%   \item 使用新的 |\XC@ifxcase| 命令简化颜色转换代码；
%   \item 内部宏一些小改动。
%   \end{itemize}
% \end{itemize}
%
% \subsubsection*{2005/06/06 v2.03}
% \begin{itemize}
% \item 新功能：
%   \begin{itemize}
%   \item \Option{fixpdftex} 选项加载 \Package{pdfcolmk} 宏包，以改善 pdf\TeX 分页时的颜色行为。
%   \end{itemize}
% \item 变更：
%   \begin{itemize}
%   \item 内部宏一些小改动。
%   \end{itemize}
% \item 错误修复：
%   \begin{itemize}
%   \item 由于 |\XC@info| 内不正确的 |\if| 语句，|\colorlet| 的第二个参数以两个相同字母开头时会出问题，例如 |\colorlet{rab}{oof}|；
%   \item |\XC@getcolor| 的参数处理导致与 \Package{msc} 宏包不兼容；
%   \item \Option{prologue} 选项导致与 \Package{preview} 宏包不兼容。
%   \end{itemize}
% \end{itemize}
%
% \subsubsection*{2005/03/24 v2.02}
% \begin{itemize}
% \item 新功能：
%   \begin{itemize}
%   \item |\aftergroupedef| 命令，重现 v2.01 之前 |\aftergroupdef| 的行为；
%   \item \XC 主页 \XCW 现也提供可直接运行的 TDS 兼容归档，包含所有必需文件。
%   \end{itemize}
% \item 变更：
%   \begin{itemize}
%   \item |\rowcolors| 及相关命令仅由 \Option{table} 选项启用；
%   \item |\@ifxempty| 改回 v2.00 更健壮的版本。
%   \end{itemize}
% \item 错误修复：
%   \begin{itemize}
%   \item |\psset{linecolor=\ifcase\foo| |red\or green\or blue\fi}| 在 \Package{pstricks} 下不工作（v2.01 引入的错误）。
%   \end{itemize}
% \end{itemize}
%
% \subsubsection*{2005/03/15 v2.01}
% \begin{itemize}
% \item 新功能：
%   \begin{itemize}
%   \item \Option{prologue} 选项，配合 \Option{dvips} 提供全面的 `named' 颜色支持：即时生成包含所有颜色定义的 \Index{PostScript} 序言文件，供 \Program{dvips} 插入和/或用设备特定参数后处理（例如专色）；
%   \item \Program{dvips} 序言文件 \File{xcolor.pro}，支持额外的 `named' 颜色；
%   \item |\colorlet| 现在也可用于从任意颜色表达式创建命名颜色；
%   \item 增强的颜色定义语法，允许指定目标模型特定颜色参数，例如 |\definecolor| |{red}{rgb/cmyk}{1,0,0/0,1,1,0}|，便于为显示器和打印机分别定制颜色；
%   \item 颜色的"延迟定义"：|\preparecolor| 和 |\definecolors| 支持颜色规范与控制序列生成解耦，对于大型颜色列表（其中仅少数名称实际被使用）特别有用（即节省内存）；
%   \item \Option{dvipsnames*} 和 \Option{svgnames*} 选项支持延迟定义。
%   \end{itemize}
% \item 变更：
%   \begin{itemize}
%   \item 更高精度：大多数补色计算现在对所有 5 位小数都是精确的；
%   \item |\rangeRGB| 及类似变量现在可在文档任意位置修改；
%   \item |\aftergroupdef| 现在仅对其代码参数做一级展开；
%   \item |\XCfileversion| 及类似内部常量从 \File{.sty} 和 \File{.def} 文件中移除；
%   \item 改进内存管理（减少 |\@ifundefined| 测试产生的"多字母控制序列"）；
%   \item 若干内部宏改进和/或重命名。
%   \end{itemize}
% \item 错误修复：
%   \begin{itemize}
%   \item 在 |pspicture| 环境（\Package{pstricks}）内使用 |\psset| 时，|\XC@getcolor| 可能产生多余空格；
%   \item 颜色参数中小数位数过多时可能发生算术溢出，例如 \Package{fp} 计算的结果。
%   \end{itemize}
% \end{itemize}
%
% \subsubsection*{2004/07/04 v2.00}
% \begin{itemize}
% \item 新功能：
%   \begin{itemize}
%   \item 扩展颜色表达式功能：像画家一样混合颜色；
%   \item 支持颜色混合：指定颜色混合表达式，与每个显示的颜色混合；
%   \item |\xglobal| 命令，用于选择性控制颜色定义、混合和蒙版的全局性；
%   \item 颜色系列中的多步操作（例如 |\color{foo!!+++}|）和访问单个成员（例如 |\color{foo!![7]}|）；
%   \item |\providecolor| 命令，仅定义不存在的颜色；
%   \item |\definecolorset| 和 |\providecolorset| 命令，便于构建具有共同底层颜色模型的颜色集合；
%   \item 根据 \Index{SVG}~1.1 规范新增 147 个预定义颜色名；
%   \item \Key{xpdfborder} 键，在使用 \Program{dvips} 时以更驱动无关的方式设置超链接边框宽度。
%   \end{itemize}
% \item 变更：
%   \begin{itemize}
%   \item \LCP 现已完全整合进 \XC；
%   \item \oOption{override}、\oOption{usenames}、\oOption{nodvipsnames} 选项和 |\xdefinecolor| 命令不再需要；
%   \item \Option{dvips} 和 \Option{dvipsnames} 选项现在相互独立；
%   \item |\tracingcolors| 的行为改为更通用，并在标准情况下减小日志文件大小；
%   \item |\rdivide| 的语法更灵活（可除以数字和/或维度）；
%   \item 代码重构，部分内部命令重命名；
%   \item 文档重新编排并增强。
%   \end{itemize}
% \item 错误修复：
%   \begin{itemize}
%   \item |\definecolor{foo}{named}{bar}| 不工作（v1.11 引入的错误）；
%   \item \Package{pstricks} 键值内条件命令更健壮。
%   \end{itemize}
% \end{itemize}
%
% \subsubsection*{2004/05/09 v1.11}
% \begin{itemize}
% \item 新功能：
%   \begin{itemize}
%   \item 开关 |\ifglobalcolors|，控制颜色定义是全局还是局部；
%   \item \Option{hyperref} 选项为超链接边框颜色提供颜色表达式支持，例如 |\hypersetup| |{xurlbordercolor=red!50!yellow}|；
%   \item 内部钩子 |\XC@bcolor|、|\XC@mcolor| 和 |\XC@ecolor|，用于在当前颜色显示前后立即执行额外代码。
%   \end{itemize}
% \item 变更：
%   \begin{itemize}
%   \item |\XC@logcolor| 重命名为 |\XC@display|，现为核心颜色显示命令；
%   \item 改进与 \Package{pstricks} 的接口。
%   \end{itemize}
% \end{itemize}
%
% \subsubsection*{2004/03/27 v1.10}
% \begin{itemize}
% \item 新功能：
%   \begin{itemize}
%   \item 支持 \nModel{named} 模型；
%   \item 支持 \Program{dvips} 颜色（现可在颜色表达式中使用）；
%   \item `普通' 和 `named' 颜色的内部表示合并为统一数据结构；
%   \item 允许颜色表达式开头出现多个 `-'.
%   \end{itemize}
% \item 错误修复：
%   \begin{itemize}
%   \item 颜色蒙版或目标模型转换激活时，|\color[named]{foo}| 等命令会报错；
%   \item 与 \Package{soul} 宏包不兼容：|\hl|、|\ul| 等命令可能产生意外结果。
%   \end{itemize}
% \item 文档：
%   \begin{itemize}
%   \item 添加一般颜色表达式的公式；
%   \item 增强文本和索引；
%   \item 移除索引生成对本地配置文件的依赖。
%   \end{itemize}
% \end{itemize}
%
% \subsubsection*{2004/02/16 v1.09}
% \begin{itemize}
% \item 新功能：
%   \begin{itemize}
%   \item \Model{HTML} 颜色模型，24 位十六进制 \Model{RGB} 变体；允许指定如 |\color[HTML]{AFFE90}| 的颜色；
%   \item 预定义颜色集合中添加颜色名 \Color{orange}、\Color{violet}、\Color{purple} 和 \Color{brown}。
%   \end{itemize}
% \item \XC 新主页：\XCW
% \item 错误修复：|\xdefinecolor| 有时不对其参数规范化。
% \item 变更：
%   \begin{itemize}
%   \item 文档小幅改进；
%   \item 示例文件 |xcolor1.tex| 重新组织并精简。
%   \end{itemize}
% \end{itemize}
%
% \subsubsection*{2004/02/04 v1.08}
% \begin{itemize}
% \item 新命令：
%   \begin{itemize}
%   \item |\selectcolormodel| 在文档内更改目标模型；
%   \item |\adjustUCRBG| 在转换到 \Model{cmyk} 时微调底色去除和黑版生成。
%   \end{itemize}
% \item 错误修复：颜色表达式与活动 `!' 字符配合使用时不正确，例如在 |\usepackage[frenchb]babel}| 的情况下。
% \item 代码重组：
%   \begin{itemize}
%   \item |\XC@xdefinecolor| 合并入 |\xdefinecolor|，使前者过时；
%   \item 若干内部命令改进/精简。
%   \end{itemize}
% \end{itemize}
%
% \subsubsection*{2004/01/20 v1.07}
% \begin{itemize}
% \item 新功能：支持颜色蒙版和颜色分色。
% \item 新命令：
%   \begin{itemize}
%   \item |\rmultiply| 将维度寄存器乘以实数；
%   \item |\xcolorcmd| 传递宏包末尾要执行的命令。
%   \end{itemize}
% \item 变更：
%   \begin{itemize}
%   \item 更一致的颜色处理：扩展颜色现在始终优先于标准颜色；
%   \item 若干命令通过使用 \LaTeX 内核代码改进。
%   \end{itemize}
% \item 文档：一些小改动。
% \item 示例文件：额外 \Package{pstricks} 示例（文件 \href{xcolor2.pdf}{\File{xcolor2.tex}}）。
% \end{itemize}
%
% \subsubsection*{2003/12/15 v1.06}
% \begin{itemize}
% \item 新功能：扩展颜色表达式，允许级联混合操作，例如 |\color{red!30!green!40!blue}|。
% \item 文档：新增颜色表达式章节。
% \item 错误修复：~在不显示的命令如 |\extractcolorspec{foo!!+}| 中，颜色系列步进不正确（此错误在 v1.05 引入）。
% \item 重命名命令：~|\ukfileversion| 及类似内部常量重命名为 |\XCfileversion| 等。
% \item 移除命令：~|\ifXCpst| 和 |\ifXCtable| 通过简单技巧变得过时。
% \end{itemize}
%
% \subsubsection*{2003/11/21 v1.05}
% \begin{itemize}
% \item 错误修复：
%   \begin{itemize}
%   \item 宏包选项 \Option{hideerrors} 现在应按预期工作；
%   \item 文档中第一个颜色表达式使用 `.' 因初始化不正确而报错。
%   \end{itemize}
% \item 代码重组：|\extractcolorspec| 现使用 |\XC@splitcolor|，使 |\XC@extract| 过时。
% \end{itemize}
%
% \subsubsection*{2003/11/09 v1.04}
% \begin{itemize}
% \item 新功能：在颜色表达式中轻松访问当前颜色。
% \item 新选项：\oOption{override}，用 |\xdefinecolor| 替换 |\definecolor|。
% \item 新命令：|\tracingcolors|，用于记录颜色相关信息。
% \end{itemize}
%
% \subsubsection*{2003/09/21 v1.03}
% \begin{itemize}
% \item 变更：绕过某些驱动的奇怪行为。
% \item 新功能：与 \Package{hyperref} 共享驱动。
% \end{itemize}
%
% \subsubsection*{2003/09/19 v1.02}
% \begin{itemize}
% \item 变更：|\extractcolorspec| 和 |\colorlet| 现也接受颜色系列作为参数。
% \end{itemize}
%
% \subsubsection*{2003/09/15 v1.01}
% \begin{itemize}
% \item 新功能：|\definecolorseries| 及相关命令。
% \item 文档：移除一些 \Package{doc} 相关副作用。
% \item 代码重组：所有计算相关工具集中到一处。
% \item 错误修复：
%   \begin{itemize}
%   \item |\@rdivide|：添加 |\relax| 修复负分子问题；
%   \item |\rowc@l@rs|：用 |\@ifxempty| 替换 |\@ifempty|。
%   \end{itemize}
% \end{itemize}
%
% \subsubsection*{2003/09/09 v1.00}
% \begin{itemize}
% \item 首次公开发布。
% \end{itemize}
%
% \end{multicols}
% \end{adjustwidth}
%
% \StopEventually{}
% \clearpage
%
% \section{The Implementation}
%
% Some commands have been copied unchanged from \cite{color} (Copyright (C) 1994--1999 \People{David P.}{Carlisle}), which is part of the Standard \LaTeX{} `Graphics Bundle'.
% Additionally, also some parts of the code documentation were copied from \cite{color}.
%
%
%
% \subsection{Preparation}
%    \begin{macrocode}
%<*package>
%    \end{macrocode}
%
% \subsubsection{Release handling}
%
%    \begin{macrocode}
\providecommand\DeclareRelease[3]{}
\providecommand\DeclareCurrentRelease[2]{}
\DeclareRelease{v2}{0000-00-00}{xcolor-2022-06-12.sty}
\DeclareCurrentRelease{}{2023-11-11}
%    \end{macrocode}
%
% \subsubsection{File identification}
%
% \renewcommand*\MacroFont{\footnotesize\ttfamily}
% \renewcommand*\AltMacroFont{\footnotesize\ttfamily\slshape}
%
%
% \fbox{\MacroFont\string\ProvidesPackage\string{xcolor\string}[\XCfiledate\space\XCfileversion\space LaTeX color extensions (UK)]}
%
%    \begin{macrocode}
\NeedsTeXFormat{LaTeX2e}
\providecommand\IfFormatAtLeastTF{\@ifl@t@r\fmtversion}
\IfFormatAtLeastTF{2022/06/01}{}
 {\RequirePackageWithOptions{xcolor-2022-06-12}}
\IfFormatAtLeastTF{2022/06/01}{}{\endinput}
%    \end{macrocode}
%
%
% \subsubsection{Safety first}
%
% \begin{macro}{\XC@catcodes}
% Save catcodes for later restore:
%    \begin{macrocode}
\expandafter\edef\csname XC@catcodes\endcsname
 {\catcode33 \the\catcode33 % !
  \catcode34 \the\catcode34 % "
  \catcode39 \the\catcode39 % '
  \catcode42 \the\catcode42 % *
  \catcode44 \the\catcode44 % ,
  \catcode45 \the\catcode45 % -
  \catcode46 \the\catcode46 % .
  \catcode47 \the\catcode47 % /
  \catcode58 \the\catcode58 % :
  \catcode59 \the\catcode59 % ;
  \catcode60 \the\catcode60 % <
  \catcode61 \the\catcode61 % =
  \catcode62 \the\catcode62 % >
  \catcode63 \the\catcode63 % ?
  \catcode64 \the\catcode64 % @
  \catcode94 \the\catcode94 % ^
  \catcode96 \the\catcode96 % `
  \catcode`\noexpand\^^A\the\catcode`\^^A\relax}
%    \end{macrocode}
% \end{macro}
%
%    \begin{macrocode}
\catcode64 11 %
\@makeother\! \@makeother\" \@makeother\' \@makeother\* \@makeother\,
\@makeother\- \@makeother\. \@makeother\/ \@makeother\: \@makeother\;
\@makeother\< \@makeother\= \@makeother\> \@makeother\? \@makeother\`
\catcode94 7 %
\catcode`\^^A=14 %
%    \end{macrocode}
%
%
% \subsubsection{Constants}
%
%  \begin{macro}{\XC@@}
% Macro consisting of the macro |\\XC@@|, serves as a delimiter.
%    \begin{macrocode}
\edef\XC@@{\expandafter\noexpand\csname\@backslashchar XC@@\endcsname}
%    \end{macrocode}
%  \end{macro}
%
%
% \subsubsection{Switches}
%
%  \begin{macro}{\ifcolors@}
% This boolean can be tested by higher level macros that may want to alter their behaviour if a monochrome driver is being used.
%    \begin{macrocode}
\newif\ifcolors@ \colors@true
%    \end{macrocode}
%  \end{macro}
%
%  \begin{macro}{\ifcolorerrors@}
% Controls whether missing color models/names result in error or warning messages.
%    \begin{macrocode}
\newif\ifcolorerrors@
%    \end{macrocode}
%  \end{macro}
%
% \begin{macro}{\ifglobalcolors}
% Controls whether |\definecolor| and |\colorlet| generate global or local color definitions.
%    \begin{macrocode}
\newif\ifglobalcolors \globalcolorsfalse
%    \end{macrocode}
% \end{macro}
%
% \begin{macro}{\ifdefinecolors}
% Controls whether |\definecolorset| should \emph{define} or \emph{prepare} colors.
%    \begin{macrocode}
\newif\ifdefinecolors \definecolorstrue
%    \end{macrocode}
% \end{macro}
%
% \begin{macro}{\ifconvertcolorsD}
% \begin{macro}{\ifconvertcolorsU}
% Control whether conversion to target models take place at \emph{definition} time (|\ifconvertcolorsD|) and/or \emph{usage} (|\ifconvertcolorsU|).
%    \begin{macrocode}
\newif\ifconvertcolorsD
\newif\ifconvertcolorsU
%    \end{macrocode}
% \end{macro}
% \end{macro}
%
% \begin{macro}{\ifblendcolors}
% \begin{macro}{\ifmaskcolors}
% Control whether colors are to be blended and/or masked at display time.
%    \begin{macrocode}
\newif\ifblendcolors
\newif\ifmaskcolors
%    \end{macrocode}
% \end{macro}
% \end{macro}
%
% \begin{macro}{\ifXC@@fbox}
% \begin{macro}{\ifXC@@tst}
% Auxiliary test variables.
%    \begin{macrocode}
\newif\ifXC@@fbox \XC@@fboxtrue
\newif\ifXC@@tst
%    \end{macrocode}
% \end{macro}
% \end{macro}
%
%
% \subsubsection{Error handling}
%
% \begin{macro}{\c@lor@error}
% Multi-purpose error handler (issues warnings instead of errors if the \Option{hideerrors} option is used).
% This command is used already in the options section.
%    \begin{macrocode}
\def\c@lor@error#1%
 {\ifcolorerrors@\PackageError{xcolor}{Undefined color #1}\@ehc
  \else\PackageWarning{xcolor}{Undefined color #1}\fi}
%    \end{macrocode}
% \end{macro}
%
%
% \subsubsection{Macro concatenation}
%
% \begin{macro}{\XC@append}
%   \marg{cmd}\marg{token list}\\
% Append \Meta{token list} to the definition of \Meta{cmd}.
% This command is used already in the options section.
%    \begin{macrocode}
\def\XC@append#1#2%
 {\ifx#1\@undefined\def#1{#2}\else\ifx#1\relax\def#1{#2}\else
  \toks@\expandafter{#1#2}\edef#1{\the\toks@}\fi\fi}
%    \end{macrocode}
% \end{macro}
%
%
% \subsubsection{Macro definition}
%
% The following |\XC@let@..| commands mimic \LaTeX3's coding syntax for the type of arguments used:
% the `c' type is for brace delimited arguments which turn into a control sequence inside |\csname ... \endcsname|, whereas `N' is just a single token.
% Thanks to \People{Morten}{H\o gholm} for pointing me to that.
%
% \begin{macro}{\XC@let@cc}
%   \marg{target character string}\marg{source character string}\\
% Performs |\let\|\Meta{target character string}|=\|\Meta{source character string}.
%    \begin{macrocode}
\def\XC@let@cc#1{\expandafter\XC@let@Nc\csname#1\endcsname}
%    \end{macrocode}
% \end{macro}
%
% \begin{macro}{\@namelet}
% This is provided for compatibility reasons only, although deprecated because of a name clash with \Package{memoir}.
%    \begin{macrocode}
\providecommand*\@namelet[1]{\expandafter\XC@let@Nc\csname#1\endcsname}
%    \end{macrocode}
% \end{macro}
%
% \begin{macro}{\XC@let@Nc}
%   \marg{target name}\marg{source character string}\\
% Performs |\let|\Meta{target name}|=\|\Meta{source character string}.
% Formerly called |\@@namelet|.
%    \begin{macrocode}
\def\XC@let@Nc#1#2{\expandafter\let\expandafter#1\csname#2\endcsname}
%    \end{macrocode}
% \end{macro}
%
% \begin{macro}{\XC@let@cN}
%   \marg{target character string}\marg{source name}\\
% Performs |\let\|\Meta{target character string}|=|\Meta{source name}.
%    \begin{macrocode}
\def\XC@let@cN#1{\expandafter\let\csname#1\endcsname}
%    \end{macrocode}
% \end{macro}
%
% \begin{macro}{\@namexdef}
%   \marg{character string}\marg{code}\\
% Performs |\xdef\|\Meta{character string}|{|\Meta{code}|}|.
% Similar to the \LaTeX{} internal |\@namedef| command.
%    \begin{macrocode}
\def\@namexdef#1{\expandafter\xdef\csname #1\endcsname}
%    \end{macrocode}
% \end{macro}
%
% \begin{macro}{\aftergroupdef}
% \begin{macro}{\aftergroupedef}
%   \marg{cmd}\marg{code}\\
% Defines \meta{cmd} with the first-level or complete expansion of \meta{code} after the end of the current group.
%    \begin{macrocode}
\def\aftergroupdef#1#2%
 {\expandafter\endgroup\expandafter\def\expandafter#1\expandafter{#2}}
\def\aftergroupedef#1#2%
 {\edef\@@tmp{\def\noexpand#1{#2}}\expandafter\endgroup\@@tmp}
%    \end{macrocode}
% \end{macro}
% \end{macro}
%
% \begin{macro}{\XC@edef}
% \begin{macro}{\XC@mdef}
% \begin{macro}{\XC@sdef}
%   \marg{cmd}\marg{code}\\
% Defines \Meta{cmd} as a safe expansion of \Meta{code}; needed for color expressions in certain languages.
% The second form is a streamlined version that takes only care of active `|/|' and `|:|' characters that may occur in \Meta{model-list} and \Meta{spec-list} arguments.
% The third form defines a sanitized expansion without considering active characters.
% It is faster since it needs no grouping operation.
%    \begin{macrocode}
\begingroup
\catcode`\!=13 \catcode`\:=13 \catcode`\-=13 \catcode`\+=13
\catcode`\;=13 \catcode`\/=13 \catcode`\"=13 \catcode`\>=13
\gdef\XC@edef#1#2%
 {\begingroup
  \ifnum\catcode`\!=13 \edef!{\string!}\fi
  \ifnum\catcode`\:=13 \edef:{\string:}\fi
  \ifnum\catcode`\-=13 \edef-{\string-}\fi
  \ifnum\catcode`\+=13 \edef+{\string+}\fi
  \ifnum\catcode`\;=13 \edef;{\string;}\fi
  \ifnum\catcode`\"=13 \edef"{\string"}\fi
  \ifnum\catcode`\>=13 \edef>{\string>}\fi
  \edef#1{#2}\@onelevel@sanitize#1\aftergroupdef#1#1}
\gdef\XC@mdef#1#2%
 {\begingroup
  \ifnum\catcode`\/=13 \edef/{\string/}\fi
  \ifnum\catcode`\:=13 \edef:{\string:}\fi
  \edef#1{#2}\@onelevel@sanitize#1\aftergroupdef#1#1}
\endgroup
\def\XC@sdef#1#2{\edef#1{#2}\@onelevel@sanitize#1}
%    \end{macrocode}
% \end{macro}
% \end{macro}
% \end{macro}
%
%
% \subsubsection{Conditional processing}
%
% \begin{macro}{\@ifxempty}
%   \marg{arg}\marg{empty}\marg{non-empty}\\
% Tests without expanding, whether the argument \Meta{arg} is empty and executes the following code accordingly; \Meta{arg} must not start with the token |\XC@@|.
% Can also be used within |\edef|.
%    \begin{macrocode}
\def\@ifxempty#1{\@@ifxempty#1\@@ifxempty\XC@@}
\def\@@ifxempty#1#2\XC@@
 {\ifx#1\@@ifxempty
  \expandafter\@firstoftwo\else\expandafter\@secondoftwo\fi}
%    \end{macrocode}
% \end{macro}
%
% \begin{macro}{\XC@ifxcase}
%   \marg{cmd}|{{|\Meta[1]{cmd}|}{|\Meta[1]{code}|}...{|\Meta[n]{cmd}|}{|\Meta[n]{code}|}}|\marg{default code}\\
% Performs the test |\ifx|\Meta{cmd}\Meta[k]{cmd} for $k=1,\dots,n$.
% If for some $k$ the test is `true', then the attached code \Meta[k]{code} is executed.
% If all tests are `false', then \Meta{default code} is executed.
%    \begin{macrocode}
\long\def\XC@ifxcase#1#2#3%
 {\long\def\XC@if@##1##2%
   {\ifx#1##1%
      \toks@{##2}\expandafter\remove@to@nnil\else\expandafter\XC@if@\fi}%
  \XC@if@#2#1{#3}\@nnil\the\toks@}
%    \end{macrocode}
% \end{macro}
%
% \begin{macro}{\XC@ifcase}
%   \marg{str}|{{|\Meta[1]{str-list}|}{|\Meta[1]{code}|}...{|\Meta[n]{str-list}|}{|\Meta[n]{code}|}}|\marg{default code}\\
% Tests whether \Meta{str} appears in one of the comma-separated lists \Meta[k]{str-list}.
% Performs the test |\if|\Meta{str}\Meta[kl]{str} for $k=1,\dots,n$ and $l=1,\ldots,m_n$.
% If for some $k,l$ the test is `true', then the attached code \Meta[k]{code} is executed.
% If all tests are `false', then \Meta{default code} is executed.
% \changes{3.0}{2022-08-10}{moved the patch from hycolor to correct Gray model, issue \#23}
%    \begin{macrocode}
\long\def\XC@ifcase#1#2#3%
 {\long\def\XC@if@##1##2%
   {\@expandtwoargs\in@{,#1,}{,##1,}\ifin@
      \toks@{##2}\expandafter\remove@to@nnil\else\expandafter\XC@if@\fi}%
  \XC@if@#2{#1}{#3}\@nnil\the\toks@}
%    \end{macrocode}
% \end{macro}
%
%
%
% \subsubsection{Color model selection}
%
% \begin{macro}{\selectcolormodel}
%   \marg{num model}\\
% Determines the target model for color definitions with |\definecolor| and friends.
%    \begin{macrocode}
\def\selectcolormodel#1%
 {\@ifundefined{XC@mod@#1}%
   {\c@lor@error{target model `#1'}}%
   {\expandafter\ifx\csname XC@mod@#1\endcsname\XC@mod@natural
      \def\XC@tgt@mod##1{##1}\convertcolorsDfalse\convertcolorsUfalse
    \else\expandafter\ifx\csname XC@mod@#1\endcsname\XC@mod@named
      \c@lor@error{target model `#1'}
    \else
      \edef\XC@tgt@mod##1{#1}\convertcolorsDtrue\convertcolorsUtrue
    \fi\fi}}
%    \end{macrocode}
% \end{macro}
%
% \begin{macro}{\XC@mod@Hsb}
% \begin{macro}{\XC@mod@tHsb}
% \begin{macro}{\XC@mod@ignore}
% \begin{macro}{\XC@mod@named}
% \begin{macro}{\XC@mod@ps}
% \begin{macro}{\XC@mod@wave}
% \begin{macro}{\XC@mod@natural}
% \begin{macro}{\XC@mod@rgb}
% \begin{macro}{\XC@mod@cmy}
% \begin{macro}{\XC@mod@cmyk}
% \begin{macro}{\XC@mod@hsb}
% \begin{macro}{\XC@mod@gray}
% \begin{macro}{\XC@mod@RGB}
% \begin{macro}{\XC@mod@HTML}
% \begin{macro}{\XC@mod@HSB}
% \begin{macro}{\XC@mod@Gray}
% Put model names into commands and define package options
% \Option{natural}, \Option{rgb}, \Option{cmy}, \Option{cmyk}, \Option{hsb}, \Option{gray}, \Option{RGB}, \Option{HTML}, \Option{HSB}, \Option{Gray}.
%    \begin{macrocode}
\@tfor\@@tmp:={Hsb}{tHsb}{ignore}{named}{ps}{wave}\do
 {\expandafter\XC@sdef\csname XC@mod@\@@tmp\endcsname{\@@tmp}}
%    \end{macrocode}
%
%    \begin{macrocode}
\@tfor\@@tmp:={natural}{rgb}{cmy}{cmyk}{hsb}{gray}{RGB}{HTML}{HSB}{Gray}\do
 {\expandafter\XC@sdef\csname XC@mod@\@@tmp\endcsname{\@@tmp}%
  \edef\@@tmp
   {\noexpand\DeclareKeys{\@@tmp .code= {\noexpand\selectcolormodel{\@@tmp}}}}%
  \@@tmp}
%    \end{macrocode}
% \end{macro}
% \end{macro}
% \end{macro}
% \end{macro}
% \end{macro}
% \end{macro}
% \end{macro}
% \end{macro}
% \end{macro}
% \end{macro}
% \end{macro}
% \end{macro}
% \end{macro}
% \end{macro}
% \end{macro}
% \end{macro}
%
% The \Option{monochrome} option should be used \emph{in addition} to one of the standard driver options.
%    \begin{macrocode}
\DeclareKeys
 {
    monochrome .code =
    {
      \colorerrors@false
      \AtEndOfPackage
       {\colors@false
        \def\set@color{\c@lor@special\m@ne
          {color push \current@color}\aftergroup\reset@color}%
        \def\reset@color{\c@lor@special\m@ne{color pop}}%
        \def\set@page@color{\c@lor@special\m@ne{background \current@color}}}
    }
   ,monochrome .usage = load
 }
%    \end{macrocode}
%
%
% \subsubsection{Sets of predefined colors}
%
% \begin{macro}{\XC@@names}
% Initialise the macro to hold the names of files that contain predefined colors; these files are loaded near to the end of this package.
%    \begin{macrocode}
\let\XC@@names\relax
%    \end{macrocode}
% \end{macro}
%
% \begin{macro}{\colornameprefix}
% \begin{macro}{\XC@@nameprefix}
% Initialise the macro to hold the string that is prefixed to driver-specific color names.
% Save the value for later restore.
%    \begin{macrocode}
\providecommand*\colornameprefix{XC@}
\let\XC@@nameprefix\colornameprefix
%    \end{macrocode}
% \end{macro}
% \end{macro}
%
% \begin{macro}{\XC@declarenames}
%   \marg{base name}\marg{file name}\marg{initial commands}\\
% Construct package options for color names.
%    \begin{macrocode}
\def\XC@declarenames#1#2#3%
 {%
  \DeclareKeys
   {
    #1  .code =
       {
        \AddToHook{package/xcolor/after}
          {#3\definecolorstrue\input{#2.def}\definecolorstrue
           \let\colornameprefix\XC@@nameprefix}%
        \DeclareKeys[xcolor]{#1 .code = {},#1* .code={}}
       }
   ,#1  .usage = preamble
   ,#1* .code =
       {
         \AddToHook{package/xcolor/after}
           {#3\definecolorsfalse\input{#2.def}\definecolorstrue
            \let\colornameprefix\XC@@nameprefix}
         \DeclareKeys[xcolor]{#1 .code = {},#1* .code={}}
       }
   ,#1* .usage = preamble
   }%
 }
%    \end{macrocode}
% \end{macro}
%
%    \begin{macrocode}
\XC@declarenames{dvipsnames}{dvipsnam}{\let\colornameprefix\@empty}
\XC@declarenames{svgnames}{svgnam}{}
\XC@declarenames{x11names}{x11nam}{}
\XC@declarenames{hksnames}%
 {hks\ifcase\paperquality k\or n\or k\or z\else n\fi nam}{}
\XC@declarenames{pmsnames}{pmsnam}{}
\XC@declarenames{ralnames}{ralnam}{}
%    \end{macrocode}
%
%
% \subsubsection{Support for other packages}
%
% Use this option to load \People{David P.}{Carlisle}'s \Package{colortbl} package:
%    \begin{macrocode}
\DeclareKeys
 {
    table .code  ={\AddToHook{package/xcolor/after}{\RequirePackage{colortbl}}}
   ,table .usage = preamble
 }
%    \end{macrocode}
%
% \subsubsection{Driver options}
%
% \begin{macro}{\Gin@driver}
% Initialise the macro to hold the driver file name.
%    \begin{macrocode}
\providecommand*\Gin@driver{}
%    \end{macrocode}
%
% Drivers that come with their individual definition files:
%    \begin{macrocode}
\@tfor\@@tmp:={dvipdf}{dvipdfmx}{dvips}{dvipsone}{dviwin}%
              {emtex}{luatex}{pctex32}{pctexhp}{pctexps}{pctexwin}{pdftex}%
              {tcidvi}{textures}{truetex}{vtex}{xetex}{dvisvgm}\do
 {\edef\@@tmp{%
   \noexpand\DeclareKeys
     {
      \@@tmp .code = {\def\noexpand\Gin@driver{\@@tmp.def}},
      \@@tmp .usage=load
     }}%
  \@@tmp}
%    \end{macrocode}
%
% Drivers that are derivatives of others:
%    \begin{macrocode}
\DeclareKeys
 {
  ,dviwindo .code   = {\SetKeys{dvipsone}}
  ,dviwindo .usage  = load
  ,oztex    .code   = {\SetKeys{dvips}}
  ,oztex    .usage  = load
  ,xdvi     .code   = {\SetKeys{dvips,monochrome}}
  ,xdvi     .usage  = load
  ,dvipdfm  .code   = {\SetKeys{dvipdfmx}}
  ,dvipdfm  .usage  = load
 }
%    \end{macrocode}
%
% Options used by (some) drivers.
%    \begin{macrocode}
\DeclareKeys
 {
   ,xcdraw   .code    = {\def\XC@@xcd@{1}}
   ,xcdraw   .usage   = load %TODO check
   ,noxcdraw .code    = {\def\XC@@xcd@{-1}}
   ,noxcdraw .usage   = load %TODO check
   ,fixinclude .code  = {\def\XC@@xcf@{1}}
   ,fixinclude .usage = load %TODO check
   ,prologue .code    = {\def\XC@@xcp@{1}}
   ,prologue .usage   = load %TODO check
%    \end{macrocode}
% \end{macro}
%
%
% \subsubsection{Other options}
%
% Switch to original kernel |\fbox| command instead of new \XC{} version.
%    \begin{macrocode}
   ,kernelfbox .code    = {\XC@@fboxfalse}
   ,kernelfbox .usage   = load %TODO check
%    \end{macrocode}
%
% These options control how non-existing colors are treated:
%    \begin{macrocode}
   ,showerrors .code   = {\colorerrors@true}
   ,hideerrors .code   = {\colorerrors@false}
%    \end{macrocode}
%
% The \Option{debugshow} option turns on debugging info (perhaps).
%    \begin{macrocode}
   ,debugshow .code    = {\catcode`\^^A=9 \let\GDebug\typeout}
  }
%    \end{macrocode}
%
% \subsubsection{Obsolete options}
% The following options are obsolete but perhaps still around in the wild,
% so we generate a warning.
%    \begin{macrocode}
\DeclareKeys
 {
   ,usenames   .code    =
     {\PackageWarning{xcolor}{Package option `usenames' is obsolete and ignored}}
   ,usenames    .usage   = preamble
   ,hyperref   .code    =
     {\PackageWarning{xcolor}{Package option `hyperref' is obsolete and ignored}}
   ,hyperref   .usage   = preamble
   ,fixpdftex  .code    =
     {\PackageWarning{xcolor}{Package option `fixpdftex' is obsolete and ignored}}
   ,fixpdftex  .usage   = preamble
 }
%    \end{macrocode}
%
% All other options (including \LC's \oOption{nodvipsnames}) error.
%
% Choose default options:
%    \begin{macrocode}
\def\XC@@xcd@{-1}
\def\XC@@xcf@{-1}
\def\XC@@xcp@{-1}
\SetKeys{natural,showerrors}
%    \end{macrocode}
%
% \subsubsection{Processing options}
%
% A local configuration file may declare more options.
% It should also make one driver option the default, by calling |\ExecuteOptions| with the appropriate option.
%    \begin{macrocode}
\let\XC@ExecuteOptions\ExecuteOptions
\ExplSyntaxOn
\def\ExecuteOptions{\keys_set_known:nn {xcolor}}
\ExplSyntaxOff
\InputIfFileExists{color.cfg}{}{}
\let\ExecuteOptions\XC@ExecuteOptions
%    \end{macrocode}
%
% After the options are processed, load the appropriate driver file.
% If a site wants a default driver (e.g., \Option{dvips}) it just needs
% to put |\ExecuteOptions{dvips}| in a \File{color.cfg} file.
%
% Note |\relax| added here to ensure that |\ProcessKeyOptions| not the following
% |\edef| shows in error context lines for unknown options.
% \changes{3.02}{2024-03-13}{Added \cs{relax}}
%    \begin{macrocode}
\ProcessKeyOptions\relax
%    \end{macrocode}
%
% \begin{macro}{\XC@@drv}
% Holds the name of the loaded driver (without \File{.def} extension).
%    \begin{macrocode}
\edef\Gin@driver{\Gin@driver}
\if!\Gin@driver!
  \PackageError{xcolor}
    {No driver specified}
    {You should make a default driver option in a file \MessageBreak
     color.cfg\MessageBreak
     eg: \protect\ExecuteOptions{dvips}}
\else
  \PackageInfo{xcolor}{Driver file: \Gin@driver}
  \@ifundefined{ver@\Gin@driver}{\input{\Gin@driver}}{}
  \def\@@tmp#1.def{\def\XC@@drv{#1}}
  \expandafter\@@tmp\Gin@driver
\fi
%    \end{macrocode}
% \end{macro}
%
% \begin{macro}{\XC@@xcd@}
% \begin{macro}{\XC@@xcf@}
% \begin{macro}{\XC@@xcp@}
% These are used firstly by the \Option{xcdraw} and \Option{prologue} options, indicating whether either option is switched on or off.
% Secondly, the macros serve as test variables for |\ifcase| expressions, storing the intended type of output: $-1=\text{dvi}, 0=\text{ps}, 1=\text{pdf}$.
%    \begin{macrocode}
\count@\XC@@xcd@ \advance\count@\XC@@xcf@ \advance\count@\XC@@xcp@\relax
\ifnum\count@>-3
%    \end{macrocode}
%
% \begin{macro}{\XC@special}
%   \meta{code}\\
% Driver-specific variant of |\special|.
% Even if the intended type of output is the same (e.g., pdf), the |\special| commands may differ.
%    \begin{macrocode}
  \XC@ifcase\XC@@drv
   {{dvips,pctex32}%
    {\def\@tempa{0}
     \def\XC@special#1{\special{ps:XC!b}\special{"#1}\special{ps:XC!e}}}%
    {pdftex}%
    {\def\@tempa{1}
     \def\XC@special#1{\pdfliteral{ q #1 Q }}}%
    {luatex}%
    {\def\@tempa{1}
     \ifnum\luatexversion<87%
       \def\XC@special#1{\pdfliteral{ q #1 Q }}%
     \else%
       \def\XC@special#1{\pdfextension literal{ q #1 Q }}%
     \fi}%
    {dvipdfm}%
    {\def\@tempa{1}
     \def\XC@special#1{\special{pdf: content #1}}}%
   }{\def\@tempa{-1}}
%    \end{macrocode}
% \end{macro}
%
% Since the options are independent of each other, we have to adjust the affiliated test variables independently as well.
%    \begin{macrocode}
  \ifnum\XC@@xcd@>\m@ne\let\XC@@xcd@\@tempa\fi
  \ifnum\XC@@xcf@>\m@ne\let\XC@@xcf@\@tempa\fi
  \ifnum\XC@@xcp@>\m@ne\let\XC@@xcp@\@tempa\fi
%    \end{macrocode}
%
%    \begin{macrocode}
  \ifnum\@tempa=0 \AtBeginDvi{\special{header=xcolor.pro}}\fi
  \ifcase\XC@@xcf@
    \AtBeginDvi{\special{! XCfixnormalscale}}
  \fi
  \ifcase\XC@@xcp@
    \AtBeginDvi{\special{header=\jobname.xcp}}
    \newwrite\XC@@xcp\immediate\openout\XC@@xcp=\jobname.xcp
    \immediate\write\XC@@xcp{\@percentchar!}
    \immediate\write\XC@@xcp{TeXDict begin}
    \AtEndDocument
     {\immediate\write\XC@@xcp{end}\immediate\closeout\XC@@xcp}
  \fi
%    \end{macrocode}
%
% End of driver-specific part.
%    \begin{macrocode}
\fi
%    \end{macrocode}
% \end{macro}
% \end{macro}
% \end{macro}
%
%
% \subsubsection{Required packages}
%
% We pretent that color.sty has already been loaded.
%
%    \begin{macrocode}
\@namedef{ver@color.sty}{1999/02/16}
%    \end{macrocode}
%
%
% \subsection{Macro tools}
%
% \subsubsection{General tools}
%
% \begin{macro}{\@secondoffive}
% \begin{macro}{\@thirdoffive}
% Argument grabbing, similar to \LaTeX{} kernel.
%    \begin{macrocode}
\long\def\@secondoffive#1#2#3#4#5{#2}
\long\def\@thirdoffive#1#2#3#4#5{#3}
%    \end{macrocode}
% \end{macro}
% \end{macro}
%
% \begin{macro}{\xglobal}
% \begin{macro}{\xglobal@test}
% If |\foo| occurs in the token list |\xglobal@@|, then the command |\xglobal\foo| will set the switch |\xglobal@true| which can be used inside |\foo| to determine whether certain definitions are to be made explicitly global.
% |\foo| is responsible for resetting |\xglobal@false| in order to avoid side effects.
% To include |\foo| in the magic list, simply say |\XC@append\xglobal@list{\foo}|.
% If |\foo| is not in the list, |\xglobal\foo| will behave like |\global\foo|, thus |\xglobal\let| will be like |\global\let| etc.
%    \begin{macrocode}
\def\xglobal{\xglobal@stop\begingroup\futurelet\@@tok\xglobal@test}
%    \end{macrocode}
%
%    \begin{macrocode}
\def\xglobal@test
 {\expandafter\@tfor\expandafter\@@tmp
   \expandafter:\expandafter=\xglobal@list\do
    {\expandafter\ifx\@@tmp\@@tok
       \aftergroup\let\aftergroup\xglobal@\@break@tfor\fi}%
  \endgroup\global}
%    \end{macrocode}
% \end{macro}
% \end{macro}
%
% \begin{macro}{\xglobal@stop}
% Resetting |\xglobal@|.
%    \begin{macrocode}
\def\xglobal@stop{\let\xglobal@\@empty}
\xglobal@stop
%    \end{macrocode}
% \end{macro}
%
% \begin{macro}{\xglobal@list}
% Initial list of commands that may be prefixed by |\xglobal|.
%    \begin{macrocode}
\def\xglobal@list
 {\definecolor\definecolors\definecolorset\colorlet\providecolor
  \providecolors\providecolorset\blendcolors\maskcolors
  \substitutecolormodel}
%    \end{macrocode}
% \end{macro}
%
% \begin{macro}{\xglobalize@}
% Globalisation command: makes a local definition global.
%    \begin{macrocode}
\def\xglobalize@#1{\ifx\xglobal@\global\global\let#1#1\fi}
%    \end{macrocode}
% \end{macro}
%
% \begin{macro}{\GetGinDriver}
% Transforms the chosen driver option ready for \Package{hyperref}.
%    \begin{macrocode}
\def\GetGinDriver
 {\@tempswafalse
  \@tfor\@@tmp:={dvips}{luatex}{pdftex}{dvipdf}{dvipdfm}{dvipsone}{vtex}{textures}%
  \do{\ifx\XC@@drv\@@tmp\@tempswatrue\@break@tfor\fi}%
  \edef\GinDriver{\if@tempswa\XC@@drv\else hypertex\fi}}
%    \end{macrocode}
% \end{macro}
%
%
% \subsubsection{String manipulation}
%
% \begin{macro}{\@hex@Hex}
%   \Meta{hex-string}|\relax|\\
% Converts \Meta{hex-string} to a string where all letters a--f are changed into their uppercase variants.
% All other letters are left unchanged.
% May be used inside |\edef|.
%    \begin{macrocode}
\def\@hex@Hex#1{\ifx#1\relax\else\expandafter\@hex@@Hex\expandafter#1\fi}
\def\@hex@@Hex#1%
 {\if a#1A\else \if b#1B\else \if c#1C\else \if d#1D\else
  \if e#1E\else \if f#1F\else #1\fi\fi\fi\fi\fi\fi \@hex@Hex}
%    \end{macrocode}
% \end{macro}
%
% \begin{macro}{\XC@strip@comma}
%   \Meta[1]{string}|,|\Meta[2]{string}|,...,|\Meta[n]{string}|,,\@nnil|\\
% Converts a comma-separated list into a space-separated list.
% The evaluation stops at the first empty argument, skipping all remaining tokens until (and including) |\@nnil|.
% Thus, the list to be converted should always contain at least one `|,,|' sequence.
% May be used inside |\edef|.
%    \begin{macrocode}
\def\XC@strip@comma#1,#2%
 {\ifx,#2%
    #1\expandafter\remove@to@nnil\else#1 \expandafter\XC@strip@comma\fi
  #2}
%    \end{macrocode}
% \end{macro}
%
% \begin{macro}{\XC@replace}
%   \marg{cmd}\marg{search}\marg{replace}\\
% Replace all occurrences of \meta{search} within the first-level expansion of \meta{cmd} by \meta{replace} and save the result in \meta{cmd}.
% The replacement is recursive, so don't put the search pattern into the replacement text!
% Note that this macro is incapable of seeing `into' braces.
% The code and its explanation is taken from \People{Donald}{Arseneau}'s \Package{url} package \cite{url}, with only minor changes and renamings applied.
% The original name of the game is |\Url@Edit|.
% Thanks to \People{Donald}{Arseneau} for giving me valuable hints.
%    \begin{macrocode}
{\catcode`Q=3
 \gdef\XC@replace#1#2#3%
  {\begingroup
%    \end{macrocode}
%
% If finished, then remove |\@empty|'s and redefine macro else, re-iterate.
%    \begin{macrocode}
   \def\XC@repl@ce##1#2##2Q##3%
    {\@ifxempty{##2}{\XC@r@pl@ce##1Q}{\XC@repl@ce##1##3##2Q{##3}}}%
   \def\XC@r@pl@ce##1\@empty Q%
    {\expandafter\endgroup\expandafter\def\expandafter#1\expandafter{##1}}%
%    \end{macrocode}
%
% The leading |\@empty| preserves braces, as does the trailing pattern.
%    \begin{macrocode}
   \expandafter\XC@repl@ce\expandafter\@empty #1\@empty#2Q{#3}}
}
%    \end{macrocode}
% \end{macro}
%
% \begin{macro}{\XC@clean}
%   \marg{cmd}\\
% Makes a proper comma-separated list from \meta{cmd}'s contents and stores it again in \meta{cmd}.
% Firstly, all commas are changed into spaces, secondly, `|. |' strings are changed to spaces.
% Thirdly, non-trailing spaces are turned into commas (using a technique similar to \LaTeX's |\zap@space|).
% Finally, the two leading commas are gobbled.
% For example, `| 1 2. ,  3 |' is mapped to `|1,2,3|'.
%    \begin{macrocode}
\def\XC@clean#1%
 {\edef#1{ #1 }%
  \XC@cle@n#1,{ }\XC@cle@n#1{. }{ }\XC@cle@n#1{ },%
  \edef#1{\expandafter\@gobbletwo#1}}
\def\XC@cle@n#1#2#3%
 {\def\XC@cl@@n##1#2##2{#3##1\ifx##2\@empty\else\expandafter\XC@cl@@n\fi##2}%
  \edef#1{\expandafter\XC@cl@@n#1#2\@empty}}%
%    \end{macrocode}
% \end{macro}
%
%
% \subsubsection{Scalar arithmetic}
%
% \begin{macro}{\rdivide}
%   \marg{dimension register}\marg{div}\\
% Divides \Meta{dimension register} by a real number \Meta{div}, which may also be buried in a dimension (register), and stores the resulting rounded dimension in \Meta{dimension register}.
% The denominator \Meta{div} is assumed to be non-zero.
%    \begin{macrocode}
\def\rdivide#1#2%
 {\begingroup\toks@{#1}\edef\@@tmp{\noexpand\@tempdima\the#1}%
  \@defaultunits\@tempdimb=#2pt\relax\@nnil \@@tmp\relax
  \ifdim\@tempdima<\z@ \@tempdima-\@tempdima\def\@sign{-}\else
    \def\@sign{}\fi
  \ifdim\@tempdimb<\z@ \@tempdimb-\@tempdimb\edef\@sign{\@sign-}\fi
  \ifdim\@tempdima<.1\maxdimen \ifdim\@tempdimb<.01\maxdimen
    \lshift\@tempdima \lshift\@tempdimb \fi\fi
  \@tempcnta=\@tempdima \@tempcntb=\@tempdimb
  \count@=\@tempcnta \divide\count@\@tempcntb
  \edef\@@tmp{\number\count@.}%
  \@tempswatrue  \@rdivide \@rdivide \@rdivide \@rdivide \@rdivide
  \@tempswafalse \@rdivide
  \dimen@\@@tmp\p@ \ifnum\count@>4\advance\dimen@1sp\fi
  \dimen@\@sign\dimen@
  \edef\@@tmp{\the\toks@\the\dimen@}%
  \expandafter\endgroup\@@tmp}
%    \end{macrocode}
%
%    \begin{macrocode}
\def\@rdivide
 {\multiply\count@\@tempcntb \advance\@tempcnta-\count@
  \multiply\@tempcnta10 \count@=\@tempcnta \divide\count@\@tempcntb
  \if@tempswa\edef\@@tmp{\@@tmp\number\count@}\fi}
%    \end{macrocode}
% \end{macro}
%
% \begin{macro}{\rmultiply}
%   \marg{dimension register}\marg{dec}\\
% Multiplies \Meta{dimension register} by \Meta{dec} and stores the resulting dimension in \Meta{dimension register}.
%    \begin{macrocode}
\def\rmultiply#1#2%
 {\begingroup
  \edef\@@mul{#2}\@tempdima\the#1\relax
  \afterassignment\@tempdimb\@tempcntb=\@@mul\p@
  \ifdim\@tempdima<\z@ \@tempdima-\@tempdima\def\@sign{-}\else
    \def\@sign{}\fi
  \ifnum\@tempcntb<\z@ \@tempcntb-\@tempcntb\edef\@sign{\@sign-}\fi
  \lshift\@tempdima \dimen@\z@
  \expandafter\@rmultiply\@@mul.0.\@nil
  \advance\dimen@\@tempcntb\@tempdima \rshift\dimen@
  \dimen@\@sign\dimen@
  \edef\@@tmp{#1\the\dimen@}%
  \expandafter\endgroup\@@tmp}
%    \end{macrocode}
%
%    \begin{macrocode}
\def\@rmultiply#1.#2.#3\@nil
 {\def\@@tmp{}%
  \@tfor\@@mul:=#2\do
   {\edef\@@tmp
     {\advance\dimen@\@@mul\@tempdima \divide\dimen@10 \@@tmp}}%
  \@@tmp}
%    \end{macrocode}
% \end{macro}
%
% \begin{macro}{\lshift}
% \begin{macro}{\llshift}
%   \marg{dimension register}\\
% Multiplies \Meta{dimension register} safely by 10 resp. 100.
%    \begin{macrocode}
\def\lshift#1{#1\expandafter\lshift@\the#1}
\def\llshift#1{\lshift#1\lshift#1}
%    \end{macrocode}
% \end{macro}
% \end{macro}
%
% \begin{macro}{\lshiftnum}
% \begin{macro}{\llshiftnum}
%   \marg{decimal number}\\
% Multiplies \Meta{decimal number} by 10 resp. 100.
%    \begin{macrocode}
\def\lshiftnum#1{\expandafter\lshiftnum@#1.00.:}
\def\lshiftnum@#1.#2#3.#4:{#1#2\ifnum0#3>\z@.#3\fi}
\def\llshiftnum#1{\expandafter\llshiftnum@#1\@empty.00.:}
\def\llshiftnum@#1.#2#3#4.#5:%
 {#1#2\ifnum0#3>\z@#3\else0\fi\ifnum0#4>\z@.#4\fi}
%    \end{macrocode}
% \end{macro}
% \end{macro}
%
% \begin{macro}{\lshiftset}
% \begin{macro}{\llshiftset}
%   \marg{dimension register}\marg{decimal number}\\
% Multiplies \Meta{decimal number} by 10 resp. 100 and puts it into \Meta{dimension register}.
%    \begin{macrocode}
\def\lshiftset#1#2{\edef\@@XC{\noexpand\lshiftnum{#2}}#1\@@XC\p@}
\def\llshiftset#1#2{\edef\@@XC{\noexpand\llshiftnum{#2}}#1\@@XC\p@}
%    \end{macrocode}
% \end{macro}
% \end{macro}
%
% \begin{macro}{\rshift}
% \begin{macro}{\rrshift}
%   \marg{dimension register}\\
% Divides \Meta{dimension register} safely by 10 resp. 100.
%    \begin{macrocode}
\def\rshift#1{#1\expandafter\rshift@\the#1}
\def\rrshift#1{\rshift#1\rshift#1}
%    \end{macrocode}
% \end{macro}
% \end{macro}
%
% \begin{macro}{\lshift}
% \begin{macro}{\llshift}
%   \marg{dec}\\
% Multiplies \Meta{dec}, which may be either a number or something like |\the\dimen@|, safely by 10 resp. 100.
% The following definitions are similar to the way how |\strip@pt| is defined in the \LaTeX{} kernel.
%    \begin{macrocode}
\begingroup
\catcode`P=12
\catcode`T=12
\lowercase
 {\def\@@tmp{\def\lshift@##1.##2##3PT{##1##2\ifnum0##3>\z@.##3\fi\p@}
             \def\rshift@##1.##2PT{\rshift@@##1\relax##2\p@}}}
\expandafter\endgroup\@@tmp
%    \end{macrocode}
% \end{macro}
% \end{macro}
%
%    \begin{macrocode}
\def\rshift@@#1#2%
 {\ifx#2\relax.#1\else#1\expandafter\rshift@@\expandafter#2\fi}
%    \end{macrocode}
%
%
% \subsubsection{Vector arithmetic}
%
% \begin{macro}{\XC@inflate}
%   \marg{core model}\marg{first}\marg{second}\marg{cmd}\\
% Constructs an 8-dimensional vector based on two vectors \Meta{first}, \Meta{second} (of common dimension $1,3,4$) and an underlying color model (stored in the macro \Meta{core model}).
% The resulting vector is stored in \Meta{cmd}.
%    \begin{macrocode}
\def\XC@inflate#1#2#3#4%
 {\edef#4{#2,\ifx#1\XC@mod@cmyk #3\else
             \ifx#1\XC@mod@gray ,,,#3,,,\else ,#3,\fi\fi}}
%    \end{macrocode}
% \end{macro}
%
% \begin{macro}{\XC@vec}
%   \Meta{a}|,|\Meta{b}|,|\Meta{c}|,|\Meta{d}|,|\Meta{a'}|,|\Meta{b'}|,|\Meta{c'}|,|\Meta{d'}|,|\marg{cmd}\\
% Generic vector calculation for two vectors (of common dimension $1,3,4$), puts the result into \Meta{cmd}.
%    \begin{macrocode}
\def\XC@vec#1,#2,#3,#4,#5,#6,#7,#8,#9%
 {\def\@@tmp{}%
  \XC@vec@{#1}{#5}%
  \ifx,#2,\else
  \XC@vec@{#2}{#6}\XC@vec@{#3}{#7}\ifx,#4,\else\XC@vec@{#4}{#8}\fi\fi
  \aftergroupdef#9\@@tmp}
%    \end{macrocode}
%
%    \begin{macrocode}
\def\XC@vec@@{\edef\@@tmp{\ifx\@@tmp\@empty\else\@@tmp,\fi\strip@pt\dimen@}}
%    \end{macrocode}
% \end{macro}
%
% \begin{macro}{\XC@add}
%   \Meta{a}|,|\Meta{b}|,|\Meta{c}|,|\Meta{d}|,|\Meta{a'}|,|\Meta{b'}|,|\Meta{c'}|,|\Meta{d'}|,|\marg{cmd}\\
% Calculates the sum of two vectors (of common dimension $1,3,4$) and puts it into \Meta{cmd}.
%    \begin{macrocode}
\def\XC@add{\begingroup \let\XC@vec@\XC@add@ \XC@vec}
%    \end{macrocode}
%
%    \begin{macrocode}
\def\XC@add@#1#2%
 {\llshiftset\dimen@{#1}\llshiftset\@tempdima{#2}%
  \advance\dimen@\@tempdima \rrshift\dimen@ \XC@vec@@}
%    \end{macrocode}
% \end{macro}
%
% \begin{macro}{\XC@sub}
%   \Meta{a}|,|\Meta{b}|,|\Meta{c}|,|\Meta{d}|,|\Meta{a'}|,|\Meta{b'}|,|\Meta{c'}|,|\Meta{d'}|,|\marg{cmd}\\
% Calculates the difference of two vectors (of common dimension $1,3,4$) and puts it into \Meta{cmd}.
%    \begin{macrocode}
\def\XC@sub{\begingroup \let\XC@vec@\XC@sub@ \XC@vec}
%    \end{macrocode}
%
%    \begin{macrocode}
\def\XC@sub@#1#2{\XC@add@{#1}{-#2}}
%    \end{macrocode}
% \end{macro}
%
% \begin{macro}{\XC@mul}
%   \Meta{a}|,|\Meta{b}|,|\Meta{c}|,|\Meta{d}|,|\Meta{a'}|,|\Meta{b'}|,|\Meta{c'}|,|\Meta{d'}|,|\marg{cmd}\\
% Calculates the (component-wise) product of two vectors (of common dimension $1,3,4$) and puts it into \Meta{cmd}.
%    \begin{macrocode}
\def\XC@mul{\begingroup \let\XC@vec@\XC@mul@ \XC@vec}
%    \end{macrocode}
%
%    \begin{macrocode}
\def\XC@mul@#1#2%
 {\llshiftset\dimen@{#1}\rmultiply\dimen@{#2}\rrshift\dimen@ \XC@vec@@}
%    \end{macrocode}
% \end{macro}
%
% \begin{macro}{\XC@mix}
%   \Meta{a}|,|\Meta{b}|,|\Meta{c}|,|\Meta{d}|,|\Meta{a'}|,|\Meta{b'}|,|\Meta{c'}|,|\Meta{d'}|,|\marg{cmd}\\
% Calculates the linear interpolation of two color parameter vectors and puts it into \Meta{cmd}.
% The percentage parameter $p$ is input via register |\dimen@|.
%    \begin{macrocode}
\def\XC@mix%
 {\begingroup
  \@tempdima\dimen@ \@tempdimb100\p@ \advance\@tempdimb-\@tempdima
  \let\XC@vec@\XC@mix@ \XC@vec}
%    \end{macrocode}
%
%    \begin{macrocode}
\def\XC@mix@#1#2%
 {\dimen@#1\@tempdima \advance\dimen@#2\@tempdimb \rrshift\dimen@ \XC@vec@@}
%    \end{macrocode}
% \end{macro}
%
%
% \begin{macro}{\XC@calc@}
%   \Meta{x}|,|\Meta{y}|,|\Meta{z}|,|\Meta{w}|,|\Meta{dummy}|:|\marg{function}\\
% Performs, dependent on the number of non-empty parameters, |\edef\@@tmp{|$f(x)$|}| or |\edef\@@tmp{|$f(x),f(y),f(z)$|}| or |\edef\@@tmp{|$f(x),f(y),f(z),f(w)$|}|, where the function $f$ is determined by the parameter \Meta{function}.
%    \begin{macrocode}
\def\XC@calc@#1,#2,#3,#4,#5:#6%
 {\XC@let@Nc\@@cmd{XC@calc#6}\@@cmd{#1}\@@tmp
  \ifx,#2,\else
    \@@cmd{#2}\@tempa\@@cmd{#3}\@tempb
    \ifx,#4,%
      \edef\@@tmp{\@@tmp,\@tempa,\@tempb}\else
      \@@cmd{#4}\@tempc\edef\@@tmp{\@@tmp,\@tempa,\@tempb,\@tempc}\fi\fi}
%    \end{macrocode}
% \end{macro}
%
% \begin{macro}{\XC@calcC}
%   \marg{x}\marg{cmd}\\
% \textbf Complement of $x$: performs |\def|\Meta{cmd}|{|$1-x$|}|.
% Yields exact results for all numbers within $[0,1]$ with 5 decimal digits.
%    \begin{macrocode}
\def\XC@calcC#1#2%
 {\llshiftset\dimen@{-#1}\advance\dimen@100\p@
  \rrshift\dimen@ \XC@calcN{\strip@pt\dimen@}#2}
%    \end{macrocode}
% \end{macro}
%
% \begin{macro}{\XC@calcD}
%   \marg{x}\marg{cmd}\\
% \textbf Division of $x$ by $s$: performs |\def|\Meta{cmd}|{|$x/s$|}|.
% The parameter $s$ is input via |\@@scl|.
%    \begin{macrocode}
\def\XC@calcD{\XC@c@lcD\@@scl{}}
%    \end{macrocode}
% \end{macro}
%
% \begin{macro}{\XC@c@lcD}
%   \marg{s}\marg{string}\marg{x}\marg{cmd}\\
% \textbf Division of $x$ by $s$, appending \Meta{string}: performs |\def|\Meta{cmd}|{|$x/s$\Meta{string}|}|.
%    \begin{macrocode}
\def\XC@c@lcD#1#2#3#4%
 {\dimen@#3\p@ \rdivide\dimen@{#1}\edef#4{\strip@pt\dimen@#2}}
%    \end{macrocode}
% \end{macro}
%
% \begin{macro}{\XC@calcF}
%   \marg{x}\marg{cmd}\\
% \textbf Format $x$ with 5 decimals, no range checks.
%    \begin{macrocode}
\def\XC@calcF#1#2%
 {\edef#2{#1}\expandafter\XC@c@lcF#2.0.:#2}
%    \end{macrocode}
%
% Special treatment of integer part because of so many possible variations ($---.0$, $7$, \dots).
%    \begin{macrocode}
\def\XC@c@lcF#1.#2.#3:#4%
 {\edef#4{\XC@c@@cR#200000:}%
  \dimen@#1.5\p@ \afterassignment\dimen@\count@=\the\dimen@
  \edef#4%
   {\ifnum#4>0 \ifnum\count@=0 \ifdim#1.5\p@<\z@ -\fi\fi\fi
    \number\count@.#4}}
%    \end{macrocode}
% \end{macro}
%
% \begin{macro}{\XC@calcH}
%   \marg{x}\marg{cmd}\\
% \textbf Hexadecimal transformation of $x\in\{0,1,\dots,255\}$: performs |\def|\Meta{cmd}|{|$x_{16}$|}|.
%    \begin{macrocode}
\def\XC@calcH#1#2%
 {\count@#1 \divide\count@\sixt@@n \edef\XC@@tmp{\hexnumber@\count@}%
  \multiply\count@-\sixt@@n \advance\count@#1
  \edef#2{\XC@@tmp\hexnumber@\count@}}
%    \end{macrocode}
% \end{macro}
%
% \begin{macro}{\XC@calcM}
%   \marg{x}\marg{cmd}\\
% \textbf Multiplication of $x$ by $s$ plus rounding: performs |\def|\Meta{cmd}|{|$[s\cdot x+0.5]$|}|.
% The parameter $s$ is input via |\@@scl|.
%    \begin{macrocode}
\def\XC@calcM#1#2%
 {\dimen@#1\p@ \dimen@\@@scl\dimen@ \advance\dimen@.5\p@
  \afterassignment\dimen@\count@=\the\dimen@ \edef#2{\number\count@}}
%    \end{macrocode}
% \end{macro}
%
% \begin{macro}{\XC@calcN}
%   \marg{x}\marg{cmd}\\
% \textbf Normalisation of $x$: similar to |\XC@calcR|, but trailing zeros are eliminated.
%    \begin{macrocode}
\def\XC@calcN#1#2%
 {\XC@calcR{#1}#2\edef#2{\expandafter\XC@c@lcN#200000:}}
%    \end{macrocode}
%
% Gobble trailing zeros.
%    \begin{macrocode}
\def\XC@c@lcN#1.#200000#3:{#1\ifnum0#2=0 \else.#2\fi}
%    \end{macrocode}
% \end{macro}
%
% \begin{macro}{\XC@calcR}
%   \marg{x}\marg{cmd}\\
% \textbf Reduction of $x$ to the unit interval: performs |\def|\Meta{cmd}|{|$\min\{1,\max\{0,x\}\}$|}|.
% Numbers are output with 5 decimals.
% This could also be done using |\XC@calcF|, followed by a range check, but perhaps less efficiently.
%    \begin{macrocode}
\def\XC@calcR#1#2{\edef#2{#1}\edef#2{\expandafter\XC@c@lcR#2.0.:}}
%    \end{macrocode}
%
% Check interval.
%    \begin{macrocode}
\def\XC@c@lcR#1.#2.#3:%
 {\ifnum#10>0 1.00000\else 0.\ifnum#11<0 00000\else
  \XC@c@@cR#200000:\fi\fi}
%    \end{macrocode}
%
% Normalize to 5 digits.
%    \begin{macrocode}
\def\XC@c@@cR#1#2#3#4#5#6:{#1#2#3#4#5}
%    \end{macrocode}
% \end{macro}
%
% \begin{macro}{\XC@calcS}
%   \marg{x}\marg{cmd}\\
% \textbf Scaling of $x$ by $s$: performs |\def|\Meta{cmd}|{|$s\cdot x$|}|.
% The parameter $s$ is input via |\@@scl|.
%    \begin{macrocode}
\def\XC@calcS{\XC@c@lcS\@@scl{}}
%    \end{macrocode}
% \end{macro}
%
% \begin{macro}{\XC@c@lcS}
%   \marg{s}\marg{string}\marg{x}\marg{cmd}\\
% \textbf Scaling of $x$ by $s$, appending \Meta{string}: performs |\def|\Meta{cmd}|{|$s\cdot x$\Meta{string}|}|.
%    \begin{macrocode}
\def\XC@c@lcS#1#2#3#4%
 {\lshiftset\dimen@{#3}%
  \ifdim#1\p@<100\p@
    \dimen@\lshiftnum#1\dimen@ \rrshift\dimen@ \else
    \dimen@#1\dimen@ \rshift\dimen@ \fi
  \edef#4{\strip@pt\dimen@#2}}
%    \end{macrocode}
% \end{macro}
%
% \begin{macro}{\XC@calcT}
%   \marg{x}\marg{cmd}\\
% \textbf Translate $x$ by $a$ and cut to unit interval: performs |\def|\Meta{cmd}|{|$\min\{1,\max\{0,x+a\}\}$|}|.
% The parameter $a$ is input via |\@@arg|.
%    \begin{macrocode}
\def\XC@calcT#1#2%
 {\lshiftset\dimen@{#1}\advance\dimen@\lshiftnum\@@arg\p@ \rshift\dimen@
  \edef#2%
   {\ifdim\dimen@>\p@1\else\ifdim\dimen@<\z@0\else\strip@pt\dimen@\fi\fi}}
%    \end{macrocode}
% \end{macro}
%
% \begin{macro}{\XC@calcU}
%   \marg{x}\marg{cmd}\\
% \textbf Unit interval mapping of $x$: performs |\def|\Meta{cmd}|{|$x-[x]$|}| if $x\neq1$, $1$ otherwise.
%    \begin{macrocode}
\def\XC@calcU#1#2%
 {\dimen@#1\p@
  \ifdim\dimen@=\p@\else
    \afterassignment\dimen@\count@=\the\dimen@
    \relax\ifdim\dimen@=\z@\else
      \ifdim#1\p@<\z@ \dimen@-\dimen@ \advance\dimen@\p@ \fi
    \fi
  \fi
  \edef#2{\strip@pt\dimen@}}
%    \end{macrocode}
% \end{macro}
%
%
% \subsection{Color definition}
%
% \subsubsection{Constants}
%
% \begin{macro}{\rangeRGB}
% \begin{macro}{\rangeHSB}
% \begin{macro}{\rangeHsb}
% \begin{macro}{\rangetHsb}
% \begin{macro}{\rangeGray}
% Scaling ranges for `integer' and other special models:
%    \begin{macrocode}
\providecommand*\rangeRGB{255}
\providecommand*\rangeHSB{240}
\providecommand*\rangeHsb{360}
\providecommand*\rangetHsb{60,30;120,60;180,120;210,180;240,240}
\providecommand*\rangeGray{15}
%    \end{macrocode}
% \end{macro}
% \end{macro}
% \end{macro}
% \end{macro}
% \end{macro}
%
% \begin{macro}{\adjustUCRBG}
% Parameters for undercolor-removal (UCR) and black-generation (BG) during \Model{cmy} to \Model{cmyk} conversion:
%    \begin{macrocode}
\providecommand*\adjustUCRBG{1,1,1,1}
%    \end{macrocode}
% \end{macro}
%
% \begin{macro}{\paperquality}
% Parameter for choosing different color sets according to the paper quality; yet supported only for the HKS colors.
%    \begin{macrocode}
\providecommand*\paperquality{1}
%    \end{macrocode}
% \end{macro}
%
% \begin{macro}{\colorseriescycle}
% Default cycle length for color series:
%    \begin{macrocode}
\providecommand*\colorseriescycle{16}
%    \end{macrocode}
% \end{macro}
%
% Model-dependent base colors:
%    \begin{macrocode}
\def\XC@clr@rgb@white{1,1,1}
\def\XC@clr@cmy@white{0,0,0}
\def\XC@clr@cmyk@white{0,0,0,0}
\def\XC@clr@hsb@white{\@@hue,0,1}
\def\XC@clr@hsb@gray{\@@hue,0,0.5}
\def\XC@clr@hsb@black{\@@hue,0,0}
\def\XC@clr@gray@white{1}
%    \end{macrocode}
%
% \begin{macro}{\XC@met@grad}
% \begin{macro}{\XC@met@last}
% \begin{macro}{\XC@met@step}
% Methods for color series:
%    \begin{macrocode}
\@tfor\@@tmp:={grad}{last}{step}\do
 {\XC@let@cN{XC@met@\@@tmp}\@@tmp}
%    \end{macrocode}
% \end{macro}
% \end{macro}
% \end{macro}
%
%
% \subsubsection{Colors}\label{sec.definition}
%
% \begin{macro}{\definecolor}
%   \oarg{type}\marg{name}\oarg{prefix}\marg{model-list}\marg{spec-list}\\
% Extends the facilities of \LC's |\definecolor|.
% The optional argument \Meta{type} has currently to be either empty (= normal color definition) or `named' or `ps'.
% \Meta{prefix} is an argument used for the internal driver-dependent representation.
% |\definecolor| creates a new command |\\color@|\Meta{name} that contains the \emph{internal representation}
% |\xcolor@|\marg{type-spec}\marg{driver-spec}\marg{n-model}\marg{n-color-spec} of the color, where
% \begin{itemize}
% \item \Meta{type-spec} is either empty (if \Meta{type} was empty), or a driver-dependent macro or expression (if \Meta{type} was `named', `ps');
% \item \Meta{driver-spec} is a driver-dependent representation of the color, which also depends on \Meta{type};
% \item \Meta{n-model} is either \Meta{model} or the name of another color model, dependent on the desired transformation(s) at definition time;
% \item \Meta{n-color-spec} are the comma-separated normalised parameter values of the color with respect to \Meta{n-model}.
% \end{itemize}
% See table \vref{tab.internal} for some examples.
%    \begin{macrocode}
\protected\def\definecolor{\@testopt{\XC@definecolor}{}}
\def\XC@definecolor[#1]#2%
 {\@testopt{\XC@definec@lor[#1]{#2}}\colornameprefix}
%    \end{macrocode}
%    \begin{macrocode}
\def\XC@definec@lor[#1]#2[#3]#4#5%
 {%
  \begingroup
  \XC@sdef\@@cls{#1}\XC@edef\@@nam{#2}\edef\colornameprefix{#3}%
  \XC@logdef\@@nam
  \XC@getmodclr02{#4}{#5}{\let\@@tmp\relax}%
   {\ifx\@@mod\XC@mod@named
      \XC@c@l@rlet\@@nam\@@clr
    \else
      \ifx\@@cls\XC@mod@ps
        \edef\@@drv{\@@mod\space\@@clr}\def\@@hue{0}%
        \edef\@@clr{\@nameuse{XC@clr@\@@mod @white}}%
      \else
        \ifconvertcolorsD
          \let\@@tmp\@@mod\XC@sdef\@@mod{\XC@tgt@mod\@@tmp}%
          \convertcolorspec\@@tmp\@@clr\@@mod\@@clr
        \fi
        \XC@coremodel\@@mod\@@clr
%    \end{macrocode}
%
% Determine driver-dependent color definition:
%    \begin{macrocode}
        \csname color@\@@mod\expandafter\endcsname
          \expandafter\@@drv\expandafter{\@@clr}%
%    \end{macrocode}
%
% Handle prologue and `named' colors:
%    \begin{macrocode}
        \ifcase\XC@@xcp@
          \immediate\write\XC@@xcp{/\colornameprefix\@@nam{\expandafter
            \XC@strip@comma\@@clr,,\@nnil}XC\@@mod}%
          \let\@@cls\XC@mod@named
        \fi
        \ifx\@@cls\@empty\else
          \let\@@tmp\@@cls
          \@nameuse{define@color@\@@cls}\@@nam\@@drv
          \csname color@\@@tmp\expandafter\endcsname
            \expandafter\@@drv\expandafter{\@@nam}%
        \fi
      \fi
%    \end{macrocode}
%
% Finally write all information into macro |\\color@|\Meta{name}:
%    \begin{macrocode}
      \toks@\expandafter{\@@drv}%
      \edef\@@tmp
       {\ifglobalcolors\global\else\xglobal@\fi
        \noexpand\@namedef{\@backslashchar color@\@@nam}%
         {\noexpand\xcolor@{\@@cls}{\the\toks@}{\@@mod}{\@@clr}}}%
    \fi}%
  \expandafter\endgroup\@@tmp\xglobal@stop}
%    \end{macrocode}
% \end{macro}
%
% \begin{macro}{\XC@getmodclr}
%   \marg{mod-opt}\marg{clr-opt}\marg{model-list}\marg{spec-list}\marg{undefined}\marg{defined}\\
% Depending on \Meta{mod-opt}, different tasks are performed:
% \begin{itemize}
% \item[0:] Retrieve color model and target model from \Meta{model-list} if |\XC@@getmod| is `true'.
% \item[1:] Retrieve color model and target model from \Meta{model-list} always.
% \end{itemize}
% The color model is stored in |\@@mod|, its position in |\XC@@pos|, a possible target model in |\@@Mod|.
% Depending on \Meta{clr-opt}, different tasks are performed:
% \begin{itemize}
% \item[0:] Don't pick color specification from \Meta{spec-list}.
% \item[1:] Pick color specification from \Meta{spec-list}, convert it from model |\@@mod| to |\@@Mod|, unless the latter is `ignore' or |\@@cls| is `ps'.
% \item[2:] Like the former case, additionally |\let\@@mod\@@Mod|.
% \end{itemize}
% The color specification is stored in |\@@clr|.
% Color specification picking relies on the position |\XC@@pos|, conversion relies on the model |\@@mod|, and the target model |\@@Mod|.
% If |\@@mod| represents an undefined color model, then \Meta{undefined} is executed at the end of the macro, else \Meta{defined}.
%    \begin{macrocode}
\def\XC@getmodclr#1#2#3#4%
 {\ifnum\ifcase#1 \ifXC@@getmod1\else0 \fi \or1 \fi>0
    \XC@mdef\@@mod{#3}\expandafter\XC@getmod\@@mod////////\XC@@
    \@ifundefinedmodel\@@Mod{\let\@@Mod\XC@mod@ignore}{}\fi
  \@ifundefined{XC@mod@\@@mod}%
   {\c@lor@error{model `\@@mod'}\@firstoftwo}%
   {\ifcase#2\else
      \XC@mdef\@@clr{#4}%
      \edef\@@clr{\expandafter\XC@getclr\@@clr////////\XC@@}%
      \ifx\@@cls\XC@mod@ps\else\XC@clean\@@clr\fi
      \ifx\@@Mod\XC@mod@ignore\else
        \convertcolorspec\@@mod\@@clr\@@Mod\@@clr
        \ifnum#2>1 \let\@@mod\@@Mod\fi
      \fi
    \fi \@secondoftwo}}
%    \end{macrocode}
%
% \begin{macro}{\XC@getmod}
% Retrieve model position |\XC@@pos|, model |\@@mod|, and target model |\@@Mod| from multiple definitions (up to~8), separated by `|/|'.
%    \begin{macrocode}
\def\XC@getmod#1/#2/#3/#4/#5/#6/#7/#8/#9\XC@@
 {\def\XC@dflt@@mod{#1}%
  \XC@g@tm@d0#1::\XC@@ \XC@getm@d1{#2}\XC@getm@d2{#3}%
  \XC@getm@d3{#4}\XC@getm@d4{#5}\XC@getm@d5{#6}\XC@getm@d6{#7}%
  \XC@getm@d7{#8}\def\XC@@pos{0}\let\@@mod\XC@dflt@@mod\@gobble\@nnil}%
%    \end{macrocode}
%
%    \begin{macrocode}
\def\XC@getm@d#1#2%
 {\edef\XC@@pos{#2}\XC@sdef\@@mod{\XC@tgt@mod\XC@@pos}%
  \ifx\XC@@pos\@@mod\def\XC@@pos{#1}\expandafter\remove@to@nnil\fi}%
%    \end{macrocode}
%
%    \begin{macrocode}
\def\XC@g@tm@d#1#2:#3:#4\XC@@
 {\@ifxempty{#3}{\let\@@Mod\XC@mod@ignore\XC@getm@d0{#2}}%
    {\def\@@Mod{#2}%
     \def\XC@tgt@mod##1{#2}%
     \def\XC@dflt@@mod{#3}%
     \XC@getm@d0{#3}}}%
%    \end{macrocode}
% \end{macro}
%
% \begin{macro}{\XC@getclr}
% Pick color specification on position |\XC@@pos| from multiple definitions (up to~8), separated by `|/|'.
% Expandable.
%    \begin{macrocode}
\def\XC@getclr#1/#2/#3/#4/#5/#6/#7/#8/#9\XC@@
 {\ifcase\XC@@pos\space #1\or#2\or#3\or#4\or#5\or#6\or#7\or#8\fi}
%    \end{macrocode}
% \end{macro}
%
% \begin{macro}{\ifXC@@getmod}
% Switch used to determine whether the model has to be retrieved (e.g., ordinary |\definecolor|) or not (e.g., |\definecolor| acting as subroutine to |\definecolorset|).
%    \begin{macrocode}
\newif\ifXC@@getmod \XC@@getmodtrue
%    \end{macrocode}
% \end{macro}
% \end{macro}
%
% \begin{macro}{\xdefinecolor}
% This command is identical with |\definecolor|, it is kept only for compatibility reasons.
%    \begin{macrocode}
\let\xdefinecolor\definecolor
%    \end{macrocode}
% \end{macro}
%
% \begin{macro}{\providecolor}
%   \oarg{type}\marg{name}\marg{model-list}\marg{spec-list}\\
% Similar to |\definecolor|, but the color \Meta{name} is only defined if it does not exist already.
%    \begin{macrocode}
\def\providecolor{\@testopt{\XC@providecolor}{}}
\def\XC@providecolor[#1]#2%
 {\@testopt{\XC@providec@lor[#1]{#2}}\colornameprefix}
%    \end{macrocode}
%
%    \begin{macrocode}
\def\XC@providec@lor[#1]#2[#3]%
 {\XC@edef\XC@@tmp{#2}%
  \@ifundefinedcolor\XC@@tmp{\XC@definecolor[#1]{#2}[#3]}%
   {\ifnum\XC@tracing>0
      \PackageInfo{xcolor}{Ignoring color definition `\XC@@tmp'}\fi
    \@gobbletwo}}
%    \end{macrocode}
% \end{macro}
%
% \begin{macro}{\xcolor@}
%   \marg{type-spec}\marg{driver-spec}\marg{model}\marg{spec}\\
% Simply grabs the driver-specific part of the color definition.
% Therefore, the command |\\color@|\Meta{name} behaves like in the original case, when first-level expanded.
%    \begin{macrocode}
\def\xcolor@#1#2#3#4{#2}
%    \end{macrocode}
% \end{macro}
%
% \begin{macro}{\colorlet}
%   \oarg{type}\marg{name}\oarg{num model}\marg{color}\\
% Create a color `name' consisting of \Meta{color}, converted to \Meta{num model} if non-empty.
%    \begin{macrocode}
\def\colorlet{\@testopt{\XC@colorlet}{}}
\def\XC@colorlet[#1]#2{\@testopt{\XC@col@rlet[#1]{#2}}{}}
\def\XC@col@rlet[#1]#2[#3]#4%
 {\begingroup
  \edef\@@cls{#1}\XC@edef\@@nam{#2}\XC@sdef\@@mod{#3}\XC@edef\@@clr{#4}%
%    \end{macrocode}
%
% Check whether |#4| is an explicit name, |#3| is empty and |#1| is either empty or equal to the corresponding parameter of |#4|: in this case, a simple |\let| command will do the job.
%    \begin{macrocode}
  \XC@info\@@clr\@@tmp\XC@@tstfalse
  \ifnum\@@tmp=\@ne\ifx\@@mod\@empty
    \ifx\@@cls\@empty
      \XC@@tsttrue
    \else
      \edef\@@tmp
       {\expandafter\expandafter\expandafter\@secondoffive
         \csname\@backslashchar color@\@@clr\endcsname}%
      \ifx\@@cls\@@tmp\XC@@tsttrue\fi
    \fi
  \fi\fi
  \ifXC@@tst
    \XC@logdef\@@nam\XC@c@l@rlet\@@nam\@@clr
  \else
    \extractcolorspec\@@clr\@@clr
    \ifx\@@mod\@empty\else
      \expandafter\convertcolorspec\@@clr\@@mod\@@clr
      \edef\@@clr{{\@@mod}{\@@clr}}\fi
    \edef\@@tmp{\noexpand\XC@definecolor[\@@cls]{\@@nam}\@@clr}%
  \fi
  \expandafter\endgroup\@@tmp\xglobal@stop}
%    \end{macrocode}
%
%    \begin{macrocode}
\def\XC@c@l@rlet#1#2%
 {\@ifundefinedcolor{#2}%
   {\c@lor@error{`#2'}\let\@@tmp\@empty}%
   {\edef\@@tmp
     {\ifglobalcolors\global\else\xglobal@\fi
      \noexpand\XC@let@cc{\@backslashchar color@#1}{\@backslashchar color@#2}}}}
%    \end{macrocode}
% \end{macro}
%
% \begin{macro}{\DefineNamedColor}
%   \marg{type}\marg{name}\marg{model-list}\marg{spec-list}
%    \begin{macrocode}
\protected\def\DefineNamedColor#1%
 {\@nameuse{XC@\ifdefinecolors defin\else prepar\fi ecolor}[#1]}
%    \end{macrocode}
%  \end{macro}
%
%
% \subsubsection{Deferred definition}
%
% \begin{macro}{\preparecolor}
%   \oarg{type}\marg{name}\oarg{prefix}\marg{model-list}\marg{spec-list}\\
%    \begin{macrocode}
\def\preparecolor{\@testopt{\XC@preparecolor}{}}
\def\XC@preparecolor[#1]#2%
 {\@testopt{\XC@preparec@lor[#1]{#2}}\colornameprefix}
%    \end{macrocode}
%
%    \begin{macrocode}
\def\XC@preparec@lor[#1]#2[#3]#4#5%
 {\XC@sdef\@@cls{#1}\XC@edef\@@nam{#2}%
  \XC@getmodclr01{#4}{#5}{}%
   {\xdef\XC@@stack
     {\ifx\XC@@stack\@empty\else\XC@@stack,\fi{#1}{\@@nam}{#3}%
      {\ifx\@@Mod\XC@mod@ignore\@@mod\else\@@Mod\fi}{\@@clr}}}}
%    \end{macrocode}
%
%    \begin{macrocode}
\let\XC@@stack\@empty
%    \end{macrocode}
%  \end{macro}
%
% \begin{macro}{\definecolors}
%   \marg{id-list}
%    \begin{macrocode}
\def\definecolors#1%
 {\let\XC@@cmd\definecolor\let\XC@@glb\xglobal@
  \expandafter\XC@definecolors#1,,}
%    \end{macrocode}
%
%    \begin{macrocode}
\def\XC@definecolors#1,%
 {\@ifxempty{#1}{}{\expandafter\XC@definec@lors#1=#1=:\XC@definecolors}}
%    \end{macrocode}
%
%    \begin{macrocode}
\def\XC@definec@lors#1=#2=#3:%
 {\XC@edef\@@nam{#1}\XC@edef\@@arg{#2}%
  \let\next\XC@definec@l@rs\expandafter\next\XC@@stack,,\@nnil}
%    \end{macrocode}
%
%    \begin{macrocode}
\def\XC@definec@l@rs#1,%
 {\ifx,#1,\let\next\remove@to@nnil\else\XC@defin@c@l@rs#1\fi\next}
%    \end{macrocode}
%
%    \begin{macrocode}
\def\XC@defin@c@l@rs#1#2#3#4#5%
 {\def\@@tmp{#2}%
  \ifx\@@tmp\@@arg
    \let\xglobal@\XC@@glb\XC@@cmd[#1]{\@@nam}[#3]{#4}{#5}%
    \let\next\remove@to@nnil
  \fi}
%    \end{macrocode}
%  \end{macro}
%
% \begin{macro}{\providecolors}
%   \marg{id-list}\\
% Similar to |\definecolor|, but based on |\providecolor|, thus the individual colors are defined only if they do not exist already.
%    \begin{macrocode}
\def\providecolors#1%
 {\let\XC@@cmd\providecolor\let\XC@@glb\xglobal@
  \expandafter\XC@definecolors#1,,}
%    \end{macrocode}
%  \end{macro}
%
%
% \subsubsection{Sets of colors}
%
% \begin{macro}{\definecolorset}
%   \oarg{type}\marg{model-list}\marg{head}\marg{tail}\marg{set spec}\\
% This command facilitates the construction of \emph{color sets} with common underlying \Meta{model-list} and \Meta{type}.
% Here, \Meta{set spec} = \Meta[1]{name},\Meta[1]{spec-list};\dots;\Meta[k]{name},\Meta[k]{spec-list} ($k\ge1$ name/specification-list pairs).
% Individual colors are being constructed by |\definecolor|\oarg{type}|{|\Meta{head}\Meta[1]{name}\Meta{tail}|}|\marg{model-list}|{|\Meta[1]{spec-list}|}| etc.
%    \begin{macrocode}
\def\definecolorset
 {\let\XC@@cmd\definecolor\@testopt{\XC@definecolorset}{}}
%    \end{macrocode}
%
%    \begin{macrocode}
\def\XC@definecolorset[#1]#2#3#4#5%
 {\XC@getmodclr10{#2}{}{}%
   {\XC@@getmodfalse\XC@edef\XC@@tmp{#5}\let\XC@@glb\xglobal@
    \def\XC@definec@lorset##1,##2;%
     {\@ifxempty{##1}{}%
       {\let\xglobal@\XC@@glb
        \XC@@cmd[#1]{#3##1#4}{\@@mod}{##2}\XC@definec@lorset}}%
    \expandafter\XC@definec@lorset\XC@@tmp;,;\XC@@getmodtrue}}
%    \end{macrocode}
% \end{macro}
%
% \begin{macro}{\providecolorset}
%   \oarg{type}\marg{model-list}\marg{name-prefix}\marg{name-postfix}\marg{set spec}\\
% Similar to |\definecolorset|, but based on |\providecolor|, thus the individual colors are defined only if they do not exist already.
%    \begin{macrocode}
\def\providecolorset
 {\let\XC@@cmd\providecolor\@testopt{\XC@definecolorset}{}}
%    \end{macrocode}
% \end{macro}
%
% \begin{macro}{\preparecolorset}
%   \oarg{type}\marg{model-list}\marg{name-prefix}\marg{name-postfix}\marg{set spec}\\
% Similar to |\definecolorset|, but dependent on |\ifdefinecolors|, either a |\definecolor| or a |\preparecolor| command will be executed for each single set member.
%    \begin{macrocode}
\def\preparecolorset
 {\ifdefinecolors\let\XC@@cmd\definecolor\else\let\XC@@cmd\preparecolor\fi
  \@testopt{\XC@definecolorset}{}}
%    \end{macrocode}
% \end{macro}
%
%
% \subsubsection{Color series}
%
% \begin{macro}{\definecolorseries}
%   \marg{name}\marg{core model}\marg{method}\oarg{b-model}\marg{base color}\oarg{l-model}\marg{delta}
%    \begin{macrocode}
\def\definecolorseries#1#2#3%
 {\@testopt{\XC@definecolorseries{#1}{#2}{#3}}{}}
%    \end{macrocode}
%
%    \begin{macrocode}
\def\XC@definecolorseries#1#2#3[#4]#5%
 {\@testopt{\XC@definec@lorseries{#1}{#2}{#3}[#4]{#5}}{}}
%    \end{macrocode}
%
%    \begin{macrocode}
\def\XC@definec@lorseries#1#2#3[#4]#5[#6]#7%
 {\begingroup
  \XC@sdef\@@mod{#2}%
  \@ifxempty{#4}{\extractcolorspec{#5}\@@clr}{\edef\@@clr{{#4}{#5}}}%
  \expandafter\convertcolorspec\@@clr\@@mod\@@clr
  \@namexdef{\@backslashchar colorB@#1}{{\@@mod}{\@@clr}}%
  \edef\@@met{#3}%
  \ifx\@@met\XC@met@last
    \@ifxempty{#6}{\extractcolorspec{#7}\@@tmp}{\edef\@@tmp{{#6}{#7}}}%
    \expandafter\convertcolorspec\@@tmp\@@mod\@@tmp
    \XC@inflate\@@mod\@@tmp\@@clr\@@tmp
    \expandafter\XC@sub\@@tmp,\@@tmp
  \else
    \edef\@@tmp{#7}%
  \fi
  \@namexdef{\@backslashchar colorD@#1}{{\@@met}{\@@tmp}}%
  \endgroup}
%    \end{macrocode}
% \end{macro}
%
% \begin{macro}{\resetcolorseries}
%   \oarg{div}\marg{name}
%    \begin{macrocode}
\def\resetcolorseries{\@testopt{\XC@resetcolorseries}{}}
\def\XC@resetcolorseries[#1]#2%
 {\begingroup
  \@namexdef
   {\@backslashchar color@#2}{\noexpand\xcolor@{}{}\@nameuse{\@backslashchar colorB@#2}}%
  \XC@let@Nc\@@tmp{\@backslashchar colorD@#2}%
  \edef\@@met{\expandafter\@firstoftwo\@@tmp}%
  \edef\@@tmp{\expandafter\@secondoftwo\@@tmp}%
  \ifx\@@met\XC@met@step\else
    \edef\@@scl{\@ifxempty{#1}\colorseriescycle{#1}}%
    \expandafter\XC@calc@\@@tmp,,,,:D%
  \fi
  \@namexdef{\@backslashchar colorS@#2}{\@@tmp}%
  \endgroup}
%    \end{macrocode}
% \end{macro}
%
% \begin{macro}{\XC@step}
% \begin{macro}{\XC@stepabs}
%   \marg{name}\\
% Perform a relative or absolute (= direct access) step within a color series \Meta{name}.
% In the relative case, the (non-negative) number of single steps is supposed to be stored in |\@tempcnta|.
%    \begin{macrocode}
\def\XC@step#1%
 {\let\xcolor@\@gobbletwo\edef\@@tmp{\@nameuse{\@backslashchar color@#1}}%
  \expandafter\XC@st@p\@@tmp{#1}%
  \@namexdef{\@backslashchar color@#1}{\noexpand\xcolor@{}{}{\@@mod}{\@@tmp}}}
%    \end{macrocode}
%
%    \begin{macrocode}
\def\XC@st@p#1#2#3%
 {\def\@@mod{#1}\def\@@tmp{#2}\edef\@@vec{\@nameuse{\@backslashchar colorS@#3}}%
  \@whilenum\@tempcnta>\z@\do
   {\XC@inflate\@@mod\@@tmp\@@vec\@@tmp
    \expandafter\XC@add\@@tmp,\@@tmp \expandafter\XC@calc@\@@tmp,,,,:U%
    \advance\@tempcnta\m@ne}}
%    \end{macrocode}
%
%    \begin{macrocode}
\def\XC@stepabs#1%
 {\edef\@@tmp{\@nameuse{\@backslashchar colorB@#1}}%
  \expandafter\XC@st@p\@@tmp{#1}%
  \edef\@@tmp
   {\def\noexpand\@@mod{\@@mod}\def\noexpand\@@clr{\@@tmp}}}
%    \end{macrocode}
% \end{macro}
% \end{macro}
%
%
% \subsubsection{Color blending}
%
% \begin{macro}{\blendcolors}
% \begin{macro}{\blendcolors*}
%   \marg{mix expr}\\
% Initialises all necessary parameters for color blending.
% The actual (completed) color blend expression is stored in |\colorblend|.
% In the starred version, the argument will be appended to a previously defined blend expression.
% An empty \Meta{mix expr} argument will switch blending off.
%    \begin{macrocode}
\def\blendcolors{\@ifstar{\XC@blendcolors\colorblend}{\XC@blendcolors{}}}
\def\XC@blendcolors#1#2%
 {\edef\colorblend{#1}\XC@edef\XC@@tmp{#2}%
  \ifx\XC@@tmp\@empty
    \xglobal@\blendcolorsfalse
  \else
    \xglobal@\blendcolorstrue
    \expandafter\XC@blendc@lors\XC@@tmp!!\@nil
  \fi
  \xglobalize@\colorblend \xglobal@stop}
%    \end{macrocode}
%
%    \begin{macrocode}
\def\XC@blendc@lors!#1!#2!#3\@nil
 {\@ifxempty{#1}{}%
   {\@ifxempty{#2}%
    {\edef\colorblend{\colorblend!#1!white}}%
    {\edef\colorblend{\colorblend!#1!#2}\XC@blendc@lors!#3!\@nil}}}
%    \end{macrocode}
%
%    \begin{macrocode}
\blendcolors{}
%    \end{macrocode}
% \end{macro}
% \end{macro}
%
%
% \subsubsection{Color masks and separation}
%
% \begin{macro}{\maskcolors}
%   \oarg{core model}\marg{color}\\
% Initialises all necessary parameters for color masking.
% The actual color mask expression is stored in |\colormask|.
% An empty \Meta{color} argument will switch masking off.
% The optional \Meta{core model} argument may be used to determine the underlying color model.
%    \begin{macrocode}
\def\maskcolors{\@testopt{\XC@maskcolors}{}}
\def\XC@maskcolors[#1]#2%
 {\XC@edef\colormask{#2}%
  \ifx\colormask\@empty
    \xglobal@\maskcolorsfalse \def\XC@mmod{}\def\XC@mclr{}%
  \else
    \xglobal@\maskcolorstrue \extractcolorspecs\colormask\XC@mmod\XC@mclr
    \@ifxempty{#1}{}%
     {\convertcolorspec\XC@mmod\XC@mclr{#1}\XC@mclr \XC@sdef\XC@mmod{#1}}%
    \XC@coremodel\XC@mmod\XC@mclr
  \fi
  \xglobalize@\colormask \xglobalize@\XC@mmod \xglobalize@\XC@mclr
  \xglobal@stop}
%    \end{macrocode}
%
%    \begin{macrocode}
\maskcolors{}
%    \end{macrocode}
% \end{macro}
%
%
%
% \subsection{Color application}
%
% \subsubsection{Standard color commands}
%
%  \begin{macro}{\color}
%   \marg{color}\\
%   \oarg{model-list}\marg{spec-list}\\
% Switches to the color given either by name/expression or by model/specification.
% This color will stay in effect until the end of the current \TeX\ group.
%    \begin{macrocode}
\DeclareRobustCommand\color
 {\@ifnextchar[\@undeclaredcolor\@declaredcolor}
%    \end{macrocode}
%  \end{macro}
%
% \begin{macro}{\@undeclaredcolor}
%   \oarg{model-list}\marg{spec-list}\\
% Display a color that is given by \Meta{model} name and numerical \Meta{spec}.
%    \begin{macrocode}
\def\@undeclaredcolor[#1]#2%
 {\begingroup
  \let\@@cls\@empty
  \XC@getmodclr12{#1}{#2}{}%
   {\ifblendcolors
      \ifx\colorblend\@empty\else
        \edef\@@mix{\expandafter\@gobble\colorblend}\@tempswafalse
        \XC@coremodel\@@mod\@@clr\@xcolor@{}{}\@@mod\@@clr
      \fi
    \fi
    \ifconvertcolorsU
      \edef\@@tmp{\XC@tgt@mod{\@@mod}}%
      \convertcolorspec\@@mod\@@clr\@@tmp\@@clr \let\@@mod\@@tmp
    \fi
    \ifmaskcolors
      \convertcolorspec\@@mod\@@clr\XC@mmod\@@clr
      \let\@@mod\XC@mmod
      \XC@inflate\@@mod\XC@mclr\@@clr\@@tmp
      \expandafter\XC@mul\@@tmp,\@@clr
    \fi
    \edef\@@tmp{\noexpand\XC@undeclaredcolor{\@@mod}{\@@clr}}%
    \expandafter\endgroup\@@tmp}}
%    \end{macrocode}
%
%    \begin{macrocode}
\def\XC@undeclaredcolor#1#2%
 {\@ifundefinedmodel{#1}%
   {\c@lor@error{model `#1'}}%
   {\@nameuse{color@#1}\current@color{#2}%
    \edef\XC@current@color{\noexpand\xcolor@{}{\current@color}{#1}{#2}}%
    \XC@display}}
%    \end{macrocode}
% \end{macro}
%
% \begin{macro}{\@declaredcolor}
%   \marg{color}\\
% Display a color that is given by an arbitrary \Meta{color}.
%    \begin{macrocode}
\def\@declaredcolor#1%
 {\XC@edef\XC@@tmp{#1}\XC@@tstfalse
  \ifcase\XC@type\XC@@tmp\relax \XC@@tsttrue\or \relax\else
    \ifblendcolors\XC@@tsttrue\fi \ifmaskcolors\XC@@tsttrue\fi
    \ifconvertcolorsU\XC@@tsttrue\fi
  \fi
  \ifXC@@tst
    \expandafter\XC@declaredcolor\else
    \expandafter\XC@declaredc@lor\fi}
%    \end{macrocode}
%
%    \begin{macrocode}
\def\XC@declaredcolor
 {\XC@split\XC@@tmp
  \edef\@@tmp{\noexpand\@undeclaredcolor[\@@mod]{\@@clr}}%
  \expandafter\endgroup\@@tmp}
%    \end{macrocode}
%
%    \begin{macrocode}
\def\XC@declaredc@lor
 {\XC@let@Nc\XC@current@color{\@backslashchar color@\XC@@tmp}%
  \edef\current@color{\XC@current@color}%
  \XC@display}
%    \end{macrocode}
% \end{macro}
%
%  \begin{macro}{\textcolor}
%   \marg{color}\marg{text}\\
%   \oarg{model-list}\marg{spec-list}\marg{text}\\
% are just alternative syntax for |\color|, in which the groups are added implicitly.
% Thus \emph{text} appears in the specified color, but then the color reverts to its previous value.
% The naming is by analogy with |\textrm| (as opposed to |\rm| and |\rmfamily|) although it is slightly a misnomer as the command also works in math-mode.
% Additionally, it calls |\leavevmode| to ensure the start of horizontal mode.
%    \begin{macrocode}
\protected\def\textcolor#1#{\@textcolor{#1}}
\def\@textcolor#1#2#3{\protect\leavevmode{\color#1{#2}#3}}
%    \end{macrocode}
%  \end{macro}
%
% \begin{macro}{\pagecolor}
%   \marg{color}\\
%   \oarg{model-list}\marg{spec-list}\\
% Specifies the background color for the current, and all following, pages.
% It is a global declaration which does not respect \TeX\ groups.
%    \begin{macrocode}
\protected\def\pagecolor
 {\begingroup
  \let\ignorespaces\endgroup \let\set@color\set@page@color
  \color}
%    \end{macrocode}
%  \end{macro}
%
% \begin{macro}{\nopagecolor}
% Removes any currently specified page color returning to the default transparent background.
% It is not yet supported by all driver options and so generates a warning if there is no definition in the driver file.
%    \begin{macrocode}
\protected\def\nopagecolor
 {\@ifundefined{no@page@color}%
   {\PackageInfo{xcolor}{\@backslashchar nopagecolor\space is not supported}}%
   {\no@page@color}}
%    \end{macrocode}
%  \end{macro}
%
% \subsubsection{mathcolor}
% We load the new mathcolor support from the kernel if it exists. The documentation
% is in \texttt{mathcolor.pdf}
%
\InputIfFileExists{mathcolor.ltx}{}{}
%
% \subsubsection{Colored boxes}
%
% \begin{macro}{\colorbox}
%   \marg{color}\marg{text}\\
%   \oarg{model}\marg{spec}\marg{text}\\
% Takes the same argument forms as |\textcolor|, but the color specifies the \emph{background} color of the box.
%    \begin{macrocode}
\protected\def\colorbox#1#{\color@box{#1}}
\def\color@box#1#2{\color@b@x\relax{\color#1{#2}}}
%    \end{macrocode}
% \end{macro}
%
% \begin{macro}{\fcolorbox}
%   \marg{frame color}\marg{background color}\marg{text}\\
%   \oarg{model}\marg{frame spec}\marg{background spec}\marg{text}\\
%   \oarg{frame model}\marg{frame spec}\oarg{background model}\marg{background spec}\marg{text}\\
%   \marg{frame color}\oarg{background model}\marg{background spec}\marg{text}\\
% Puts a frame of the first color around a box with a background specified by the second color.
% If only the first optional argument is given, it specifies the color model for both colors.
%    \begin{macrocode}
\protected\def\fcolorbox#1#{\color@fbox{#1}}
\def\color@fbox#1#2#3#{\protect\color@fb@x{#1}{#2}{#3}}
\def\color@fb@x#1#2#3#4%
 {\color@b@x{\fboxsep\z@\color#1{#2}\XC@fbox}%
   {\@ifxempty{#3}{\color#1}{\color#3}{#4}}}
%    \end{macrocode}
% \end{macro}
%
%  \begin{macro}{\color@b@x}
% Internal macro for |\colorbox| and |\fcolorbox|.
%    \begin{macrocode}
\long\def\color@b@x#1#2#3%
 {\leavevmode
  \setbox\z@\hbox{\kern\fboxsep{\set@color#3}\kern\fboxsep}%
  \dimen@\ht\z@\advance\dimen@\fboxsep\ht\z@\dimen@
  \dimen@\dp\z@\advance\dimen@\fboxsep\dp\z@\dimen@
  {#1{#2\color@block{\wd\z@}{\ht\z@}{\dp\z@}\box\z@}}}
%    \end{macrocode}
%  \end{macro}
%
%  \begin{macro}{\XC@frameb@x}
% Common part of |\framebox| and |\fbox|.
% |#1| is a negative kern in the |\framebox| case so that the vertical rules do not add to the width of the box.
% Optimised by \People{Donald}{Arseneau}.
%    \begin{macrocode}
\def\XC@frameb@x#1%
 {\hbox
   {\@tempdima\fboxrule \advance\@tempdima\fboxsep
    \advance\@tempdima\dp\@tempboxa
    \setbox\@tempboxa\hbox
     {\lower\@tempdima\hbox
       {\vbox
         {\kern\fboxrule
          \hbox{\kern\fboxrule#1%
                \vbox{\kern\fboxsep\box\@tempboxa\kern\fboxsep}%
                #1\kern\fboxrule}%
          \kern\fboxrule}}}%
    \@tempdima\wd\@tempboxa
    \edef\@tempa{{\the\@tempdima}{\the\ht\@tempboxa}{\the\dp\@tempboxa}}%
    \box\@tempboxa\kern-\@tempdima\expandafter\boxframe\@tempa}}
%    \end{macrocode}
%  \end{macro}
%
% \begin{macro}{\XC@fbox}
% \begin{macro}{\@frameb@x}
% If no \Option{kernelfbox} option was specified, switch to the new method in general, else only for |\fcolorbox|.
%    \begin{macrocode}
\ifXC@@fbox
  \let\@frameb@x\XC@frameb@x \let\XC@fbox\fbox
\else
  \def\XC@fbox{\let\@frameb@x\XC@frameb@x\fbox}
\fi
%    \end{macrocode}
%  \end{macro}
%  \end{macro}
%
%
% \subsubsection{Driver-specific and generic drawing}
%
% \begin{macro}{\color@block}
%   \marg{width}\marg{height}\marg{depth}\\
% Should take up no space for \TeX, but produce a block in the current color of the specified size.
% It is mainly used for producing box backgrounds.
%  \begin{macro}{\boxframe}
%   \marg{width}\marg{height}\marg{depth}\\
% Draw a frame with a linewidth of |\fboxrule|.
% Return a |\hbox| with outer dimensions \meta{width}, \meta{height}, \meta{depth}.

%  \begin{macro}{\XC@boxframe}
%   \marg{width}\marg{height}\marg{depth}\marg{frame construction}\\
% Wrapper for |\boxframe| commands that return zero-dimension frames.
%    \begin{macrocode}
\ifnum\XC@@xcd@>\m@ne
\def\XC@boxframe#1#2#3#4%
 {\hbox{\sbox\z@{\ifdim\fboxrule>\z@ #4\fi}%
    \wd\z@#1\relax\ht\z@#2\relax\dp\z@#3\relax\box\z@}}
\fi
%    \end{macrocode}
% \end{macro}
%
% \textbf{\Option{dvips}-optimised code}
%    \begin{macrocode}
\ifcase\XC@@xcd@
%    \end{macrocode}
%
%    \begin{macrocode}
\providecommand*\color@block[3]%
 {{\set@color\hbox
    {\dimen@#3\edef\@tempa{ \number\dimen@\space}%
     \dimen@#2\edef\@tempa{ \number\dimen@\@tempa}%
     \dimen@#1\relax
     \XC@special{\number\dimen@\@tempa XCcolorblock}}}}
%    \end{macrocode}
%
%    \begin{macrocode}
\providecommand*\boxframe[3]%
 {\XC@boxframe{#1}{#2}{#3}%
   {\dimen@#3\edef\@tempa{ \number\dimen@\space}%
    \dimen@#2\edef\@tempa{ \number\dimen@\@tempa}%
    \dimen@#1\relax
    \XC@special{\number\fboxrule\space\number\dimen@\@tempa XCboxframe}}}
%    \end{macrocode}
%
% \textbf{\Option{pdftex}-optimised code }
% There seems to be no improvement by adding a specific |\color@block| command, since the PDF format doesn't offer higher accuracy than \TeX.
% For |\boxframe| we want to achieve this PDF pseudo-code (besides initial scaling):\\[1ex]
% \centerline{\ttfamily $l$ w [] $0$ d $2$ J $0$ j $4$ M
% $1$ $0$ $0$ $1$ $\frac12\,l$ $\frac12\,l$ cm
% $0$ $-d$ $(w-l)$ $(h+d-l)$ re s}
%    \begin{macrocode}
\or
%    \end{macrocode}
%
%    \begin{macrocode}
\providecommand*\boxframe[3]%
 {\XC@boxframe{#1}{#2}{#3}%
   {\dimen@#3\edef\@tempa{ 0 -\strip@pt\dimen@\space}%
    \edef\@tempb{\advance\dimen@\the\dimen@}%
    \dimen@#2\advance\dimen@-\fboxrule\@tempb
    \edef\@tempb{ \strip@pt\dimen@\space re s}%
    \dimen@#1\advance\dimen@-\fboxrule
    \edef\@tempb{ cm\@tempa\strip@pt\dimen@\@tempb}%
    \dimen@.5\fboxrule\edef\@tempa{ \strip@pt\dimen@}%
    \XC@special
     {0.996264 0 0 0.996264 0 0 cm
      \strip@pt\fboxrule\space w [] 0 d 2 J 0 j 4 M
      1 0 0 1\@tempa\@tempa\@tempb}}}
%    \end{macrocode}
%
%    \begin{macrocode}
\fi
%    \end{macrocode}
%
% \textbf{Generic code }
% The definition here works by selecting a color, and then drawing a \TeX\ rule (unless |\ifcolors@false|).
%    \begin{macrocode}
\providecommand*\color@block[3]%
 {{\set@color\rlap{\ifcolors@\vrule\@width#1\@height#2\@depth#3\fi}}}
%    \end{macrocode}
%
% The following definition was optimised by \People{Donald}{Arseneau}.
%    \begin{macrocode}
\providecommand*\boxframe[3]%
 {\hbox
   {\dimen@#2\advance\dimen@#3\relax
    \lower#3\vbox
     {\hrule\@height\fboxrule
      \@tempdima-0.5\fboxrule \ifodd\fboxrule\advance\@tempdima\m@ne sp\fi
      \kern\@tempdima
      \hbox
       {\advance\dimen@-\fboxrule
        \vrule\@width\fboxrule\@height\dimen@\@depth\z@
        \@tempdima#1\advance\@tempdima-\tw@\fboxrule \kern\@tempdima
        \vrule\@width\fboxrule\@height\dimen@\@depth\z@}%
      \kern-0.5\fboxrule
      \hrule\@height\fboxrule}}}
%    \end{macrocode}
% \end{macro}
% \end{macro}
%
%
% \subsubsection{Color argument processing}
%
% \begin{macro}{\@xcolor@}
%   \marg{type-spec}\marg{driver-spec}\marg{model}\marg{spec}\\
% This is the `work-horse' that does all necessary conversions and calculations if a color (expression) has to be displayed.
% \Meta{type-spec}  and \Meta{driver-spec} are ignored.
% |\@@mix| may contain a mix expression without leading `!'.
%    \begin{macrocode}
\def\@xcolor@#1#2#3#4%
 {\XC@sdef\@@mod{#3}\edef\@@clr{#4}%
  \ifx\@@mix\@empty\else\expandafter\@xcolor@@\@@mix!!!!\@nil\fi
  \if@tempswa
    \XC@let@Nc\@@tmp{XC@cnv@\@@mod @compl}%
    \ifx\@@tmp\relax \expandafter\XC@calc@\@@clr,,,,:C\else
                     \expandafter\@@tmp\@@clr,\fi
    \let\@@clr\@@tmp
  \fi}
%    \end{macrocode}
%
%    \begin{macrocode}
\def\@xcolor@@#1!#2!#3!#4!#5\@nil
 {\dimen@\@ifxempty{#1}{100}{#1}\p@
  \edef\@@mix{\@ifxempty{#2}{white}{#2}}%
  \ifdim\dimen@=100\p@
    \expandafter\XC@postfix\@@mix]\relax
  \else\ifdim\dimen@=\z@
    \extractcolorspecs\@@mix\@@mod\@@clr
  \else
    \XC@let@Nc\@@tmp{XC@clr@\@@mod @\@@mix}%
    \ifx\@@tmp\relax
      \extractcolorspec\@@mix\@@tmp
      \ifx\@@mod\XC@mod@gray
        \edef\@@mod{\expandafter\@firstoftwo\@@tmp}%
        \edef\@@tmp{\expandafter\@secondoftwo\@@tmp}%
        \convertcolorspec\XC@mod@gray\@@clr\@@mod\@@clr
      \else
        \expandafter\convertcolorspec\@@tmp\@@mod\@@tmp
      \fi
    \else
      \ifx\@@mod\XC@mod@hsb\expandafter\XC@extract@hue\@@clr:\@@hue\fi
    \fi
    \XC@inflate\@@mod\@@clr\@@tmp\@@tmp
    \expandafter\XC@mix\@@tmp,\@@clr
  \fi\fi
  \@ifxempty{#3}%
   {\@ifxempty{#4}{}{\@xcolor@@!#4!#5\@nil}}%
   {\@xcolor@@#3!#4!#5\@nil}}
%    \end{macrocode}
% \end{macro}
%
% \begin{macro}{\XC@split}
%   \marg{ext expr}\\
% Split \Meta{ext expr} into all of its components, do the necessary calculations, and save the model of the resulting color in |\@@mod| and the parameters in |\@@clr|.
% Note: |\XC@split| opens a group which has to be closed afterwards by |\endgroup|.
%    \begin{macrocode}
\def\XC@split#1%
 {\begingroup\XC@edef\@@tmp{#1}\expandafter\XC@split@i\@@tmp>>\@nnil}
%    \end{macrocode}
%
% \begin{macro}{\XC@split@i}
% Separate first part of expression chain and apply function(s) afterwards as requested.
%    \begin{macrocode}
\def\XC@split@i#1>%
 {\def\@@tmp{#1}\XC@split@ii#1:\@nil \XC@func}
%    \end{macrocode}
% \end{macro}
%
% \begin{macro}{\XC@split@ii}
% Check whether it is an \emph{extended} or \emph{standard} color expression.
%    \begin{macrocode}
\def\XC@split@ii#1:#2\@nil
 {\@ifxempty{#2}{\XC@split@vi}{\XC@split@iii#1,,:#2}}
%    \end{macrocode}
% \end{macro}
%
% \begin{macro}{\XC@split@iii}
% Initialisation for extended color expressions.
% The magic \Meta{div} is either given as second argument, or calculated as the sum of all weights.
%    \begin{macrocode}
\def\XC@split@iii#1,#2,#3:#4:%
 {\edef\XC@@mod{#1}\let\@tempa\@empty
  \ifx,#2,\XC@@tsttrue\@tempdima\z@\else\XC@@tstfalse\@tempdima#2\p@\fi
  \XC@split@iv#4;,;}
%    \end{macrocode}
% \end{macro}
%
% \begin{macro}{\XC@split@iv}
%   \marg{expr}\marg{dec}\\
% Step through the list of \Meta{expr},\Meta{dec} pairs.
% Add \Meta{dec} to the sum of weights (in order to calculate the ultimate \Meta{div}) and append pair information to |\@tempa| command.
% At the end of the list, execute |\@tempa| and map resulting parameter vector into unit cube.
% Finally, |\@@mod| and |\@@clr| contain the desired information about color model and parameters.
%    \begin{macrocode}
\def\XC@split@iv#1,#2;%
 {\@ifxempty{#1}%
   {\let\XC@@clr\@empty \@tempa \expandafter\XC@calc@\XC@@clr,,,,:U%
    \let\@@mod\XC@@mod \let\@@clr\@@tmp}
   {\XC@append\@tempa{\XC@split@v{#1}{#2}}%
    \ifXC@@tst\advance\@tempdima#2\p@\fi
    \XC@split@iv}}
%    \end{macrocode}
% \end{macro}
%
% \begin{macro}{\XC@split@v}
%   \marg{expr}\marg{dec}\\
% Extract color specification from \Meta{expr}, convert to target model, multiply the resulting vector by the quotient of \Meta{dec} and \Meta{div}, add that vector to the overall parameter vector.
%    \begin{macrocode}
\def\XC@split@v#1#2%
 {\def\@@tmp{#1}\XC@split@vi
  \ifx\@@mod\XC@@mod\else
    \convertcolorspec\@@mod\@@clr\XC@@mod\@@clr
  \fi
  \dimen@#2\p@ \rdivide\dimen@\@tempdima
  \edef\@@scl{\strip@pt\dimen@}%
  \expandafter\XC@calc@\@@clr,,,,:S%
  \ifx\XC@@clr\@empty\let\XC@@clr\@@tmp\else
    \XC@inflate\XC@@mod\XC@@clr\@@tmp\@@tmp
    \expandafter\XC@add\@@tmp,\XC@@clr
  \fi}
%    \end{macrocode}
% \end{macro}
%
% \begin{macro}{\XC@split@vi}
% Start splitting process for standard color expression.
%    \begin{macrocode}
\def\XC@split@vi{\count@\z@ \expandafter\XC@split@vii\@@tmp!\@nil}
%    \end{macrocode}
% \end{macro}
%
% \begin{macro}{\XC@split@vii}
% Count number of `-' signs at the beginning of the string.
%    \begin{macrocode}
\def\XC@split@vii#1%
 {\ifx#1-\advance\count@\@ne\expandafter\XC@split@vii
  \else\expandafter\XC@split@viii\expandafter#1\fi}
%    \end{macrocode}
% \end{macro}
%
% \begin{macro}{\XC@split@viii}
% Extract \Meta{name} and \Meta{mix expr} part of color expression.
% Delegate control to |\@xcolor@|.
%    \begin{macrocode}
\def\XC@split@viii#1!#2\@nil
 {\edef\@@nam{#1}\edef\@@mix{#2}%
  \ifodd\count@\@tempswatrue\else\@tempswafalse\fi
  \@ifundefinedcolor\@@nam
   {\c@lor@error{`\@@nam'}\def\@@nam{black}}{}%
  \let\xcolor@\@xcolor@\@nameuse{\@backslashchar color@\@@nam}}
%    \end{macrocode}
% \end{macro}
% \end{macro}
%
% \begin{macro}{\XC@postfix}
% Examine the postfix of a color expression and derive the actions to be taken for color series.
%    \begin{macrocode}
\def\XC@postfix#1#2]#3\relax
 {\begingroup
  \if[#1\@tempcnta#2\relax\XC@stepabs\@@nam
  \else\@tempcnta\z@
    \@tfor\@@tok:=#1#2\do
     {\if+\@@tok\advance\@tempcnta\@ne
      \else\if-\@@tok\advance\@tempcnta\m@ne
      \else\@tempcnta\z@\expandafter\@break@tfor\fi\fi}%
    \ifnum\@tempcnta=\z@\relax\else\XC@step\@@nam\fi
    \let\@@tmp\@empty
  \fi\expandafter\endgroup\@@tmp}
%    \end{macrocode}
% \end{macro}
%
% \begin{macro}{\XC@func}
%   \marg{function}|,{|\Meta[1]{arg}|,...,{|\Meta[n]{arg}|}>...\@nnil>|\\
% General macro to handle color functions.
% The number of arguments depends on the function.
% Multiple consecutive functions may be applied.
% Color model and specifications are input and output via |\@@mod| and |\@@clr|.
%    \begin{macrocode}
\def\XC@func#1>{\@ifxempty{#1}{\remove@to@nnil}{\XC@fun@#1,>\XC@func}}
\def\XC@fun@#1,#2>{\@nameuse{XC@func@#1}#2,,>}
%    \end{macrocode}
% \end{macro}
%
% \begin{macro}{\XC@func@wheel}
% \begin{macro}{\XC@func@twheel}
%   \marg{angle}|,|\marg{full circle}|,>|\\
% Calculate color wheel position according to \Model{Hsb} or \Model{tHsb} model: \emph{hue} is rotated by \meta{angle} relative to \meta{full circle}.
% The latter value is taken from |\rangeHsb| if left empty.
% Finally, |\@@mod| is changed to \Model{hsb}.
%    \begin{macrocode}
\def\XC@func@wheel#1,#2,#3>{\XC@func@wh@el{#1}{#2}0}
%    \end{macrocode}
%
%    \begin{macrocode}
\def\XC@func@twheel#1,#2,#3>{\XC@func@wh@el{#1}{#2}1}
%    \end{macrocode}
%
%    \begin{macrocode}
\def\XC@func@wh@el#1#2#3%
 {\convertcolorspec\@@mod\@@clr
   {\ifcase#3 \XC@mod@Hsb\or\XC@mod@tHsb\fi}\@@clr
  \expandafter\XC@extract@hue\@@clr:\@@hue \dimen@#1\p@
  \@ifxempty{#2}{}%
   {\rdivide\dimen@{#2}\lshift\dimen@
    \dimen@\rangeHsb\dimen@\rshift\dimen@}%
  \advance\dimen@\@@hue\p@ \rdivide\dimen@\rangeHsb
  \edef\@@hue{\strip@pt\dimen@}\XC@calcU\@@hue\@@hue
  \ifcase#3
    \expandafter\XC@replace@hue\@@clr:\@@hue,\@@clr
  \or
    \XC@c@lcS\rangeHsb{}\@@hue\@@hue
    \expandafter\XC@replace@hue\@@clr:\@@hue,\@@clr
    \convertcolorspec\XC@mod@tHsb\@@clr\XC@mod@hsb\@@clr
  \fi
  \let\@@mod\XC@mod@hsb}
%    \end{macrocode}
% \end{macro}
% \end{macro}
%
% \begin{macro}{\XC@display}
% Write information to log file and display the current color.
%    \begin{macrocode}
\def\XC@display
 {\XC@bcolor\XC@let@cN{\@backslashchar color@.}\XC@current@color
  \ifnum\XC@type{.}=1
    \PackageWarning{xcolor}{Incompatible color definition}\else
  \ifnum\XC@tracing>2
    \begingroup \let\xcolor@\@empty \def\@nil{\@backslashchar @nil}%
    \PackageInfo{xcolor}{Setting color `\XC@current@color'}%
    \endgroup
  \fi\fi
  \XC@mcolor\set@color\XC@ecolor\ignorespaces}
%    \end{macrocode}
% \end{macro}
%
% \begin{macro}{\XC@bcolor}
% \begin{macro}{\XC@mcolor}
% \begin{macro}{\XC@ecolor}
% Hooks that may be used to execute additional code within |\XC@display| immediately before/after a color is being displayed.
% Note that |\XC@bcolor| can use the current color information given by |\current@color| and |\XC@current@color| if further manipulation of the color is necessary.
% These hooks are intended for class/package writers and not used by \XC{} itself.
% In order to prevent unwanted side effects, please make sure that all internal \XC{} variables are restored properly at the end of |\XC@display|.
%    \begin{macrocode}
\let\XC@bcolor\relax
\let\XC@mcolor\relax
\let\XC@ecolor\relax
%    \end{macrocode}
% \end{macro}
% \end{macro}
% \end{macro}
%
%
% \subsubsection{Border colors for hyperlinks}
%
% \begin{macro}{\XC@bordercolor}
%   \marg{type}\marg{color}\\
% Basic macro to transform \Meta{color} into the format required by \Package{hyperref}.
% Here, \Meta{type} corresponds to the key `x\Meta{type}bordercolor'.
%    \begin{macrocode}
\def\XC@bordercolor#1#2%
 {\extractcolorspec{#2}\XC@@tmp
  \expandafter\convertcolorspec\XC@@tmp\XC@mod@rgb\XC@@tmp
  \expandafter\edef\csname @#1bordercolor\endcsname
   {\expandafter\XC@strip@comma\XC@@tmp,,\@nnil}}
%    \end{macrocode}
% \end{macro}
%
% \begin{macro}{\XC@pdfborder}
%   \Meta[1]{number}| |\Meta[2]{number}| |\Meta[3]{number}|\relax|\\
% Basic macro to transform a pdfborder expression into a suitable format for \Program{dvips} with a resolution of 8000~dpi.
%    \begin{macrocode}
\def\XC@pdfborder#1 #2 #3\relax
 {\def\@@scl{80}\XC@calc@#1,#2,#3,,:S%
  \edef\@pdfborder{\expandafter\XC@strip@comma\@@tmp,,\@nnil}}
%    \end{macrocode}
% \end{macro}
%
%
%
%
% \subsubsection{Color information}
%
% \begin{macro}{\extractcolorspec}
%   \marg{color}\marg{cmd}\\
% Extracts the color specification of \Meta{color} and puts it into \Meta{cmd}; equivalent to
% |\def|\Meta{cmd}|{|\marg{model}\marg{spec}|}|.
%    \begin{macrocode}
\def\extractcolorspec#1#2%
 {\XC@split{#1}\edef\@@tmp{{\@@mod}{\@@clr}}\aftergroupdef#2\@@tmp}
%    \end{macrocode}
% \end{macro}
%
% \begin{macro}{\extractcolorspecs}
%   \marg{color}\marg{model-cmd}\marg{color-cmd}\\
% Extracts the color specification of \Meta{color} and puts it into \Meta{model-cmd} and \Meta{color-cmd}, respectively.
%    \begin{macrocode}
\def\extractcolorspecs#1#2#3%
 {\extractcolorspec{#1}#3%
  \edef#2{\expandafter\@firstoftwo#3}%
  \edef#3{\expandafter\@secondoftwo#3}}
%    \end{macrocode}
% \end{macro}
%
% \begin{macro}{\XC@extract@hue}
% Extract first entry of comma-separated value list.
%    \begin{macrocode}
\def\XC@extract@hue#1,#2:#3{\edef#3{#1}}
%    \end{macrocode}
% \end{macro}
%
% \begin{macro}{\XC@extract@hue}
% Change first entry of comma-separated value list.
%    \begin{macrocode}
\def\XC@replace@hue#1,#2:#3,#4{\edef#4{#3,#2}}
%    \end{macrocode}
% \end{macro}
%
% \begin{macro}{\@ifundefinedcolor}
%   \marg{name}\marg{undefined}\marg{defined}\\
% Checks whether \Meta{name} is a defined color name and executes the respective code.
%    \begin{macrocode}
\def\@ifundefinedcolor#1{\@ifundefined{\@backslashchar color@#1}}
%    \end{macrocode}
% \end{macro}
%
% \begin{macro}{\@ifundefinedmodel}
%   \marg{model}\marg{undefined}\marg{defined}\\
% Checks whether \Meta{model} is a defined color model and executes the respective code.
%    \begin{macrocode}
\def\@ifundefinedmodel#1{\@ifundefined{color@#1}}
%    \end{macrocode}
% \end{macro}
%
% \begin{macro}{\XC@type}
%   \marg{name}\\
% Determines the \emph{type} of color \Meta{name}, expressed as a natural number (here, $t_1,t_2,\dots$ denote the tokens of |\\color@|\Meta{name}):
% \begin{center}
% \begin{tabular}{cll}
% \em type& \em description& \em conditions\\
% 0 & undefined\\
% 1 & unidentified& $t_1\neq|\xcolor@|$\\
% 2 & named & $t_1=|\xcolor@|\quad t_2\neq\emptyset$\\
% 3 & declared & $t_1=|\xcolor@|\quad t_2=\emptyset\quad t_3\neq\emptyset$\\
% 4 & undeclared & $t_1=|\xcolor@|\quad t_2=\emptyset\quad t_3=\emptyset$\\
% \end{tabular}
% \end{center}
%    \begin{macrocode}
\def\XC@type#1%
 {\expandafter\expandafter\expandafter\XC@typ@
  \csname\@backslashchar color@#1\endcsname\@empty\@empty\@empty\XC@@}
%    \end{macrocode}
%
%    \begin{macrocode}
\def\XC@typ@#1#2#3#4\XC@@
 {\ifx#1\relax 0\else
    \ifx#1\xcolor@
      \ifx$#2$%
        \ifx$#3$4\else3\fi\@gobbletwo
      \else2\fi\@gobbletwo
    \else1\fi
  \fi}
%    \end{macrocode}
% \end{macro}
%
% \begin{macro}{\XC@info}
%   \marg{color}\marg{cmd}\\
% Determines the \emph{type} of color \Meta{color}, expressed as a natural number; the result is stored in \Meta{cmd}.
% \begin{center}
% \begin{tabular}{cll}
% \em type& \em description& \em examples\\
% 0 & implicit name& `|.|'\\
% 1 & explicit name& `|red|', `|green|'\\
% 2 & standard color expression& `|-red|', `|green!50|'\\
% 3 & extended color expression& `|rgb:-red,1;green!50,2|'\\
% 4 & color function expression& `|rgb:-red,1;green!50,2>wheel,30|'\\
% \end{tabular}
% \end{center}
%    \begin{macrocode}
\def\XC@info#1#2%
 {\XC@edef#2{#1}%
  \expandafter\in@\expandafter>\expandafter{#2}%
  \ifin@\def#2{4}\else\expandafter\in@\expandafter:\expandafter{#2}%
  \ifin@\def#2{3}\else\expandafter\in@\expandafter!\expandafter{#2}%
  \ifin@\def#2{2}\else % now there's no >,:,!
  \if-\expandafter\@car#2\relax\@nil\def#2{2}\else % catches -red
  \edef#2{\noexpand\in@{/./}{/#2/}}#2%
  \edef#2{\ifin@0\else1\fi}\fi\fi\fi\fi}
%    \end{macrocode}
% \end{macro}
%
% \begin{macro}{\tracingcolors}
% Put the integer following |\tracingcolors| into the macro |\XC@tracing|.
%    \begin{macrocode}
\def\tracingcolors{\begingroup\afterassignment\XC@tracingcolors\count@}
\def\XC@tracingcolors{\aftergroupdef\XC@tracing{\the\count@}}
%    \end{macrocode}
%
%    \begin{macrocode}
\tracingcolors=0
%    \end{macrocode}
% \end{macro}
%
% \begin{macro}{\XC@logdef}
%   \marg{name}\\
% Issue a log file information if color \Meta{name} is being (re)defined.
%    \begin{macrocode}
\def\XC@logdef#1%
 {\ifnum\XC@tracing>1
    \@ifundefinedcolor{#1}%
     {\ifnum\XC@tracing>2
        \PackageInfo{xcolor}{Defining color `#1'}%
      \fi}%
     {\PackageInfo{xcolor}{Redefining color `#1'}}%
  \fi}
%    \end{macrocode}
% \end{macro}
%
%
% \subsubsection{Color testing}
%
% \begin{environment}{testcolors}
%   \oarg{num models}\\
% This is a simple tabular environment in order to test (display) colors in different models, showing both the visual result and the model-specific parameters.
% The optional \Meta{num models} argument is a comma-separated list of \emph{numerical} color models.
%    \begin{macrocode}
\newenvironment*{testcolors}[1][rgb,cmyk,hsb,HTML]%
 {\let\@@nam\@empty\count@\z@
  \@for\@@tmp:=#1\do
    {\advance\count@\@ne
     \XC@sdef\@@tmp{\@@tmp}\edef\@@nam{\@@nam{\@@tmp}}}%
  \edef\@@num{\the\count@}%
%    \end{macrocode}
%
% Check whether the actual font belongs to a set of OT1-encoded monospaced fonts in order to display |>| instead of |\textgreater|.
% \begin{macro}{\XC@@gt}
%    \begin{macrocode}
  \def\XC@@gt{\textgreater}\def\@@tmp{OT1}%
  \ifx\f@encoding\@@tmp
    \@expandtwoargs\in@{,\f@family,}{,cmtt,pcr,}%
    \ifin@\def\XC@@gt{>}\fi
  \fi
%    \end{macrocode}
%  \end{macro}
%
% Disable prologue and tracing commands for temporary color definitions.
%    \begin{macrocode}
  \def\XC@@xcp@{-1}\ifnum\XC@tracing>1 \def\XC@tracing{1}\fi
%    \end{macrocode}
%
% \begin{macro}{\testcolor}
%   \marg{color}\\
%   \oarg{model-list}\marg{spec-list}\\
% Each |\testcolor| command generates a table row.
% If the column-model matches the model of the color in question, its parameters are underlined.
%    \begin{macrocode}
  \def\testcolor{\@testopt{\@testcolor}{}}%
  \def\@testcolor[##1]##2%
   {\XC@mdef\@@mod{##1}\XC@edef\@@clr{##2}%
    \ifx\@@mod\@empty
      \let\@@arg\@@clr\XC@replace\@@arg>\XC@@gt\else
      \edef\@@arg{[\@@mod]{\@@clr}}\XC@definecolor[]{*}\@@mod\@@clr
      \def\@@clr{*}\fi
    \XC@append\@@arg{&}\extractcolorspecs\@@clr\@@mod\@@clr
    \@testc@lor}%
%    \end{macrocode}
%  \end{macro}
%
% Common loop macro for table head and |\testcolor| rows.
% Note that this command is only available within the |testcolors| environment.
%    \begin{macrocode}
  \def\@testc@lor
   {\count@\z@
    \expandafter\@tfor\expandafter\@@tmp\expandafter:\expandafter=\@@nam\do
     {\ifx\@@clr\@empty
        \edef\@@cmd{\noexpand\textbf{\@@tmp}}%
      \else
        \convertcolorspec\@@mod\@@clr\@@tmp\@@cmd
        \edef\@@cmd
         {\noexpand\@testc@l@r{\@@tmp}{\@@cmd}%
          \ifx\@@mod\@@tmp\noexpand\underline\fi
          {\expandafter\XC@strip@comma\@@cmd,,\@nnil}}%
      \fi
      \expandafter\XC@append\expandafter\@@arg\expandafter{\@@cmd}%
      \advance\count@\@ne
      \ifnum\count@=\@@num\XC@append\@@arg{\\}\else\XC@append\@@arg{&}\fi}%
    \@@arg}%
%    \end{macrocode}
%
% Display color sample box:
%    \begin{macrocode}
  \def\@testc@l@r##1##2%
   {\fboxsep\z@\fbox{\colorbox[##1]{##2}{\phantom{XX}}} }%
%    \end{macrocode}
%
% Now start the table:
%    \begin{macrocode}
  \tabular{@{}l*{\@@num}{l}@{}}%
  \def\@@arg{\textbf{color}& }\let\@@clr\@empty\@testc@lor}%
%    \end{macrocode}
%
% End of environment:
%    \begin{macrocode}
 {\endtabular\ignorespacesafterend}
%    \end{macrocode}
%  \end{environment}
%
%
% \subsubsection{Extra groups}
%
% Like |\normalcolor|, the following five commands are defined in the kernel. 
% This means that they can be used to make macros in packages `color safe'.
% The extra groups have been made permanent in the \LaTeX{} kernel 2021 so
% they are no longer redefined here.
% The |\color@setgroup| will not have any effect unless a user also uses 
% a color package, when the `active definition' here will take effect. 
%
% \begin{macro}{\color@setgroup}
% This is to be used in contexts (e.g., `lrbox') where text is to be saved and used after some other, unknown, text that may contain color commands.
% A matching |\color@endgroup| should be used at the end of the text.
%    \begin{macrocode}
\def\color@setgroup{\begingroup\set@color}
%    \end{macrocode}
%  \end{macro}
%
%  \begin{macro}{\color@begingroup}
% This is to be used at the start of contexts that may contain color commands, but where it is not necessary to save the current color.
% Examples of this are in the box commands of \File{ltboxes.dtx} where user-supplied text is saved internally in a box between |\color@begingroup|, |\color@endgroup|, but the box is used before any other color commands could intervene.
% A matching |\color@endgroup| should be used at the end of the text.
%  \end{macro}
%
%  \begin{macro}{\color@endgroup}
% To be used to close the `group' started by one of the above two commands.
% The |\endgraf| in its definition is required in the case of groups of text in vertical `par' mode, but doesn't do any harm in horizontal `LR' contexts.
%  \end{macro}
%
%  \begin{macro}{\color@hbox}
% To be used to open a `colored hbox'
%  \end{macro}
%
%  \begin{macro}{\color@vbox}
% To be used to open a `colored hbox'
%  \end{macro}
%
%  \begin{macro}{\color@endbox}
% To be used to close a `colored (h/v)box'
%  \end{macro}
%
%
%
% \subsection{Color calculation}
%
% \subsubsection{Color conversion}
%
% \begin{macro}{\convertcolorspec}
%   \marg{source model}\marg{spec}\marg{target model}\marg{cmd}\\
% General color conversion command.
%    \begin{macrocode}
\def\convertcolorspec#1#2#3#4%
 {\begingroup
  \XC@sdef\sm{#1}\XC@sdef\tm{#3}\edef\@@tmp{#2}%
  \ifx\sm\XC@mod@named\extractcolorspecs\@@tmp\sm\@@tmp\fi
  \ifx\sm\tm\else\csname XC@cnv@\sm\expandafter\endcsname\@@tmp,\fi
  \aftergroupdef#4\@@tmp}
%    \end{macrocode}
% \end{macro}
%
% \begin{macro}{\XC@coremodel}
%   \marg{model-cmd}\marg{color-cmd}\\
% Convert color, given by \marg{model-cmd} and \marg{color-cmd}, into parametrisation in one of the \emph{core} models.
% The result is again saved in \marg{model-cmd} and \marg{color-cmd}.
%    \begin{macrocode}
\def\XC@coremodel#1#2%
 {\XC@ifxcase#1%
   {\XC@mod@RGB {\convertcolorspec#1#2\XC@mod@rgb#2\let#1\XC@mod@rgb}%
    \XC@mod@HTML{\convertcolorspec#1#2\XC@mod@rgb#2\let#1\XC@mod@rgb}%
    \XC@mod@Hsb {\convertcolorspec#1#2\XC@mod@hsb#2\let#1\XC@mod@hsb}%
    \XC@mod@tHsb{\convertcolorspec#1#2\XC@mod@hsb#2\let#1\XC@mod@hsb}%
    \XC@mod@HSB {\convertcolorspec#1#2\XC@mod@hsb#2\let#1\XC@mod@hsb}%
    \XC@mod@Gray{\convertcolorspec#1#2\XC@mod@gray#2\let#1\XC@mod@gray}%
    \XC@mod@wave{\convertcolorspec#1#2\XC@mod@hsb#2\let#1\XC@mod@hsb}%
    \XC@mod@named{}%
    \XC@mod@ps{}}%
   {\expandafter\XC@calc@#2,,,,:N\let#2\@@tmp}}
%    \end{macrocode}
% \end{macro}
%
% \begin{macro}{\XC@cnv@rgb}
%   \Meta{red}|,|\Meta{green}|,|\Meta{blue}|,|
%    \begin{macrocode}
\def\XC@cnv@rgb#1,#2,#3,%
 {\XC@ifxcase\tm
   {\XC@mod@cmy {\XC@calc@#1,#2,#3,,:C}%
    \XC@mod@cmyk
     {\XC@calc@#1,#2,#3,,:C\expandafter\XC@cnv@cmy@cmyk\@@tmp,}%
    \XC@mod@RGB {\edef\@@scl{\rangeRGB}\XC@calc@#1,#2,#3,,:M}%
    \XC@mod@HTML{\edef\@@scl{\@cclv}\XC@calc@#1,#2,#3,,:M%
      \expandafter\XC@calc@\@@tmp,,:H%
      \expandafter\XC@cnv@rgb@HTML\@@tmp,}%
    \XC@mod@gray{\XC@cnv@rgb@gray#1,#2,#3,}%
    \XC@mod@Gray{\XC@cnv@rgb@gray#1,#2,#3,\XC@cnv@gray\@@tmp,}}%
   {\XC@cnv@rgb@hsb#1,#2,#3,\XC@cnv@hsb@}}
%    \end{macrocode}
%
% \begin{macro}{\XC@cnv@rgb@}
%    \begin{macrocode}
\def\XC@cnv@rgb@
 {\ifx\tm\XC@mod@rgb\else\expandafter\XC@cnv@rgb\@@tmp,\fi}
%    \end{macrocode}
% \end{macro}
% \end{macro}
%
% \begin{macro}{\XC@cnv@rgb@hsb}
%   \Meta{red}|,|\Meta{green}|,|\Meta{blue}|,|
%    \begin{macrocode}
\def\XC@cnv@rgb@hsb#1,#2,#3,%
 {\count@\z@
  \ifdim#2\p@>#1\p@ \else\advance\count@4 \fi
  \ifdim#3\p@>#2\p@ \else\advance\count@2 \fi
  \ifdim#1\p@>#3\p@ \else\advance\count@1 \fi
  \ifcase\count@\or
  \XC@c@v@rgb@hsb{#3}{#2}{#1}3{}\or
  \XC@c@v@rgb@hsb{#2}{#1}{#3}1{}\or
  \XC@c@v@rgb@hsb{#2}{#3}{#1}3-\or
  \XC@c@v@rgb@hsb{#1}{#3}{#2}5{}\or
  \XC@c@v@rgb@hsb{#3}{#1}{#2}5-\or
  \XC@c@v@rgb@hsb{#1}{#2}{#3}1-\or
  \def\@tempa{0}\def\@tempb{0}\XC@calcN{#3}\@tempc
  \fi
  \edef\@@tmp{\@tempa,\@tempb,\@tempc}}
%    \end{macrocode}
% \end{macro}
%
% \begin{macro}{\XC@c@v@rgb@hsb}
%   \marg{max}\marg{med}\marg{min}\marg{const}\marg{sign}
%    \begin{macrocode}
\def\XC@c@v@rgb@hsb#1#2#3#4#5%
 {\XC@calcN{#1}\@tempc \llshiftset\dimen@{#1}%
  \llshiftset\@tempdimb{-#3}\advance\@tempdimb\dimen@
  \llshiftset\@tempdima{-#5#2}\advance\@tempdima#5\dimen@
  \advance\@tempdima#4\@tempdimb \rdivide\@tempdima{6\@tempdimb}%
  \edef\@tempa{\ifdim\@tempdima>\p@1\else\strip@pt\@tempdima\fi}%
  \rdivide\@tempdimb\dimen@
  \edef\@tempb{\strip@pt\@tempdimb}}
%    \end{macrocode}
% \end{macro}
%
% \begin{macro}{\XC@cnv@rgb@gray}
%   \Meta{red}|,|\Meta{green}|,|\Meta{blue}|,|
%    \begin{macrocode}
\def\XC@cnv@rgb@gray#1,#2,#3,%
 {\llshiftset\dimen@{#1}\@tempdima30\dimen@
  \llshiftset\dimen@{#2}\advance\@tempdima59\dimen@
  \llshiftset\dimen@{#3}\advance\@tempdima11\dimen@
  \rrshift\@tempdima \rdivide\@tempdima{100}%
  \edef\@@tmp{\strip@pt\@tempdima}}
%    \end{macrocode}
% \end{macro}
%
% \begin{macro}{\XC@cnv@rgb@HTML}
%   \Meta{red}|,|\Meta{green}|,|\Meta{blue}|,|
%    \begin{macrocode}
\def\XC@cnv@rgb@HTML#1,#2,#3,%
 {\edef\@@tmp{#1#2#3}}
%    \end{macrocode}
% \end{macro}
%
%
% \begin{macro}{\XC@cnv@cmy}
%   \Meta{cyan}|,|\Meta{magenta}|,|\Meta{yellow}|,|
%    \begin{macrocode}
\def\XC@cnv@cmy#1,#2,#3,%
 {\XC@ifxcase\tm
   {\XC@mod@cmyk{\XC@cnv@cmy@cmyk#1,#2,#3,}%
    \XC@mod@gray{\XC@cnv@rgb@gray#1,#2,#3,\XC@calcC\@@tmp\@@tmp}%
    \XC@mod@Gray{\XC@cnv@rgb@gray#1,#2,#3,\XC@calcC\@@tmp\@@tmp
      \XC@cnv@gray\@@tmp,}}%
   {\XC@calc@#1,#2,#3,,:C\XC@cnv@rgb@}}
%    \end{macrocode}
% \end{macro}
%
% \begin{macro}{\XC@cnv@cmy@cmyk}
%   \Meta{cyan}|,|\Meta{magenta}|,|\Meta{yellow}|,|\\
% This transformation includes linear \emph{undercolor-removal} and \emph{black-generation}, as specified by the 4-dimensional parameter vector |\adjustUCRBG|.
%    \begin{macrocode}
\def\XC@cnv@cmy@cmyk#1,#2,#3,%
 {\edef\@@scl
   {\ifdim#1\p@<#2\p@
    \ifdim#1\p@<#3\p@ #1\else #3\fi \else
    \ifdim#2\p@<#3\p@ #2\else #3\fi \fi}%
  \expandafter\XC@calc@\adjustUCRBG,:S%
  \expandafter\XC@c@v@cmy@cmyk\@@tmp,#1,#2,#3,}
%    \end{macrocode}
%
%    \begin{macrocode}
\def\XC@c@v@cmy@cmyk#1,#2,#3,#4,#5,#6,#7,%
 {\XC@add#5,#6,#7,,-#1,-#2,-#3,,\@@tmp \edef\@@tmp{\@@tmp,#4}%
  \expandafter\XC@calc@\@@tmp,:N}
%    \end{macrocode}
% \end{macro}
%
%
% \begin{macro}{\XC@cnv@cmyk}
%   \Meta{cyan}|,|\Meta{magenta}|,|\Meta{yellow}|,|\Meta{black}|,|
%    \begin{macrocode}
\def\XC@cnv@cmyk#1,#2,#3,#4,%
 {\XC@ifxcase\tm
   {\XC@mod@gray{\XC@cnv@rgb@gray#1,#2,#3,\edef\@@arg{#4}%
      \XC@calcT\@@tmp\@@tmp \XC@calcC\@@tmp\@@tmp}%
    \XC@mod@Gray{\XC@cnv@rgb@gray#1,#2,#3,\edef\@@arg{#4}%
      \XC@calcT\@@tmp\@@tmp \XC@calcC\@@tmp\@@tmp \XC@cnv@gray\@@tmp,}}%
   {\XC@cnv@cmyk@cmy#1,#2,#3,#4,%
    \ifx\tm\XC@mod@cmy\else\expandafter\XC@cnv@cmy\@@tmp,\fi}}
%    \end{macrocode}
% \end{macro}
%
% \begin{macro}{\XC@cnv@cmyk@cmy}
%   \Meta{cyan}|,|\Meta{magenta}|,|\Meta{yellow}|,|\Meta{black}|,|
%    \begin{macrocode}
\def\XC@cnv@cmyk@cmy#1,#2,#3,#4,%
 {\edef\@@arg{#4}\XC@calc@#1,#2,#3,,:T}
%    \end{macrocode}
% \end{macro}
%
%
% \begin{macro}{\XC@cnv@hsb}
%   \Meta{hue}|,|\Meta{saturation}|,|\Meta{brightness}|,|
%    \begin{macrocode}
\def\XC@cnv@hsb#1,#2,#3,%
 {\XC@ifxcase\tm
   {\XC@mod@Hsb {\XC@cnv@hsb@Hsb#1,#2,#3,}%
    \XC@mod@tHsb{\XC@cnv@hsb@Hsb#1,#2,#3,\expandafter\XC@cnv@Hsb\@@tmp,}%
    \XC@mod@HSB {\edef\@@scl{\rangeHSB}\XC@calc@#1,#2,#3,,:M}}%
   {\XC@cnv@hsb@rgb#1,#2,#3,\XC@cnv@rgb@}}
%    \end{macrocode}
%
% \begin{macro}{\XC@cnv@hsb@}
%    \begin{macrocode}
\def\XC@cnv@hsb@
 {\ifx\tm\XC@mod@hsb\else\expandafter\XC@cnv@hsb\@@tmp,\fi}
%    \end{macrocode}
% \end{macro}
% \end{macro}
%
% \begin{macro}{\XC@cnv@hsb@rgb}
%   \Meta{hue}|,|\Meta{saturation}|,|\Meta{brightness}|,|
%    \begin{macrocode}
\def\XC@cnv@hsb@rgb#1,#2,#3,%
 {\llshiftset\dimen@{#1}\dimen@6\dimen@ \rrshift\dimen@
  \afterassignment\dimen@\count@=\the\dimen@
  \edef\@@tmp{\strip@pt\dimen@}%
  \ifodd\count@\else\XC@calcC\@@tmp\@@tmp\fi
  \edef\@@tmp
   {\ifcase\count@ 0,\@@tmp,1\or \@@tmp,0,1\or 1,0,\@@tmp\or
    1,\@@tmp,0\or \@@tmp,1,0\or 0,1,\@@tmp\or 0,\@@tmp,1\fi}%
  \edef\@@scl{#2}\expandafter\XC@calc@\@@tmp,,:S%
  \expandafter\XC@calc@\@@tmp,,:C%
  \edef\@@scl{#3}\expandafter\XC@calc@\@@tmp,,:S%
  \expandafter\XC@calc@\@@tmp,,:N}
%    \end{macrocode}
% \end{macro}
%
% \begin{macro}{\XC@cnv@hsb@Hsb}
%   \Meta{hue}|,|\Meta{saturation}|,|\Meta{brightness}|,|
%    \begin{macrocode}
\def\XC@cnv@hsb@Hsb#1,#2,#3,%
 {\XC@c@lcS\rangeHsb{,#2,#3}{#1}\@@tmp}
%    \end{macrocode}
% \end{macro}
%
% \begin{macro}{\XC@cnv@gray}
%   \Meta{gray}|,|
% \changes{3.0}{2022-08-10}{moved the patch from hycolor to correct Gray model, issue \#23}w
%    \begin{macrocode}
\def\XC@cnv@gray#1,%
 {\XC@ifxcase\tm
   {\XC@mod@rgb {\XC@calcN{#1}\@@tmp \edef\@@tmp{\@@tmp,\@@tmp,\@@tmp}}%
    \XC@mod@gray{}%
    \XC@mod@cmy {\XC@calcC{#1}\@@tmp \edef\@@tmp{\@@tmp,\@@tmp,\@@tmp}}%
    \XC@mod@cmyk{\XC@calcC{#1}\@@tmp \edef\@@tmp{0,0,0,\@@tmp}}%
    \XC@mod@RGB {\edef\@@scl{\rangeRGB}\XC@calcM{#1}\@@tmp
      \edef\@@tmp{\@@tmp,\@@tmp,\@@tmp}}%
    \XC@mod@HTML{\edef\@@scl{\@cclv}\XC@calcM{#1}\@@tmp
      \XC@calcH\@@tmp\@@tmp \edef\@@tmp{\@@tmp\@@tmp\@@tmp}}%
    \XC@mod@HSB {\edef\@@scl{\rangeHSB}\XC@calcM{#1}\@@tmp
      \edef\@@tmp{0,0,\@@tmp}}%
    \XC@mod@Gray{\edef\@@scl{\rangeGray}\XC@calcM{#1}\@@tmp}}%
   {\XC@calcN{#1}\@@tmp \edef\@@tmp{0,0,\@@tmp}}}
%    \end{macrocode}
% \end{macro}
%
%
% \begin{macro}{\XC@cnv@RGB}
%   \Meta{Red}|,|\Meta{Green}|,|\Meta{Blue}|,|
%    \begin{macrocode}
\def\XC@cnv@RGB#1,#2,#3,%
 {\edef\@@scl{\rangeRGB}\XC@calc@#1,#2,#3,,:D\XC@cnv@rgb@}
%    \end{macrocode}
% \end{macro}
%
%
% \begin{macro}{\XC@cnv@HTML}
%   \Meta{RRGGBB}|,|
%    \begin{macrocode}
\def\XC@cnv@HTML#1,%
 {\edef\@@tmp{\expandafter\@hex@Hex#1\relax}\expandafter\XC@cnv@HTML@\@@tmp
  \edef\@@scl{\@cclv}\expandafter\XC@calc@\@@tmp,,:D\XC@cnv@rgb@}
%    \end{macrocode}
%
%    \begin{macrocode}
\def\XC@cnv@HTML@#1#2#3#4#5#6%
 {\edef\@@tmp{\number"#1#2,\number"#3#4,\number"#5#6}}
%    \end{macrocode}
% \end{macro}
%
%
% \begin{macro}{\XC@cnv@Hsb}
%   \Meta{hue\Circ}|,|\Meta{saturation}|,|\Meta{brightness}|,|
%    \begin{macrocode}
\def\XC@cnv@Hsb#1,#2,#3,%
 {\ifx\tm\XC@mod@tHsb\XC@cnv@@Hsb#1,#2,#3,\XC@cnv@Hsb@
  \else\XC@cnv@Hsb@hsb#1,#2,#3,\XC@cnv@hsb@\fi}
%    \end{macrocode}
% \end{macro}
%
% \begin{macro}{\XC@cnv@Hsb@hsb}
%   \Meta{hue\Circ}|,|\Meta{saturation}|,|\Meta{brightness}|,|
%    \begin{macrocode}
\def\XC@cnv@Hsb@hsb#1,#2,#3,%
 {\XC@c@lcD\rangeHsb{,#2,#3}{#1}\@@tmp}
%    \end{macrocode}
% \end{macro}
%
% \begin{macro}{\XC@cnv@@Hsb}
%   \Meta{hue\Circ}|,|\Meta{saturation}|,|\Meta{brightness}|,|\marg{cmd}
%    \begin{macrocode}
\def\XC@cnv@@Hsb#1,#2,#3,#4%
 {\dimen@\rangeHsb\p@ \advance\dimen@\p@ \edef\@@tmp{\strip@pt\dimen@}%
  \edef\@@tmp{0,0;\rangetHsb;\rangeHsb,\rangeHsb;\@@tmp,\@@tmp;}%
  \dimen@#1\p@ \expandafter#4\@@tmp\@nnil
  \edef\@@tmp{\strip@pt\dimen@,#2,#3}}
%    \end{macrocode}
% \end{macro}
%
% \begin{macro}{\XC@cnv@Hsb@}
%   \Meta{x}|,|\Meta{y}|;|\Meta{x'}|,|\Meta{y'}|;|
%    \begin{macrocode}
\def\XC@cnv@Hsb@#1,#2;#3,#4;%
 {\ifdim\dimen@<#4\p@
    \@tempdima#3\p@ \advance\@tempdima-#1\p@
    \@tempdimb#4\p@ \advance\@tempdimb-#2\p@
    \rdivide\@tempdima\@tempdimb
    \advance\dimen@-#2\p@ \rmultiply\dimen@{\strip@pt\@tempdima}%
    \advance\dimen@#1\p@ \expandafter\remove@to@nnil
  \else
    \def\@@tmp{\XC@cnv@Hsb@#3,#4;}\expandafter\@@tmp
  \fi}
%    \end{macrocode}
% \end{macro}
%
%
% \begin{macro}{\XC@cnv@tHsb}
%   \Meta{hue\Circ}|,|\Meta{saturation}|,|\Meta{brightness}|,|
%    \begin{macrocode}
\def\XC@cnv@tHsb#1,#2,#3,%
 {\XC@cnv@@Hsb#1,#2,#3,\XC@cnv@tHsb@
  \ifx\tm\XC@mod@Hsb\else\expandafter\XC@cnv@Hsb\@@tmp,\fi}
%    \end{macrocode}
% \end{macro}
%
% \begin{macro}{\XC@cnv@tHsb@}
%   \Meta{x}|,|\Meta{y}|;|\Meta{x'}|,|\Meta{y'}|;|
%    \begin{macrocode}
\def\XC@cnv@tHsb@#1,#2;#3,#4;%
 {\ifdim\dimen@<#3\p@
    \@tempdima#4\p@ \advance\@tempdima-#2\p@
    \@tempdimb#3\p@ \advance\@tempdimb-#1\p@
    \rdivide\@tempdima\@tempdimb
    \advance\dimen@-#1\p@ \rmultiply\dimen@{\strip@pt\@tempdima}%
    \advance\dimen@#2\p@ \expandafter\remove@to@nnil
  \else
    \def\@@tmp{\XC@cnv@tHsb@#3,#4;}\expandafter\@@tmp
  \fi}
%    \end{macrocode}
% \end{macro}
%
%
% \begin{macro}{\XC@cnv@HSB}
%   \Meta{Hue}|,|\Meta{Saturation}|,|\Meta{Brightness}|,|
%    \begin{macrocode}
\def\XC@cnv@HSB#1,#2,#3,%
 {\edef\@@scl{\rangeHSB}\XC@calc@#1,#2,#3,,:D\XC@cnv@hsb@}
%    \end{macrocode}
% \end{macro}
%
%
% \begin{macro}{\XC@cnv@Gray}
%   \Meta{Gray}|,|
%    \begin{macrocode}
\def\XC@cnv@Gray#1,%
 {\XC@c@lcD\rangeGray{}{#1}\@@tmp \XC@cnv@gray\@@tmp,}
%    \end{macrocode}
% \end{macro}
%
%
% \begin{macro}{\XC@cnv@wave}
%   \Meta{lambda}|,|
%    \begin{macrocode}
\def\XC@cnv@wave#1,%
 {\XC@cnv@wave@hsb#1,\XC@cnv@hsb@}
%    \end{macrocode}
% \end{macro}
%
% \begin{macro}{\XC@cnv@wave@hsb}
%   \Meta{lambda}|,|
%    \begin{macrocode}
\def\XC@cnv@wave@hsb#1,%
 {\@tempdima#1\p@
%    \end{macrocode}
%
% Calculate hue:
%    \begin{macrocode}
  \ifdim\@tempdima<440\p@
    \XC@c@v@wave@hsb{-.0166667}{7.33333}4{}\else
  \ifdim\@tempdima<490\p@
    \XC@c@v@wave@hsb{.02}{-8.8}4-\else
  \ifdim\@tempdima<510\p@
    \XC@c@v@wave@hsb{-.05}{25.5}2{}\else
  \ifdim\@tempdima<580\p@
    \XC@c@v@wave@hsb{.0142857}{-7.28571}2-\else
  \ifdim\@tempdima<645\p@
    \XC@c@v@wave@hsb{-.0153846}{9.92308}0{}\else
  \dimen@\z@ \fi\fi\fi\fi\fi
  \divide\dimen@6\relax
  \edef\@@tmp{\strip@pt\dimen@,1,}%
%    \end{macrocode}
%
% Calculate brightness:
%    \begin{macrocode}
  \ifdim\@tempdima<420\p@
    \XC@c@v@wave@hsb{.0175}{-6.35}0{}\else
  \ifdim\@tempdima>700\p@
    \XC@c@v@wave@hsb{-.00875}{7.125}0{}\else
  \dimen@\p@ \fi\fi
  \edef\@@tmp{\@@tmp\strip@pt\dimen@}}
%    \end{macrocode}
% \end{macro}
%
% \begin{macro}{\XC@c@v@wave@hsb}
%   \marg{factor}\marg{const}\marg{offset}\marg{sign}\\
% Calculates $\Meta{offset} + \Meta{sign}\cdot\varrho\bigl(\Meta{const} + \Meta{factor}\cdot\lambda\bigr)$, where $\varrho$ is the reduction function.
% The result is stored in |\dimen@|.
%    \begin{macrocode}
\def\XC@c@v@wave@hsb#1#2#3#4%
 {\dimen@#1\@tempdima \advance\dimen@#2\p@
  \XC@calcR{\strip@pt\dimen@}\@tempa
  \dimen@#4\@tempa\p@ \advance\dimen@#3\p@}
%    \end{macrocode}
% \end{macro}
%
%
% \subsubsection{Color complement}
%
% \begin{macro}{\XC@cnv@cmyk@compl}
%   \Meta{cyan}|,|\Meta{magenta}|,|\Meta{yellow}|,|\Meta{black}|,|
%    \begin{macrocode}
\def\XC@cnv@cmyk@compl#1,#2,#3,#4,%
 {\XC@cnv@cmyk@cmy#1,#2,#3,#4,%
  \expandafter\XC@calc@\@@tmp,,:C\expandafter\XC@cnv@cmy@cmyk\@@tmp,}
%    \end{macrocode}
% \end{macro}
%
% \begin{macro}{\XC@cnv@hsb@compl}
%   \Meta{hue}|,|\Meta{saturation}|,|\Meta{brightness}|,|
%    \begin{macrocode}
\def\XC@cnv@hsb@compl#1,#2,#3,%
 {\lshiftset\dimen@{#1}\advance\dimen@\ifdim\dimen@<5\p@-\fi-5\p@
  \rshift\dimen@ \edef\@@tmp{\strip@pt\dimen@}%
  \lshiftset\dimen@{#3}\@tempdima#2\dimen@
  \@tempdimb\@tempdima \advance\@tempdimb10\p@ \advance\@tempdimb-\dimen@
  \ifdim\@tempdimb=\z@
    \edef\@@tmp{\@@tmp,0,0}%
  \else
    \rdivide\@tempdima\@tempdimb \rshift\@tempdimb
    \edef\@@tmp{\@@tmp,\strip@pt\@tempdima,\strip@pt\@tempdimb}%
  \fi}
%    \end{macrocode}
% \end{macro}
%
%
%
% \subsection{Color initialisation}
%
% \subsubsection{Support for missing color models}
%
% These macros provide generic support for missing color models in specific drivers.
% For example, \File{pdftex.def} does not support the \Model{hsb} model.
% It is assumed that at least one of the models \Model{rgb}, \Model{RGB}, \Model{HTML}, \Model{hsb}, \Model{Hsb}, \Model{HSB}, \Model{cmy}, \Model{cmyk}, \Model{gray}, \Model{Gray} exists.
%
% \begin{macro}{\substitutecolormodel}
%   \marg{source model}\marg{target model-list}\\
% Substitute \Meta{source model} by the first actually present model from \Meta{target model-list}:
% each reference to the driver-specific command |\color@|\Meta{source model} is replaced by a color model conversion, followed by the driver-specific command |\color@|\Meta{target model}.
%    \begin{macrocode}
\def\substitutecolormodel#1#2%
 {\XC@sdef\XC@@tmp{\expandafter\XC@substitutecolormodel#2/ignore/\@nnil}%
  \ifx\XC@@tmp\XC@mod@ignore
    \XC@let@cN{color@#1}\color@ignore
  \else
    \expandafter\edef\csname color@#1\endcsname##1##2%
     {\noexpand\convertcolorspec
      \expandafter\noexpand\csname XC@mod@#1\endcsname{##2}%
      \expandafter\noexpand\csname XC@mod@\XC@@tmp\endcsname
      \noexpand\XC@@tmp\noexpand\expandafter
      \expandafter\noexpand\csname color@\XC@@tmp\endcsname
      \noexpand\expandafter##1\noexpand\expandafter{\noexpand\XC@@tmp}}%
  \fi
  \expandafter\xglobalize@\csname color@#1\endcsname \xglobal@stop
  \PackageInfo{xcolor}{Model `#1' substituted by `\XC@@tmp'}}
%    \end{macrocode}
%
%    \begin{macrocode}
\def\XC@substitutecolormodel#1/%
 {\@ifundefinedmodel{#1}{\XC@substitutecolormodel}{#1\remove@to@nnil}}
%    \end{macrocode}
% \end{macro}
%
% \begin{macro}{\color@ignore}
%   \marg{cmd}\marg{spec}\\
% Empty color model: no driver-specific representation at all.
%    \begin{macrocode}
\def\color@ignore#1#2{\def#1{}}
%    \end{macrocode}
% \end{macro}
%
% \begin{macro}{\color@rgb}
%   \marg{cmd}\marg{spec}\\
% Takes the color specification \Meta{spec}, turns it into a driver-specific representation and puts the result into the macro \Meta{cmd}.
% For example, if \File{dvips.def} is loaded, then |\color@rgb\foo{1,0,0}| will be equivalent to |\def\foo{rgb 1 0 0}|.
% In the case of \File{pdftex.def}, it will act like |\def\foo{1 0 0 rg 1 0 0 RG}| instead.
%    \begin{macrocode}
\@ifundefinedmodel{rgb}%
 {\substitutecolormodel{rgb}{RGB/HTML/hsb/Hsb/HSB/cmy/cmyk/gray/Gray}}{}
%    \end{macrocode}
% \end{macro}
%
% \begin{macro}{\color@cmyk}
%   \marg{cmd}\marg{spec}
%    \begin{macrocode}
\@ifundefinedmodel{cmyk}{\substitutecolormodel{cmyk}{cmy/rgb}}{}
%    \end{macrocode}
% \end{macro}
%
% \begin{macro}{\color@cmy}
%   \marg{cmd}\marg{spec}
%    \begin{macrocode}
\@ifundefinedmodel{cmy}%
 {\PackageInfo{xcolor}{Model `cmy' substituted by `cmy0'}
  \def\color@cmy#1#2{\color@cmyk#1{#2,0}}}{}
%    \end{macrocode}
% \end{macro}
%
% \begin{macro}{\color@hsb}
%   \marg{cmd}\marg{spec}\\
% Since the \Option{dvipdfm} and \Option{textures} drivers only pretend to support \Model{hsb}, we have to check for these drivers.
%    \begin{macrocode}
\@ifundefinedmodel{hsb}{\@tempswatrue}%
 {\XC@ifcase\XC@@drv{{dvipdfm,textures}\@tempswatrue}\@tempswafalse}
\if@tempswa
  \substitutecolormodel{hsb}{Hsb/HSB/rgb}
\fi
%    \end{macrocode}
% \end{macro}
%
% \begin{macro}{\color@RGB}
%   \marg{cmd}\marg{spec}\\
% If the model is supported and |\rangeRGB| equals 255, we stick to the algorithm provided by the driver.
%    \begin{macrocode}
\@ifundefinedmodel{RGB}%
 {\substitutecolormodel{RGB}{HTML/rgb}}
 {\PackageInfo{xcolor}{Model `RGB' extended}
  \let\XC@color@RGB\color@RGB
  \def\color@RGB#1#2%
   {\relax\ifnum\rangeRGB=255
      \XC@color@RGB#1{#2}%
    \else
      \convertcolorspec\XC@mod@RGB{#2}\XC@mod@rgb\XC@@tmp
      \expandafter\color@rgb\expandafter#1\expandafter{\XC@@tmp}%
    \fi}}
%    \end{macrocode}
% \end{macro}
%
% \begin{macro}{\color@gray}
% \begin{macro}{\color@HTML}
% \begin{macro}{\color@Hsb}
% \begin{macro}{\color@tHsb}
% \begin{macro}{\color@HSB}
% \begin{macro}{\color@Gray}
% \begin{macro}{\color@wave}
%   \marg{cmd}\marg{spec}
%    \begin{macrocode}
\@ifundefinedmodel{gray}{\substitutecolormodel{gray}{Gray/rgb}}{}
\@ifundefinedmodel{HTML}{\substitutecolormodel{HTML}{rgb}}{}
\@ifundefinedmodel{Hsb}{\substitutecolormodel{Hsb}{hsb}}{}
\@ifundefinedmodel{tHsb}{\substitutecolormodel{tHsb}{hsb}}{}
\@ifundefinedmodel{HSB}{\substitutecolormodel{HSB}{hsb}}{}
\@ifundefinedmodel{Gray}{\substitutecolormodel{Gray}{gray}}{}
\@ifundefinedmodel{wave}{\substitutecolormodel{wave}{hsb}}{}
%    \end{macrocode}
% \end{macro}
% \end{macro}
% \end{macro}
% \end{macro}
% \end{macro}
% \end{macro}
% \end{macro}
%
%
% \subsubsection{\Index{PostScript}-related stuff}
%
% \begin{macro}{\XC@getcolor}
%   \marg{color}\marg{cmd}\\
%   \marg{color command}\marg{cmd}\\
%   |{|\oarg{model}\marg{spec}|}|\marg{cmd}\\
% Transformation of a color, given by the first argument, especially needed for \Package{pstricks}.
% The result will always be an internal color representation (as described in section \vref{sec.definition}) and is stored in \Meta{cmd}.
% There are different ways of specifying the color:
% \begin{itemize}
% \item \Meta{color} is a color name, expression, etc., e.g., |\XC@getcolor{green!50}\foo|;
% \item \Meta{color command} is a macro whose \emph{first level} expansion is an internal color representation, e.g., |\XC@getcolor\pslinecolor\foo|;
% \item \Meta{model} and \Meta{spec} provide, as usual, explicit parameter specifications.
% \end{itemize}
%    \begin{macrocode}
\def\XC@getcolor#1#2%
 {\begingroup\toks@{#1}\XC@getc@lor#1\XC@@\aftergroupdef#2{\@@tmp}}
%    \end{macrocode}
%
% Check whether the argument consists of one or more tokens.
% In the first case, expand by one level and go to next check.
% In the second case, assume it's a \Meta{color} and simulate its display, but without performing the real |\set@color| command (which would insert |\special|s, possibly causing trouble with blank spaces).
% This simulation ensures that we get exactly the same color --- including all blending, masking --- as if we had executed an explicit |\color| command at that point.
%    \begin{macrocode}
\def\XC@getc@lor#1#2\XC@@
 {\@ifxempty{#2}%
   {\expandafter\XC@getc@l@r#1\XC@@}%
   {\let\set@color\relax
    \ifx#1[%
      \expandafter\@undeclaredcolor\the\toks@\else
      \expandafter\@declaredcolor\expandafter{\the\toks@}\fi
    \let\@@tmp\XC@current@color}}
%    \end{macrocode}
%
% Check whether the first token is |\xcolor@|.
% In that case, |#1#2| is assumed to be a complete internal color representation.
% Otherwise, perform a simulation as described before.
%    \begin{macrocode}
\def\XC@getc@l@r#1#2\XC@@
 {\ifx#1\xcolor@
    \def\@@tmp{\xcolor@#2}%
  \else
    \let\set@color\relax
    \expandafter\@declaredcolor\expandafter{\the\toks@}%
    \let\@@tmp\XC@current@color
  \fi}
%    \end{macrocode}
% \end{macro}
%
% \begin{macro}{\XC@usecolor}
%   \marg{name}\\
%   \marg{color command}\\
% Transformation of a color into \Index{PostScript} code.
% The argument may be either a command whose \emph{first-level} expansion is an internal color representation (e.g., the outcome of an |\XC@getcolor| command: |\XC@usecolor\pslinecolor|), or a \Meta{name} (e.g., |\XC@usecolor{green}|).
% Note that non-trivial color expressions like |green!50| are forbidden and that |\XC@usecolor| is fully expandable, for example, |\edef\foo{\XC@usecolor{green}}| is valid.
%    \begin{macrocode}
\def\XC@usecolor#1{\expandafter\XC@usec@lor#1\XC@@}
\def\XC@usec@lor#1#2\XC@@
 {\ifx#1\xcolor@
    \expandafter\expandafter\expandafter\c@lor@to@ps#1#2\@@
  \else
    \expandafter\expandafter\expandafter\expandafter\expandafter\expandafter
    \expandafter\c@lor@to@ps\csname\@backslashchar color@#1#2\endcsname\@@
  \fi
  \space}
%    \end{macrocode}
%
% For compatibility reasons, we need the following command temporarily; however, |\XC@c@lor@to@ps| will disappear!
%    \begin{macrocode}
\let\XC@c@lor@to@ps\XC@usecolor
%    \end{macrocode}
% \end{macro}
%
%    \begin{macrocode}
\let\pst@getcolor\XC@getcolor
\let\pst@usecolor\XC@usecolor
%    \end{macrocode}
%
%
% \subsubsection{Whatsit \ldots}
%
% \begin{macro}{\c@lor@special}
% Some drivers do not support all the features of this package.
% They should always put a \Meta{whatsit} in the current list though.
% The following macro has most of the features of |\special|, but does not put anything into the \File{.dvi} file.
% It does write to the |log| file or the terminal (depending on the value of |#1|).
%    \begin{macrocode}
\def\c@lor@special#1#2{\edef\@tempa{\write#1{#2}}\@tempa}
%    \end{macrocode}
%  \end{macro}
%
%
% \subsubsection{Default color}
%
% \begin{macro}{\normalcolor}
% |\normalcolor| is defined (to |\relax|) in the \LaTeX\ kernel, so it is safe to use this in macros that may possibly be used in conjunction with color.
% It will have no effect until the \XCP{} is also loaded.
%    \begin{macrocode}
\protected\def\normalcolor{%
  \let\current@color\default@color
  \let\XC@current@color\XC@default@color\XC@display}
\AtBeginDocument{\let\XC@default@color\XC@current@color}
%    \end{macrocode}
% \end{macro}
%
%  \begin{macro}{\default@color}
% Internal macro to store the `default' color used by |\normalcolor|.
%    \begin{macrocode}
\AtBeginDocument{\let\default@color\current@color}
%    \end{macrocode}
% \end{macro}
% |\current@color| (which is the expansion of |\XC@current@color|) contains an internal representation of the color at this point in the document.
% This can only be an approximation to the truth as the `macro layer' of \TeX\ does not know where the output routine is going to re-insert floats.
% This is why drivers must maintain their own stack of colors in order to fully support these commands.
%
%
% \subsubsection{Driver-specific definitions for named colors}
%
% \begin{macro}{\c@lor@@named}
% \begin{macro}{\define@color@named}
% We have to modify a few driver-specific commands in order to avoid generation of superfluous macros.
%    \begin{macrocode}
\XC@ifcase\XC@@drv
 {{dvipdf,dvipdfm,dvipdfmx,dvips,pctex32}%
  {\def\c@lor@@named#1,#2,#3\@@#4%
    {\edef#4{\expandafter\expandafter\expandafter\@thirdoffive
             \csname\@backslashchar color@#1\endcsname}}
   \def\define@color@named#1#2%
    {\@namedef{\@backslashchar color@#1}{{}{}{ \colornameprefix#1}{}{}}}}%
%    \end{macrocode}
%
%    \begin{macrocode}
  {dvipsone,luatex,pctexps,pdftex,tcidvi,truetex,vtex,xetex}%
  {\def\c@lor@@named#1,#2,#3\@@#4%
    {\edef#4{\expandafter\expandafter\expandafter\@thirdoffive
             \csname\@backslashchar color@#1\endcsname}}
   \def\define@color@named#1#2%
    {\let\@@cls\@empty\@namedef{\@backslashchar color@#1}{{}{}{#2}{}{}}}}%
%    \end{macrocode}
%
%    \begin{macrocode}
  {textures}%
  {\def\c@lor@@named#1,#2,#3\@@#4%
    {\edef\@tempa{\color@P@custom{#1}{\if!#2!1\else#2\fi}}%
     \expandafter\def\expandafter#4\@tempa}
   \def\define@color@named#1#2%
    {\edef\@tempa{{#1}{#2}}\expandafter\AddCustomColor\@tempa}
   \def\c@lor@ps@#1 #2\@@
    {\expandafter\expandafter\expandafter\expandafter\expandafter
     \expandafter\expandafter\c@lor@to@ps@
     \csname\@backslashchar color@#1\expandafter\endcsname\space
     #2. \@@{#1}}}%
%    \end{macrocode}
%
%    \begin{macrocode}
 }{}
%    \end{macrocode}
% \end{macro}
% \end{macro}
%
%
% \subsubsection{Basic color definitions}
%
% \Model{rgb}-minded colors.\footnote{Generic definitions:
% \Color{brown} = \Color{violet}|!|50|!|\Color{yellow},
% \Color{lime} = \Color{green}|!|25|!|\Color{yellow},
% \Color{orange} = \Color{red}|!|50|!|\Color{yellow},
% \Color{pink} = \Color{red}|!|25,
% \Color{purple} = \Color{violet}|!|50|!|\Color{red},
% \Color{teal} = \Color{green}|!|50|!|\Color{blue},
% \Color{violet} = \Color{blue}|!|50|!|\Color{red}.}
%    \begin{macrocode}
\definecolorset{rgb/hsb/cmyk/gray}{}{}%
 {red,1,0,0/0,1,1/0,1,1,0/.3;%
  green,0,1,0/.33333,1,1/1,0,1,0/.59;%
  blue,0,0,1/.66667,1,1/1,1,0,0/.11;%
  brown,.75,.5,.25/.083333,.66667,.75/0,.25,.5,.25/.5475;%
  lime,.75,1,0/.20833,1,1/.25,0,1,0/.815;%
  orange,1,.5,0/.083333,1,1/0,.5,1,0/.595;%
  pink,1,.75,.75/0,.25,1/0,.25,.25,0/.825;%
  purple,.75,0,.25/.94444,1,.75/0,.75,.5,.25/.2525;%
  teal,0,.5,.5/.5,1,.5/.5,0,0,.5/.35;%
  violet,.5,0,.5/.83333,1,.5/0,.5,0,.5/.205}%
%    \end{macrocode}
%
% \Model{cmyk}-minded colors.\footnote{Generic definitions:
% \Color{olive} = \Color{red}|!|50|!|\Color{green}.}
%    \begin{macrocode}
\definecolorset{cmyk/rgb/hsb/gray}{}{}%
 {cyan,1,0,0,0/0,1,1/.5,1,1/.7;%
  magenta,0,1,0,0/1,0,1/.83333,1,1/.41;%
  yellow,0,0,1,0/1,1,0/.16667,1,1/.89;%
  olive,0,0,1,.5/.5,.5,0/.16667,1,.5/.39}
%    \end{macrocode}
%
% \Model{gray}-minded colors.
%    \begin{macrocode}
\definecolorset{gray/rgb/hsb/cmyk}{}{}%
 {black,0/0,0,0/0,0,0/0,0,0,1;%
  darkgray,.25/.25,.25,.25/0,0,.25/0,0,0,.75;%
  gray,.5/.5,.5,.5/0,0,.5/0,0,0,.5;%
  lightgray,.75/.75,.75,.75/0,0,.75/0,0,0,.25;%
  white,1/1,1,1/0,0,1/0,0,0,0}
%    \end{macrocode}
%
% Load predefined sets of colors (e.g., \File{dvipsnam.def}).
%    \begin{macrocode}
\XC@@names
%    \end{macrocode}
%
%
% \subsubsection{Finale}
%
% Clear color stack and restore catcodes.
%    \begin{macrocode}
\AtBeginDocument{\let\XC@@stack\@empty}
\let\colornameprefix\XC@@nameprefix
\XC@catcodes
%    \end{macrocode}
%
%    \begin{macrocode}
\csname xcolorcmd\endcsname
\color{black}
%</package>
%    \end{macrocode}
%
%
% \iffalse
%-------------------------------------------------------------------------------
% \fi
%
% \subsection{Auxiliary definitions}
%
% \subsubsection{\Index{PostScript} prologue file \File{xcolor.pro}}
%
%    \begin{macrocode}
%<*dvips>
userdict begin
/setcmycolor{0 setcmykcolor}def
/XCbd{bind def}bind def
/XCed{exch def}XCbd
%    \end{macrocode}
%
% \Index{PostScript} version of |\providecommand|.
%    \begin{macrocode}
/XCdef{exch dup TeXDict exch known{pop pop}{XCed}ifelse}XCbd
%    \end{macrocode}
%
% \Index{PostScript} version of |\providecolor|.
%    \begin{macrocode}
/XCsetcolor{cvx counttomark array astore cvx exch pop XCdef}XCbd
%    \end{macrocode}
%
% Now define model-dependent color definition commands.
%    \begin{macrocode}
/XCcmy{mark exch exec/setcmycolor XCsetcolor}XCbd
/XCcmyk{mark exch exec/setcmykcolor XCsetcolor}XCbd
/XCgray{mark exch exec/setgray XCsetcolor}XCbd
/XChsb{mark exch exec/sethsbcolor XCsetcolor}XCbd
/XCrgb{mark exch exec/setrgbcolor XCsetcolor}XCbd
%    \end{macrocode}
%
% This is a fix for \File{special.pro}'s behaviour to perform a |0 setgray| during |normalscale|.
%    \begin{macrocode}
/XCfixnormalscale{/normalscale{Resolution 72 div VResolution 72 div neg
  scale magscale{DVImag dup scale}if}def}def
%    \end{macrocode}
%
% Grouping operators.
% We have to call |XCfixnormalscale| since |normalscale| is invoked by |@beginspecial| as defined by \File{special.pro}.
%    \begin{macrocode}
/XC!b{save XCfixnormalscale}XCbd
/XC!e{restore}XCbd
%    \end{macrocode}
%
% Scale `sp' coordinates back to `bp' units.
%    \begin{macrocode}
/XCsp2bp{1 65781.76 div dup scale}XCbd
%    \end{macrocode}
%
% \begin{macro}{XCsetline}
%   \meta{linewidth} |XCsetline|\\
% Common line definition.
%    \begin{macrocode}
/XCsetline
 {setlinewidth [] 0 setdash 2 setlinecap 0 setlinejoin 4 setmiterlimit}XCbd
%    \end{macrocode}
% \end{macro}
%
% \begin{macro}{XCboxframe}
%   \meta{linewidth} \meta{width} \meta{height} \meta{depth} |XCboxframe|\\
% Draw a frame in `sp' coordinates using LanguageLevel~2.
%    \begin{macrocode}
/XCboxframe{/d XCed/h XCed/w XCed/l XCed w 0 gt{h d neg gt{l 0 gt
 {XCsp2bp l XCsetline
  l 2 div dup translate
  0 d neg w l sub h l sub d add rectstroke}if}if}if}XCbd
%    \end{macrocode}
% \end{macro}
%
% \begin{macro}{XCcolorblock}
%   \meta{width} \meta{height} \meta{depth} |XCcolorblock|\\
% Draw a colored box in `sp' coordinates using LanguageLevel~2.
%    \begin{macrocode}
/XCcolorblock{/d XCed/h XCed/w XCed w 0 gt{h d neg gt
 {XCsp2bp 0 XCsetline
  0 d neg w h d add rectfill}if}if}XCbd
%    \end{macrocode}
% \end{macro}
%
%    \begin{macrocode}
end
%</dvips>
%    \end{macrocode}
%
%
% \clearpage
% \Finale
% \endinput
%
%
% \iffalse
%-------------------------------------------------------------------------------
%<*svgnames|x11names|hksknames|hksnnames|hksznames|pmsnames|ralnames>
\def\colornameprefix{XC@}
%</svgnames|x11names|hksknames|hksnnames|hksznames|pmsnames|ralnames>
%
%<*svgnames>
\preparecolorset{rgb}{}{}{%
AliceBlue,.94,.972,1;%
AntiqueWhite,.98,.92,.844;%
Aqua,0,1,1;%
Aquamarine,.498,1,.83;%
Azure,.94,1,1;%
Beige,.96,.96,.864;%
Bisque,1,.894,.77;%
Black,0,0,0;%
BlanchedAlmond,1,.92,.804;%
Blue,0,0,1;%
BlueViolet,.54,.17,.888;%
Brown,.648,.165,.165;%
BurlyWood,.87,.72,.53;%
CadetBlue,.372,.62,.628;%
Chartreuse,.498,1,0;%
Chocolate,.824,.41,.116;%
Coral,1,.498,.312;%
CornflowerBlue,.392,.585,.93;%
Cornsilk,1,.972,.864;%
Crimson,.864,.08,.235;%
Cyan,0,1,1;%
DarkBlue,0,0,.545;%
DarkCyan,0,.545,.545;%
DarkGoldenrod,.72,.525,.044;%
DarkGray,.664,.664,.664;%
DarkGreen,0,.392,0;%
DarkGrey,.664,.664,.664;%
DarkKhaki,.74,.716,.42;%
DarkMagenta,.545,0,.545;%
DarkOliveGreen,.332,.42,.185;%
DarkOrange,1,.55,0;%
DarkOrchid,.6,.196,.8;%
DarkRed,.545,0,0;%
DarkSalmon,.912,.59,.48;%
DarkSeaGreen,.56,.736,.56;%
DarkSlateBlue,.284,.24,.545;%
DarkSlateGray,.185,.31,.31;%
DarkSlateGrey,.185,.31,.31;%
DarkTurquoise,0,.808,.82;%
DarkViolet,.58,0,.828;%
DeepPink,1,.08,.576;%
DeepSkyBlue,0,.75,1;%
DimGray,.41,.41,.41;%
DimGrey,.41,.41,.41;%
DodgerBlue,.116,.565,1;%
FireBrick,.698,.132,.132;%
FloralWhite,1,.98,.94;%
ForestGreen,.132,.545,.132;%
Fuchsia,1,0,1;%
Gainsboro,.864,.864,.864;%
GhostWhite,.972,.972,1;%
Gold,1,.844,0;%
Goldenrod,.855,.648,.125;%
Gray,.5,.5,.5;%
Green,0,.5,0;%
GreenYellow,.68,1,.185;%
Grey,.5,.5,.5;%
Honeydew,.94,1,.94;%
HotPink,1,.41,.705;%
IndianRed,.804,.36,.36;%
Indigo,.294,0,.51;%
Ivory,1,1,.94;%
Khaki,.94,.9,.55;%
Lavender,.9,.9,.98;%
LavenderBlush,1,.94,.96;%
LawnGreen,.488,.99,0;%
LemonChiffon,1,.98,.804;%
LightBlue,.68,.848,.9;%
LightCoral,.94,.5,.5;%
LightCyan,.88,1,1;%
LightGoldenrod,.933,.867,.51;%
LightGoldenrodYellow,.98,.98,.824;%
LightGray,.828,.828,.828;%
LightGreen,.565,.932,.565;%
LightGrey,.828,.828,.828;%
LightPink,1,.712,.756;%
LightSalmon,1,.628,.48;%
LightSeaGreen,.125,.698,.668;%
LightSkyBlue,.53,.808,.98;%
LightSlateBlue,.518,.44,1;%
LightSlateGray,.468,.532,.6;%
LightSlateGrey,.468,.532,.6;%
LightSteelBlue,.69,.77,.87;%
LightYellow,1,1,.88;%
Lime,0,1,0;%
LimeGreen,.196,.804,.196;%
Linen,.98,.94,.9;%
Magenta,1,0,1;%
Maroon,.5,0,0;%
MediumAquamarine,.4,.804,.668;%
MediumBlue,0,0,.804;%
MediumOrchid,.73,.332,.828;%
MediumPurple,.576,.44,.86;%
MediumSeaGreen,.235,.7,.444;%
MediumSlateBlue,.484,.408,.932;%
MediumSpringGreen,0,.98,.604;%
MediumTurquoise,.284,.82,.8;%
MediumVioletRed,.78,.084,.52;%
MidnightBlue,.098,.098,.44;%
MintCream,.96,1,.98;%
MistyRose,1,.894,.884;%
Moccasin,1,.894,.71;%
NavajoWhite,1,.87,.68;%
Navy,0,0,.5;%
NavyBlue,0,0,.5;%
OldLace,.992,.96,.9;%
Olive,.5,.5,0;%
OliveDrab,.42,.556,.136;%
Orange,1,.648,0;%
OrangeRed,1,.27,0;%
Orchid,.855,.44,.84;%
PaleGoldenrod,.932,.91,.668;%
PaleGreen,.596,.985,.596;%
PaleTurquoise,.688,.932,.932;%
PaleVioletRed,.86,.44,.576;%
PapayaWhip,1,.936,.835;%
PeachPuff,1,.855,.725;%
Peru,.804,.52,.248;%
Pink,1,.752,.796;%
Plum,.868,.628,.868;%
PowderBlue,.69,.88,.9;%
Purple,.5,0,.5;%
Red,1,0,0;%
RosyBrown,.736,.56,.56;%
RoyalBlue,.255,.41,.884;%
SaddleBrown,.545,.27,.075;%
Salmon,.98,.5,.448;%
SandyBrown,.956,.644,.376;%
SeaGreen,.18,.545,.34;%
Seashell,1,.96,.932;%
Sienna,.628,.32,.176;%
Silver,.752,.752,.752;%
SkyBlue,.53,.808,.92;%
SlateBlue,.415,.352,.804;%
SlateGray,.44,.5,.565;%
SlateGrey,.44,.5,.565;%
Snow,1,.98,.98;%
SpringGreen,0,1,.498;%
SteelBlue,.275,.51,.705;%
Tan,.824,.705,.55;%
Teal,0,.5,.5;%
Thistle,.848,.75,.848;%
Tomato,1,.39,.28;%
Turquoise,.25,.88,.815;%
Violet,.932,.51,.932;%
VioletRed,.816,.125,.565;%
Wheat,.96,.87,.7;%
White,1,1,1;%
WhiteSmoke,.96,.96,.96;%
Yellow,1,1,0;%
YellowGreen,.604,.804,.196}
%</svgnames>
%
%<*x11names>
\preparecolorset{rgb}{}{}{%
AntiqueWhite1,1,.936,.86;%
AntiqueWhite2,.932,.875,.8;%
AntiqueWhite3,.804,.752,.69;%
AntiqueWhite4,.545,.512,.47;%
Aquamarine1,.498,1,.83;%
Aquamarine2,.464,.932,.776;%
Aquamarine3,.4,.804,.668;%
Aquamarine4,.27,.545,.455;%
Azure1,.94,1,1;%
Azure2,.88,.932,.932;%
Azure3,.756,.804,.804;%
Azure4,.512,.545,.545;%
Bisque1,1,.894,.77;%
Bisque2,.932,.835,.716;%
Bisque3,.804,.716,.62;%
Bisque4,.545,.49,.42;%
Blue1,0,0,1;%
Blue2,0,0,.932;%
Blue3,0,0,.804;%
Blue4,0,0,.545;%
Brown1,1,.25,.25;%
Brown2,.932,.23,.23;%
Brown3,.804,.2,.2;%
Brown4,.545,.136,.136;%
Burlywood1,1,.828,.608;%
Burlywood2,.932,.772,.57;%
Burlywood3,.804,.668,.49;%
Burlywood4,.545,.45,.332;%
CadetBlue1,.596,.96,1;%
CadetBlue2,.556,.898,.932;%
CadetBlue3,.48,.772,.804;%
CadetBlue4,.325,.525,.545;%
Chartreuse1,.498,1,0;%
Chartreuse2,.464,.932,0;%
Chartreuse3,.4,.804,0;%
Chartreuse4,.27,.545,0;%
Chocolate1,1,.498,.14;%
Chocolate2,.932,.464,.13;%
Chocolate3,.804,.4,.112;%
Chocolate4,.545,.27,.075;%
Coral1,1,.448,.336;%
Coral2,.932,.415,.312;%
Coral3,.804,.356,.27;%
Coral4,.545,.244,.185;%
Cornsilk1,1,.972,.864;%
Cornsilk2,.932,.91,.804;%
Cornsilk3,.804,.785,.694;%
Cornsilk4,.545,.532,.47;%
Cyan1,0,1,1;%
Cyan2,0,.932,.932;%
Cyan3,0,.804,.804;%
Cyan4,0,.545,.545;%
DarkGoldenrod1,1,.725,.06;%
DarkGoldenrod2,.932,.68,.055;%
DarkGoldenrod3,.804,.585,.048;%
DarkGoldenrod4,.545,.396,.03;%
DarkOliveGreen1,.792,1,.44;%
DarkOliveGreen2,.736,.932,.408;%
DarkOliveGreen3,.635,.804,.352;%
DarkOliveGreen4,.43,.545,.24;%
DarkOrange1,1,.498,0;%
DarkOrange2,.932,.464,0;%
DarkOrange3,.804,.4,0;%
DarkOrange4,.545,.27,0;%
DarkOrchid1,.75,.244,1;%
DarkOrchid2,.698,.228,.932;%
DarkOrchid3,.604,.196,.804;%
DarkOrchid4,.408,.132,.545;%
DarkSeaGreen1,.756,1,.756;%
DarkSeaGreen2,.705,.932,.705;%
DarkSeaGreen3,.608,.804,.608;%
DarkSeaGreen4,.41,.545,.41;%
DarkSlateGray1,.592,1,1;%
DarkSlateGray2,.552,.932,.932;%
DarkSlateGray3,.475,.804,.804;%
DarkSlateGray4,.32,.545,.545;%
DeepPink1,1,.08,.576;%
DeepPink2,.932,.07,.536;%
DeepPink3,.804,.064,.464;%
DeepPink4,.545,.04,.312;%
DeepSkyBlue1,0,.75,1;%
DeepSkyBlue2,0,.698,.932;%
DeepSkyBlue3,0,.604,.804;%
DeepSkyBlue4,0,.408,.545;%
DodgerBlue1,.116,.565,1;%
DodgerBlue2,.11,.525,.932;%
DodgerBlue3,.094,.455,.804;%
DodgerBlue4,.064,.305,.545;%
Firebrick1,1,.19,.19;%
Firebrick2,.932,.172,.172;%
Firebrick3,.804,.15,.15;%
Firebrick4,.545,.1,.1;%
Gold1,1,.844,0;%
Gold2,.932,.79,0;%
Gold3,.804,.68,0;%
Gold4,.545,.46,0;%
Goldenrod1,1,.756,.145;%
Goldenrod2,.932,.705,.132;%
Goldenrod3,.804,.608,.112;%
Goldenrod4,.545,.41,.08;%
Green1,0,1,0;%
Green2,0,.932,0;%
Green3,0,.804,0;%
Green4,0,.545,0;%
Honeydew1,.94,1,.94;%
Honeydew2,.88,.932,.88;%
Honeydew3,.756,.804,.756;%
Honeydew4,.512,.545,.512;%
HotPink1,1,.43,.705;%
HotPink2,.932,.415,.655;%
HotPink3,.804,.376,.565;%
HotPink4,.545,.228,.385;%
IndianRed1,1,.415,.415;%
IndianRed2,.932,.39,.39;%
IndianRed3,.804,.332,.332;%
IndianRed4,.545,.228,.228;%
Ivory1,1,1,.94;%
Ivory2,.932,.932,.88;%
Ivory3,.804,.804,.756;%
Ivory4,.545,.545,.512;%
Khaki1,1,.965,.56;%
Khaki2,.932,.9,.52;%
Khaki3,.804,.776,.45;%
Khaki4,.545,.525,.305;%
LavenderBlush1,1,.94,.96;%
LavenderBlush2,.932,.88,.898;%
LavenderBlush3,.804,.756,.772;%
LavenderBlush4,.545,.512,.525;%
LemonChiffon1,1,.98,.804;%
LemonChiffon2,.932,.912,.75;%
LemonChiffon3,.804,.79,.648;%
LemonChiffon4,.545,.536,.44;%
LightBlue1,.75,.936,1;%
LightBlue2,.698,.875,.932;%
LightBlue3,.604,.752,.804;%
LightBlue4,.408,.512,.545;%
LightCyan1,.88,1,1;%
LightCyan2,.82,.932,.932;%
LightCyan3,.705,.804,.804;%
LightCyan4,.48,.545,.545;%
LightGoldenrod1,1,.925,.545;%
LightGoldenrod2,.932,.864,.51;%
LightGoldenrod3,.804,.745,.44;%
LightGoldenrod4,.545,.505,.298;%
LightPink1,1,.684,.725;%
LightPink2,.932,.635,.68;%
LightPink3,.804,.55,.585;%
LightPink4,.545,.372,.396;%
LightSalmon1,1,.628,.48;%
LightSalmon2,.932,.585,.448;%
LightSalmon3,.804,.505,.385;%
LightSalmon4,.545,.34,.26;%
LightSkyBlue1,.69,.888,1;%
LightSkyBlue2,.644,.828,.932;%
LightSkyBlue3,.552,.712,.804;%
LightSkyBlue4,.376,.484,.545;%
LightSteelBlue1,.792,.884,1;%
LightSteelBlue2,.736,.824,.932;%
LightSteelBlue3,.635,.71,.804;%
LightSteelBlue4,.43,.484,.545;%
LightYellow1,1,1,.88;%
LightYellow2,.932,.932,.82;%
LightYellow3,.804,.804,.705;%
LightYellow4,.545,.545,.48;%
Magenta1,1,0,1;%
Magenta2,.932,0,.932;%
Magenta3,.804,0,.804;%
Magenta4,.545,0,.545;%
Maroon1,1,.204,.7;%
Maroon2,.932,.19,.655;%
Maroon3,.804,.16,.565;%
Maroon4,.545,.11,.385;%
MediumOrchid1,.88,.4,1;%
MediumOrchid2,.82,.372,.932;%
MediumOrchid3,.705,.32,.804;%
MediumOrchid4,.48,.215,.545;%
MediumPurple1,.67,.51,1;%
MediumPurple2,.624,.475,.932;%
MediumPurple3,.536,.408,.804;%
MediumPurple4,.365,.28,.545;%
MistyRose1,1,.894,.884;%
MistyRose2,.932,.835,.824;%
MistyRose3,.804,.716,.71;%
MistyRose4,.545,.49,.484;%
NavajoWhite1,1,.87,.68;%
NavajoWhite2,.932,.81,.63;%
NavajoWhite3,.804,.7,.545;%
NavajoWhite4,.545,.475,.37;%
OliveDrab1,.752,1,.244;%
OliveDrab2,.7,.932,.228;%
OliveDrab3,.604,.804,.196;%
OliveDrab4,.41,.545,.132;%
Orange1,1,.648,0;%
Orange2,.932,.604,0;%
Orange3,.804,.52,0;%
Orange4,.545,.352,0;%
OrangeRed1,1,.27,0;%
OrangeRed2,.932,.25,0;%
OrangeRed3,.804,.215,0;%
OrangeRed4,.545,.145,0;%
Orchid1,1,.512,.98;%
Orchid2,.932,.48,.912;%
Orchid3,.804,.41,.79;%
Orchid4,.545,.28,.536;%
PaleGreen1,.604,1,.604;%
PaleGreen2,.565,.932,.565;%
PaleGreen3,.488,.804,.488;%
PaleGreen4,.33,.545,.33;%
PaleTurquoise1,.732,1,1;%
PaleTurquoise2,.684,.932,.932;%
PaleTurquoise3,.59,.804,.804;%
PaleTurquoise4,.4,.545,.545;%
PaleVioletRed1,1,.51,.67;%
PaleVioletRed2,.932,.475,.624;%
PaleVioletRed3,.804,.408,.536;%
PaleVioletRed4,.545,.28,.365;%
PeachPuff1,1,.855,.725;%
PeachPuff2,.932,.796,.68;%
PeachPuff3,.804,.688,.585;%
PeachPuff4,.545,.468,.396;%
Pink1,1,.71,.772;%
Pink2,.932,.664,.72;%
Pink3,.804,.57,.62;%
Pink4,.545,.39,.424;%
Plum1,1,.732,1;%
Plum2,.932,.684,.932;%
Plum3,.804,.59,.804;%
Plum4,.545,.4,.545;%
Purple1,.608,.19,1;%
Purple2,.57,.172,.932;%
Purple3,.49,.15,.804;%
Purple4,.332,.1,.545;%
Red1,1,0,0;%
Red2,.932,0,0;%
Red3,.804,0,0;%
Red4,.545,0,0;%
RosyBrown1,1,.756,.756;%
RosyBrown2,.932,.705,.705;%
RosyBrown3,.804,.608,.608;%
RosyBrown4,.545,.41,.41;%
RoyalBlue1,.284,.464,1;%
RoyalBlue2,.264,.43,.932;%
RoyalBlue3,.228,.372,.804;%
RoyalBlue4,.152,.25,.545;%
Salmon1,1,.55,.41;%
Salmon2,.932,.51,.385;%
Salmon3,.804,.44,.33;%
Salmon4,.545,.298,.224;%
SeaGreen1,.33,1,.624;%
SeaGreen2,.305,.932,.58;%
SeaGreen3,.264,.804,.5;%
SeaGreen4,.18,.545,.34;%
Seashell1,1,.96,.932;%
Seashell2,.932,.898,.87;%
Seashell3,.804,.772,.75;%
Seashell4,.545,.525,.51;%
Sienna1,1,.51,.28;%
Sienna2,.932,.475,.26;%
Sienna3,.804,.408,.224;%
Sienna4,.545,.28,.15;%
SkyBlue1,.53,.808,1;%
SkyBlue2,.494,.752,.932;%
SkyBlue3,.424,.65,.804;%
SkyBlue4,.29,.44,.545;%
SlateBlue1,.512,.435,1;%
SlateBlue2,.48,.404,.932;%
SlateBlue3,.41,.35,.804;%
SlateBlue4,.28,.235,.545;%
SlateGray1,.776,.888,1;%
SlateGray2,.725,.828,.932;%
SlateGray3,.624,.712,.804;%
SlateGray4,.424,.484,.545;%
Snow1,1,.98,.98;%
Snow2,.932,.912,.912;%
Snow3,.804,.79,.79;%
Snow4,.545,.536,.536;%
SpringGreen1,0,1,.498;%
SpringGreen2,0,.932,.464;%
SpringGreen3,0,.804,.4;%
SpringGreen4,0,.545,.27;%
SteelBlue1,.39,.72,1;%
SteelBlue2,.36,.675,.932;%
SteelBlue3,.31,.58,.804;%
SteelBlue4,.21,.392,.545;%
Tan1,1,.648,.31;%
Tan2,.932,.604,.288;%
Tan3,.804,.52,.248;%
Tan4,.545,.352,.17;%
Thistle1,1,.884,1;%
Thistle2,.932,.824,.932;%
Thistle3,.804,.71,.804;%
Thistle4,.545,.484,.545;%
Tomato1,1,.39,.28;%
Tomato2,.932,.36,.26;%
Tomato3,.804,.31,.224;%
Tomato4,.545,.21,.15;%
Turquoise1,0,.96,1;%
Turquoise2,0,.898,.932;%
Turquoise3,0,.772,.804;%
Turquoise4,0,.525,.545;%
VioletRed1,1,.244,.59;%
VioletRed2,.932,.228,.55;%
VioletRed3,.804,.196,.47;%
VioletRed4,.545,.132,.32;%
Wheat1,1,.905,.73;%
Wheat2,.932,.848,.684;%
Wheat3,.804,.73,.59;%
Wheat4,.545,.494,.4;%
Yellow1,1,1,0;%
Yellow2,.932,.932,0;%
Yellow3,.804,.804,0;%
Yellow4,.545,.545,0;%
Gray0,.745,.745,.745;%
Green0,0,1,0;%
Grey0,.745,.745,.745;%
Maroon0,.69,.19,.376;%
Purple0,.628,.125,.94}
%</x11names>
%
%<*test1>
\documentclass[a4paper]{article}

\usepackage{array,multicol,multido,textcomp}
\usepackage{xcolor}[2005/03/24]
\usepackage[hmargin={1.25cm,.75cm},vmargin=1.25cm,footskip=.5cm,nohead]{geometry}

\newdimen\Argwd
\newdimen\Errsum
\newdimen\Errmax
\newcount\Colors
\newcount\Errors
\newcount\Loops
\newif\ifColor

\makeatletter

\def\DisplayColor#1#2#3%
 {\global\advance\Loops\@ne
  \@tempdima=#1\p@ \edef\@tempa{\the\@tempdima}%
  \edef\Colori{\strip@pt\@tempdima}%
  \ifColor
    \@tempdimb=#2\p@ \edef\@tempb{\the\@tempdimb}%
    \@tempdimc=#3\p@ \edef\@tempc{\the\@tempdimc}%
    \edef\Colori{\Colori,\strip@pt\@tempdimb,\strip@pt\@tempdimc}%
  \fi
  \Make12\Make23\Make34\Make45\Make56\Make67%
  \begin{tabular}{@{}ll*4{@{ }p{\Argwd}}@{}}
  \Row1\Row2\Row3\Row4\Row5\Row6\noalign{\medskip}
  \end{tabular}
  \global\advance\Colors1 \relax\ifnum\Colors=30 \newpage\global\Colors=0 \fi}

\def\Make#1#2%
 {\edef\@@tmp
   {\noexpand\definecolor{Color\romannumeral#1}%
     {\@nameuse{Model\romannumeral#1}}{\@nameuse{Color\romannumeral#1}}%
   }\@@tmp
  \relax\ifnum#1<6
    \convertcolorspec
     {\@nameuse{Model\romannumeral#1}}{\@nameuse{Color\romannumeral#1}}%
     {\@nameuse{Model\romannumeral#2}}\@@tmp
    \XC@let@cN{Color\romannumeral#2}\@@tmp
  \fi}
\def\Row#1%
 {\edef\@@tmp{\@nameuse{Color\romannumeral#1}\@nameuse{Row\romannumeral#1}}%
  \expandafter\@Row\@@tmp}

\def\@Row#1,#2,#3,#4,#5#6%
 {\fcolorbox{black}{Color\romannumeral#5}{\phantom{000}}&
  \csname Model\romannumeral#5\endcsname &
  \ifodd#6 \@Error{#1}\@tempa \else #1\fi &
  \ifColor\ifodd#6 \@Error{#2}\@tempb \else #2\fi\else #2\fi &
  \ifColor\ifodd#6 \@Error{#3}\@tempc \else #3\fi\else #3\fi &
  #4\\}

\def\defModel#1#2#3#4% number, name, commas, check
 {\@namedef{Model\romannumeral#1}{#2}\@namedef{Row\romannumeral#1}{#3#1#4}}

\def\@Error#1#2%
 {\ifdim#1\p@=#2\relax #1\else
    \textcolor{red}{#1}\global\advance\Errors\@ne
    \dimen@#2\advance\dimen@-#1\p@ \ifdim\dimen@<\z@\dimen@-\dimen@\fi
    \global\advance\Errsum\dimen@
    \ifdim\dimen@>\Errmax \global\Errmax\dimen@
      \dimen@\@tempa \edef\@@tempa{\strip@pt\dimen@}%
      \ifColor
        \dimen@\@tempb \edef\@@tempb{\strip@pt\dimen@}%
        \dimen@\@tempc \edef\@@tempc{\strip@pt\dimen@}%
        \xdef\Errcase{\@@tempa\noexpand~\@@tempb\noexpand~\@@tempc}%
      \else
        \xdef\Errcase{\@@tempa}%
      \fi
    \fi
  \fi}

\newenvironment{ColorTest}[1]%
 {\addcontentsline{toc}{section}{Color conversion chain: #1}
  \Colors\z@ \Errors\z@ \Loops\z@ \Errsum\z@ \Errmax\z@
  \begin{multicols}{3}\centering}%
 {\end{multicols}\bigskip
  \begin{tabular}{@{}l@{ = }l@{}}
  \multicolumn2{@{}l@{}}{Rounding error statistics of color conversion test}\\
  \multicolumn2{@{}l@{}}%
   {$\textbf\Modeli\rightarrow \textbf\Modelii\rightarrow \textbf\Modeliii\rightarrow
     \textbf\Modeliv\rightarrow \textbf\Modelv\rightarrow \textbf\Modelvi$:}\\[1ex]
  number of examples& \number\Loops\\
  number of errors& \number\Errors\\
  sum of absolute errors& \strip@pt\Errsum\\
  \ifnum\Errors>\z@
    average error& \Errsum100\Errsum \dimen@\Errors\p@ \divide\dimen@\@m
    \rdivide\Errsum\dimen@ \strip@pt\Errsum\textperiodcentered10\textsuperscript{-5}\\
    maximum error& \strip@pt\Errmax{} @ \Modeli~\Errcase\\
  \fi
  \end{tabular}\newpage}

\def\thepage{\sffamily\footnotesize\@nameuse{@arabic}\c@page}

\makeatother

\parindent0pt
\setcounter{secnumdepth}{-1}

\begin{document}

\title{Color conversion with the \textsf{xcolor} package --- Examples}
\author{Dr. Uwe Kern\thanks{This file is part of the \textsf{xcolor} distribution which can be downloaded from the CTAN mirrors (\texttt{macros/latex/contrib/xcolor/}) or the homepage \texttt{www.ukern.de/tex/xcolor.html}. Please send error reports and suggestions for improvements to \texttt{https://github.com/latex3/xcolor/issues}.}}
\date{\XCfileversion{} (\XCfiledate)}
\maketitle

\begin{center}
\begin{minipage}{.75\textwidth}
\tableofcontents
\end{minipage}
\end{center}
\newpage

\sffamily\footnotesize
\fboxrule.25pt
\settowidth\Argwd{0.00000}

\defModel1{rgb}{,,}0
\defModel2{cmy}{,,}0
\defModel3{cmyk}{,}0
\defModel4{hsb}{,,}0
\defModel5{rgb}{,,}1
\defModel6{gray}{,,,,}0
\Colortrue

\begin{ColorTest}{rgb, cmy, cmyk, hsb, rgb, gray}
\multido{\nR=1+-.1}{11}%
 {\multido{\nG=1+-.1}{11}%
  {\multido{\nB=1+-.1}{11}%
   {\DisplayColor{\nR}{\nG}{\nB}}}}%
\end{ColorTest}

\defModel1{gray}{,,,,}0
\defModel2{rgb}{,,}0
\defModel3{cmy}{,,}0
\defModel4{cmyk}{,}0
\defModel5{hsb}{,,}0
\defModel6{gray}{,,,,}1
\Colorfalse

\begin{ColorTest}{gray, rgb, cmy, cmyk, hsb, gray}
\multido{\nG=1+-.01}{101}{\DisplayColor{\nG}{}{}}%
\end{ColorTest}

\end{document}
%</test1>
%
%<*test2>
\listfiles
\documentclass{article}

\usepackage{pstricks,pst-tree,pst-slpe,multido}
\makeatletter
\makeatother
\usepackage{fancyvrb-ex}

\def\textfraction{0.01}
\def\topfraction{0.99}
\def\bottomfraction{0.99}
\def\floatpagefraction{0.70}
\setlength\abovecaptionskip{12pt}
\setlength\belowcaptionskip{10pt}
\parindent0pt

\begin{document}
\title{Color extensions with the \textsf{xcolor} package --- \textsf{pstricks} examples}
\author{Dr. Uwe Kern\thanks{This file is part of the \textsf{xcolor} distribution which can be downloaded from the CTAN mirrors (\texttt{macros/latex/contrib/xcolor/}) or the homepage \texttt{www.ukern.de/tex/xcolor.html}. Please send error reports and suggestions for improvements to \texttt{https://github.com/latex3/xcolor/issues}.}}
\date{\XCfileversion{} (\XCfiledate)}
\maketitle

\centerline{\fcolorbox{red}{yellow!50}{This document is not suitable for pdf\LaTeX! Please use \LaTeX{} + \texttt{dvips} etc.}}

\begin{figure}[htb]\caption{Modified version of an example from the \textsf{pstricks} manual; requires \textsf{pst-tree}}
\begin{SideBySideExample}[fontsize=\footnotesize,xrightmargin=.45\textwidth]
  \begin{pspicture}(-0.5,-2)(3,2)
  \cnode(0,0){.5cm}{root}
  \cnode*[linecolor=red](3,1.5){4pt}{A}
  \cnode*[linecolor=red!72.5375!blue](3,0){4pt}{B}
  \cnode*[linecolor=-red](3,-1.5){4pt}{C}
  \psset{nodesep=3pt}
  \ncline[linecolor=green!50!red]{root}{A}
  \ncline[linecolor=blue]{root}{B}
  \ncline[linecolor=-green!50!red]{root}{C}
  \end{pspicture}
\end{SideBySideExample}
\end{figure}

\begin{figure}[htb]\caption{Moving from one color to its complement; requires \textsf{pst-slpe}}
\begin{SideBySideExample}[fontsize=\footnotesize,xrightmargin=.45\textwidth]
  \begin{pspicture}(4,1)
  \psframe[fillstyle=slope,
           slopeangle=30,
           slopebegin=red!72.5375!blue,
           slopeend=-red!72.5375!blue](4,1)
  \end{pspicture}
\end{SideBySideExample}
\end{figure}

\begin{figure}[htb]\caption{Explicit color specification via a loop command; requires \textsf{multido}}
\begin{SideBySideExample}[fontsize=\footnotesize,xrightmargin=.45\textwidth]
  \psset{unit=1.75}%
  \begin{pspicture}(0,-1)(2,1)
  \multido{\rHue=0.00+0.01}{100}%
   {\pscircle[linewidth=0.01,
              linecolor={[hsb]{\rHue,1,1}}](1,0){\rHue}}
  \end{pspicture}
\end{SideBySideExample}
\end{figure}

\begin{figure}[htb]\caption{Color series --- modified version of an example from the \textsf{pst-fill} manual; note that the \texttt{\string\multirput} command does not give the desired result here}
\begin{SideBySideExample}[fontsize=\footnotesize,xrightmargin=.45\textwidth]
  \newcommand*\Sheep
   {\begin{pspicture}(3,1.5)
    \pscustom[liftpen=2,fillstyle=solid,fillcolor=sheep!!+]%
     {\pscurve(0.5,-0.2)(0.6,0.5)(0.2,1.3)(0,1.5)(0,1.5)
       (0.4,1.3)(0.8,1.5)(2.2,1.9)(3,1.5)(3,1.5)(3.2,1.3)
       (3.6,0.5)(3.4,-0.3)(3,0)(2.2,0.4)(0.5,-0.2)}
    \pscircle*(2.65,1.25){0.12\psunit}% Eye
    \psccurve*(3.5,0.3)(3.35,0.45)(3.5,0.6)(3.6,0.4)% Muzzle
    \pscurve(3,0.35)(3.3,0.1)(3.6,0.05)% Mouth
    \pscurve(2.3,1.3)(2.1,1.5)(2.15,1.7)
    \pscurve(2.1,1.7)(2.35,1.6)(2.45,1.4)% Ear
    \end{pspicture}}
  \definecolorseries
   {sheep}{rgb}{step}[rgb]{.95,.85,.55}{.17,.47,.37}
  \resetcolorseries{sheep}
  \psset{unit=0.4}
  \begin{pspicture}(-3,-6)(0,7.5)
  \Multido{\ry=6.0+-1.5}{5}%
   {\rput(0,\ry){\Multido{}{5}{\Sheep}}}
  \resetcolorseries{sheep}%
  \multirput(-6,-6)(3,0)5{\Sheep}
  \end{pspicture}
\end{SideBySideExample}
\end{figure}

\begin{figure}[htb]\caption{Interaction with native PostScript code --- $\gamma$-corrected wavelengths}
\begin{Example}[fontsize=\footnotesize,xrightmargin=.45\textwidth]
  \newcount\WL \unitlength.75pt
  \def\WaveToPS#1%
   {\definecolor{tmp}{rgb:wave}{#1}\extractcolorspecs{tmp}\tmpm\tmpc
    \expandafter\WaveToPSi\tmpc,}
  \def\WaveToPSi#1,#2,#3,{\pstVerb{/Red{#1}def /Green{#2}def /Blue{#3}def}}
  \def\DisplayBar#1#2%
   {\linethickness{1.25\unitlength}\WL=360
    \pstVerb{/Gamma{#1}def}%
    \multiput(360,#2)(1,0){456}%
     {\WaveToPS{\the\WL}{\color{lambda}\line(0,1){50}}\global\advance\WL1}%
    \linethickness{0.25\unitlength}\WL=360
    \multiput(360,#2)(20,0){23}%
     {\picture(0,0)
      \line(0,-1){5}\multiput(5,0)(5,0){3}{\line(0,-1){2.5}}%
      \put(0,-10){\makebox(0,0){\the\WL}}\global\advance\WL20
      \endpicture}%
    \put(350,#2){\makebox(0,50)[r]{\small$\gamma$\,=\,#1}}}
  \pstVerb{/Corr{dup 0 gt {Gamma exp}if}def}
  \definecolor[ps]{lambda}{rgb}{Red Corr Green Corr Blue Corr}%
  \begin{picture}(510,345)(310,-10)
  \sffamily\tiny
  \DisplayBar{0.4}{0}%
  \DisplayBar{0.6}{70}%
  \DisplayBar{0.8}{140}%
  \DisplayBar{1.0}{210}%
  \DisplayBar{1.2}{280}%
  \end{picture}
\end{Example}
\end{figure}

\end{document}
%</test2>
%
%<*test3>
\listfiles
\documentclass[a4paper]{article}
\usepackage[showerrors,table,dvipsnames*,hyperref]{xcolor}[2005/12/21]
\usepackage[margin=2.25cm]{geometry}
\usepackage
 [\GinDriver,hyperindex=false,bookmarks,bookmarksopen,bookmarksopenlevel=1,%
  pdftitle={xcolor3 \XCfileversion{} (\XCfiledate)},pdfauthor={Dr. Uwe Kern},%
  pdfsubject={Color extensions for LaTeX and pdfLaTeX},%
  pdfkeywords={xcolor,color,colour,model,tint,tone,shade,harmony,spot,latex,pdftex,dvips,%
   conversion,blend,mix,mask,separation,rgb,cmy,cmyk,hsb,gray,html,wave,thsb,wheel}]{hyperref}

\tracingcolors=4
%%\tracingcolors=3
%%\tracingcolors=2
%%\tracingcolors=1
%%\tracingcolors=0

\parindent0pt
\pagecolor{gray!25}

\definecolors{JungleGreen,DarkOrchid}

\begin{document}
\title{Color extensions with the \textsf{xcolor} package --- various examples}
\author{\href{mailto:xcolor@ukern.de}{\fboxrule0pt\fboxsep2pt\fbox{Dr. Uwe Kern}}}
\date{\XCfileversion{} (\XCfiledate)
\thanks{This file (\texttt{\jobname.tex}) is part of the \textsf{xcolor} distribution which can be downloaded from the CTAN mirrors \texttt{\href{http://www.ctan.org/tex-archive/macros/latex/contrib/xcolor/}{CTAN/macros/latex/contrib/xcolor/}} or the homepage \texttt{\href{http://www.ukern.de/tex/xcolor.html}{www.ukern.de/tex/xcolor.html}}. Please send error reports and suggestions for improvements to \url{https://github.com/latex3/xcolor/issues}.}}
\maketitle

The purpose of this file is to demonstrate a variety of capabilities including the logging facilities of the \textsf{xcolor} package.
By playing around with different values of \texttt{\string\tracingcolors}, one can observe the different behavior in the \texttt{log} file.

\section{Predefined colors}

\begingroup
\small\sffamily
\rowcolors1{}{}
\begin{testcolors}[rgb,cmyk,hsb,HTML,gray]
\testcolor{red}
\testcolor{green}
\testcolor{blue}
\testcolor{cyan}
\testcolor{magenta}
\testcolor{yellow}
\testcolor{orange}
\testcolor{violet}
\testcolor{purple}
\testcolor{brown}
\testcolor{pink}
\testcolor{olive}
\testcolor{black}
\testcolor{darkgray}
\testcolor{gray}
\testcolor{lightgray}
\testcolor{white}
\noalign{\medskip}\hline\noalign{\medskip}
\testcolor{-red}
\testcolor{-green}
\testcolor{-blue}
\testcolor{-cyan}
\testcolor{-magenta}
\testcolor{-yellow}
\testcolor{-orange}
\testcolor{-violet}
\testcolor{-purple}
\testcolor{-brown}
\testcolor{-pink}
\testcolor{-olive}
\testcolor{-black}
\testcolor{-darkgray}
\testcolor{-gray}
\testcolor{-lightgray}
\testcolor{-white}
\noalign{\medskip}\hline\noalign{\medskip}
\testcolor{JungleGreen}
\testcolor{DarkOrchid}
\noalign{\medskip}\hline\noalign{\medskip}
\testcolor{-JungleGreen}
\testcolor{-DarkOrchid}
\end{testcolors}
\endgroup

\vfill

\clearpage
\nopagecolor
Switch off \texttt{\string\pagecolor} by \texttt{\string\nopagecolor}.
The driver may or may not obey this command \dots

\clearpage
\pagecolor{white}

\section{Color definition and application}

\providecolor{dummy}{rgb}{.6,.5,.4}
\definecolor{dummy}{rgb}{.6,.5,.4}
\providecolor{dummy}{rgb}{.6,.5,.4}
\hbox{\textcolor{dummy}{Test with \texttt{\string\definecolor}}}

\bigskip

Comma-separated and space-separated definitions:

\definecolor{c1}{rgb}{.7,.6,.5}
\definecolor{c2}{rgb}{.7 .6 .5}
\colorlet{c1a}{c1}
\colorlet{c2a}{c2}

\textcolor{c1}{identical} =
\textcolor{c2}{identical} =
\textcolor{c1a}{identical} =
\textcolor{c2a}{identical} =
\textcolor[rgb]{.7,.6,.5}{identical} =
\textcolor[rgb]{.7 .6 .5}{identical} =
\textcolor{rgb,10:red,7;green,6;blue,5}{identical}
\textcolor{rgb,15:red,10.5;green,9;blue,7.5}{identical}

\medskip

\begingroup
\sffamily
\begin{testcolors}
\testcolor{c1}
\testcolor{c2}
\testcolor{c1a}
\testcolor{c2a}
\testcolor[rgb]{.7,.6,.5}
\testcolor[rgb]{.7 .6 .5}
\testcolor{rgb,10:red,7;green,6;blue,5}
\testcolor{rgb,15:red,10.5;green,9;blue,7.5}
\end{testcolors}
\endgroup

\bigskip

\textcolor{rgb:red!50,4;green!25,2}{Another extended color expression (rgb:red!50,4;green!25,2)}.

\bigskip

\begingroup
\color{black}
Test with named colors:\par
\color{blue}
Test: \textcolor[named]{JungleGreen}{JungleGreen};
Test: \textcolor{JungleGreen}{JungleGreen};
Test: \textcolor{JungleGreen!50!DarkOrchid}{JungleGreen!50!DarkOrchid};
Test: \textcolor{green!50!red}{green!50!red}.
\endgroup

\bigskip

{\color[rgb]{.4,.5,.6}Test with \texttt{\string\color}}

\bigskip
Current color application:\par
\def\test{current, \textcolor{.!50}{50\%}, \textcolor{-.}{complement},
          \textcolor{yellow!50!.}{mix}}
\textcolor{blue}{\test} and \textcolor{red}{\test},\par
\def\Test{\color{.!80}Test}
\textcolor{blue}{\Test\Test\Test\Test\Test} and
\textcolor{red}{\Test\Test\Test\Test\Test}.

\bigskip
Current color test with \texttt{\string\definecolorseries}:\par
\begingroup
\color{blue}
\definecolorseries{foo}{rgb}{last}{.}{-.}
\resetcolorseries[5]{foo}
\def\test{\hbox to 1em{{\color{foo!!+}\vrule width 1em height 1.5ex}}}
Test\test\test\test\test\test\test Test

\resetcolorseries[5]{foo}
\def\test{\hbox to 1em{{\color{foo!!++}\vrule width 1em height 1.5ex}}}
Test\test\test\test\test\test\test Test

\resetcolorseries[5]{foo}
\def\test{\hbox to 1em{{\color{foo!![2]}\vrule width 1em height 1.5ex}}}
Test\test\test\test\test\test\test Test

\endgroup

\section{Color in tables}

\rowcolors[\hline]{1}{green!25}{yellow!50}
\begin{tabular}{ll}
test & row \number\rownum\\
test & row \number\rownum\\
\rowcolor{blue!25}
test & row \number\rownum\\
test & row \number\rownum\\
\hiderowcolors
test & row \number\rownum\\
test & row \number\rownum\\
\showrowcolors
test & row \number\rownum\\
test & row \number\rownum\\
\multicolumn{1}%
 {>{\columncolor{red!12}}l}{test} & row \number\rownum\\
\end{tabular}

\section{Color information}

Type test:
\makeatletter
\@namedef{\string\color@foo1}{foo1{}{}{}{}}\edef\tempa{\XC@type{foo1}}\tempa
\@namedef{\string\color@foo2}{\xcolor@{foo2}{}{}{}}\edef\tempb{\XC@type{foo2}}\tempb
\@namedef{\string\color@foo3}{\xcolor@{}{foo3}{}{}}\edef\tempc{\XC@type{foo3}}\tempc
\@namedef{\string\color@foo4}{\xcolor@{}{}{foo4}{}}\edef\tempd{\XC@type{foo4}}\tempd
\makeatother

\end{document}
%</test3>
%
%<*test4>
\listfiles
\documentclass[a4paper]{article}
%%\PassOptionsToPackage{dvipdf}{xcolor}
%%\PassOptionsToPackage{dvipdfm}{xcolor}
%%\PassOptionsToPackage{dvipdfmx}{xcolor}
%%\PassOptionsToPackage{dvips}{xcolor}
%%\PassOptionsToPackage{dvipsone}{xcolor}
%%\PassOptionsToPackage{dviwin}{xcolor}
%%\PassOptionsToPackage{dviwindo}{xcolor}
%%\PassOptionsToPackage{emtex}{xcolor}
%%\PassOptionsToPackage{oztex}{xcolor}
%%\PassOptionsToPackage{luatex}{xcolor}
%%\PassOptionsToPackage{pctex32}{xcolor}
%%\PassOptionsToPackage{pctexhp}{xcolor}
%%\PassOptionsToPackage{pctexps}{xcolor}
%%\PassOptionsToPackage{pctexwin}{xcolor}
%%\PassOptionsToPackage{pdftex}{xcolor}
%%\PassOptionsToPackage{tcidvi}{xcolor}
%%\PassOptionsToPackage{textures}{xcolor}
%%\PassOptionsToPackage{truetex}{xcolor}
%%\PassOptionsToPackage{vtex}{xcolor}
%%\PassOptionsToPackage{xdvi}{xcolor}
%%\PassOptionsToPackage{xetex}{xcolor}
%%\PassOptionsToPackage{prologue}{xcolor}
%%\PassOptionsToPackage{kernelfbox}{xcolor}
%%\PassOptionsToPackage{xcdraw}{xcolor}
\usepackage[emtex]{graphics}
\usepackage{xcolor}
\usepackage[margin=1.25cm,noheadfoot]{geometry}
\pagestyle{empty}
\parindent0pt
\nofiles
\let\familydefault\sfdefault
\makeatletter
%% special code for driver list construction via
%%   latex \def\xcdriverlist{1}\input{xcolor4}
%% writes list in file xcolor4.tmp
\@ifundefined{xcdriverlist}{}{%
%% optional * (skips test whether #1.def exists)
%% #1 = driver
%% #2 = additional options (stored in \options@#2)
%% #3 = processor 1
%% #4 = output format 1
%% #5 = processor 2
%% #6 = output format 2
%% #7 = processor 3
%% #8 = output format 3
\def\0{\@ifstar
  {\def\@@tmp####1{\@firstoftwo}\1}{\def\@@tmp####1{\IfFileExists{####1.def}}\1}}
\def\1#1#2#3#4#5#6#7#8{\@@tmp{#1}%
  {\XC@let@Nc\@@tmp{options@#2}
   \ifx\@@tmp\@empty\else
     \edef\@@tmp{,\@@tmp}\XC@replace\@@tmp,{;#1-}\XC@replace\@@tmp;,
   \fi
   \edef\@@tmp{:=#1\@@tmp}
   \expandafter\@for\expandafter\@tempa\@@tmp\do
    {\let\@@tmp\@tempa\XC@replace\@@tmp-,
     \edef\@@tmp{\@tempa\space\string\PassOptionsToPackage{\@@tmp}{xcolor}}
     \@namexdef{list@#4}{\@nameuse{list@#4},#3 \@@tmp}
     \@namexdef{list@#6}{\@nameuse{list@#6},#5 \@@tmp}
     \@namexdef{list@#8}{\@nameuse{list@#8},#7 \@@tmp}}}%
  {\typeout{No file `#1.def' found.}}}
\def\list@{}
\def\list@dvi{}
\def\list@ps{}
\def\list@pdf{}
\@namedef{options@0}{}
\@namedef{options@1}{kernelfbox,prologue,xcdraw,kernelfbox-xcdraw,prologue-xcdraw}
\0{dvips}1{latex}{dvi}{dvips}{ps}{ps2pdf}{pdf}
\0{pctex32}0{latex}{dvi}{dvips}{ps}{ps2pdf}{pdf}
\0{dvipdf}0{latex}{dvi}{dvipdf}{pdf}{}{}
\0{dvipdfm}1{latex}{dvi}{dvipdfm}{pdf}{}{}
\0{dvipdfmx}0{latex}{dvi}{dvipdfmx}{pdf}{}{}
\0{dvipsone}0{latex}{dvi}{}{}{}{}
\0{dviwin}0{latex}{dvi}{}{}{}{}
\0*{dviwindo}0{latex}{dvi}{}{}{}{}
\0{emtex}0{latex}{dvi}{}{}{}{}
\0*{oztex}0{latex}{dvi}{}{}{}{}
\0{pctexhp}0{latex}{dvi}{}{}{}{}
\0{pctexps}0{latex}{dvi}{}{}{}{}
\0{pctexwin}0{latex}{dvi}{}{}{}{}
\0{tcidvi}0{latex}{dvi}{}{}{}{}
\0{textures}0{latex}{dvi}{}{}{}{}
\0{truetex}0{latex}{dvi}{}{}{}{}
\0*{xdvi}0{latex}{dvi}{}{}{}{}
\0{pdftex}1{pdflatex}{pdf}{}{}{}{}
\0{vtex}0{vlatex}{pdf}{}{}{}{}
\0{xetex}0{xelatex}{pdf}{}{}{}{}
\newwrite\@tempa\immediate\openout\@tempa=\jobname.tmp
\@tfor\@tempb:={dvi}{ps}{pdf}\do
 {\edef\@@tmp{\@nameuse{list@\@tempb}}\edef\@@tmp{:=\expandafter\@gobble\@@tmp}
  \expandafter\@for\expandafter\@@tmp\@@tmp\do
   {\immediate\write\@tempa{\@@tmp}}}
\immediate\closeout\@tempa
\@@end}% end of driver list construction code
\edef\Options{}
\@tfor\@@tmp:={dvipdf}{dvipdfm}{dvipdfmx}{dvips}{dvipsone}{dviwin}{dviwindo}%
              {emtex}{oztex}{pctex32}{pctexhp}{pctexps}{pctexwin}{pdftex}%
              {tcidvi}{textures}{truetex}{vtex}{xdvi}{xetex}%
              {natural}{table}{fixpdftex}{hyperref}{showerrors}{hideerrors}%
              {override}{usenames}{nodvipsnames}{debugshow}\do
 {\expandafter\XC@replace\expandafter\Options\expandafter{\@@tmp,}{}}
\edef\Options{\XC@@drv,\Options}
\def\@@tmp#1,#2%
 {\ifx,#2#1\expandafter\remove@to@nnil\else#1', `\expandafter\@@tmp\fi#2}
\edef\Options{`\expandafter\@@tmp\Options,,\@nnil'}
\makeatother
\begin{document}
\vbox to\textheight{\vss\centering
{\Large\textbf{xcolor} driver test: \Options}
\vskip0.5cm
\begin{testcolors}[rgb,hsb,cmyk,gray,HTML]
\testcolor{magenta}
\testcolor[rgb]{1,0,1}
\testcolor[hsb]{0.83333,1,1}
\testcolor[cmyk]{0,1,0,0}
\testcolor[gray]{0.41}
\noalign{\medskip}
\testcolor{blue}
\testcolor[rgb]{0,0,1}
\testcolor[hsb]{0.66667,1,1}
\testcolor[cmyk]{1,1,0,0}
\testcolor[gray]{0.11}
\noalign{\medskip}
\testcolor{cyan}
\testcolor[rgb]{0,1,1}
\testcolor[hsb]{0.5,1,1}
\testcolor[cmyk]{1,0,0,0}
\testcolor[gray]{0.7}
\noalign{\medskip}
\testcolor{green}
\testcolor[rgb]{0,1,0}
\testcolor[hsb]{0.33333,1,1}
\testcolor[cmyk]{1,0,1,0}
\testcolor[gray]{0.59}
\noalign{\medskip}
\testcolor{yellow}
\testcolor[rgb]{1,1,0}
\testcolor[hsb]{0.16667,1,1}
\testcolor[cmyk]{0,0,1,0}
\testcolor[gray]{0.89}
\noalign{\medskip}
\testcolor{red}
\testcolor[rgb]{1,0,0}
\testcolor[hsb]{0,1,1}
\testcolor[cmyk]{0,1,1,0}
\testcolor[gray]{0.3}
\noalign{\medskip}
\testcolor{black}
\testcolor[rgb]{0,0,0}
\testcolor[hsb]{0,0,0}
\testcolor[cmyk]{0,0,0,1}
\testcolor[gray]{0}
\noalign{\medskip}
\testcolor{gray}
\testcolor[rgb]{0.5,0.5,0.5}
\testcolor[hsb]{0,0,0.5}
\testcolor[cmyk]{0,0,0,0.5}
\testcolor[gray]{0.5}
\noalign{\medskip}
\testcolor{white}
\testcolor[rgb]{1,1,1}
\testcolor[hsb]{0,0,1}
\testcolor[cmyk]{0,0,0,0}
\testcolor[gray]{1}
\noalign{\medskip\smallskip}
f(color)box &
\multicolumn5{c}{\def\i{\fbox{\a}}\def\ii{\fcolorbox{red}{yellow}{\a}}\def\iii{\scalebox{3}}%
\def\a{{\setbox0=\hbox to3ex{\hss\textcolor{blue}{\vrule width .75ex height1.5ex depth1.5ex}\hss}\ht0=0pt\dp0=0pt\box0}}%
\fboxsep0.333ex \i\quad\ii\quad\iii\i\quad\iii\ii\quad\fboxrule.75ex \iii\i\quad\iii\ii}
\end{testcolors}
\vskip0.5cm
{\footnotesize This table was generated with \textbf{xcolor4.tex} which is part of the \textbf{xcolor} package \XCfileversion{} (\XCfiledate) by Uwe Kern.}
\vss}
\end{document}
%</test4>
%
% \fi
\endinput
%%
%% End of file `xcolor.dtx'.
